% La Longue Marche — Volume 140-4 — Flash-Lite + mateo-canonical
% Model: Gemini 3.1 Flash-Lite (thinking=low)
% Prompt: mateo-canonical (notation + structure block, Part I few-shot pool)
% Pages transcribed: 280/280
% Pages overlaid from diagram-tikzcd branch: 58
% Rebuilt: 2026-04-19

%% ===== Page 1 =====

UNIVERSITÉ MONTPELLIER

IMAG
INSTITUT MONTPELLIÉRAIN
ALEXANDER GROTHENDIECK

La "Longue Marche" à travers la théorie de Galois

[suite à La "Longue Marche" à travers la théorie de Galois..., pages 650 bis à 787, dont
table des matières] : notes manuscrites (s.d.).

Cote n° 140-4, 279 pages

Alexander GROTHENDIECK

s.d.

\newpage

%% ===== Page 2 =====

\newpage
\thispagestyle{empty}

Pages 650--787 + table des matières

NB. La page 650 figure deux fois (erreur de numérotation)

\newpage

%% ===== Page 3 =====

\chapter*{Table des matières}
\addcontentsline{toc}{chapter}{Table des matières}

\noindent \textbf{50. Enveloppes schématiques des groupes discrets : cas anabélien} \hfill 650
\begin{enumerate}
\item[1)] Sections aux points \hfill 650
\item[2)] Essai de caractérisation des sections \hfill 651
\item[3)] Schémas ou groupes associés à des groupes discrets (noyau ob.) \hfill 651
\item[4)] Enveloppe schématique d'un groupe métabélien \hfill 651
\item[5)] Caractérisation universelle de l'enveloppe schématique \hfill 661
\end{enumerate}

\noindent \textbf{51. Exemple : La tour de Fermat.} \hfill 669
\begin{enumerate}
\item[1)] Le groupe $\mathcal{G} = \\operatorname{Sl}(2, \mathbf{Z}) / [\dots]$ \hfill 669
\item[2)] $Z_{p,\sigma}$ : version provisoire \hfill 673
\item[3)] Les groupes $\mathcal{G}_N$ et variantes \hfill 673
\item[4)] La suite $1 \to \mathcal{M}_{p,\sigma}^+ \to \mathcal{M}_{p,\sigma} \to \mathcal{D}_{12} \to 1$, et relations avec $\\operatorname{Sl}(2, \mathbf{Z}/ \dots)$ \hfill 677
\item[5)] $\mathcal{M}_{p,\sigma}, \mathcal{M}'_{p,\sigma}$ via $\mathcal{M}_{p,\sigma}$ \hfill 682
\item[6)] La suite $1 \to \mathcal{G}_\sigma \to Z_{p,\sigma} \to \mathcal{D}_{12} \to 1$, et $1 \to \mathcal{G}_{p,\sigma} \to \mathcal{D}_{12} \to 1$ \hfill 684
\item[7)] $\mathcal{N}_p, \mathcal{N}'_p$ etc. \hfill 686
\item[8)] $\mathcal{N}_p, \mathcal{N}'_p$ etc. \hfill 689
\item[9)] Retour sur $\mathcal{M}_{p,\sigma}, \mathcal{M}'_{p,\sigma}$ etc. \hfill 690
\item[10)] $\mathcal{M}_{p,\sigma}(0), \mathcal{M}_{p,\sigma}(1/2), \mathcal{M}_{p,\sigma}(-j)$ etc. \hfill 691
\item[11)] L'homomorphisme $\mathcal{M} \to \mathcal{M}_{p,\sigma} = \mathcal{M}_{p,\sigma}(\mathbf{Z})^*$, et relations avec la tour des courbes de Fermat : position du problème. \hfill 695
\item[12)] Scindages de l'extension $\mathcal{G}_\sigma / \mathcal{M}_\sigma$ de $\mathbf{Z}^* \times \mathfrak{S}_3$ par $\hat{\mathcal{M}}$. \hfill 700
\item[13)] Sections de la tour de Fermat, et relation de compatibilité avec $\mathcal{M} \to \mathcal{G}_{p,\sigma}$. \hfill 707
\item[14)] Relation aux jacobiennes généralisées. \hfill 710
\item[15)] Interprétation de $\mathbf{Q} \to \mathcal{N} \to \hat{\mathbf{Z}}$ \hfill 725
\item[16)] Réflexions sur le théorème de Fermat \hfill 733
\end{enumerate}

\newpage

%% ===== Page 4 =====

\section*{52. Enveloppes schématiques des groupes et \text{Teichm\"{u}ller}-Teichm\"{u}ller}

\begin{enumerate}
\item[1)] Notion de la \text{Teichm\"{u}ller}-groupe : cas discret \hfill 737
\item[2)] Cas profini \hfill 744
\item[3)] Description des lois de système de Galois-\text{Teichm\"{u}ller} \hfill 748
\item[4)] Apparents en \text{Teichm\"{u}ller}-droits par une apparition de la \text{Teichm\"{u}ller}-géométrie \hfill 758
\end{enumerate}

\section*{53. Le groupe de \text{Teichm\"{u}ller} spécial $\mathcal{SL}_{1,1} \simeq \mathcal{SL}(2, \mathbf{Z})$, et le groupe $\mathcal{S}_{1,1}$} \hfill 760

\begin{enumerate}
\item[1)] Action définie de la \text{Teichm\"{u}ller}-spécial (cabinet provisoire) \hfill 765
\item[2)] Le groupe $\mathcal{SL}(2, \mathbf{R}) \simeq \mathcal{SL}(2, \mathbf{R})$ \hfill 761
\item[3)] Le groupe $\mathcal{SL}(2, \mathbf{Z})$ \hfill 763
\item[4)] L'isomorphisme $\mathcal{GL}(2, \mathbf{Z}) \simeq \mathcal{S}_{1,1}$ \hfill 765
\item[5)] Digression sur les rev\^{e}tements des genres d'espaces, et relations entre genres typiques et genres d'efficacités des \text{Teichm\"{u}ller}-morphismes \hfill 768
\end{enumerate}

\hfill 691

\newpage

%% ===== Page 5 =====

50. Enveloppes schématiques de groupes discrets. (Cas anabélien).

1. Soit $G$ un schéma en groupes localement de type fini sur un schéma de base $S$ (on songe à $S = \operatorname{Spec} \mathbf{Z}$). Une section $g$ de $G$ sur $S$ est dite unipotente si, sur tout corps géométrique, liée $k \in \mathcal{G}(k^s)$ (où $k = \operatorname{kér}$)\marginpar{Sections unipotentes} est unipotente, i.e.\ s'il existe un plongement
$$
G_{\bar{k}} \hookrightarrow \mathbf{G}_{a, \bar{k}} \quad (k = \operatorname{kér}) \leqno{(1)}
$$
tel que
$$
g_{\bar{k}} = g_a(1) \leqno{(2)}
$$

On dit que $g$ est quasi-unipotente si pour tout $H \in S$, il existe $n \in \mathbf{N}^*$ tel que $g^n$ soit unipotente. Si le corps de base $k$ est pro, alors $g$ est d'ordre fini. On voit que si $S$ est quasi-compact, les $g$ quasi-unipotentes sont telles qu'il existe $n \in \mathbf{N}^*$ tel que $g^n$ sont unipotentes.

\textbf{Théorème.} Soit $G$ un groupe profini sur $S$ (quasi-compact). Pour que $g$ soit unipotente, il faut et il suffit qu'il existe un homomorphisme
$$
\mathbf{p}: \mathbf{G}_a \to G \leqno{(3)}
$$

\newpage

%% ===== Page 6 =====

tel que
$$
\varphi(1) = g \leqno{(4)}
$$
Je me demande si c'est aussi nécessaire, dans des "bons" cas. Prenons déjà le cas $S = \Spec K$. En ce cas, lorsque $G$ on a un résultat d'unicité pour $\varphi$, que pouvons-nous dire que les "bons" cas (...) :
Si l'on considère $S' \to S$ fidèle, il existe $\varphi'$ ayant la propriété, lorsque $\varphi$ existe (sur $S$ bien sûr). Dans ce cas on comprend le cas. C'est OK. En ce qui concerne l'effectivité précise des faits que j'ai dits $G_0 \to G_1$: des liens des groupes ($x \mapsto x^{\operatorname{fix}}$) qui à un élément ... fixe sur 1. On s'en tire.

Obligeons la descente (les quotients) qui sont mous. Sont-ils des sous-groupes de $G$?
Remarquons: $G_0$ est-il un groupe qui a une section ... triviale?
Est-ce dans un main?

Je n'ai pas envie de ... question (très intéressante!), et parfois

\newpage

%% ===== Page 7 =====

\noindent une question : on peut parler de cela
- que l'on ait, $S = \text{Spec } \mathbf{Z}$, un groupe $S$ multiplicatif avec des points génériques ?
En fait. Cela s'explique déjà d'après les
résultats sur un corps de base, qu'il
existe un ouvert $U$ et
$$
U \subset S \quad \text{et} \quad \mathcal{G}|U \to \mathcal{G}|U \leqno{(5)}
$$
des groupes que l'on ait $(4)$ sur $U$ et
d'après ce qui est unique. Par l'unicité,
on sait dire qu'il y a plus qu'un
le $\pi_1$, et la question ouverte :
nécessairement $U = S$.

Si on sait $U \neq S$, quel est le localisé
en un point maximal de $\mathcal{B}(U)$, pour se
remettre à examiner le cas $S$ local,
de point fermé $s$ - et : peut-on faire
dans ce cas, la preuve que sur $S$
l'action $\mathcal{G}$ est : une contradiction ... ). De
ce qui est intéressant, des règles de
détail, dans ce contexte on trouve un
trait. Je vais me borner à ce cas, initial.

\newpage

%% ===== Page 8 =====

sans doute rien aux évidences en général, et qui suffit pour mon propos.

Soit $y$ le point générique de $S$, $s$ un point fermé. Considérons le sous-groupe type $G_{\text{local}} \subset G_y$, soit $L$ un sous-schéma dans $G$, qui est un sous-schéma, n'étant pas un sous-groupe de $G$. Soit $L^\circ = L \setminus (\bigcup L_{s_i})$ pour prouver que $L^\circ \cong G_{\text{loc}}$. En fait, ... je pense que $(L^\circ)^\circ = (L_s)^\circ$ est n'étant pas...

Mais pour que $g : G_s \to G$ soit... pas nécessairement... sans doute que $L^\circ \cong G_s$. Car il est possible que $g$ ne soit pas mono, i.e. que $g_s$ ne soit pas mono. Il pourrait être l'homomorphisme trivial ($g_s = 1$) ! Mais même si $g_s \neq 1$, il n'est pas clair que $g_s$ soit mono.

Prenons l'exemple :
$$
G = \\operatorname{Aut}(M) \simeq \\operatorname{Aut}(S) \leqno{(6)}
$$
$M$ un modèle lisse de ... sur $S$, donc avec ...

$$
\varphi_g(\lambda) = \sum x_i \lambda^i \leqno{(7)}
$$
$x \in \text{End}(M)$, $K = k(y)$

\newpage

%% ===== Page 9 =====

est un endomorphisme nilpotent de $M_K$.

$$
x \in \End_A(M) \otimes K, \quad x^n = 0 \leqno{(8)}
$$

(A... de $S$, qui est une jauge). Donc par hypothèse
$$
\phi_y(1) = \sum \frac{x^i}{i!} \in \\operatorname{Aut}_A(M) \leqno{}
$$
i.e.
$$
\sum \frac{x^i}{i!}, \sum \frac{(-1)^i x^i}{i!} \in \End_A(M) \leqno{(9)}
$$
Et la question si $\phi_y$ provient d'un homomorphisme de groupes $\mathcal{G}_s \to G$
$$
\frac{x^i}{i!} \in \End_A(M) \quad (1 \le i \le n-1) \leqno{(10)} \marginpar{c'est une condition, c'est trivial}
$$

Pour $n=2$ c'est trivial, pour $n=3$ (dans (9) il y a un seul terme non trivial).
$$
\phi(x) = 1 + x + \frac{x^2}{2} \quad (X \in M_2(K), X^2 = 0)
$$
Pour $n=3$
$$
\phi_y(x) = 1 + x + \frac{x^2}{2} + \frac{x^3}{6} \quad (X^3 = 0, X \in M_3(K))
$$
... par hyp. $\phi_y(1)$ et $\phi_y(-1)$ dans $\\operatorname{Aut}_A(M)$
i.e.
$$
\frac{x^2}{2} \pm \frac{x^3}{6} \in M_4(A) \leqno{(11)}
$$
...
$$
\frac{x^2}{2}, \frac{x^3}{6} \in M_4(A) \quad ? \leqno{(12)}
$$
Si $2$ est inv. dans $A$ c'est clair, on tire de (11) que $\frac{x^2}{2} \in M_4(A)$ et $\frac{x^3}{6} \in M_4(A)$, OK. D'où le résultat.

\newpage

%% ===== Page 10 =====

En présence de corps résiduels 2. (Pour simplification par 2) Supposons un élément de $A$, un $X \in M_4(A)$
$$ X = \frac{1}{2} X_0 \quad (X_0 \in M_4(A), X_0 \notin 2 M_4(A)) $$
La donnée de $X$ équivaut à la donnée de $X_0 = (a_{ij})_{1 \le i, j \le 4}$, les $a_{ij} \in A$ pas tous divisibles par 2.
$$ X^4 = 0 \leqno{(11 \text{ bis})} $$
et la question est s'il résulte que
$$ \frac{X^2}{2}, \frac{X^3}{6} \in M_4(A) \quad ? \leqno{(12 \text{ bis})} $$

Bien sûr, la réponse est, si $X^3 \neq 0$, i.e. $g_y$ est ``nilpotent principal''.

\marginpar{Essai de construction des sections indépendantes?}
Pour en avoir le cœur net, il faudrait un contre-exemple dans les cas précédents — comme une matrice nilpotente d'ordre 4 — théorème ! Mais pour... il semble qu'il ne soit pas utile d'ouvrir une solution de ce pb.

Par contre, je vais énoncer le théorème (!) : Soit $G$ un schéma en groupes localement de présentation finie sur un schéma $S$ quasi-compact, et soit $g$ une section de $G$. Pour que $g$ soit quasi-nilpotent, il

\newpage

%% ===== Page 11 =====

Il faut et il suffit qu'il existe un entier $n > 0$, et une homomorphisme de groupes $\varphi: \mathcal{G} \to \mathcal{G}$, tel que soit
$$
\varphi(1) = g^n \leqno{(13)}
$$
Le cas $\mathcal{G}_{res}$ étant traité, il s'agit d'examiner la nécessité. Pour simplifier, on suppose $S$ schéma de base — un corps $S$ supposé noethérien, je pense que la réduction au cas noethérien doit se faire par des d'ossages faibles, que je n'ai pas envie de faire, faute de besoin essentiel! Je vais examiner le cas $S$ schéma, plus généralement :

a) $\mathcal{G}_{y} \Rightarrow \infty$. Alors l'hypothèse sur $g_y$ est que $g_y$ est d'ordre fini — donc $\exists n > 0$ tel que $g_y^n = 1$, donc $g^n = 1$. On prendra $\varphi = 0$, et c'est l'homomorphisme trivial.

b) $\mathcal{G}_{y}$ de $\mathcal{G}_{y}$ nul. Soit $\varphi$ :
$$
\varphi_y: \mathcal{G}_{y} \to \mathcal{G}_{y} \leqno{(14)}
$$
défini par la condition
$$
g_y = \varphi_y(1)
$$
sur $U_y$, suivant les définitions dans $S$.
Soit $\varphi_y(\lambda) = \lambda^n = \varphi_y(n \lambda) : \mathcal{G}_{y} \to \mathcal{G}_{y}$ et sur $U_y$ suivant les définitions. On a
$$
\varphi_y(1) = g^n \text{ sur } U_y \leqno{(15)}
$$

\newpage

%% ===== Page 12 =====

$$ n | n' \implies U_n \subset U_{n'} \leqno{(16)} $$

Il faut prouver que la réunion des $U_n$ est $S$. On a nommé $S$ local, et $\varphi_1$ est déjà défini dans
$$ U_1 = S \setminus \{s\}, \quad s \text{ le point fermé de } S \leqno{(17)} $$
et cela vient au

\textbf{Lemme (?)} Soient $S$ un schéma intègre, $G$ un schéma en groupes profini, lisse (pas fini) sur $S$, $U$ un ouvert non vide de $S$, $\varphi: G_{\bar{u}} \to G_u$ un homomorphisme, $\dots$ on suppose que $\varphi_1(1) = g^n$ et prolonge une section $g$ de $G$ sur $S$ entière (c'est-à-dire importante), d'après des résultats assez généraux de SGA3). Alors il existe un entier $n > 0$ tel que $\varphi_n: \lambda \mapsto \varphi_u(\lambda^n)$ se prolonge en un homomorphisme $\varphi: G_{\bar{u}} \to G$.

Je ne vais pas essayer de les démontrer dans ce cadre général des relèvements en général, et vais me limiter, à titre de test décisif (pour imposer les conditions que je demande) le cas où $G = G(\bar{S})$, comme tantôt, et celui-ci $S$ de plus est régulier et donc.

Donc on est ramené à la situation d'une jauge $A$, et de $X \in M_n(K)$ ($K$ corps de fonctions) [unclear: faisant].

\newpage

%% ===== Page 13 =====

$$
X^n = 0, \quad \sum_{i \ge 0} \frac{X^i}{i!} \in M_n(A), \quad \exp(X) \in M_n(A) \leqno{(18)}
$$
\marginpar{\raggedright Le cas de car. nulle se résout par l'affirmative à l'aide de la fonction logarithme : Posons $\log(1+Y) = \sum (-1)^{k-1} \frac{Y^k}{k}$...}
et on sent vraiment qu'il existe un entier $d > 0$
$$
\frac{d! X^i}{i!} \in M_n(A) \quad (1 \leq i \leq n-1) \leqno{(19)}
$$
Si $A$ est de caractéristique $0$, i.e.\ s'il est une $\mathbf{Q}$-algèbre, ceci implique d'ailleurs : $X \in M_n(A)$. Il faudrait que dans le cas
en $0$, c'est que lui qui puisse définir
... $S$ est entier. Le premier
on a -t-on si $X \in M_5(A)$, $X^5=0$
$$
\left\{ \frac{X^2}{2} \pm \frac{X^4}{24} + \dots \in M_5(A) \right\} \leqno{(20)}
$$
Par contre, dans le cas envisagé — au précédent
($n=4$, car.\ résiduelle $2$ — le seul qui pose des
difficultés) il est clair qu'il suffit de prendre
pour $d$ une puissance de $2$ assez grande.
Cette remarque s'applique aux autres
constituants. On devrait venir au fait.

Ceci
Corollaire au lemme de densité : si $S$ est local,
de car.\ résiduelle $p > 0$, et si $\varphi$ est définie
sur un ouvert qui contient $S - V(p)$
(i.e.\ $\varphi$ est complémentaire de $U$ et de
car.\ $p > 0$, il est inutile de dire sans que $S$ soit
local) alors $\exists q = p^n$ tel que $\varphi$ se prolonge.

\newpage

%% ===== Page 14 =====

Quand $G$ est affine, ceci devrait être trivial. On se ramène à

Proposition. Soit $S$ un préschéma [unclear: fin. qc?], $n$ un entier $\ge 1$, $S_0 = V(n) \subset S$ le sous-schéma fermé défini par $n.1_S$ (donc [unclear: crossed out text]), $U = S \setminus S_0$, $G$ un schéma en groupes affine sur $S$, et

(21) $$ \varphi_U : G_{a|U} \to G|_U $$

un hom. de groupes. On suppose $U$ schématiquement dense dans $S$. Alors il existe une puiss. $m = n^r$ de $n$, et un hom. de schémas en groupes

$$ \psi_m : G_a \to G , $$

tels qu'on ait

(22) $$ \psi_m|_U = (\varphi_U^m : \lambda \mapsto \varphi_U(m\lambda) = \varphi_U(\lambda)^m) $$

Considérons en effet l'endo.

$$ u_m : G_a \to G_a \qquad \lambda \mapsto m\lambda $$

Alors la formule (22) signifie qu'on a commutativité

\begin{tikzcd}
G_{a|U} \ar[r, "(u_m)_U"] \ar[dr, "\psi_m|_U"'] & G_{a|U} \ar[d, "\varphi_U"] \\
& G_U
\end{tikzcd}

et l'existence de $\psi_m$ pour un $m$ grand, signifie que pour un $m$ grand le hom. composé

$$ \varphi_U \circ u_m : G_{a|U} \to G_U $$

se prolonge en un hom de groupes $G_a \to G$.

\newpage

%% ===== Page 15 =====

Pour raisons de densité schématique et numérique, il suffit de poser que $\varphi_n$ se prolonge en une borne de schémas en groupes sur
$$
\text{Sch} \times \operatorname{Sch} \times \operatorname{Sch} \text{ en groupes}.
$$
Soit $A$ un anneau, et soit $B$ celui de $\mathcal{G}$. Donc l'anneau de $A$ est $A_n = A[T]$, celui de $\mathcal{G}$ est $B$ et $A_n = B_n$, par hypothèse, via un homomorphisme
$$
B_n \xrightarrow{\varphi_n} A_n[T] \leqno{(23)}
$$
tandis que l'homomorphisme correspondant $\varphi_n$ sera le produit des points
$$
A[T] \xrightarrow{\varphi_n^*} A_n[T] \leqno{(24)}
$$
$T \mapsto T$ de $T \mapsto m_i T^i$.

Soit $x = a_0 + a_1 T + \dots + a_n T^n \in A_n[T]$ dans $A_n$,
$$
\varphi_n^*(x) = \sum a_i m_i T^i = \sum (a_i m_i) T^i
$$
d'où si $m_n = 1$ en général,
$$
m_i \in A \text{ i.e. } m_i \in A[T] \text{ fix } \implies \varphi_n^*(x) \equiv a_0 \mod A[T]
$$
ce qui entraîne
$$
\varphi_n(x) \in A[T] \text{ polynomial } \iff a_0 \in A \leqno{(25)}
$$

\newpage

%% ===== Page 16 =====

Appliquons ceci aux 1-gerbes d'étale de $B$ qui engendrent $B$ comme $A$-algèbre...

Corollaire $S$ étant, $U$ étant comme des $S$-pro-finis, et $G$ étale un schéma affine (et de type fini $S$ étale-fini quasi-structure de groupes), et qui $E^U \to G$ un homomorphisme de schémas, pour qu'il existe (précisément un $S$-schéma $U$ en groupe tel que) donné par (24) se prolonge, il faut et il suffit que $G$ se prolonge en une rétraction de $S$.

Remarquons les arguments précédents trouvent un écho dans la "Théorie" de la déstabilité p. 6621 dans la on a un $\dots$ sur [unclear: la base] de $S$ avec [unclear: notation] de [unclear: base] $1$ (le cas intéressant pour nos états [unclear: être] de la géométrie $S$ de [unclear: G.A.G.A.]).

\marginpar{Schémas en groupes associés à des groupes discrets}
Soit maintenant $M$ un $\mathbf{Z}$-module libre de type fini, on va construire un pré-schéma en groupe sur $S$, associé au système projectif des sous-groupes $\mathbf{Z}$-modules d'indice fini de $M$. Pour faire la construction, je vais développer...

\newpage

%% ===== Page 17 =====

656

[MARGIN: Y a une complication sur l'existence d'une construction de [unclear] $M_{S_0}$ posera $\underline{G} = \underline{G}_{S_0} \times M_{S_0}$]

Lemme préliminaire, relatif à une extension

(26) $$1 \to M \to G \to H \to 1$$

suite exacte de groupes discrets, avec $M$ commutatif — le cas qui nous intéressera pour la suite étant celui où $M$ est un $\mathbb{Z}$-module libre de t. f. En tous cas $M$ est un $\mathbb{Z}$-module, et à ce titre définit un Groupe sur $\operatorname{Spec}\ \mathbb{Z}$, savoir 
$$A \longmapsto M \otimes_{\mathbb{Z}} A$$
($A$, $\mathbb{Z}$-algèbre commutative) — qu'on désigne par $\underline{M}$. En fait, $\underline{M}$ est représentable, i.e. un schéma en groupes sur $\mathbb{Z}$ si $M$ est libre de type fini. J'ai envie de définir aussi des schémas en Groupes sur $S_0$ (qui seront des schémas en groupes si $M$ l'est) $\underline{G}, \underline{H}$, et une suite exacte

(27) $$1 \to \underline{M} \to \underline{G} \to \underline{H} \to 1$$

de Groupes sur $S_0$, et un isomorphisme de suites exactes

(28)
\begin{tikzcd}
1 \ar[r] & M \ar[r] \ar[d, equal] & G \ar[r] \ar[d, "\simeq"'] & H \ar[r] \ar[d, "\simeq"'] & 1 \\
1 \ar[r] & \underline{M}(\mathbb{Z}) \ar[r] & \underline{G}(\mathbb{Z}) \ar[r] & \underline{H}(\mathbb{Z}) \ar[r] & 1
\end{tikzcd}

On prendra

(29) $$\underline{H} = H_{S_0} \text{ Groupe constant de val. } H$$

et les flèches verticales extrêmes de (28)

\newpage

%% ===== Page 18 =====

sont les isomorphismes canoniques.

On peut aussi définir $\underline{\mathfrak{S}}$ comme foncteur (en $A \in \mathbf{Z}$-algèbres) au-dessus de $S/S_0$, par
$$
\underline{\mathfrak{S}}(A) = \underline{\mathcal{G}}_M(A) \cong M \otimes_{\mathbf{Z}} A \leqno{(30)}
$$
$\underline{\mathcal{G}}_M(A)$ est donc défini comme l'cat. de $H$ par $M(A)$ s'induit de l'cat. $\underline{\mathfrak{S}}$ de $H$ par $M$ par l'hom.
$$
M \to M \otimes_{\mathbf{Z}} A \leqno{(31)}
$$
pour ainsi dire, et ainsi, au lieu de $H$-$\mathbf{Z}$-modules, on fait opérer $H$ sur $M \otimes_{\mathbf{Z}} A$ via son opération sur $M$.

La petite difficulté est que le foncteur ainsi défini par (30) n'est pas lui-même représentable — il ne donne l'équivalence correcte de $\underline{\mathfrak{S}}(A)$ que si $A$ est connexe — il correspond plutôt à un foncteur $S \mapsto \underline{\mathfrak{S}}(A) = \underline{\mathfrak{S}}(S)$ ($S = \operatorname{Spec} A$) faisant en plus des $\underline{\mathfrak{S}}$ sur $S$ qui
$$
S \to H_S \leqno{(32)}
$$
(qui est comme une section loc. constante sur $S$ au-dessus de $H$)

\newpage

%% ===== Page 19 =====

L'expression correcte s'obtient en partant des «foncteurs invariants»
$$
\underline{\mathcal{G}} : A \longmapsto \underline{\mathcal{G}}_M(M \otimes_{\mathcal{M}} A)
$$
ou mieux, des contrefoncteurs
$$
S \longmapsto \underline{\mathcal{G}}_M(S, M \otimes_{\mathcal{M}} \mathcal{O}_S),
$$
et en prenant le faisceau Zariski-noerien. En fait on a
$$
\underline{\mathcal{G}}(S) = \varprojlim_{\substack{\mathcal{U} = (U_i)_{i \in I} \\ (\mathcal{U} \in I \text{ partition} \\ \text{de } S \text{ schémas}) \\ S \cong \coprod U_i}} \prod_{i \in I} \underline{\mathcal{G}}(U_i) \leqno{(33)}
$$

Voici une autre façon de présenter la construction de $\underline{\mathcal{G}}$, en utilisant la construction générale du §..., dans le contexte faisceautique. On se donne un système
$$
M, \mathcal{G}_0 \rightrightarrows \mathcal{M}_0, \quad \text{Op. de } \mathcal{G}_0 \text{ sur } M \leqno{(34)}
$$
(avec en $S_0$, $\mathcal{G}_0$ (considéré) en $S_0$, $\mathcal{M}_0$ homomorphisme $\mathcal{G}_0 \to \mathcal{M}$)

Si l'op. de $\mathcal{G}_0$ sur $M$, qui induit un gr. de $\mathcal{G}$ sur $M$, est celle du l'op. de $\mathcal{G}$ sur $M$, par fonctorialité, on demande alors la compatibilité avec les loc. cit. Par la construction de loc. cit., on construit

\newpage

%% ===== Page 20 =====

À l'aide de données, un système
$$ \underline{\mathfrak{G}} \supset \underline{M}, \text{ et } \underline{\varphi} : \underline{\mathfrak{G}}_S \to \underline{\mathfrak{G}} \leqno{(35)} $$
Groupe sur $S_0$, $\underline{M}$ sous-groupe distingué.

Tel que (34) se reconstitue à partir de (35) en prenant $\underline{M}, \underline{\mathfrak{G}}_S$ dans (35), $M_S$ l'image inverse de $\underline{M} \subset \underline{\mathfrak{G}}$ par $\underline{\varphi}$, qui induit par $\underline{\varphi}$, et $\Theta$ l'opération induite de $\underline{\varphi}$. De plus, $\underline{\varphi}$ induit un isomorphisme.
$$ \underline{\mathfrak{G}}_S / M_S \isom \underline{\mathfrak{G}} / \underline{M} $$
$$ H_S $$
... et ... lie une description "axiomatique" des systèmes (35) en termes des (39). La suite exacte de groupes sur $S$\footnote{L'hypothèse que $H$ soit un [unclear: ...] de fait, [unclear: l'hypothèse ...] les $M$-torsions (plus générales) cas $S_0$ définis [unclear: ...] images [unclear: ...] triviaux.}
$$ 1 \to \underline{M} \to \underline{\mathfrak{G}} \to \underline{\mathfrak{G}}/\underline{M} \to 1 \leqno{(36)} $$
$$ H_S $$

$$ 1 \to \underline{M}(S) \to \underline{\mathfrak{G}}(A) \to H_S(S) \to 1 $$
$$ \qquad M \otimes_{\mathbf{Z}} A \qquad \qquad H_S $$

\newpage

%% ===== Page 21 =====

En fait, on peut [unclear: construire] un diagr.
commut. de suites exactes

(38)
\begin{tikzcd}
1 \ar[r] & \underline{M}_{S_0} \ar[r] \ar[d] & \underline{\mathfrak{S}}_{S_0} \ar[r] \ar[d] & \underline{H}_{S_0} \ar[r] \ar[d] & 1 \\
1 \ar[r] & \underline{M} \ar[r] & \underline{\mathfrak{S}} \ar[r] & \underline{H}_{S_0} \ar[r] & 1
\end{tikzcd}

qui par passage aux sections donne
pour $S$ quelconque affine (disons) connexe
un diagr. comm.

[MARGIN: NB En [unclear: déduisant] de (39)
$\underline{\mathfrak{S}}(A)$ est isomorphe
au produit semi-direct
de $H$ par
$M \otimes_{\mathbb{Z}} A$
[unclear]
[unclear]]

(39)
\begin{tikzcd}
1 \ar[r] & M \ar[r] \ar[d] & \mathfrak{S} \ar[r] \ar[d] & H \ar[r] \ar[d, equal] & 1 \\
1 \ar[r] & M \otimes_\mathbb{Z} A \ar[r] & \underline{\mathfrak{S}}(A) \ar[r] & H \ar[r] & 1
\end{tikzcd}

pour $S=S_0$ on retrouve l'iso (28)

- . - . - .

Cette construction du $S_0$-schéma en
groupes $\underline{\mathfrak{S}}$, à partir d'une
suite exacte (27), est [unclear: manifestement]
fonctorielle par r. à cette dernière : un
hom. de suites exactes

(40)
\begin{tikzcd}
1 \ar[r] & M \ar[r] \ar[d] & \mathfrak{S} \ar[r] \ar[d] & H \ar[r] \ar[d] & 1 \\
1 \ar[r] & M' \ar[r] & \mathfrak{S}' \ar[r] & H' \ar[r] & 1
\end{tikzcd}

donne naissance à un diagr. comm.
de suites exactes

\newpage

%% ===== Page 22 =====

(41)
\begin{tikzcd}[row sep=1.5em, column sep=2em]
& 1 \ar[r] & M_{S_0} \ar[r] \ar[dd] \ar[dl] & \underline{G}_{S_0} \ar[r] \ar[dd] \ar[dl] & H_{S_0} \ar[r] \ar[dd] \ar[dl] & 1 \\
1 \ar[r] & M'_{S_0} \ar[r] \ar[dd] & \underline{G}'_{S_0} \ar[r] \ar[dd] & H'_{S_0} \ar[r] \ar[dd] & 1 & \\
& 1 \ar[r] & \underline{M} \ar[r] \ar[dl] & \underline{\underline{G}} \ar[r] \ar[dl] & H_{S_0} \ar[r] \ar[dl] & 1 \\
1 \ar[r] & \underline{M}' \ar[r] & \underline{\underline{G}}' \ar[r] & H' \ar[r] & 1 &
\end{tikzcd}

[MARGIN: Exemple 4. réduction discrète des groupes métabéliens]

Reprenons donc un $\mathbb{Z}$-module libre
de type fini $M$, et pour $H$ un
groupe d'ordre fini, considérons
la suite exacte

(42) $$1 \longrightarrow M_i \longrightarrow M \longrightarrow M/M_i \longrightarrow 1$$

Plus gén., si on part d'une suite
exacte (27), pour un sous-groupe $M_i$
d'indice fini de $M$, stable par
$\underline{\underline{G}}$ i.e. stable par $H$, on a une
suite exacte

(43) $$1 \longrightarrow M_i \longrightarrow \underline{\underline{G}} \longrightarrow \underline{\underline{G}}/M_i \longrightarrow 1$$

qui montre que $H_i = \underline{\underline{G}}/M_i$
est l'extension de

(44) $$1 \longrightarrow M/M_i \longrightarrow H_i \longrightarrow H \longrightarrow 1$$

déduite de (27) par
$M \longrightarrow M/M_i$.

\newpage

%% ===== Page 23 =====

(51)
On note que $H$ est un groupe profini de $\mathcal{G}$, pour $H$ fini, les $M_i \subset M$ d'indice fini de $M$ stables par $H$ se finissent\footnote{Ou plutôt: "forment une base de voisinages de 0", en corrigeant: "prendre les $\lim M$!".} de tous les $M_i$ d'indice fini de $M$. Considérons la suite exacte de schémas de groupes associée
$$
1 \to M_i \to \mathfrak{S}_i \to H_i \to 1 \leqno{(45)}
$$
Pour $i$ variant, elles forment un système projectif de extensions de groupes sur $\mathbf{Z}$.

La procédure.

On considère le pro-objet $(\mathfrak{S}_i)$ comme pro-objet de $\mathcal{C}$ catégories des schémas en groupes (lions et $\dots$) sur $S_0$. On note que les morphismes de transition
$$
\mathfrak{S}_j \to \mathfrak{S}_i
$$
sont définis par quotient et, étant donnés l'homomorphisme\marginpar{ceci est une extension $E(M_i, M_j) \to E(M_i, M_k)$ avec $M_j \supset M_k$, ext. de $(M_i/M_j)_{S_0}$ par $M_j$} le lien entre les extensions de $H_i$ par $M_i$, les extensions de $H_j$ par $M_j$.

\newpage

%% ===== Page 24 =====

Or le premier est affine, comme somme finie de schémas affines, le deuxième est affine — donc le morphisme est affine.

Ceci étant, la limite projective des $\mathcal{G}_i$ dans la catégorie des schémas existe (plus précisément, le foncteur limite projective des $\mathcal{G}_i$ est représentable) — il n'est pas bon de dire : la limite existe, et de travailler dans la catégorie des schémas — pour le t.f.

Remarque. Le pro-groupe $(\mathcal{G}_i)$ sur $S_0 = \operatorname{Spec} \mathcal{G}$ est donné par des groupes discrets $\mathcal{G}_i$, munis d'une donnée de liasons permettant de travailler avec ces schémas. Il ne faut pas perdre de vue l'hypothèse faite que des schémas de type fini... [unclear: il y a que] (le paramètre assez petit, il sera libre de t.f.). Si $M_i$ est schéma d'indice fini dans $\mathcal{G}_i$ et $H_i$ sont finis, ...

\newpage

%% ===== Page 25 =====

alors ce produit ne dépend que de $\mathcal{G}$ lui-même — à condition qu'il existe un sous-groupe d'indice fini de $\mathcal{G}$ qui soit un groupe commutatif, et t.f. sur $\mathbf{Z}$ modulée. (On dit lors que $\mathcal{G}$ est "métabélien".)

Regardons la typique :
Ainsi pour $\mathcal{G}$ groupe discret métabélien, on construit un système projectif de schémas (affine), en groupes, que nous notons
$$
\operatorname{Ens}(\mathcal{G}) = (\mathcal{G}_i) \leqno{(46)}
$$
(l'image d'indice étant formée des sous-groupes distingués de $\mathcal{G}$ qui sont sans torsion). Ce produit dépend fonctoriellement de $\mathcal{G}$. Il a des propriétés d'exactitude très intéressantes et, il est multiplicatif
$$
\operatorname{Ens}(\mathcal{G} \times \mathcal{G}') \isom \operatorname{Ens}(\mathcal{G}) \times_{S_0} \operatorname{Ens}(\mathcal{G}') \leqno{(47)}
$$
Si on a une suite exacte
$$
1 \to \mathcal{G}'' \to \mathcal{G} \to \mathcal{G}' \to 1 \leqno{(48)}
$$
de groupes métabéliens, on a une suite exacte pour les morphismes
$$
1 \to \operatorname{Ens}(\mathcal{G}'') \to \operatorname{Ens}(\mathcal{G}) \to \operatorname{Ens}(\mathcal{G}') \leqno{(49)}
$$

\newpage

%% ===== Page 26 =====

un sous-groupe qui convient de préciser :
l'occasion — pour les schémas en groupe
associés (abéliens ou pas, à la limite) —
bel et bien un homomorphisme fidèlement plat de $S$.
La façon est de dire...

Le cas typique — que le cas général devra de façon opérer immédiatement, ce
celui-ci $\mathcal{G} = \mathcal{M} = \mathbf{Z}$. On — pour
l'entier $n \ge 1$ la suite exacte
$$ 1 \to \mathbf{Z} \xrightarrow{n} \mathbf{Z} \to \mathbf{Z}/n\mathbf{Z} \to 1 \leqno{(50)} $$
Je vais encore expliciter la construction
de l'extension correspondante de
$\mathbf{Z}/n\mathbf{Z}$ par $\underline{\mathbf{G}} = \mathbf{G}_m$. Soit $\mathcal{G}_n$.
C'est un schéma affine, dont l'anneau
de coordonnées est l'anneau — de façon par
bien connus — un produit de
$n$ copies de celui de $\underline{\mathbf{G}} = \mathbf{G}_m$, savoir
$\mathbf{Z}[T]$ — un terme avec les isomorphismes
en divisant aux sections assujetties à
$\mathbf{Z}$ sur $\mathbf{Z}/n\mathbf{Z}$, par exemple par l'une
des action $\mathbf{Z}_n[0, n-1]$. Les liens de...

\newpage

%% ===== Page 27 =====

transition entre les $\mathfrak{S}_n$ n'est pas plat bien sûr. P. ex. : la limite projective

$$
(\hat{\mathbf{Z}})_{so} = \varprojlim_n (\mathbf{Z}/n\mathbf{Z})_{so}
$$

— schéma — groupe profini vu comme la limite projective des schémas
$\mathfrak{S}_n$, s'envoyant dans la limite.
Un des faits de 20) — c'est le spectre
de l'anneau limite inductive des
$$
\mathbf{Z}[T_n] \subset \mathbf{Q}[T]
$$
$$
T_n = \frac{1}{n} T
$$
qui est, donc etc., la $\mathbf{Z}$-algèbre
de $\mathbf{Q}[T]$ formé des polynômes
dont le terme constant est non nul.

On a une description [unclear: claire] de
la limite projective de l'anneau
scindé $\mathbf{Z}^\wedge$ — c'est une
extension de $(\hat{\mathbf{Z}})_{so}$ [unclear: avec] des composantes,
qui est le spectre du sous-anneau
$\mathbf{Q}[T_n, \dots, T_m]$ formé des polynômes dont
le terme constant ...

\newpage

%% ===== Page 28 =====

C'est dans une situation où [les] groupes sont [unclear: scindés] de toutes les fois qu'il existe un groupe fini, i.e. le plus générique, qui est $\Gamma_{\mathbf{Q}}$.

$\bullet$ Dans le cas d'un groupe arbitraire $\mathcal{G}$, l'image de $\mathcal{G}$ dans le complété profini $\hat{\mathcal{G}}$ de $\mathcal{G}$, [qui est] le groupe profini sur $\mathbf{Z}$.
$$ \mathcal{G}_{\sigma} = \varprojlim_i (\Pi_i) \leqno (50) $$
On aura alors une suite exacte canonique
$$ 1 \to \operatorname{Eno}^0(\mathcal{G}) \to \varprojlim \operatorname{Eno}(\mathcal{G}) \to (\hat{\mathcal{G}}) \to 1 \leqno (51) $$
$\operatorname{Eno}^0(\mathcal{G})$ désigne le complété profini des $\varprojlim \operatorname{Eno}(\mathcal{G})$, i.e.
$$ \operatorname{Eno}^0(\mathcal{G}) \defeq \varprojlim_i M_i \simeq \varprojlim_M N_M \leqno (52) $$
où les $M_i$ sont les sous-groupes d'indice fini dans $\mathcal{G}$ qui sont de type libre fini. [...] pour un $M$ plus haut, $\dots$
$$ \underline{M} = \check{V}(M) = \operatorname{Spec} \operatorname{Sym}(\check{M}) \leqno (53) $$
désigne la fibre vectorielle [?] $M$ qui définit la...

\newpage

%% ===== Page 29 =====

$M_0$ est un sous-groupe de $M_i$, dans la limite projective (cf. le calcul des limites des $M_i$) des $M_i$. Le choix d'une base de $M_0$ sur $\mathbf{Z}$ donne un isomorphisme
$$
\operatorname{End}^0(\mathcal{G}) \isom \operatorname{End}^0(M_0) \isom (\operatorname{End}^0(\mathbf{Z}))^n \isom (\mathcal{E}')^n \leqno{(59)}
$$
(où $\mathcal{E}' \defeq \operatorname{End}^0(\mathbf{Z}) \isom \operatorname{Spec}(\mathbf{Z}[T, T^{-1}])$ \marginpar{où $\mathcal{E}'$ représente le schéma en groupes à lacets, qui représente le foncteur-groupes $A \mapsto \mathcal{E}'(A) = \{x \in A^* \mid x_n = x^m, \forall m, n \in \mathbf{N}\}$}) est le schéma en groupes à lacets, qui représente le foncteur-groupes
$$
A \mapsto \mathcal{E}'(A) = \{x \in A^* \mid x^n = x^m, \forall n, m \in \mathbf{N}\} \leqno{(56)}
$$
[Si $A$ est tel qu'il existe $n \in \mathbf{N}$ avec $x^n = x^m$ pour tout $x \in A^*$, pour les composantes de $A$ divisées par $n$, il existe des éléments de $A$ tels que les composantes de $A$ soient telles que pour $x \in A^*$, $\mathcal{E}'(A) = \{1\}$ --- si l'on restreint $A$ à $0$, $\mathcal{E}'(A) = A$.]

\newpage

%% ===== Page 30 =====

5. Construction universelle de l'enveloppe pro-relativiste d'un groupe semi-bilinéaire.

Soit $G$ un schéma en groupes lisse et pris fidèle sur une base $S$. On dit qu'un schéma en groupes $H \to G$ est pro-unipotent (resp. quasi-unipotent) si l'image de $H$ est un idéal unipotent (resp. pro-unipotent) de $G$. Ce qui signifie faire jouer les homomorphismes d'un groupe discret
$$ \mathfrak{G} \to \Gamma(G, S) $$
et dit unipotent (resp. quasi-unipotent) si l'image de l'homomorphisme de schémas en groupes $\mathfrak{G} \to G$ est unipotent (resp. quasi-unipotent).

\newpage

%% ===== Page 31 =====

Notons que d'après les propriétés des limites de systèmes projectifs filtrants (à morphismes de transition affines) entre schémas, si $\mathcal{G}$ est arbitraire, d'où un système projectif de schémas en groupes affines sur $S$, disons $E_{\text{aff}}(\mathcal{G})$, et une limite projective
$$
\text{LEns}(\mathcal{G}) = \text{LEns}(\mathcal{G}) \times_{S_0} S \simeq \varprojlim \operatorname{Ens}(\mathcal{G}) \leqno{(57)}
$$
--- dès que $\mathcal{G}$ est local de présentation finie
$$
\\operatorname{Hom}_{\text{progr}/S}(\operatorname{Ens}(\mathcal{G}), \mathcal{G}) \isommap \\operatorname{Hom}(\operatorname{Ens}(\mathcal{G}), \mathcal{G}) \defeq \varprojlim \\operatorname{Hom}(\mathcal{G}_i, \mathcal{G}) \leqno{(58)}
$$

Notons d'autre part qu'il existe un homomorphisme canonique
$$
\mathcal{G}_S \to \operatorname{Ens}(\mathcal{G}) \leqno{(59)}
$$
déduit par passage à la limite de la base $S_0$ à la base $S$
$$
\mathcal{G}_{S_0} \to \operatorname{Ens}(\mathcal{G}) \leqno{(59 \text{ bis})}
$$
provenant du système de limites canoniques.

\newpage

%% ===== Page 32 =====

\section*{}
$$
\mathfrak{C}_{S} \longrightarrow \mathfrak{C}
$$
Dans, via (59), on trouve
$$
\\operatorname{Hom}_{\text{pro-gr}}(\varprojlim \mathfrak{E}\text{ns}(\mathfrak{C}), G) \xrightarrow{\sim} \\operatorname{Hom}_{\text{gr}}(\mathfrak{C}_{S}, G) \cong \\operatorname{Hom}_{\text{gr}}(\mathfrak{C}, G(S)) \leqno{(60)}
$$
Composant (58) et (60), on trouve
$$
\\operatorname{Hom}_{\text{pro-gr}}(\mathfrak{E}\text{ns}(\mathfrak{C}), G) \cong \\operatorname{Hom}_{\text{gr}}(\varprojlim \mathfrak{E}\text{ns}(\mathfrak{C}), G) \longrightarrow \\operatorname{Hom}_{\text{gr}}(\mathfrak{C}, G) \cong \\operatorname{Hom}_{\text{gr}}(\mathfrak{C}, G(S)) \leqno{(61)}
$$
Ceci posé, on a le résultat suivant:

\textbf{Théorème.} Sous les conditions précédentes ($\mathfrak{C}$ catégorie discrète multiplicative, $G$ schéma en groupes localement de type fini sur $S$), l'application canonique est injective, et comme image... \leqno{(62)}
plus, pour des tels foncteurs $\mathfrak{C} \to G(S)$ dans une variété de $\mathfrak{C} \to G(S)$, des vrais foncteurs... pour voir simplement l'existence, que $S$ est multiplicative et $G$ affine sur $S$.

\newpage

%% ===== Page 33 =====

\section*{50. Exemple : La ``tour'' de Fermat.}

\subsection*{1. Les sous-groupes $\mathcal{B}, \mathcal{C}$.}

Je m'intéresse au quotient de $\mathcal{G} \defeq \Gamma_0(2)$ obtenu en prenant le quotient par $[\Pi_0, \Pi_0]$ on \dots un \dots
\marginpar{On engendre $\Pi_0$ par les $\lambda_i = \begin{cases} \lambda_0 = \begin{pmatrix} 1 & -2 \\ 0 & 1 \end{pmatrix} \\ \lambda_1 = \begin{pmatrix} 1 & 0 \\ 2 & 1 \end{pmatrix} \end{cases}$}
$\mathcal{G}$ de $\mathcal{G}/\Pi_0$ par le groupe
$$
M \defeq \Pi_0^{\mathrm{ab}} \isom (\hat{\mathbf{Z}}(0) + \hat{\mathbf{Z}}(1) + \hat{\mathbf{Z}}(\infty)) / \text{diagonale} \leqno (1)
$$
Ainsi on a une suite exacte de groupes profinis
$$
1 \to M \to \mathcal{G} \to \mathcal{G}/\Pi_0 \to 1 \leqno (2)
$$
où $\mathcal{G}/\Pi_0 \isom \\operatorname{Sl}(2, \mathbf{Z}/4\mathbf{Z}) / \langle \lambda_0, \lambda_1 \rangle$.

Par générateurs et relations \dots

$$
\mathcal{G} = \left\{ \tau, \rho, \sigma \mid \begin{array}{l} \rho^6 = \sigma^4 = \tau^2 = 1 \\ \rho^3 = \sigma^2, \tau(\rho) = \rho^{-1}, \tau(\sigma) = \sigma^{-1} \\ [\lambda_0, \lambda_1] = 1, \dots \end{array} \right\} \leqno (3)
$$
$$
\left\{ \begin{array}{l} \lambda_0 = \rho = \sigma \rho \sigma = \sigma \rho \sigma^{-1} \\ \lambda_1 = \dots = \rho \sigma \rho = \rho \sigma \rho^{-1} \end{array} \right.
$$
d'où la relation de commutateurs $[\lambda_0, \lambda_1] = 1$ s'écrit aussi
$$
\sigma \rho \sigma \rho \sigma \rho = \rho \sigma \rho^3 \rho \leqno (4)
$$
$$
\sigma \rho \sigma \rho^{-1} \sigma \rho = \rho \sigma^{-1} \rho \sigma \leqno (5)
$$
On le sens des règles \dots \dots \dots les formules des générateurs et relations en supposant des \dots \dots \dots et en remplaçant les \dots \dots (\dots \dots des $\sigma$ de la formule \dots des $\sigma^{-1}$. Par exemple, de \dots

\newpage

%% ===== Page 34 =====

\section*{50. Exemple : La «tour» de Fermat (suite)}
\addcontentsline{toc}{section}{50. Exemple : La «tour» de Fermat}

\noindent
(6) $\lambda_\infty = \rho \sigma \rho^{-1} \sigma \rho = \rho (\sigma \rho^{-1} \sigma \rho) = -\rho \sigma \rho^{-1} \rho = \rho \sigma \rho^{-1} \sigma$
\vspace{1em}
\noindent
(7) $[\sigma, \lambda_\infty] = 1, \quad \begin{cases} [\lambda_0, \lambda_1] = (\sigma \rho^{-1})^3 \\ [\lambda_0, \lambda_\infty] = (\sigma \rho^{-1})^3 = \text{int}(\rho^{-1} \sigma \rho) [\lambda_0, \lambda_1] \end{cases}$ \footnote{NB On trouve $[\lambda_0, \lambda_1] = (\sigma \rho^{-1})^3 = [\sigma, \lambda_\infty] = \dots$}

On définit un schéma de groupes $\widehat{M}$ sur $\mathbf{Z}$, des triplets sans pb des $\widehat{\mathfrak{S}}^1$. On pose
$$
M = \widehat{\mathfrak{S}}^1 = \mathfrak{S}^{0,3} \text{ / diagonale} \leqno{(8)}
$$
\marginpar{Il y a un sous-groupe $M$ dont on calcule le quotient.}

On va faire opérer $\mathfrak{G}$.
$$
(9) \quad \mathfrak{S}_{0,3} \isom \\operatorname{Sl}(2, \mathbf{Z}) / [\pi_0, \pi_0]
$$
\noindent
On décrit ainsi un sous-groupe $M$,
$$
\mathfrak{G} \times \mu_2 \isom \mathfrak{G}/M \isom \mathfrak{S}_{0,3}/\pi \isom \mathfrak{T}_{0,3} \leqno{(5)}
$$
Désignons par $\lambda_0, \lambda_1, \lambda_\infty$ les images des $\lambda_0, \lambda_1, \lambda_\infty = -\lambda_i$ dans $\mathfrak{G}$, identifiés à des éléments de $M = \widehat{\mathfrak{S}}^1 = \widehat{\mathfrak{S}}^1(\hat{\mathbf{Z}})$.

\noindent
(6) $\lambda_0, \lambda_1, \lambda_\infty \in M(\mathbf{Z}) \subset M(\hat{\mathbf{Z}}) = M \otimes \hat{\mathbf{Z}}$. On a les additions des groupes, d'où on déduit des actions sur $M = \widehat{\mathfrak{S}}^1$ \marginpar{Alors $\mathfrak{G} \dots$}.

\noindent
(7) $\sigma(\lambda_0) = \lambda_1, \quad \sigma(\lambda_1) = \lambda_\infty, \quad \sigma(\lambda_\infty) = \lambda_0$

\noindent
(8) $\tau(\lambda_0) = \lambda_1, \quad \tau(\lambda_1) = \lambda_0, \quad \tau(\lambda_\infty) = \lambda_\infty$ \footnote{On a utilisé $\sigma$ pour "$\tau$" et $\tau$ pour "$\tau'$"} \footnote{ce sont des éléments de $\mathfrak{G}$ par définition.}

\newpage

%% ===== Page 35 =====

L'extension $\mathfrak{S}$ annoncée est bien sûr celle de \S~30, développé entre-temps (depuis les premières lignes du présent \S, constatant justement que j'étais gêné aux entournures pour expliciter la construction) :
$$
1 \to M \xrightarrow{\parallel} \mathfrak{S} \to H_{S_0} \to 1 \leqno{(9)}
$$
$$
\check{V}(M) \qquad \text{où } S_0 = \Spec \mathbf{Z}.
$$
Pour tout schéma $S$, on pose
$$
\mathfrak{S}^\Delta(S) \defeq \mathfrak{S}(S) \leqno{(10)}
$$
(invariance de la "diagramme" $\Gamma(H_S/S)$ (i.e. diagramme à $H \to \Gamma(H,S)$) par l'homomorphisme canonique $\mathfrak{S}(S) \to H_S(S)$), de sorte que par la construction des loc.\@ cit., on a une isomorphisme canonique
$$
\mathfrak{S}^\Delta(S) \isom \mathfrak{S}_M \Gamma(S, \mathbf{Q} \otimes_{\mathbf{Z}} M) \leqno{(11)}
$$
et une suite exacte canonique
$$
1 \to \Gamma(S, \mathbf{Q} \otimes_{\mathbf{Z}} M) \to \mathfrak{S}^\Delta(S) \to H \to 1 \leqno{(12)}
$$
Posant $A \defeq \hat{\mathbf{Z}}$, on trouve en particulier (cela est une fait général sur les schémas en groupes associés aux extensions de groupes discrets à noyau abélien)
$$
\mathfrak{S}^\Delta \isom \mathfrak{S}^\Delta(\hat{\mathbf{Z}}) \leqno{(13)}
$$

\newpage

%% ===== Page 36 =====

2. Le groupe $\mathcal{Z}_{p,\sigma}$ : version provisoire.

Je vais maintenant expliciter pour $\hat{\mathcal{G}}$ et $\hat{\mathcal{G}} \to \mathcal{G}$ les constructions générales de $§ 49$ II, pour $\hat{\mathfrak{S}}$.
$$
\hat{\mathcal{G}} \defeq M = \hat{\mathcal{G}}^0 \quad (\text{composée des éléments de } \hat{\mathcal{G}}) \leqno{(14)}
$$

Dans un cas comme celui-ci, $\mathcal{Z}_{p,\sigma}$ justement le $\pi$-groupe formé des $(p,\sigma)$-groupes ($\mathcal{Z}_{p,\sigma} \subset \hat{\mathcal{G}}$ ici) engendrés par $\lambda_0, \lambda_1, \dots$ et disposés dans les relations — là il n'y a pas d'ambiguïté — d'où un bon $\mathcal{Z}_{p,\sigma}$.

\marginpar{Je ferai ici, $\mathcal{Z}_p$ "fixe" bien sûr, compte tenu des conditions locales $\alpha, \beta \dots$}

Le pas essentiel est de définir le groupe $\mathcal{Z}_{p,\sigma}$.

$$
\mathcal{Z}_{p,\sigma} \defeq \{ u \in \\operatorname{Aut}(\hat{\mathcal{G}}^+) \mid 
\begin{cases}
a) \; \exists \alpha \in \hat{\mathcal{G}}, \; u(p) = \alpha p \\
b) \; \exists \beta \in \dots, \; u(\sigma) = \beta \sigma \\
c) \; u(\mathcal{L}_0) = \mathcal{L}_0 \\
d) \; (\text{automatique}) \; u(\hat{\mathcal{G}}) \subset \hat{\mathcal{G}}
\end{cases}
\} \leqno{(15)}
$$

Notons que $u$ est déterminé par les valeurs
$$
\{ u(p) = \rho, \quad u(\sigma) = \sigma \eta \} \leqno{(16)}
$$
avec $\rho, \eta \in \Pi$ (où $\Pi = M$, $M$ sous-groupe d'indice 2 dans $\Pi$).

\newpage

%% ===== Page 37 =====

a) $\xi$ (localement) conjugue $\rho$ via $\mathcal{G} \supset M$
b) $\dots$ $\sigma \dots$
c) égalité des $\mathcal{G}$-orbites
$$
\xi = \eta \lambda_1 \quad \text{où } \eta \in \Gamma \text{ convenable (tel que } H \ni \dots \subset \Gamma_M) \leqno{(17)}
$$
Dans ce groupe $\Gamma_N$, il faudra...

Remarque : Il n'est plus possible ici de désigner par $-1$ l'élément de centre de $\mathcal{G}^+$, et de désigner par $-x$ pour $x \in \mathcal{G}^+$ le produit de $-1$ par celui-ci, à cause de la confusion avec la notation $-x$ (groupe, selon $M$) pour l'image de $x$ dans le groupe quotient $M$ pour l'action de $x$, mais l'élément $\omega x \neq -x$, puisque $\omega \notin M, -x \notin M \dots$

On pose bien sûr
$$
\varepsilon_{\mathcal{G}}(M) = \varepsilon_{\mathcal{G}}(\alpha) \ , \quad \varepsilon_{\mathcal{G}}(u) = \varepsilon_{\mathcal{G}}(\beta) \leqno{(18)}
$$
Bien définis, ces $\mathcal{G}$-orbites sont $\mathcal{G}$-conjuguées sur $\dots$ fibre, i.e. pour $\sigma, \sigma^{-1}$, comment les expliquer en termes de $\dots$ ?

$$
\begin{cases}
\xi = \alpha \rho \alpha^{-1} \in \Gamma_M \times \alpha \in \Gamma_M = \Gamma_\mathcal{G} \text{ i.e. } \varepsilon_{\mathcal{G}}(\alpha) = 1 \\
\quad = \alpha \tau \rho(\alpha^{-1}) \dots = \alpha \sigma (\rho(\alpha)^{-1}) \rho(\alpha) \in \Gamma_M \quad \lambda_1^{-1} \in M \\
\eta = \beta \sigma \beta^{-1} \in \Gamma_M \times \beta \in \Gamma_M = \Gamma_\mathcal{G} \text{ i.e. } \varepsilon_{\mathcal{G}}(\beta) = 1 \\
\quad = (\beta \sigma) \sigma \beta^{-1} = \omega \beta \sigma \beta^{-1} \in \omega M \subset \Gamma_M
\end{cases} \leqno{(19)}
$$

donc en tous cas
$$
\begin{cases}
\xi \in \Gamma_M \\
\eta \in \Gamma (\omega^{1/2} M)
\end{cases} \leqno{(19 \text{ bis})}
$$

\newpage

%% ===== Page 38 =====

\marginpar{675}
\begin{enumerate}
\item[a)] $\xi_{\beta}$ (localement) conjugués : $\rho$ via $\Gamma \subset \mathcal{M}$
\item[b)] $\sigma$
\item[c)] équation des lacets
\end{enumerate}

$$
\xi = \eta l_1 \quad \text{où } \eta \in \Gamma \text{ convenable (tel que } H = 2\nu + 1 \in \Gamma_{\mathcal{M}} \text{)} \leqno{(17)}
$$

Remarquez: Il n'est plus possible de distinguer par $-1$ l'élément $\omega$ tiré du centre $\rho$ de $\mathcal{G}$, et de distinguer par $-x$ pour celui-ci, à cause de la relation $\omega \cdot x \in \mathcal{M}$, le groupe cyclique $\mathcal{M}$ pour l'ensemble des $x \dots$ mais l'identité $\omega x \neq -x$, puisque $\omega x \notin \mathcal{M}$, $-x \in \mathcal{M} \dots$

On pose bien sûr:

$$
\xi_3(n) = \xi_3(n), \quad \xi_2(n) = \xi_2(n) \leqno{(18)}
$$

C'est défini, ces groupes étant conjugués sur un fibré, i.e. par $\sigma, \sigma^{-1}$, comme les expliquées en termes de $\dots$

On a:
$$
\begin{cases}
\xi = \alpha \rho \alpha^{-1} \in \Gamma_{\mathcal{M}} \text{ et } \alpha \in \Gamma_{\mathcal{M}} = \Gamma \quad \text{i.e. } \xi_3(n) = 1 \\
\eta = \beta \sigma \beta^{-1} \in \Gamma_{\mathcal{M}} \text{ et } \beta \in \Gamma_{\mathcal{M}} = \Gamma \quad \text{i.e. } \xi_3(n) = 1
\end{cases} \leqno{(19)}
$$

donc en tous cas:
$$
\begin{cases}
\xi \in \Gamma_{\mathcal{M}} \\
\eta \in \Gamma(\omega^{1/2} \mathcal{M})
\end{cases} \leqno{(19\text{ bis})}
$$

\newpage

d'où, tiré par l'équation des lacets (17)

$$
l_1 = (\omega) \in \Gamma(\omega^{1/2}(n) \mathcal{M}) \quad \text{i.e. } \omega \in \Gamma(\omega^{1/2}(n) \mathcal{M}) \leqno{(20)}
$$

On pose l'écran sous la forme

$$
\xi = 2\xi' + \xi_2 \leqno{(21)}
$$

et

$$
\begin{cases}
\xi_2, \xi_3 : \mathbf{Z}_{p^n} \to (\mathbf{Z}/2\mathbf{Z})_{S_0}
\end{cases} \leqno{(22)}
$$

et des morphismes de schéma bien définis $\dots$

par les formules (21) et

$$
\begin{cases}
\mu = 2\nu + 1 \ (= 4\nu + 2\nu_2 + 1) \\
\xi_2 = (-1)^{\xi_2}, \quad \xi_3 = (-1)^{\xi_3}
\end{cases} \leqno{(23)}
$$

On a ainsi, en outre, dans (24)

$$
l_1 = l_1^{2\nu} l_1^{\nu_2} = (\omega) l_1^{\nu_2} = \dots \leqno{(24)}
$$

Comme $l_1 = \omega$, on ...

$$
l_1 = a_1 \omega \leqno{(25)}
$$

posant ... écrit (24) ...

$$
\xi = (\lambda_0, a_1) \omega^{1/2} \dots \in \Gamma_{\mathcal{M}} \leqno{(26)}
$$

Pour expliciter les $a_i$, il faut les développements ... linéaires de $\lambda_0, \lambda_1, \lambda_n \dots$

\newpage

%% ===== Page 39 =====

$$ l_1^V = \lambda_1^V \omega^V. \leqno{(24)} $$
Posons (cf. (19)):
$$ \begin{cases} \beta = \beta_0 \tau^V \\ \eta_0 = \beta_0 \sigma(\beta_0)^{-1} \end{cases} \quad (\beta \in \mathcal{M}) $$
on a $\eta = \beta \sigma(\beta)^{-1} = \beta_0 \sigma(\beta_0)^{-1} \omega^V = \eta_0 \omega^V$.

L'équation des lacets prend la forme d'une équation dans $\mathcal{M}$ :
$$ \eta = \eta_0 \lambda_1^V, \leqno{(26)} $$
$$ \eta = \eta_0 \lambda_1^V \leqno{(26 \text{bis})} $$
(indiquant deux sections de $\mathcal{M}$ qui se déterminent mutuellement, $\eta$ donné). On a ainsi exprimé de façon commode l'équation des lacets, i.e. le condition (17), il reste à exprimer a) et b), que ces formules $\eta$ sont (avec $\eta_3=0, \eta_3=1$) :

\begin{enumerate}
\item[a)] $\eta$ locales de la forme $\alpha \sigma(\alpha)^{-1} = \alpha \sigma(\lambda_1^V \sigma(\alpha)^{-1})$ non $\in \mathcal{M}$
\item[b)] $\eta_0$ locales de la forme $\beta_0 \sigma(\beta_0)^{-1}$
\end{enumerate}

(La condition que la formulation est indépendante de $\sigma, \tau$ etc., mais dépend uniquement de $\eta_3$). Pour expliciter ces conditions, on écrit $\mathcal{M}$ additivement, les relations s'écrivent :
$$ \begin{cases} \eta = \alpha_0 - \rho(\alpha_0) \\ \eta_0 = \beta_0 - \sigma(\beta_0) \end{cases} \quad (\eta_3=0) \quad (\eta_3=1) \quad (\eta = \partial_0 - \rho(\partial_0) - \lambda_1 (\eta_3=1)) \leqno{(28)} $$
On a donc pour $\alpha = \alpha_0 \tau^V$.

\newpage

%% ===== Page 40 =====

ce qui s'exprime ainsi :
$$
\begin{cases} \zeta_0 \in \operatorname{Im}(\operatorname{Id}_M - \rho_M) & (\psi_0 = 0) \\ \eta_0 \in \operatorname{Im}(\operatorname{Id}_M - \sigma_M) \end{cases} \leqno{(29)}
$$
C'est le moment de passer à des calculs matriciels explicites, en identifiant $M$ à $\mathbf{Z}^2$ via $\lambda_0, \lambda_1$. Dans cette base, les matrices de $\rho_M, \sigma_M, \tau_M$ sont données par :
$$
\rho_M = \begin{pmatrix} 0 & 1 \\ -1 & -1 \end{pmatrix}, \quad \operatorname{id}_M - \rho_M = \begin{pmatrix} 1 & -1 \\ 1 & 2 \end{pmatrix} \leqno{(30)}
$$
$$
\sigma_M = \begin{pmatrix} 0 & 1 \\ 1 & 0 \end{pmatrix}, \quad \operatorname{id}_M - \sigma_M = \begin{pmatrix} 1 & -1 \\ -1 & 1 \end{pmatrix}
$$
$$
\tau_M = \begin{pmatrix} -1 & 0 \\ 0 & -1 \end{pmatrix}
$$
On a
$$
\det(\operatorname{id}_M - \rho_M) = 3, \quad \det(\operatorname{id}_M - \sigma_M) = 0
$$
ce qui montre que les endomorphismes $\operatorname{id}_M - \rho_M$ et $\operatorname{id}_M - \sigma_M$ de $M$ sont, i.e. des isomorphismes — et ce n'est pas faire des isomorphismes —\footnote{Note-t-il : "soit distinctes les valeurs 3 pour le premier, 0 pour le second".}

On aurait pu être tenté de prendre comme "primitives" non $\zeta_0$ et $\eta_0$, mais bien convenablement $\alpha_0$ et $\beta_0$ (au plus des "primitives" discrètes $\lambda_2, \lambda_3$ et $\lambda'$), et observer donc ce qui est possible... avec les systèmes.

\newpage

%% ===== Page 41 =====

$$
\mathfrak{Z}_{g,\nu} = \{ (\nu, \nu', \nu_2, \nu_3, \alpha_0, \beta_0) \mid \nu, \nu_2, \nu_3 \in \{0, 1\} \subset \mathbf{Z}_2, \, \alpha_0, \beta_0 \in M, \dots \leqno (31)
\nu_2 + \nu_3 = 2\nu + 1 \in \mathbf{Z}_2 \}
\alpha_0 - \rho(\alpha_0) \nu_2 \lambda_1 = \beta_0 - \sigma(\beta_0) + \nu_3 \lambda_1
$$

\marginpar{NB. L'élimination qui suit est redondante.}

L'équation des lacets, $\mathfrak{Z}_{g,\nu} \ni (\alpha_0, \beta_0)$, prend donc la forme
$$
(\rho_M - \operatorname{id}_M) \alpha_0 - (\operatorname{id}_M - \sigma_M) \beta_0 = (2\nu + 1 - \nu_2 - \nu_3) \lambda_1 \leqno (32)
$$

Remarquons que le schéma $\mathfrak{Z}_{g,\nu}$ lui-même, tel qu'il a été défini au § 49 (II), n'est pas lisse, même sur le corps de base $\mathbf{Q}$. Il semble se composer en composantes irréductibles distinctes, plus fines, qui sont celles de l'espace des lacets $(\alpha_0, \beta_0)$ qui semble bien lisse. Ce qui suggère : nous pouvons, après un peu de présentation des constructions du § 49, en utilisant la géométrie du groupe $\mathfrak{Z}_{g,\nu}$ (a-t-on égard \dots \textit{lacets} $\xi, \eta$, et en les remplaçant par un groupe qui approche assez les \textit{points} rationnels du $\dots$ de $M_{g,\nu}$ (ou $M_{g,\nu}$)). On fera une présentation des $\dots$ au sens du \S suivant, et \dots -apparemment déjà esquissé cette construction.

\marginpar{Cela doit être $\operatorname{Norm}_M(\dots)$ et $(\alpha_0, \beta_0) \mapsto \dots$ Mais TS ?}

\newpage

%% ===== Page 42 =====

D'où le groupe (schéma de groupe) $\tilde{\mathfrak{S}}$ avec une structure de groupe commutatif, et opérant sur $\mathfrak{S}^+$ via
$$
\begin{cases}
u(\rho) = (\alpha_0 - \rho(\alpha_0) + \nu_3 \lambda_1)\rho \\
u(\sigma) = (\beta_0 - \sigma(\beta_0))\nu_2^{-1}\sigma
\end{cases} \leqno{(33)}
$$

Mais il faut quand même vérifier que les formules (33) déterminent bel et bien un automorphisme de $\mathfrak{S}^+$. La vérification essentielle : faire en sorte que la relation (7)
$$
[\sigma, \lambda_0] \leqno{(34)}
$$
et respectées par $u$.

Je vais expliciter l'action de $u$ sur $\mathcal{M}$, en explicitant ses valeurs sur
$$
\begin{cases}
\lambda_0 = \omega \lambda_0 = \omega \sigma \rho \\
\lambda_1 = \omega \lambda_1 = \omega \rho \sigma \rho
\end{cases}
$$
$$
\begin{cases}
\xi = \alpha - \rho(\alpha) + \nu_3 \lambda_1 \\
y = \beta_0 - \sigma(\beta_0)
\end{cases} \in \mathcal{M} \leqno{(35)}
$$

Lis :
$$
u \rho = \xi \rho, \quad u \sigma = \omega^{1/2} y \sigma \leqno{(36)}
$$
On [unclear: détermine] dans [unclear: $\tilde{\mathfrak{S}}$] par le relèvement des lacets ($\mu = 2\nu_1 = 4\nu' + 2\nu_2 + 1$)
d'où

\newpage

%% ===== Page 43 =====

$$
u(l_0) = u(\epsilon^2) = l_0(\epsilon^2) \marginpar{(69)}
$$
et
$$
u(\lambda_0) = u(\omega l_0) = u(\omega) u(l_0) = \omega l_0(\epsilon^2) = \omega (\lambda_0 \lambda_0^{-1}) = \lambda_0 \dots
$$
comme il convient. (Pour me rassurer, j'ai vérifié que la deuxième formule (33) donne bien $u(\sigma) = (y_0 \sigma(y_0))^{\dots} = \omega \dots$)
$$
u(\omega) = \omega \leqno{(67)}
$$
On a
$$
u(\lambda_1) = u(\rho \lambda_0) = (u(\rho) u(\lambda_0)) = (\rho \lambda_0) = \lambda_1
$$
des finis, comme prévisible.
$$
u(\lambda) = \mu \lambda \text{ pour } \lambda \in M \leqno{(38)}
$$
\marginpar{$\mu = 2N+1 = 4N' + 2N_2 + 1$}
On voit bien que $u(\sigma) = (y_0 \sigma(y_0))^{\mu} \sigma$.

Tous ces calculs sont peu rigoureux, logiquement, car à la fois la justification de l'existence d'un bon relèvement dans les calculs précédents par l'identité déjà connue, dans un groupe "plus gros", à savoir le groupe "plus gros", à savoir le groupe schématique engendré par $\rho, \sigma \dots$

Une façon de les rendre rigoureux serait d'utiliser la description $\mathfrak{S}$-adique de $\mathfrak{S}^+$ :
$$
\underline{\mathfrak{S}}^+ = \mathfrak{S}^+ \wedge_{M} \mathfrak{S}_0 \leqno{(39)}
$$
de définir un homomorphisme de $\mathfrak{S}^+$ sur un sous-groupe $G$ (i.e. $G = \mathfrak{S}^+$) revient à

\newpage

%% ===== Page 44 =====

définir des homomorphismes
(40) $$M_{S_0} \xrightarrow{u_M} g^\natural, \quad \underline{\mathfrak{S}}^+_{S_0} \xrightarrow{u_{\underline{\mathfrak{S}}}} g^\natural \quad (\text{i.e. } \underline{\mathfrak{S}} \to g^{c.s.})$$
satisfaisant les deux conditions suivantes :

a) Commutativité du diagramme

(41)
\begin{tikzcd}
M_{S_0} \ar[r] \ar[d] & \underline{\mathfrak{S}}^+_{S_0} \ar[d, "u_{\underline{\mathfrak{S}}}"] \\
M_{S_0} \ar[r, "u_M"'] & g
\end{tikzcd}

et b)

(42) $u_M$ est compatible avec les actions
de $\underline{\mathfrak{S}}^+_{S_0}$ sur $M_{S_0}$ (via son action sur $M_{S_0}$)
et sur $g$ (via $u_{\underline{\mathfrak{S}}}$).

Dans le cas naturel, on aura $g = \underline{\mathfrak{S}}^+$

(43) $u_M(\lambda) = i(\mu . \lambda)$

où $i : M \longrightarrow \underline{\mathfrak{S}}^+$
est l'inclusion évidente, et

(44) $u_{\underline{\mathfrak{S}}} : \underline{\mathfrak{S}}^+_{S_0} \longrightarrow g \quad \text{défini par (33)}$

Il faudra d'abord vérifier que (33)
définit bien un hom. $\underline{\mathfrak{S}}^+_{S_0} \longrightarrow g$,
ce qui revient à $= \underline{\mathfrak{S}} \to g^{c.s.}$.

\newpage

%% ===== Page 45 =====

Pour ceux qui utilisent la présentation usuelle de $\mathfrak{S}^+$ (ici $\mathfrak{S}_{0,3}^+$)
$$
\mathfrak{S}^+ = \{ \rho, \sigma \mid \rho^6 = \sigma^4 = 1, \rho^3 = \sigma^2, [\sigma, \lambda_0] = 1 \text{ où } \lambda_0 = \sigma \rho \sigma^{-1} \rho^{-1} \} \leqno{(29)}
$$
en posant $\rho' = u_M(\rho), \sigma' = u_M(\sigma)$ (cf (33)), il faut vérifier les relations
$$
\rho'^6 = \sigma'^4 = 1, \rho'^3 = \sigma'^2 = \lambda'_\infty \leqno{(46)}
$$
On commence par vérifier que
$$
\rho'^3 = \sigma'^2 = \omega \leqno{(47)}
$$
ce qui implique $\exists \mathfrak{S}^+ \to \mathcal{G}$ ($\rho \to \rho', \sigma \to \sigma'$), d'où
$$
\rho'^6 = \sigma'^4 = 1 \leqno{(48)}
$$
Ensuite, pour vérifier que $\lambda'_\infty = u_M(\lambda_\infty)$, il suffit de faire vérifier
$$
\lambda'_\infty = \mu \lambda_\infty \leqno{(49)}
$$
Il suffit de vérifier que $\lambda'_\infty$ commute : $\sigma' = u(\sigma) = \rho_0 \sigma$ puis puisque $\sigma$ et $\rho_0 \in M$ commutent. D'où la construction $u_M$, les vérifications essentielles (47) et (49) — d'ailleurs une fois (47) posés, les calculs précédents aboutissent à (38) et justifiés (là où il faut, i.e., $u_M \in \Lambda M$).

\newpage

%% ===== Page 46 =====

donc (49) sera bien vérifié.

On a
$$
\rho^3 = (\rho \rho^2 \rho^3)^3 = (\rho \rho^2) \omega
$$
$$
\sigma^2 (y_0 \sigma y_0) \sigma^2 = (y_0 \sigma y_0) \omega
$$
d'où les formules (47) s'induisent :
$$
[\rho \rho^2 \rho] = 1, \quad [y_0 \sigma y_0] = 1 \leqno{(50)}
$$
soit, définissons
$$
[\rho \rho^2 \rho] = 0, \quad y_0 + \sigma(y_0) = 0
$$
dont la vérification est triviale, constatant que
$$
\rho_M^3 = \sigma_M^2 = \operatorname{Id}_M, \quad (1 + \rho_M + \rho_M^2) = 0_M
$$
On a donc bien construit un homomorphisme $u_G: \mathcal{G}^+ \to \hat{\mathcal{G}}(S_0)$, et il reste à vérifier (41) et (42).

Pour (41), nous avons déjà noté l'existence d'un $u_G$. En voici :
$$
u_G(\lambda_0) = \dots \lambda_0, \quad u_G(\dots) = \dots
$$
Je vais, pour le faire, la vérification, pour être sûr.
$$
\varepsilon_0 = u_G(\varepsilon_0) = u_G(\sigma) u_G(y_0) = \dots = \rho \dots = \dots \sigma = \dots
$$
$$
\sigma^\rho = y_0 \sigma w^{-1/2} \varepsilon_0
$$

\newpage

%% ===== Page 47 =====

Je dis que
$y'_0 = \varepsilon_0^{(M-1)}$ i.e. $y_0 \sigma(y) \omega^{V_2} = \underbrace{\varepsilon_0^{(\mu-1)}}_{l_0^{\nu'}} \implies \mu-1=2\nu_1$, soit
$y_0 \sigma(y) \omega^{V_2} = l_0^{\nu'}$ ; $\nu = 2\nu'$ i.e.
$y_0 \sigma(y) = l_0^{2\nu'} (\omega l_0 \lambda_0^{V_2}) = \underbrace{\lambda_0^{2\nu'+V_2}}_{\lambda_0^{2\nu'} \lambda_0^{-V_2}}$ i.e.
$y_0 \sigma(y) = \lambda_0^{\nu'}$ soit, en appliquant $\sigma$ et notant
que $\sigma(y_0) = y_0^{-1}$, $\sigma^2(y)=y$, $\sigma(\lambda_0)=\lambda_1$,
$y_0^{-1} y = \lambda_1(y)$ soit enfin $y = y_0 \lambda_1(y)$ $\sim$
$y = y_0 + \nu \lambda_1$ (écrit additivement)

ce qui est vrai par hypothèse (cf (31) ).
Reste à [unclear: s'assurer] (pour achever [unclear: l'étude] de $u_{\mathcal{E}_0}$ pour
les [unclear: éléments annulés] (pour [unclear: abréger on l'écrira simplement])
$u(\lambda_1) = u\rho(\lambda_0) = (\rho \mu)(\lambda_0) = \rho(u(\lambda_0)) = \rho(\lambda_0 \mu) = \lambda_1 \mu$
OK.

Reste à vérifier (42), ce qui [unclear: se réduit] à :
commutativité de $u_M = \mu\ id_M$ à
$(\rho_M, \rho'_M)$, $(\sigma_M, \sigma'_M)$ - ce qui signifie que
$\rho_M$ et $\rho'_M$ d'une part, $\sigma_M$ et $\sigma'_M$ d'autre part,
agissent [unclear: également], ce qui est trivialement évident.

Ainsi, le groupe $\tilde{\mathcal{E}}_{\rho, \sigma}$ est bel et
bien construit - du moins, en tant que
relevant un [unclear: sous-groupe] $Z$, muni des [unclear: morphismes]
de structure
(51)
\begin{tikzcd}
\tilde{Z}_{\rho, \sigma} \ar[r, shift left, "f_0"] \ar[r, shift right, "f_{S_0}"'] & \mathcal{E}_{0, S_0} \ar[r] & \mathcal{E}_{S_0}
\end{tikzcd}
d'où $\mu = 4\nu' + 2 + V_2 : \tilde{Z}_{\rho, \sigma} \to \mathcal{E}_{S_0}$
$\varepsilon_2, \varepsilon_5 : \tilde{Z}_{\rho, \sigma} \rightrightarrows \mu_{S_0} = \mu\ id_{S_0}$

\newpage

%% ===== Page 48 =====

et un morphisme de faisceaux sur $\widetilde{Z}_{g,\nu}$
$$
\widetilde{Z}_{g,\nu} \to \\operatorname{Aut}_{gr}(\mathfrak{S}^+) \leqno{(51)}
$$
$$
\widetilde{Z}_{g,\nu} \to \\operatorname{Aut}_{gr}(\mathfrak{S}^+) \times G_m \times \mu_{s_0} \times \mu_{s_0} \leqno{(53)}
$$
$$
u \mapsto u, \mu(u), \varepsilon_2(u), \varepsilon_3(u)
$$

Des calculs déjà faits traitant à la fois des extensions diverses donnent\marginpar{un peu compliqué car $\widetilde{Z}_{g,\nu} \to G_m \times \mu_{s_0} \times \mu_{s_0}$ n'est pas, en fait, un produit $M \times M$ (le $\mu$ dans $\mu \times \mu$ est un peu différent, voir 3-\#55).} que ce morphisme (cf. \#33) est un morphisme de groupes, pour une structure de groupes convenable sur $\widetilde{Z}_{g,\nu}$ qui est déterminée, l'examen du \#13, et conduit par voie d'une démonstration géométrique. On aura
$$
1 \to \widetilde{Z}_{g,\nu}^{+0} \to \widetilde{Z}_{g,\nu} \to G_m \times \mu_{s_0} \times \mu_{s_0} \to 1
$$
explicitant $\widetilde{Z}_{g,\nu}^{+0}$ comme sous-schéma en groupes
$$
\widetilde{Z}_{g,\nu}^{+0} = M \times M \times G_a = \{ (\alpha_0, \beta_0, \gamma) \mid \alpha_0 \delta = 0, (1+\rho_M)\alpha_0 = (1-\sigma_M)\beta_0 + 2\gamma'\lambda_1 \}
$$
On aura donc
$$
1 \to \dots \to \widetilde{Z}_{g,\nu}^{+0} \to G_a \to 1 \leqno{(55)}
$$

\newpage

%% ===== Page 49 =====

$$
\tilde{\mathfrak{T}}_{0,3}^{+0} = \{ u \in \tilde{\mathfrak{T}}_{0,3} \mid \nu_1(u) = \nu_2(u) = 0 \} \leqno{(57)}
$$
$$
= \{ (\alpha_0, \beta_0) \in \Gamma(M \times M) \mid (1-\rho_M)\alpha_0 = (1-\sigma_M)\beta_0 \}
$$

\marginpar{en cas 3 de notre section sur le lien 2-3 de la section}

Schéma en groupes $\tilde{\mathfrak{T}}_{0,3}^{+0}$ est bien en cas \#3 de la relation 2-3 de la façon conjuguée, sous le module $M$.

$$
\tilde{\mathfrak{T}}_{0,3}^{+0} \isom \mathfrak{V}(\mathfrak{T}) \leqno{(58)}
$$

P est le produit fibré.
$$
0 \to M \to M \times M \to M \to 0 \leqno{(59)}
$$
$(\alpha_0, \beta_0) \mapsto (1-\rho_M)\alpha_0 - (1-\sigma_M)\beta_0$

Le homomorphisme $M \times M \to M$ s'explicite :
$$
\begin{cases} \alpha_0 = a_0 + a_1 \lambda_1 \\ \beta_0 = b_0 + b_1 \lambda_1 \end{cases} \quad d'où \quad \begin{cases} a_0 - \rho a_0 = (a_0 + a_1) \lambda_0 + (2a_1 - a_0) \lambda_1 \\ \beta_0 - \sigma \beta_0 = (b_0 - b_1) \lambda_0 + (b_0 - b_1) \lambda_1 \end{cases} \leqno{(60)}
$$

$$
(\alpha_0, \beta_0) \mapsto (a_0 + a_1, b_1 + b_0) \lambda_0 + (2a_1 - a_0 - (b_1 - b_0)) \lambda_1 \leqno{(61)}
$$

C'est un homomorphisme surjectif ; les éléments de l'image sont $c_0 \lambda_0 + c_1 \lambda_1$ avec
$$
c_0 + c_1 \equiv 0 \pmod{3} \leqno{(62)}
$$

Remarques : Si on prend un [unclear: ] sur $\hat{\mathbf{Z}}$ provenant d'un élément de $M[0]$, alors...

\newpage

%% ===== Page 50 =====

\setcounter{page}{49}
\noindent
$\mu \in \hat{\mathbf{Z}}^*$, $\mu \equiv (-1)^{\nu_3} \equiv 1-2\nu_3 \pmod 3$, et la condition sur $(\mu, \nu_1, \nu_3)$ pour être dans l'image de $\widetilde{\mathcal{Z}}_{\sigma} \to \Aut \times \Aut \times \Aut$, qui s'écrit en général comme dans $(62)$, est
$$
2\nu_1 - \nu_2 + \nu_3 \equiv 0 \pmod 3 \leqno{(3)}
$$
disons, en multipliant par 2:
$$
4\nu_1 - 2\nu_2 + 2\nu_3 \equiv 0 \pmod 3 \leqno{(3)}
$$
i.e. $\mu \equiv 1 - 2\nu_3 \pmod 3$, et elle est automatiquement satisfaite.

C'est pourquoi j'ai mis des indices. La spécification de $\widetilde{\mathcal{Z}}_{\sigma}$ est:

\marginpar{$\mu = 1 - \nu_2 \ (4)$ \\ $\mu = 1 - 2\nu_3 \ (6)$ \\ $\mu = 1 - 2\nu_3 + 6\nu'' \ \text{i.a.}$}

$$
\widetilde{\mathcal{Z}}_{\sigma} \defeq \left\{ (\nu_1', \nu_2'', \nu_3, \alpha_0, \beta_0) \in \mathcal{G} \times \mathcal{G} \times \hat{\mathbf{Z}} \times M \times M \mid 
\begin{cases}
\mu = 4\nu_1' - 2\nu_2'' + 1 = 6\nu'' - 2\nu_3 + 1, \quad 2\nu_1' + \nu_3 - \nu_2'' = 3\nu'' \\
(1-\rho_M)(\alpha_0) - (1-\sigma_M)(\beta_0) = 3\nu'' \lambda_1
\end{cases}
\right\} \leqno{(63)}
$$

\noindent
i.e. $2c = 6\nu''$, et ça fait $c = 3\nu''$.

\newpage

%% ===== Page 51 =====

\section*{3. --- LES GROUPES $\mathfrak{T}_{0,3}$ ET $\mathcal{O}_N^*$}
\addcontentsline{toc}{section}{3. Les groupes $\mathfrak{T}_{0,3}$ et $\mathcal{O}_N^*$}
\section*{}

On va introduire une construction auxiliaire, une définition, un schéma, un groupe qui comprend le domaine de quantités $\mu \in \mathcal{G}_m, \nu', \nu'' \in \mathcal{G}_m, \nu_1, \nu_2 \in \{0, 1\}$, satisfaisant :
$$
\mu = 3\nu' - 2\nu_2 + 1 = 6\nu'' - 2\nu_3 + 1 \leqno{(63)}
$$
De façon générale, soit $N$ un entier fixé. Le cas qui nous intéresse est le cas $N=12$, et considérons le schéma (ou extension de groupes discrets) :
$$
1 \to \mathbf{Z} \to \dots \to \mathbf{Z}/N\mathbf{Z} \to 1 \leqno{(64)}
$$
qui s'insère dans une suite exacte :
$$
1 \to E' \to \mathcal{O}_N \to (\mathbf{Z}/N\mathbf{Z})_{S_0} \to 1 \leqno{(65)}
$$
Rappelons la définition simple
$$
\mathcal{O}_N = (\mathbf{Z})_{S_0} \rtimes (N\mathbf{Z}) \leqno{(66)}
$$
Mais je dis qu'il y a dans la structure d'anneau, la structure additive et la structure multiplicative, on obtient ainsi : Une section de $\mathcal{O}_N$ \dots
$$
(n + N\lambda) + (n' + N\lambda') = (n+n') + N(\lambda + \lambda') \leqno{(67)}
$$
$n + N\lambda = n' + N\lambda' \iff n = n', \lambda = \lambda'$.

\newpage

%% ===== Page 52 =====

On multiplie alors les sections par
$$
(n + N\lambda)(n' + N\lambda') = nn' + N[n\lambda' + n'\lambda + \lambda\lambda'] \leqno{(68)}
$$
on vérifie que c'est commutatif avec relations (57b), que c'est distributif, associatif, que $1+N[0]$ est unité... Autre façon de voir la structure d'anneau.

$\mathfrak{O}_N = \mathbf{Z}_{S_0} A$ une $\mathbf{Z}$-algèbre finie munie d'un idéal $J$ qui soit un $\mathbf{Z}$-module libre de type fini, contenant l'idéal $NA$ et la structure d'extension
$$
0 \to J \to A \to A/J \to 0 \leqno{(69)}
$$
d'une algèbre de groupe connexe $A[N]$, que j'obtiens comme quotient de $(A)_{S_0} \times J$ semi-direct
$$
0 \to J \to A_S \times J \to A_N \to 0 \leqno{(70)}
$$
$$
\lambda \mapsto (\lambda, -\lambda)
$$
et on trouve que $A_S \times J$ a une structure d'algèbre produit semi-direct, et que

\newpage

%% ===== Page 53 =====

On multiplie alors les sections par
$$
(u + N\lambda)(u' + N\lambda') = uu' + N[u\lambda' + u'\lambda + \lambda\lambda'] \leqno(68)
$$
on vérifie que c'est bien compatible avec les relations (57 b), que c'est bijectif, que c'est unitatif, associatif, que $1 + N\lambda$ est unité... Autre façon de voir la structure d'anneau. On sait plus généralement $A = \mathbf{Z}_{S_0} A$ une $\mathbf{Z}$-algèbre munie d'un idéal $J$ qui soit un $\mathbf{Z}$-module libre de type fini, condition annulant l'idéal $NA$ et la structure d'extension
$$
0 \to J \to A \to A/J \to 0 \leqno(69)
$$
d'où un schéma en groupes comme $A[J]$, qui s'obtient comme quotient de $A_{S_0} \times J$ (semi-direct).
$$
0 \to J \to A_{S_0} \times J \to A_N \to 0 \leqno(70)
$$
avec $\lambda \mapsto (\lambda, -\lambda)$
et on trouve que $A_{S_0} \times J$ est un anneau produit semi-direct, et que

L'image de $J \to A$ est un idéal...
On trouve bien une structure d'anneau sur $A[N]$, et on peut écrire
$$
1 \to J \to A[N] \to B_{S_0} \to 1 \leqno(71)
$$
où $A[N] \to B_{S_0}$ est un homomorphisme de schémas en anneaux. Passant aux groupes multiplicatifs, on trouve une suite exacte
$$
1 \to (1+J)^* \to A[N]^* \to (B_{S_0})^* \to 1 \leqno(72)
$$
où $A[N]^*$, $(B_{S_0})^*$ et $B^*$, le groupe des points... $\times$-groupe multiplicatif et $(1+J)^*$ a la sous-extension
$$
(1+J)^* \cong 1+J \leqno(73)
$$
$$
\{ (1+\lambda) \mid 1+\lambda \text{ inversible dans } A \}
$$
$$
\{ (\lambda, \lambda') \in J \times J \mid \lambda + \lambda' + \lambda\lambda' = 0 \} \text{ dans } J
$$
Ainsi, dans le cas particulier $A = \mathbf{Z}, J = N\mathbf{Z}$
$$
1 \to (1+N\mathbf{Z})^* \to \mathbf{Z}^* \to (\mathbf{Z}/N\mathbf{Z})^*_{S_0} \to 1 \leqno(74)
$$

\newpage

%% ===== Page 54 =====

$(1+NZ)^*$ est défini comme
$$
(1+NZ)^* \defeq \{ \lambda \in \mathbf{Z}/NZ \mid 1 + N\lambda \text{ inversible} \} \leqno{(75)}
$$
On a une tour canonique (en partie).
Les $\mathcal{O}_N$ pour $N$ variable forment un système projectif d'anneaux, les $\mathcal{O}_N^*$ forment un système projectif de schémas en groupes commutatifs.
Je n'ai pas pu regarder de plus près la structure de ces systèmes projectifs, ni comment la limite projective — question qui se pose plus généralement que pour d'un $A$ comme ci-dessus. On a explicité alors des morphismes canoniques
$$
\mathcal{O}_N^* \to \mathcal{O}_1^* = \mathcal{O}^* = \mathbb{G}_m \leqno{(76)}
$$
D'une façon qui est, en fait, une... Le dernier système des
$$
\{ 1 + N\mathcal{J} \mid 1 + N\lambda = 1 \text{ i.e. } \lambda = 0 \} \leqno{(77)}
$$
qui est donc un sous-schéma de la structure de groupes de $\mathcal{O}_N$.

\newpage

%% ===== Page 55 =====

qui est un schéma en groupes complet ($N \neq 1$) qui ne nul sur $\mathbf{Z}[1/N]$, et satisfaisant aux fibres en $p|N$ $\simeq \mathbf{G}_m$.

Je ne rends compte maintenant que c'est une modulation de celles des groupes $\mathcal{O}_N^*$, qui se fait avec des rappels : $\mathbf{G}_m$.

Je rends en particulier examine :
$$
1 \to (1+12\mathbf{Z})^* \to \mathcal{O}_{12}^* \to (\mathbf{Z}/12\mathbf{Z})^* \to 1 \leqno{(78)}
$$
Notons que
$$
(\mathbf{Z}/12\mathbf{Z})^* \simeq (\mathbf{Z}/4\mathbf{Z})^* \times (\mathbf{Z}/3\mathbf{Z})^* = \mu_2 \times \mu_2.
$$
De $(78)$ pour le faisceau, on tire une section, en fait on a une factorisation
$$
1 \to (1+12\mathbf{Z})^* \to \mathcal{O}_{12}^* \xrightarrow{\varepsilon_2 \times \varepsilon_3} \mu_2 \times \mu_2 \to 1 \leqno{(79)}
$$
Les groupes $\mu_2, \mu_2$ sont des groupes définis par
$$
\mathcal{O}_{12}^* \to (\mu_2, \mu_3) \to \{0,1\}_{S_0} \times \{0,1\}_{S_0} \leqno{(80)}
$$
De là on a une suite
$$
\mathcal{O}_{12}^* \xrightarrow{(\varepsilon_2, \varepsilon_3)} \mu_2 \times \mu_2 \to \{0,1\}_{S_0} \times \{0,1\}_{S_0} \to 1 \leqno{(81)}
$$

\newpage

%% ===== Page 56 =====

que
et on voit alors les lois de groupes, d'où
$\varepsilon_2 = (-5)^{\nu_2} = 1-2\nu_2$, on a $\{0,1\}_{\operatorname{S}_0 \varepsilon_2} = 1-2\nu_2$.
Les "groupes" de Zsigmondy de $(\mathbf{Z}/12\mathbf{Z})^*$
se donnent dans le tableau :

$$
\begin{array}{c|c|c}
\varepsilon_2 \diagdown \nu_2 & +1 \ (\text{i.e. } \nu_2=0) & -1 \ (\text{i.e. } \nu_2=1) \\
\hline
+1 \ (\text{i.e. } \nu_2=0) & 1 & +5 \\
\hline
-1 \ (\text{i.e. } \nu_2=1) & -5 & -1
\end{array} \leqno{(82)}
$$

Dans ce tableau, une section semblable
de (79) permet les sections de $\mathcal{O}_{12}^*$ de la forme
$$
\mu_0 = (-5)^{\nu_2} (+5)^{\nu_3} \equiv \begin{cases} (-1)^{\nu_3} \pmod{6} \\ (-1)^{\nu_2} \pmod{4} \end{cases} \leqno{(83)}
$$
Dans une section générique de $\mathcal{O}_{12}^*$ s'écrit
$$
\mu = (-5)^{\nu_2} (+5)^{\nu_3} + 12\mu_0]
$$
$$
(\nu_2, \nu_3) \in \Gamma \{0,1\}_{\operatorname{S}_0} \times \{0,1\}_{\operatorname{S}_0}
$$
$$
\lambda \in \Gamma(\mathbb{G}_a) \leqno{(84)}
$$

Notons que $\nu_2 = \nu_2^2$ et
$$
\mu = (1+4\nu_3) \quad 5^{\nu_3} = 1+4\nu_3
$$
$$
(-5)^{\nu_2} [(-1)^{\nu_2} 5^{\nu_2}] = (1-2\nu_2)(1+4\nu_2)
$$
$$
= 1 + 2\nu_2 - 8\nu_2^2 = 1 - 6\nu_2 \leqno{\text{NB}}
$$
d'où
$$
\mu_0 = (1+4\nu_3)(1-6\nu_2)
$$
$$
= 1 + 2\nu_2 - 8\nu_2^2 + 4\nu_3 - 24\nu_2\nu_3 - 8\nu_2^2
$$
$$
= 1 - 2\nu_2 + 4\nu_3 - 24\nu_2\nu_3 \leqno{(85)}
$$
\begin{center}
(On trouve $\mu_0 = 1 + 4\nu_3 - 6\nu_2$)
\end{center}

\newpage

%% ===== Page 57 =====

Si l'on utilise (84) et (85)
$$
\mu = 1 - 2\nu_2 + 4[\nu_3 - \nu_2 + 3\nu_2 \nu_3 + 3\nu_0] \leqno{(86)}
$$
$$
\mu = 1 - 2\nu_3 + 6[\nu_3 - \nu_2 - 2\nu_2 \nu_3 + 2\nu_0] \leqno{(87)}
$$
le calcul des diagrammes (composés) $\mathcal{O}_{12}^*$ et l'aligné de la notion $\mathcal{Q}_m = \mathcal{D}_1^*$ :
$$
2\nu_2' - \nu_2 = 2\nu_3 - 3\nu_2 - 6\nu_2 \nu_3 + 6\nu_0
$$
$$
3\nu_3'' - \nu_3 = 2\nu_3 - 3\nu_2 - 6\nu_2 \nu_3 + 6\nu_0
$$

En résumé, le schéma des groupes lisses $\mathcal{O}_{12}^*$.

\paragraph*{Proposition.} Introduisons les $\nu_2, \nu_3$ (valeurs dans $\{0,1\}_{\mathcal{G}} \times \mathfrak{T}_{\mathcal{G}}$), $\nu', \nu''$ (valeurs dans $\mathfrak{T}_{\mathcal{G}}$) et $\mu$ (valeurs dans $\mathcal{O}_{12}^*$) définis par les conditions :
$$
\begin{cases}
\mu = \operatorname{Im} \text{ de } \mu \text{ de } \mathcal{D}_1^* \text{ par l'image du groupe } \mathcal{D}_N^* \hookrightarrow \mathcal{D}_1^* \\
(-1)^{\nu_2} = \dots \in (\mathbf{Z}/2\mathbf{Z})^*_{\mathcal{G}} \\
(-1)^{\nu_3} = \dots \in (\mathbf{Z}/2\mathbf{Z})^*_{\mathcal{G}}
\end{cases} \leqno{(88)}
$$

$$
\begin{aligned}
\mu &= (-5)^{\nu_2} (+5)^{\nu_3} + 12[\nu_0] \\
&= (1 + 4\nu_3)(1 - 6\nu_2) + 12[\nu_0 - 2\nu_2 \nu_3] \\
&= 1 + 4\nu_3 - 6\nu_2 + 12[\nu_0 - 2\nu_2 \nu_3]
\end{aligned} \leqno{(89)}
$$
pour une section $\nu_0$ bien déterminée de $G=0$.

\newpage

%% ===== Page 58 =====

Posons de plus
$$
\begin{cases}
\nu' = \nu_3 - \nu_2 - 3\nu_2\nu_3 + 3\nu_0 \\
\nu'' = \nu_3 - \nu_2 - 2\nu_2\nu_3 + 2\nu_0
\end{cases} \leqno{(90)}
$$
On a alors les relations :
$$
\begin{cases}
\mu = 1 - 2\nu_2 + 4\nu' & (= \varepsilon_2 + 4\nu') \\
\mu = 1 - 2\nu_3 + 6\nu'' & (= \varepsilon_3 + 6\nu'')
\end{cases} \leqno{(91)}
$$
\hspace{2cm} $- \nu_2 + 2\nu' = - \nu_3 + 3\nu'' \quad \text{i.e. } 2\nu' - \nu_2 + \nu_3 = 3\nu''$

\marginpar{où $\varepsilon_p = (-1)^p$ \\ d'où $\varepsilon_{p-1} = -2\nu_p$ \\ ($p \in \{2,3\}$)}

NB. La condition $\mu = 1$ équivaut à la condition
$$
\nu_2 = \nu_3 = \nu_0 = 0 \leqno{(92)}
$$
et implique
$$
\nu = \nu' = \nu'' = 0 \leqno{(93)}
$$

\newpage

%% ===== Page 59 =====

4. Le groupe $\widetilde{Z}_{\beta\sigma}$ est défini par l'extension $1 \to N_{\beta\sigma} \to \widetilde{Z}_{\beta\sigma} \to \mathbb{D}_{12}^* \to 1$, et le plongé dans $\mathbb{D}_{12}^* \times M \times M$ :
$$
\widetilde{Z}_{\beta\sigma} \subset \mathbb{D}_{12}^* \times M \times M \leqno{(94)}
$$
$$
\{(\mu, \alpha_0, \beta_0) \mid (1-\rho_M)(\alpha_0) - (1-\sigma_M)(\beta_0) = 3\,\nu'(\mu)\,\lambda_1\}
$$
Notons qu'on a
$$
(1-\rho_M)(\lambda_1 - \lambda_0) = 3\lambda_1 \leqno{(95)}
$$
dans la relation des cocycles de (77), pour avoir la forme
$$
(1-\rho_M)(\alpha_0 - \nu''(\lambda_1 - \lambda_0)) - (1-\sigma_M)(\beta_0) = 0 \leqno{(96)}
$$
Je vais admettre, sans quoi on a une structure de groupe naturelle sur $\widetilde{Z}_{\beta\sigma}$, pour laquelle la projection $\widetilde{Z}_{\beta\sigma} \to \mathbb{D}_{12}^*$ est un morphisme de groupes. On trouve maintenant une suite exacte
$$
1 \to \widetilde{Z}_{\beta\sigma}^{+0} \to \widetilde{Z}_{\beta\sigma} \to \mathbb{D}_{12}^* \to 1 \leqno{(97)}
$$
$$
\widetilde{Z}_{\beta\sigma}^{+0} = \\operatorname{Ker}(\widetilde{Z}_{\beta\sigma} \to \mathbb{D}_{12}^*) \leqno{(98)}
$$
$$
\{(\alpha_0, \beta_0) \mid (1-\rho_M)(\alpha_0) - (1-\sigma_M)(\beta_0) = 0\}
$$
La définition du groupe est lisse en $n \neq 3$, mais pas en $n=3$, si on divise par $n$.

\newpage

%% ===== Page 60 =====

de le dim. relative 2 : le dim. relative 3.

Le fait d'une fibre exceptionnelle de dim. plus grande en partie, il me semble, être élucidé par une autopsie, en introduisant des définitions ad-hoc pour ``meilleur'' que $\widetilde{Z}_{\beta,\sigma}$ : il y a ``plein'' d'automorphismes de $\mathfrak{S}_k$ (pour le corps alg. des diviseurs), ayant les propriétés analogues, si $h$ est de car. 3, qui dans les... de caractéristique.

Cependant il me semble qu'on peut construire un groupe $N_{\beta,\sigma}$, qui est lisse sur $\operatorname{Spec}(\mathbf{Z})$ et coïncide avec $\widetilde{Z}_{\beta,\sigma}$ au dessus de $\operatorname{Spec} \mathbf{Z}[1/p] = S_0 \setminus \{p\}$.

Pour la définition du... le... le sous-module de $M$ correspondant (96).
$$
c_0 + c_1 \equiv 0 \pmod 3 \leqno{(3)}
$$
(cf 62) ...

\newpage

%% ===== Page 61 =====

\noindent correspond à $M_0$ de $M$

$$ M_0 = \{ M \in M \mid c_0 + c_1 = 0 \} \leqno{(99)} $$

d'où un fibré sur $M_0$ sur $\widetilde{\mathfrak{T}}_\sigma$, et en fait un homomorphisme de schémas
$$ \widetilde{\mathfrak{T}}_\sigma \longrightarrow M_0 \leqno{(100)} $$
$$ (H, \alpha_0, \beta_0) \longmapsto (1-\sigma)(\lambda_1 - \lambda_0) + (1-\sigma)\beta_0 $$

qui n'est pas explicite ainsi : on a
$$ M \times M \times \mathbf{Z} \longrightarrow M_0 \leqno{(101)} $$
$$ (\alpha_0, \beta_0, \gamma) \longmapsto (1-\sigma)(\alpha_0 \dots) + (1-\sigma)\beta_0 $$

d'où pour les schémas correspondants
$$ \underline{M} \times \underline{M} \times \frac{\mathbf{Z}}{\mathcal{G}} \longrightarrow M_0 \leqno{(102)} $$

qui se compose avec
$$ \widetilde{\mathfrak{T}}_\sigma \longrightarrow \underline{M} \times \underline{M} \times \underline{\mathbf{Z}} \leqno{(103)} $$
$$ (\alpha_0, \beta_0) \longmapsto (\alpha_0, \beta_0, \gamma(\mu)) $$

On peut donc considérer l'équation des cocycles (196), comme l'équation dans $M_0$, elle devient ainsi plus stricte.

\newpage

%% ===== Page 62 =====

d'où un sous-schéma fermé de $\widetilde{Z}_{\rho, \sigma}$,
$$
N_{\rho, \sigma} = \widetilde{Z}_{\rho, \sigma} \cap \{ (\mu, \alpha_0, \beta_0) \in \widetilde{Z}_{\rho, \sigma} \mid \text{l'équation des cocycles } \dots \text{ dans } M_0 \} \leqno{(104)}
$$
\marginpar{NB Le groupe $N_{\rho, \sigma}$ devant ainsi être noté, il est clair que ce morphisme est [unclear: sur] plutôt que $N_{\rho, \sigma} \subset \widetilde{Z}_{\rho, \sigma}$.}

Soit $p$ le fibre d'un $\dots = (\alpha_0, \beta_0) \in M \times M$ [d'après] (96) des $M_0$, $(\tilde{\alpha}_0, \beta_0)$ [unclear: est solution de] l'équation linéaire :
$$
M \times M \to M_0
$$
$$
(1 - \rho_M)(\tilde{\alpha}_0) + (1 - \sigma_M)(\beta_0) = 0 \leqno{(105)}
$$
Or l'homomorphisme de modèles :
$$
M \times M \to M
$$
$$
(\tilde{\alpha}_0, \beta_0) \mapsto (1 - \rho_M)(\tilde{\alpha}_0) - (1 - \sigma_M)(\beta_0) \leqno{(106)}
$$
[unclear: et] le morphisme $M \times M \to M_0$ correspond à qui est [unclear: un] sous-schéma. D'où $N_{\rho, \sigma}$ un sous-schéma fermé de $\widetilde{Z}_{\rho, \sigma}$.

\newpage

%% ===== Page 63 =====

d'un sous-schéma fermé de $\mathcal{T}_{g,\nu}$,

$$ N_{g,\nu} = \tilde{\mathcal{T}}_{g,\nu} \leqno{(104)} $$
$$ \{(\alpha, \beta) \in \mathcal{T}_{g,\nu} \mid \text{l'application des } \dots \} \leqno \dots $$
Condition sur le morphisme
NB. Le groupe $N_{g,\nu}$ devant ainsi noter aussi (par abus) le fibre d'un morphisme $\dots$

$$ (\alpha, \beta) = (\alpha, \beta, \dots) \in M \times M \text{ et } \dots \leqno \dots $$

$$ (1 - \rho_M)(\tilde{\alpha}) - (1 - \sigma_M)(\beta) = 0 \leqno{(105)} $$

Or l'homomorphisme $M \times M \to M$
$$ (\alpha, \beta) \mapsto (1 - \rho_M)(\alpha) - (1 - \sigma_M)(\beta) \leqno{(106)} $$

qui coïncide avec lui-dessus de $\mathcal{S}_{g,\nu}$, il s'ensuit que c'est l'adhérence schématique
$$ N_{g,\nu} = \tilde{\mathcal{T}}_{g,\nu} \mid (\operatorname{Spec} \mathbf{Z} \setminus \{3\}) \leqno{(107)} $$
(codim. schématique), en ce qui implique que c'est un sous-schéma en groupes.

C'est peut-être le moment de regarder d'un peu plus près le sous-module $M_0$ de $M$,
$$ M_0 = \mathbf{Z}(\lambda_1 - \lambda_0) + \mathbf{Z}(\delta - \lambda_1) \leqno{(108)} $$
$$ = \mathbf{Z}(\lambda_1 - \lambda_0) + \mathbf{Z}(\delta - \lambda_0) $$
$$ = \mathbf{Z}\delta \oplus \mathbf{Z}(1 - \rho_M)(\delta) $$
$$ = \mathbf{Z}\delta + \mathbf{Z}\rho_M \delta $$

Notons que $M_0$ est stable sous $\rho$ (de façon géniale par construction) pour $x \in M$,
$$ x \equiv \rho_M(x) \pmod{M_0} \quad (\text{i.e. } (1 - \rho_M)(x) \in M_0) \leqno{(109)} $$
et
$$ x \equiv \sigma_M(x) \pmod{M_0} \quad (\text{i.e. } (1 - \sigma_M)(x) \in M_0). $$

Il est bien sûr stable aussi sous $\tau = \operatorname{id}_M$.
Donc $M_0$ est un sous-module de $M$ distingué dans $\mathfrak{S}_0$. C'est un sous-groupe distingué de $\mathfrak{S}_0 = \operatorname{Aut}(\hat{\mathbf{Z}})$,

\newpage

%% ===== Page 64 =====

contenant $[\Pi_0, \Pi_0]$, d'indice $3$ dans $\Pi_0$.

Je présume qu'il contient le sous-groupe
de congruences relatif à $12$ - i.e. distingué,
qu'il correspond à un sous-groupe de
$S\ell(2, \mathbb{Z}/12\mathbb{Z})$. Il est d'indice
$3 \cdot 12 = 36$ dans $S\ell(2, \mathbb{Z})$, [unclear: rayé] $G\ell(2, \mathbb{Z})$,
tandis que $$S\ell(2, \mathbb{Z}/12\mathbb{Z}) \simeq \underset{48}{S\ell(2, \mathbb{Z}/4\mathbb{Z})} \times \underset{24}{S\ell(2, \mathbb{Z}/3\mathbb{Z})}$$
il s'agit donc d'un sous-groupe (c'est même un sous-groupe distingué)
$32$ de $S\ell(2, \mathbb{Z}/12\mathbb{Z})$, engendré par les
éléments images de
$$ (\sigma(\rho)\rho^{-1})^3 = [\lambda_0, \lambda_1], \text{ ou } [\sigma_0, \lambda_\infty] $$
cf (7)
$\delta = \lambda_1 \lambda_0^{-1}$
$\rho(\delta)$

ou encore le sous-groupe distingué engendré par
le seul élément dans $S\ell(2, \mathbb{Z}/12\mathbb{Z})$

$$ (110) \quad \left\{ \begin{array}{l} \delta = \lambda_1 \lambda_0^{-1} = l_1 l_0^{-1} = \left(\begin{array}{cc} 1 & 0 \\ 2 & 1 \end{array}\right) \left(\begin{array}{cc} 1 & +2 \\ 0 & 1 \end{array}\right) \\ = \left(\begin{array}{cc} 1 & +2 \\ 2 & 5 \end{array}\right) \end{array} \right. $$

La question est de si ce sous-groupe
distingué de $S\ell(2, \mathbb{Z}/12\mathbb{Z})$ est bel et
bien d'ordre $32$, donc d'indice $36$ - un
ordre qui pourrait être plus grand, donc l'indice
plus petit. Déjà quel est l'ordre de $\bar{\delta}$ ?

\newpage

%% ===== Page 65 =====

mod 12 ? $\delta^{60}$.
$$
\left\{ \delta^2 = \begin{pmatrix} 1 & 1 \\ 2 & 5 \end{pmatrix} \begin{pmatrix} 1 & 1 \\ 2 & 5 \end{pmatrix} = \begin{pmatrix} 5 & 12 \\ 12 & 29 \end{pmatrix} \equiv \begin{pmatrix} 5 & 0 \\ 0 & 5 \end{pmatrix} \right. \leqno{(112)}
$$
d'où $\delta^{12} \equiv 1 \pmod{12}$,
c'est déjà un bon point! $\delta \equiv \delta \pmod 4$
$$
\left\{ \begin{aligned} \delta &\equiv 1 \pmod 2 \\ \delta &\not\equiv 1 \pmod 4 \\ \delta &\not\equiv 1 \pmod 3 \end{aligned} \right. \quad \Rightarrow \quad \delta^2 \equiv -1 \pmod 3 \leqno{(112')}
$$
Considérons
$$
\operatorname{Sl}(2, \mathbf{Z}/12\mathbf{Z}) \isom \operatorname{Sl}(2, \mathbf{Z}/4\mathbf{Z}) \times \operatorname{Sl}(2, \mathbf{Z}/3\mathbf{Z}) \to \mathfrak{S}_{0,3}^+ \times \mathfrak{S}_3^+ \leqno{(113)}
$$
et
$$
u: \operatorname{Sl}(2, \mathbf{Z}/4\mathbf{Z}) \to \mathfrak{S}_{0,3}^+/\pi_0
$$
et l'homomorphisme surjectif, dont le noyau est isomorphe au groupe des trois éléments $\rho$ nul dans le sujet des \dots et $\lambda_0, \lambda_1, \lambda_\infty$ permutent entre eux par $\rho$ et $\lambda_1 \lambda_0^{-1} = \lambda_0 = \dots$ etc. les $\lambda_i$ sont permutés par $\rho$ et
$v: \operatorname{Sl}(2, \mathbf{Z}/3\mathbf{Z}) \to \mathfrak{S}_3$
et l'homomorphisme
$$
\operatorname{Sl}(2, \mathbf{Z}/3\mathbf{Z}) \to \mathfrak{S}_4 \xrightarrow{\text{pi}} \mathfrak{S}_3^+
$$
d'où le noyau $\delta$ d'ordre 8. Le noyau de $u \times v$ d'ordre $4 \cdot 8 = 32$, et il existe \dots $\delta$. C'est un sous-groupe distingué engendré par $\delta$ (dont l'ordre au premier était bien multiple de 32), et

\newpage

%% ===== Page 66 =====

ce sous-groupe image de $\Pi_0 \subset \operatorname{Sl}(2,\hat{\mathbb{Z}})$, c'est-à-dire dans $\Pi_0$,
que l'image inverse d'un sous-groupe de
l'ét. précédente de $\operatorname{Sl}(2, \mathbb{Z}/12\mathbb{Z})$.

Résumons les résultats obtenus :

Proposition. Considérons dans $\hat{\mathfrak{S}}_0^+ = \operatorname{Sl}(2, \hat{\mathbb{Z}})$
la tour de sous-groupes invariants (et dans $\hat{\mathfrak{S}}_0$)
(114)
\begin{tikzcd}
\hat{\mathfrak{S}}_0^+ \ar[r, hookleftarrow, "6"] & \Pi = \Pi[2] \ar[r, hookleftarrow, "2"] \ar[d, equal, "\text{def}"'] & \Pi_0 \ar[r, hookleftarrow, "3"] \ar[d, phantom, "\text{sous-groupe N}"'] & \Pi_{00} \ar[r, hookleftarrow, "2!"] \ar[d, phantom, "\text{sous-groupe inv.}"'] & (\Pi_0, \Pi_0) \\
& \{u \in \hat{\mathfrak{S}}_0^+ \mid u \equiv 1 (2)\} & \text{engendré par } \rho, \dots & \text{engendré par les} & \\
& & \text{[unclear]} & \lambda_1 \lambda_0^{-1} = \rho_1 \rho_0^{-1} &
\end{tikzcd}

Considérons d'autre part le groupe quotient
(115) $$\operatorname{Sl}(2, \mathbb{Z}/12\mathbb{Z}) \simeq \hat{\mathfrak{S}}_0^+ / \Pi[12]$$
$$ \simeq \operatorname{Sl}(2, \mathbb{Z}/4\mathbb{Z}) \times \operatorname{Sl}(2, \mathbb{Z}/3\mathbb{Z}) $$

et la tour de sous-groupes suivants

[MARGIN: NB $(\hat{\mathfrak{S}}_0^+)^{ab} \simeq \mathbb{Z}/12\mathbb{Z}$ (cf 42 bis)
l'identif... en $H/\Pi_{00}^\bullet$ donne un homom (116)
d'ordre 3 $-$ surj. d'un ? en un
$\operatorname{Sl}(2, \mathbb{Z}/4\mathbb{Z}) / \operatorname{Sl}(2, \mathbb{F}_2)' \times \operatorname{Sl}(2, \mathbb{Z}/3\mathbb{Z}) / Syl_{2,3} \simeq \hat{\mathfrak{S}}_0^+ / \Pi_{00}$]

\begin{tikzcd}
H = \ar[d, phantom, "\cup" description, "6" left] & \operatorname{Sl}(2, \mathbb{Z}/4\mathbb{Z}) \times \operatorname{Sl}(2, \mathbb{Z}/3\mathbb{Z}) & \text{ordre } 48 \times 24 \\
\Pi^\bullet = \ar[d, phantom, "\cup" description, "2" left] & \operatorname{Sl}(2, \mathbb{F}_2) \times \operatorname{Sl}(2, \mathbb{Z}/3\mathbb{Z}) & \text{ordre } 8 \times 24 \\
\Pi_0^\bullet = \ar[d, phantom, "\cup" description, "3" left] & \operatorname{Sl}(2, \mathbb{F}_2)' \times \operatorname{Sl}(2, \mathbb{Z}/3\mathbb{Z}) & \text{ordre } 4 \times 24 \\
\Pi_{00}^\bullet = & \operatorname{Sl}(2, \mathbb{F}_2)' \times Syl_{2,3} \ar[d, "\simeq"'] & \text{ordre } 4 \times 8 \\
& \text{[unclear]} \simeq D_4 &
\end{tikzcd}

(117)
a) $\operatorname{Sl}(2, \mathbb{F}_2) \simeq \operatorname{Ker}(\operatorname{Sl}(2, \mathbb{Z}/4\mathbb{Z}) \to \operatorname{Sl}(2, \mathbb{Z}/2\mathbb{Z}))$
b) $\operatorname{Sl}(2, \mathbb{F}_2)' \subset \operatorname{Sl}(2, \mathbb{F}_2)$, sous-groupe distingué $\simeq \mathbb{F}_2 \times \mathbb{F}_2$ dont les trois él. $\neq 0$ sont $\lambda_0, \lambda_1, \lambda_\infty$ (images dans $\operatorname{Sl}(2, \mathbb{Z}/4\mathbb{Z})$) - c'est un supplémentaire dans la [unclear] du sous-groupe multiplicatif $\mu_2 \simeq \mathbb{F}_2$ des unit. [unclear] de $\operatorname{Sl}(2, \mathbb{Z}/4\mathbb{Z})$.
c) $Syl_2 \subset \operatorname{Sl}(2, \mathbb{Z}/3\mathbb{Z})$, sous-groupe de Sylow, d'ordre 8 (qui est distingué)

\newpage

%% ===== Page 67 =====

Ceci posé, la tour de sous-groupes (114) se loge dans la tour (115).

Re-parlant de conscience, je vais reprendre l'esquisse de démonstration précédente. On note que les groupes figurant dans (114) — c'est-à-dire dans (116) — puis $\pi_1, \pi_0$ — sont d'ailleurs depuis un moment que les images inverses de $\pi_1, \pi_0$ ! (pour ceci, il suffit de considérer $\\operatorname{Sl}(2, \mathbf{Z}/4\mathbf{Z})$, puis le $\\operatorname{Sl}(2, \mathbf{Z}/2\mathbf{Z})$ ...).
Ainsi l'image inverse de $\pi_0$ dans $\mathfrak{S}^+_0$ est $\pi_{00}$. Comme les deux sont d'indice 36, ce sont deux groupes égaux, CQFD (le raisonnement pouvait s'appliquer aussi aux $\pi_0, \pi_0' \dots$).

\marginpar{NB: Mais $\mathfrak{S}^+_0 \to \\operatorname{Sl}(2, \mathbf{Z})$ en fait le quotient de $\pi_0/\pi_{00}$ serait ...}
Corollaire: L'homomorphisme canonique $\mathfrak{S}^+_0 = \\operatorname{Sl}(2, \mathbf{Z}) \to \\operatorname{Sl}(2, \mathbf{Z}/4\mathbf{Z}) / (\pi_{00} \cong \\operatorname{Sl}(2, \mathbf{F}_2) \times \operatorname{Syl}_2)$ est $\mathfrak{S} = \\operatorname{Sl}(2, \mathbf{Z}) / [\pi_0, \pi_0]$ \footnote{Le quotient de $\pi_0/\pi_{00}$ serait de toute évidence à étudier dans la prise $§ 58$, et l'image de $\pi_{00}$ n'est autre que le sous-groupe $\mathcal{M}_0$ de (99) (108).}

\newpage

%% ===== Page 68 =====

Les considérations précédentes rendent naturel la considération des revêtements d'ordre 3\marginpar{$M_{15}$} de $\mathbb{P}^1_{\overline{\mathbf{Q}}}$ correspondant à des sous-groupes $\Pi_0$ de $\Pi_0 \cong \pi_1(\mathcal{M}_{0,3}, \bar{\eta})$.\marginpar{c'est un revêtement de $\mathcal{M}_{0,3}$}
On vérifie très de suite que c'est une courbe de genre 1, si donc...

D'où de travailler avec $[\Pi_0, \Pi_0]$ au lieu de $\Pi_0$, d'où un groupe extension de $\Pi_0/[\Pi_0, \Pi_0] \cong M$, par un groupe de rang 2.

\marginpar{c'est une courbe elliptique (donc de genre 1)} que ce ``elliptique'', le sous-groupe $M$ va opérer de façon non triviale -- par la suite il en faudrait examiner la liaison et le mettre en relation avec l'action qui...
$$\\operatorname{Sl}(2, \hat{\mathbf{Z}}) \to \\operatorname{Sl}(2, \mathbf{Z}).$$

\newpage

%% ===== Page 69 =====

5) Les groupes $\underline{M}_{\rho,\sigma}$, $\underline{M}_{\rho,\sigma}^+$ (via $\underline{N}_{\rho,\sigma}$).

À tout $l \in \hat{\Pi}_0$ correspond un autom
intérieur $int(l)$, qui est l'identité sur
$\underline{M}_0$. On aura, tautologiquement

(118) $$ \begin{cases} int(l)(\rho) = \alpha(\rho) \\ int(l)(\sigma) = \beta(\sigma) \end{cases} \quad \alpha = \beta = l $$

Mais [unclear] (rel. de lacets C relative : $f_1=1$ d. $\nu=0$)
$$ (1-\rho_M)(\underset{l}{\alpha}) = (1-\sigma_M)(\underset{l}{\beta}) $$
i.e.
$$ \rho_M(l) = \sigma_M(l) \quad ? $$

Posant 
(119) $$ l = c_0 \lambda_0 + c_1 \lambda_1 $$
la relation [unclear] s'écrit
$$ c_0 \lambda_1 + c_1 \underset{-\lambda_0-\lambda_1}{\lambda_\infty} = c_0 \lambda_1 + c_1 \lambda_0 $$
i.e.
$$ - c_1 \lambda_0 + (c_0 - c_1) \lambda_1 = c_1 \lambda_0 + c_0 \lambda_1 $$
i.e. $c_1=0$ i.e. $l = c_0 \lambda_0$. Donc c'est
sans remords qu'on se restreint à
considérer des $l \in \widehat{L}_0$ ; correspondant

(120) $$ l = c_0 \lambda_0 \in \widehat{L}_0 = Im (\hat{\mathbb{Z}} \xrightarrow{u} \underline{M}) $$
$$ c_0 \mapsto c_0 \lambda_0 $$

On a un demi diagramme
(121) 
\begin{tikzcd}
\mathbb{G}_a \simeq \widehat{L}_0 \ar[r] \ar[d, hook] & \underline{N}_{\rho,\sigma} \\
\underline{G}^+
\end{tikzcd}

\newpage

%% ===== Page 70 =====

l'image de
(122) $$L_0^{\underline{M}_0} \to \underline{N}_{\rho,\sigma}$$
est invariante - l'hom (122) étant en
fait compatible avec l'action de $\underline{N}_{\rho,\sigma}$
- via $p_r$ sur $L_0^{\underline{M}_0} \simeq \tilde{G}_0$ (i.e. induits
par les opérations de $\underline{N}_{\rho,\sigma}$ sur $\tilde{G}^+ \supset \underline{M}_a \to L_0^{\underline{M}}$
et opérations de [unclear] sur $\tilde{G}^+$ déduites
de (121) sont identiques. Donc d'après
le topos général on trouve un
schéma en groupes $\underline{\underline{M}}_{\rho,\sigma}$ et un
diagramme de suites exactes df $\underline{\Gamma}_{\rho,\sigma} = \underline{N}_{\rho,\sigma}/L_0^{\underline{M}}$

[MARGIN: avec les notati. du $a$ $§ 99$ $\underline{VI}$, le gr. noté $\underline{\underline{M}}_{\rho,\sigma}$ serait sans doute noté $\underline{M}_{\rho,\sigma}/H_p$ ; ici $\underline{M}_{\rho,\sigma}$ serait $\underline{M}_{\rho,\sigma}$ (loc. cit. 8°)]

(123)
\begin{tikzcd}
1 \ar[r] & L_0^{\underline{M}} \ar[r] \ar[d] & \underline{N}_{\rho,\sigma} \ar[r] \ar[d] & \underline{\Gamma}_{\rho,\sigma} \ar[r] \ar[d, "\simeq"'] & 1 \\
1 \ar[r] & \tilde{G}^{+\Delta} \ar[r] & \underline{\underline{M}}_{\rho,\sigma} \ar[r] & \underline{\Gamma}_{\rho,\sigma} \ar[r] & 1
\end{tikzcd}

Remplaçant $\tilde{G}^+$ par $\underline{M}$ , on trouve un $\underline{M}_{\rho,\sigma}$ , et
(124)
\begin{tikzcd}
1 \ar[r] & L_0^{\underline{M}} \ar[r] \ar[d] & \underline{N}_{\rho,\sigma} \ar[r] \ar[d] & \underline{\Gamma}_{\rho,\sigma} \ar[r] \ar[d, "\simeq"'] & 1 \\
1 \ar[r] & \underline{M} \ar[r] & \underline{M}_{\rho,\sigma} \ar[r] & \underline{\Gamma}_{\rho,\sigma} \ar[r] & 1
\end{tikzcd}

D'ailleurs on a dans (124)
dans (123) (comparer $§ 99$ $\underline{VI}$ 5° ...), le
(125) $$M_{\rho,\sigma} \subset \underline{M}_{\rho,\sigma} \dots$$

\newpage

%% ===== Page 71 =====

Il faut préciser la structure de $\Gamma_{p,\sigma}$ en termes de $L_0^M \subset N_{p,\sigma}^+$, donc
$$ 1 \to \Gamma_{p,\sigma}^{0^+} \to \Gamma_{p,\sigma} \to \mathcal{G}_2^* \to 1 \leqno{(126)} $$
$$ 1 \to L_0^M \to N_{p,\sigma}^{0^+} \to \Gamma_{p,\sigma}^{0^+} \to 1 \leqno{(127)} $$
Considérons le noyau commun de $\mathbb{Z}$-modules
$$ N^{0^+} = \\operatorname{Ker}(M \times M \to M_0) \leqno{(128)} $$
avec l'application $(a_0, b_0) \mapsto (1 - \rho_M)(a_0) - (1 - \sigma_M)(b_0)$, noyau formé des $(a_0, b_0)$ tels que
$$
\begin{cases}
c_0 = a_0 + a_1 + (b_0 - b_0) = 0 \\
c_1 = 2a_1 - a_0 - (b_0 - b_0) = 0
\end{cases} \text{ (cf. (61))}
$$
Soit $N^{0^+}$ le noyau, de sorte que
$$ N_{p,\sigma}^{0^+} \isom N^{0^+} \leqno{(129)} $$
Considérons l'homomorphisme
$$
\begin{aligned}
\mathbb{Z} &\to N^{0^+} \\
n &\mapsto (n\lambda_0, n\lambda_0)
\end{aligned} \leqno{(130)}
$$
L'image est un sous-groupe fini des systèmes $(a_0, a_1, b_0, b_1)$ tels que $a_1 = b_1 = 0$, $a_0 = b_0$.
On a une suite exacte
$$ 0 \to \mathbb{Z} \to N^{0^+} \to \mathbb{Z}^3 $$
avec $(a_1, b_1, a_0 - b_0)$, d'où $N^{0^+} / \mathbb{Z} \to \dots$

\newpage

%% ===== Page 72 =====

donc $\dots$ suite exacte de $\mathbf{Z}$-modules
$$
1 \to \mathbf{Z} \xrightarrow{u \mapsto (u \lambda, u \lambda)} N^{0+} \to \Gamma^{0+} \to 0 \leqno{(131)}
$$
avec $\Gamma^{0+} \cong N^{0+}/\mathbf{Z}$ un $\mathbf{Z}$-module libre de rang $1$.

Ces $\dots$ conduisent
$$
\Gamma_{\rho, \sigma}^{0+} \cong \Gamma^{0+} (\cong \mathbf{G}_a) \leqno{(132)}
$$
Bien entendu, tous ces schémas de groupes, qui sont des constructions, sont lisses sur $\mathbf{Spec } \mathbf{Z}$.

Peut-on définir une section plus naturelle de l'extension $N_{\rho, \sigma}$ de $\Gamma_{\rho, \sigma}$ par $M$ ? \marginpar{liée à des relations}

Compte tenu de l'équation qui décrit $N_{\rho, \sigma}$, il semble qu'il existe une telle section des groupes $\dots$
$$
0 \to \mathbf{Z} \to N^{0+} \to \Gamma^{0+} \to 0 \leqno{(133)}
$$
C'est un $\mathbf{Z}$-module. Mais il faut vérifier que les sections « souhaitées » des schémas permettent de définir $\dots$ multiplicative $\dots$

\newpage

%% ===== Page 73 =====

5°) $Z_{\rho,\sigma}$ et $G_{\rho,\sigma} = \mathbb{O}_{12}^* G_{12}^+$ + $G_{\rho,\sigma}^+ \simeq \mathbb{O}_{12}^* \cdot \underline{M} = \underline{M}_{\rho,\sigma} \simeq \underline{M}$.
Une fois intr. duit $\mathbb{O}_{12}^*$ et $M_0 \subset M$, il me semble que le groupe $Z_{\rho,\sigma}$ ten lieu et place de $\tilde{Z}_{\rho,\sigma}$, dès l'équa - tion tirant compte de $\alpha, \beta$ et via ses données de $\xi, \eta$ redevient sympathique :

(134)
$$ \underline{Z}_{\rho,\sigma} = \operatorname{Aut}(\underline{\mathcal{S}}^+) \times \mathbb{O}_{12}^* $$
$$ \{ (u, \mu) \mid $$
a) $u(\rho) = \xi \rho$ (loc.) (conjugué à $\rho$, par un $\alpha \in \widehat{G_{12}^+}$ tel que $\chi(\alpha) = \varepsilon_3(\mu)$)
b) $u(\sigma) = \eta \sigma$ (loc.) (conjugué à $\sigma$, par un $\beta \in \widehat{G_{12}^+}$ tel que $\chi(\beta) = \varepsilon_2(\mu)$)
c) équation des bouts $\xi_0 = \eta_0 \nu''(\mu)$ soit $\xi_0 = \eta_0 + 3 \nu'' \lambda_1$ dans $M_0$

où $\nu'' = \nu''(\mu) \in \dots$ (cf pur. p. 685, et ci-dess.!)

(135)
$$ \begin{cases} \xi_0 = [- \nu_3 \lambda_1] \in \Gamma M_0 & \nu_3 = \nu_3(\mu) \\ \eta_0 = [\nu_2 \lambda_2] \in \Gamma M_0 & \nu_2 = \nu_2(\mu) \end{cases} $$

On doit supposer de plus que

(136)
$$ \begin{cases} \xi_0 \text{ locale dans l'im. de } (1 - \rho_M)(M) \\ \eta_0 \text{ locale dans l'im. de } (1 - \sigma_M)(M) \end{cases} $$

et on aura

(137)
$$ \underline{Z}_{\rho,\sigma} \simeq \tilde{\Sigma}_{\rho,\sigma} = \left\{ (\mu, \xi_0, \eta_0) \in \mathbb{O}_{12}^* \times M_0 \times M_0 \;\middle|\; \begin{matrix} 1) \; \xi_0 = \eta_0 + (\nu'' \lambda_1) \text{ dans } M_0 \\ \quad \text{où } \nu'' = \nu''(\mu) \\ 2) \text{ on a (136)} \end{matrix} \right\} $$

$G$

$$ (1 - \rho_M) : M \to M_0 $$
$$ (1 - \sigma_M) : M \to \mathbb{Z} = C_a \to M_0 $$
$$ (\lambda_0, \lambda_0+d_1, \lambda_1) \mapsto (b_1 - b_0) \mapsto (\lambda_1 - \lambda_0) = \lambda \delta $$

\newpage

%% ===== Page 74 =====

donc (136) dit cet est une condition sur $\xi_0$, à savoir
$$
\xi_0 \in \operatorname{Im}(G_{|M} \to M_0) \cap \\operatorname{Ker}(M_0 \to G_0) \leqno{(136 \text{bis})}
$$
(où $\lambda \mapsto \lambda \delta$ et $c_0 + c_1 \mapsto c_0 \sigma_0$).

Ceci dit, on trouve donc comme
$$
\text{Notons } \hat{\xi}_{0, \sigma} \defeq \dots
$$
$$
\begin{cases}
M_0 = \mathbf{Z}\delta \oplus \mathbf{Z}(\lambda_1) \\
\underline{M}_0 = \mathbf{Z}\delta \times \mathbf{Z}(\lambda_1)
\end{cases} \leqno{(138)}
$$
d'où, par suite, on trouve
$$
\hat{\xi}_{0, \sigma} \simeq \mathbf{D}_{12}^* \times \mathbf{Z} \delta \in G_a \leqno{(139)}
$$
$$
(\xi_0, \beta_0) \mapsto (\mu, \rho_0)
$$
car $\xi_0$ se récupère est $\mu, \rho_0$ dans l'équation des lacets
$$
\xi_0 = (1-\sigma_M)^{-1} \left( \rho_0 + \delta''(\lambda_1) \right) \in \dots
$$
(où $(1-\sigma_M)\xi = \dots$)
$= \delta'' + \dots$

NB: Si on se donne par contre $\xi_0$, on en tire par l'équation des lacets et (138) la valeur de $\rho_0$, et celle de $\delta''$ — dans $\mathbf{Z}$ —, on connaît en plus $\nu_2$ et $\nu_3$, on en déduit la valeur de $\delta_0$ dans $\mathbf{Z}$ (sous bon 2-ci section régulière de $G_0$) i.e. $\mathbf{Z}$-finie $\mathbf{Z}$. Donc c'est un...

\newpage

%% ===== Page 75 =====

\begin{equation*}
\Sigma_{g,\nu}^{+t} (S) \hookrightarrow (M, \{0,1\} \times \{0,1\}) (S) \leqno{(140)}
\end{equation*}
$$ (\xi_0, \xi_1, \xi_2) \longmapsto (\xi_0, \varepsilon_2(\xi_1), \varepsilon_3(\xi_1)) $$

L'image étant formée des systèmes $(\xi_0, \xi_1, \xi_2)$ tels que, posant
$$ \xi_0 = \delta \delta + \nu'' (3\lambda_1) \quad \delta, \nu'' \in \Gamma_{\mathbf{Q}}(S) \leqno{(140 bis)} $$
qu'en (138), on ait
$$
\begin{cases}
\nu'' + \nu_2 - \nu_3 [2\nu_2, \nu_3] \text{ divisible par } 2 \\
(-5)\nu_2 + (+5)\nu_3 + 6\nu'' + \nu_2 - \nu_3 + 2\nu_2\nu_3 \text{ invariant par } 1 - 2\nu_3 - 2\nu_2\nu_3 + 6\nu''
\end{cases} \leqno{(141)}
$$

En tout état de cause on a que $\Sigma_{g,\nu}^{+t}$ des sous-groupes $\Sigma_{g,\nu}$ est lisse de dimension relative $2$
$$ 1 \to \mathcal{G}_a \simeq \mathbf{Z} \to \Sigma_{g,\nu}^{+t} \to \Sigma_{g,\nu}^{+t} \to 1 \leqno{(142)} $$
(L'isomorphisme $\mathcal{G}_a \simeq \Sigma_{g,\nu}^{+t}$ étant donné par $\lambda \mapsto (\eta=1, \xi_0=\lambda\delta, \eta_0=\xi_0=\delta\delta)$)

On aura posé
$$ L \defeq \{ \lambda, \delta \mid \lambda \in \Gamma_{\mathbf{Q}} \} \subset M \leqno{(143)} $$
et on définit
$$ \mathcal{G} \to \Sigma_{g,\nu}^{+t} \leqno{(144)} $$
$$ \lambda \mapsto (\operatorname{int}(u), 1) \mapsto (\dots) $$

\newpage

%% ===== Page 76 =====

d'où

(145) $$L^\natural \overset{\sim}{\longrightarrow} Z_{\rho,\sigma}^{+0}$$

ce qui donne

(146) $$\Gamma_{\rho,\sigma} \overset{dfn}{=} Z_{\rho,\sigma}/L^\natural \simeq \mathbb{O}_{12}^* \quad .$$

Si on construit maintenant $G_{\rho,\sigma}$
comme

(147) $$G_{\rho,\sigma} = (Z_{\rho,\sigma} \cdot \underline{\underline{C}}^+) / L^\natural$$

on a le diagramme

(148)
\begin{tikzcd}
& & & \mathbb{O}_{12}^* \ar[d, equal] \\
1 \ar[r] & L^\natural \ar[r] \ar[d, hook] & Z_{\rho,\sigma} \ar[r] \ar[d, hook] & \Gamma_{\rho,\sigma} \ar[r] \ar[d, equal] & 1 \\
1 \ar[r] & \underline{\underline{C}}^+ \ar[r] & G_{\rho,\sigma} \ar[r] & \Gamma_{\rho,\sigma} \ar[r] & 1
\end{tikzcd}

D'ailleurs on a un scindage canonique

de $Z_{\rho,\sigma}$ sur $\Gamma_{\rho,\sigma} = \mathbb{O}_{12}^*$

(149)
\begin{tikzcd}
\mathbb{O}_{12}^* \ar[r] & Z_{\rho,\sigma} \ar[r, "\sim"] & \Sigma_{\rho,\sigma}^{1/\natural} \\
\mu \ar[r, mapsto] & \tilde{\mu} \ar[r, mapsto] & (\mu, \xi_0 = \nu''(3\lambda_1), \eta_0 = 0)
\end{tikzcd}

d'où

(150) $$ \begin{cases} Z_{\rho,\sigma} \simeq \mathbb{O}_{12}^* \cdot L^\natural \\ G_{\rho,\sigma} \simeq \mathbb{O}_{12}^* \cdot \underline{\underline{C}}^+ \end{cases} $$

L'opération de $\mathbb{O}_{12}^*$ sur $L^\natural$ se fait
via $\mathbb{O}_{12}^* \to \operatorname{Aut}(M)$ homothéties, celle
de $\mathbb{O}_{12}^*$ sur $\underline{\underline{C}}^+$ se fait par

(151) $$ \begin{cases} u_\mu(c) = c & (= \delta_1 c \text{ en écriture add. pour } C^+) \\ u_\mu(x) = \mu x & \text{si } x \in M \end{cases} $$

[MARGIN:
NB On pose
$G_{\rho,\sigma} = \mathbb{O}_{12}^* \cdot \underline{\underline{C}}^+$
$\cup$
$Z_{\rho,\sigma} = \mathbb{O}_{12}^* \cdot L^\natural$
]

\newpage

%% ===== Page 77 =====

79 Les groupes $\mathcal{N}_{\mathfrak{g}}, \mathcal{N}^{\mathfrak{g}}$ etc.

Étudions la structure de $\mathcal{N}_{\mathfrak{g}} = \mathcal{G}_{\mathfrak{g}, \sigma}$

$$ \mathcal{N}_{\mathfrak{g}} = \mathcal{G}_{\mathfrak{g}, \sigma} \defeq \{ u \in \mathcal{G}_{\mathfrak{g}, \sigma} \mid u(\rho) = \rho^{\varepsilon(u)} \} \leqno{(152)} $$
\footnote{Notation provisoire, cf plus bas p. 67}

L'homomorphisme composé $\mathcal{N}_{\mathfrak{g}} \to \mathcal{G}_{\mathfrak{g}, \sigma} \xrightarrow{\chi} \mathcal{O}_{12}^*$ est épi, car surjectif — comme on le verra —. Soit une section locale de $\mathcal{O}_{12}^*$, comme $\delta^{\nu_2}, \tau^{\nu_3}, \dots$. On a ainsi
$$ \operatorname{int}(\delta^{\nu_2} \tau^{\nu_3})^{-1} \circ u(\rho) = \rho \leqno{(153)} $$
distingue les $\dots$ i.e. $\nu=3\nu''$ on a alors
$$ \begin{cases} (\delta^{\nu_2} \tau^{\nu_3})^{-1} u = \tau^{\nu_3} \dots \\ \chi(u) = \dots \end{cases} \leqno{(154)} $$
\marginpar{NB on note $\chi(\dots) \to \mathcal{O}_{12}^*$}

Donc posant $v = (\tau^{\nu_3})^{-1} u = (\omega^{\nu_2} \tau^{\nu_3}) \dots$
$$ \chi(v) = 1, \text{ avec } v = \dots \leqno{(154)} $$

\newpage

%% ===== Page 78 =====

On a aussi la suite exacte
$$
1 \to \mathbb{Z}/2\mathbb{Z} \to \mathcal{N}_\rho \to \mathcal{G}_{\rho, \sigma}^* \to 1 \leqno{(155)}
$$
et
$$
\mathbb{Z}/2\mathbb{Z} = \operatorname{Centr}_{\mathcal{G}}(\rho)
$$
On a, si $H_{\rho, \sigma}^+ \defeq \mathfrak{S} / \mathcal{M} \simeq (\mathcal{M}^+)_\rho$,
$$
\left\{ \begin{aligned} H_{\rho, \sigma}^+ &\simeq \mathfrak{S} / \mathcal{M} \\ H^+ &= \mathfrak{S} / \mathcal{M} \end{aligned} \right. \leqno{(156)}
$$
Des isomorphismes canoniques
$$
\mathbb{Z}_\rho \to \operatorname{Centr}_{H_{\rho, \sigma}^+}(\rho) = \mathcal{L}_\rho \simeq (\mathbb{Z}/2\mathbb{Z})_\rho \leqno{(157)}
$$
et un surjectif qui est épi (puisqu'il y a une section, provenant de l'inclusion $\mathcal{L}_\rho \subset \mathbb{Z}_\rho$) et de la notation
$$
\mathbb{Z}_\rho^\circ = \mathbb{Z}_\rho \cap \mathcal{M} = \{ x \in \mathcal{M} \mid \rho_H(x) = x \quad \text{i.e. } (\operatorname{Int} \rho)(x) = x \} \leqno{(158)}
$$
Le noyau est une $\mathbb{Z}/2\mathbb{Z}$, mais de dimension 1 ou dimension 3. Pour autant, il est bien que la "bonne" définition de $\mathbb{Z}_\rho^\circ$ doit être pour une équation
$$
(\operatorname{Id} - \rho_H)(x) = 0 \quad \text{dans } \mathcal{M}_\circ \leqno{(159)}
$$

\newpage

%% ===== Page 79 =====

687

Ce qui revient en fait à remplacer $N_\rho$ défini dans (152), par des définitions schématiques et restrictives ci-dessous de §33. On trouve dans ce même groupe d'ailleurs, que désormais je note $N_\rho$, inséré dans la suite exacte
$$
1 \to L_\rho \to N_\rho \to \frac{\mathcal{D}_{12}^*}{\Upsilon_{g,\nu}} \to 1 \leqno{(159)}
$$
Posons
$$
\begin{cases}
\mathcal{G}_\rho \defeq \frac{\mathcal{D}_{12}^*}{\Upsilon_{g,\nu}} \subset \mathcal{D}_{12}^* \cdot \mathfrak{S}^+ = \mathfrak{S}_{g,\nu} \\
N_\rho^{\mathcal{G}} \defeq N_\rho \cap \mathcal{G}_\rho
\end{cases}
\leqno{(160)}
$$
On trouve alors
$$
N_\rho^{\mathcal{G}} \cap L_\rho = \mathcal{M} \cap L_\rho = \{1\} \leqno{(161)}
$$
\marginpar{donc $N_\rho^{\mathcal{G}} \isom \pi_1$}
$$
N_\rho^{\mathcal{G}} \hookrightarrow \frac{\mathcal{D}_{12}^*}{\Upsilon_{g,\nu}} \leqno{(161)}
$$
Je dis que l'image est
$$
\Upsilon_{g,\nu} = \mathcal{D}_{12}^* \defeq \{h \in \Gamma_{\mathbf{Q}} \mid \dots = 1\} \leqno{(162)}
$$
Or $\Gamma$ est contenu dans $\mathcal{D}_{12}^*$ résultat de (159), montre qu'il est contenu...
Soit $u = f^{-1} u_f$, $f \in \mathcal{G}, f \in \mathcal{D}_{12}^*$ section $\mathcal{G}_\rho = \mathcal{D}_{12}^* \cap \mathfrak{S}$, écrites qui donnent dans $N_\rho$ d'où...

\newpage

%% ===== Page 80 =====

$\operatorname{int}(f^{-1}) u_p(g) = g^{\epsilon_3}$ i.e. $\lambda_1 \rho = \operatorname{int}(f) \rho^{\epsilon_3}$.

En réduisant modulo $\mathcal{M}$ on trouve
$$ \rho = \rho^{\epsilon_3} \quad \text{dans } \mathcal{H} $$
ce qui implique $\epsilon_3 = 1$, ok. Donc
$$ \mathfrak{N}_{\rho, \sigma} \isom \mathfrak{D}_{12}^{\ast 1} = \mathfrak{D}_{12}^{\ast \epsilon_3} \leqno{(163)} $$
— ainsi $\mathfrak{N}_{\rho, \sigma}$ est une section partielle
au-dessus de $\mathfrak{D}_{12}^{\ast 1}$ de l'extension $\mathfrak{N}_{\rho, \sigma}$ de $\mathfrak{D}_{12}^{\ast}$ par $\mathcal{L}_p$.

Passons aux autres sections partielles,
$$ \mathfrak{N}_{\rho, \sigma} = \{1, \sigma'\} \leqno{(164)} $$
sections par $\sigma'$, qui est un relevé de $\sigma$. Est-ce que $\sigma'$ normalise $\mathfrak{N}_{\rho, \sigma}$? Comme il normalise $\mathfrak{N}_{\rho, \sigma}$ [unclear: dans] $\mathfrak{G}_{\rho, \sigma}$, il suffit qu'elle normalise $\mathcal{G}_{\rho, \sigma}$ (donc $\mathfrak{D}_{12}^{\ast} = \mathcal{G}_{\rho, \sigma} / \mathcal{M}$), ce qui est évident car $\sigma' = \sigma$, il suffit que $\sigma \in \operatorname{Norm}(\mathfrak{G}_{\rho, \sigma})$ normalise $\mathfrak{G}_{\rho, \sigma}$, il vient que $\sigma$ normalise $\mathfrak{G}_{\rho, \sigma}^{\ast} = \mathfrak{D}_{12}^{\ast} \mathcal{G}_{\rho, \sigma}$.
Comme $\sigma$ normalise $\mathfrak{G}_{\rho, \sigma}^{\ast} = \mathfrak{G}_{\rho, \sigma}$, il suffit de regarder l'action de $\sigma$ sur $\mathfrak{D}_{12}^{\ast} \subset \mathcal{G}_{\rho, \sigma}$ défini par les $u_p$ de $(151)$. La question est si $\sigma u_p \sigma^{-1} u_p^{-1} \in \mathcal{G}_{\rho, \sigma}$.

\newpage

%% ===== Page 81 =====

$$\sigma u \sigma^{-1} u^{-1} = \sigma u_\mu(\sigma)^{-1} = \sigma (w_{12} \sigma)^{-1} = w_{12}$$

\section*{Proposition. Pour une section $\tau$ de $\mathcal{G}_{0,\nu}$ [unclear: (i.e. $\sigma(x)$ est une section de $\mathcal{G}_{0,\nu}$)], et son image dans $\mathfrak{T}_{1,\nu}$ est égale à $\iota_{1/2}(x)$ --- donc les ex des sections de $\mathcal{G}_{1,\nu}$ [unclear: sont $f(x)$] ---}

$$ \underline{\mathfrak{D}}^{*\nu}_{12} \defeq \{ \mathcal{H} \subset \mathfrak{T}_{12}^\nu \mid \epsilon_2(\mathcal{H}) = 1 \} \leqno{(154)} $$

\emph{Corollaire.} $\mathcal{G}_{1,\nu}$ n'est pas normalisé par $\sigma$ ni $\tau$, mais
$$ \mathcal{G}_{1,\nu}^{\sigma} \defeq \{ x \in \Gamma \mid \epsilon_2(x) = 1 \} \leqno{(155)} $$
est égal à $\mathcal{G}_{1,\nu} \cap \sigma (\mathcal{G}_{1,\nu}) = \mathcal{G}_{1,\nu} \cap \tau (\mathcal{G}_{1,\nu})$ et est normalisé par $\sigma$ et $\tau$.

\emph{Corollaire.} $\mathfrak{T}_{1}^\nu$ n'est pas normalisé par $\sigma$ ni $\tau$, mais
$$ \mathfrak{T}_{1}^{\nu, \sigma} \defeq \mathfrak{T}_{1}^\nu \cap \mathcal{G}_{1,\nu}^\sigma \subset \mathfrak{T}_{1}^{\nu, \text{supp}} \leqno{(165)} $$
est normalisé par $\sigma$ et $\tau$, et contient $\mathfrak{T}_{1}^{\nu, *0} (= \mathfrak{T}_{12}^\nu \cap \mathfrak{T}_{12}^{\nu})$.

\footnote{NB On a $\mathcal{N}^{\nu, \sigma} = (\mathcal{N}^\nu)^\sigma$ [unclear: (en un sens)].}
\footnote{A-t-on $\mathcal{Z}^\nu = (\mathcal{N}^\nu)^\sigma$? Je ferais [unclear: (Corollaire (168))] pour [unclear: (le groupe)] $\mathcal{G}_{1,\nu}$, [unclear: (le groupe)] $\mathcal{G}_{1,\nu}$. [unclear: (alors)] (38) p. 622.}

Soit
$$ \mathcal{N}^{\nu, \sigma} = \{ 1, \tau, \tau^2, \dots \} \subset \mathcal{N}^\nu, $$
$$ \mathcal{N}^\nu \isom \mathfrak{T}_{1}^{\nu, *} \defeq \{ \mathcal{H} \subset \mathfrak{T}_{12}^\nu \mid (\epsilon_2 \dots)(\mathcal{H}) = \} \leqno{(169)} $$

\newpage

%% ===== Page 82 =====

Donc $N_{\mathfrak{p}}^{\prime \prime \prime}$ est un sous-groupe profini de $N_{\mathfrak{p}}$ sur $\mathcal{D}_{12}^{\ast \prime \prime \prime}$. Enfin, posons
$$
\mathcal{Z}_{\mathfrak{p}}^{\prime} \defeq \mathcal{Z}_{\mathfrak{p}}^{\sigma_{\mathfrak{p}}} \times \mu_{\mathfrak{p}} \quad (\mu = \{+1, -1\}) \leqno{(170)}
$$
alors $\mathcal{Z}_{\mathfrak{p}}^{\prime}$ est invariant par $\tau^{\prime}$, et
$$
\mathcal{N}_{\mathfrak{p}}^{\prime \prime} \defeq \operatorname{Centr}_{\mathcal{Z}_{\mathfrak{p}}^{\prime}} \subset N_{\mathfrak{p}} \leqno{(171)}
$$
de sorte que $\mathcal{Z}_{\mathfrak{p}}^{\sigma_{\mathfrak{p}}}$ divise $d$ des $N_{\mathfrak{p}}^{\sigma}$. On a alors

\textsc{Corollaire 4}. --- \textit{On a une suite exacte}
$$
1 \to \mu_{\mathfrak{p}} \xrightarrow{\operatorname{diag}} \mathcal{N}_{\mathfrak{p}}^{\prime \prime} \to \mathcal{D}_{12}^{\prime \prime} \to 1 \leqno{(172)}
$$
\marginpar{NB $N_{\mathfrak{p}}^{\prime \prime}$ est invariant par $\mathcal{D}_{12}$ et l'action de $\mathcal{D}_{12}$ sur $N_{\mathfrak{p}}^{\prime \prime}$ est via l'action sur $\mathcal{D}_{12}^{\prime \prime}$.}
\textit{où $\mathcal{N}_{\mathfrak{p}}^{\prime \prime}$ est un sous-groupe de l'extension $N_{\mathfrak{p}}$ de $\mathcal{D}_{12}^{\ast \prime \prime}$ par $\mu_{\mathfrak{p}}$.}

Il y aurait lieu de développer les groupes de sous-groupes remarquables de $\mathcal{N}_{\mathfrak{p}}$, calqués sur la digression (38) de p. 622\textit{bis}.

\newpage

%% ===== Page 83 =====

689

8) Les groupes $N_{\sigma}, N_{\sigma}^S$ ...

Posons
(173) $$ \begin{cases} N_{\sigma} \subset \Gamma_{\rho,\sigma} \\ =^{dfn} \{ u \in \Gamma_{\rho,\sigma} \mid u(\sigma) = \varepsilon_2(u)\sigma \quad (= \omega^{\nu_2(u)} \sigma) \} \end{cases} $$

(174) $$ N_{\sigma}^S =^{dfn} N_{\sigma} \cap G_{\rho,\sigma}^S $$

L'hom. composé
\begin{tikzcd}
N_{\sigma}^S \ar[r, hook] \ar[rr, bend right, "\chi_{\sigma}"'] & S_{\rho,\sigma} \ar[r, "\chi"] & \mathbb{O}_{12}^*
\end{tikzcd}
est épi. En effet, les $u_f$ sont dans $N_{\sigma}^S$ et donnent un scindage $\mathbb{O}_{12}^* \hookrightarrow N_{\sigma}^S$, d'où des décompositions en pro. $1/2$ directs
(175) $$ \begin{cases} N_{\sigma}^S \simeq \mathbb{O}_{12}^* \cdot S Z_{\sigma}^S \\ N_{\sigma} \simeq \mathbb{O}_{12}^* \cdot S Z_{\sigma} \end{cases} $$

où on a posé
(176) $$ \begin{cases} Z_{\sigma} = \operatorname{Ker}(N_{\sigma} \xrightarrow{\varepsilon_2 \circ p} \mu_{S_0}) = \operatorname{Cent}_{G_{\rho,\sigma}}(\sigma) \\ S Z_{\sigma} = Z_{\sigma} \cap S^+ = \operatorname{Cent}_{S^+_{\rho,\sigma}}(\sigma) \\ Z_{\sigma}^S = Z_{\sigma} \cap G_{\rho,\sigma}^S \\ S Z_{\sigma}^S = S Z_{\sigma} \cap G_{\rho,\sigma}^S \\ \quad = S^{\circ} = M \end{cases} $$

D'ailleurs
\begin{tikzcd}
S Z_{\sigma} \ar[r, "p_{S_0}^+"] & \operatorname{Cent}_{H_{S_0}^+}(\sigma) \ar[d, "\wr"'] \\
& L_{\sigma} \simeq (\mathbb{Z}/q\mathbb{Z})_{S_0} \ar[ul, dashed]
\end{tikzcd}
qui a une section, donc

(177) $$ S Z_{\sigma} \simeq L_{\sigma} \times S Z_{\sigma}^S $$

\newpage

%% ===== Page 84 =====

Enfin
$$
S\mathcal{Z}_{\sigma}^G = M \leqno{(178)}
$$
$$
\{x \in \Gamma \mid (1-\sigma_M)(x) = 0\} = \{a \in \Gamma \mid \sigma(a) = a\} = \Gamma_{\sigma} \cong G_{\sigma}
$$
d'où, compte tenu de (177):
$$
\begin{cases}
S\mathcal{Z}_{\sigma} \cong L_{\sigma} \times G_{\sigma} \\
S\mathcal{Z}_{\sigma}^G = (S\mathcal{Z}_{\sigma})^G \cong G_{\sigma}
\end{cases} \leqno{(179)}
$$
On a donc
$$
1 \to G_{\sigma} \to N_{\sigma}^G \to \mathcal{O}_{12}^* \to 1 \leqno{(180)}
$$
plus précisément, les isom.\ (175), (177), (179) donnent
$$
\begin{cases}
N_{\sigma}^G \cong \mathcal{O}_{12}^* \cdot G_{\sigma} \\
N_{\sigma} \cong \mathcal{O}_{12}^* \cdot (L_{\sigma} \times G_{\sigma})
\end{cases} \leqno{(180)}
$$
Dans ces formules, $\mathcal{O}_{12}^*$ opère sur $G_{\sigma}$ via $\chi: \mathcal{O}_{12}^* \to G_{\sigma}$, et sur $L_{\sigma}$ via $\epsilon_2: \mathcal{O}_{12}^* \to (\pm 1)_{S_{\sigma}}$:
$$
\begin{cases}
\mu(\lambda) = \mu.\lambda \\
\mu(\sigma) = \sigma^{-\epsilon_2}
\end{cases} \quad \text{où } \chi = \chi(\mu), \epsilon_2 = \epsilon_2(\mu) \leqno{(181)}
$$

\newpage

%% ===== Page 85 =====

9°) Retour sur $M_{\mathfrak{g}, \sigma}$.

Nous inspirant des développements du § 49 (section VI, 5), nous poserons
$$
\begin{cases}
\mathcal{S}_0 M_{\mathfrak{g}, \sigma} = \mathbf{Z}_{\mathfrak{g}, \sigma} \times \mathcal{N}_{\mathfrak{g}}^{\mathcal{G}} \times \mathcal{N}_{\mathfrak{g}} \times \mathfrak{S}_{\mathfrak{g}, \sigma} \\
\mathcal{S}_0 M_{\mathfrak{g}, \sigma}^{\mathcal{G}} = \mathbf{Z}_{\mathfrak{g}, \sigma} \times \mathcal{N}_{\mathfrak{g}}^{\mathcal{G}} \times \mathcal{N}_{\mathfrak{g}}^{\mathcal{G}} \times \mathfrak{S}_{\mathfrak{g}, \sigma}
\end{cases} \leqno{(182)}
$$
$$
\begin{cases}
M_{\mathfrak{g}, \sigma} = \mathcal{N}_{\mathfrak{g}}^{\mathcal{G}} \times \mathcal{N}_{\mathfrak{g}} \times \mathfrak{S}_{\mathfrak{g}, \sigma} \\
M_{\mathfrak{g}, \sigma}^{\mathcal{G}} = \mathcal{N}_{\mathfrak{g}}^{\mathcal{G}} \times \mathcal{N}_{\mathfrak{g}}^{\mathcal{G}} \times \mathfrak{S}_{\mathfrak{g}, \sigma}
\end{cases} \leqno{(183)}
$$
$$
\begin{cases}
\mathcal{S}_0 \Gamma_{\mathfrak{g}, \sigma} = \mathbf{Z}_{\mathfrak{g}, \sigma} \times \mathcal{N}_{\mathfrak{g}}^{\mathcal{G}} \times \mathcal{N}_{\mathfrak{g}}^{\mathcal{G}} \\
\Gamma_{\mathfrak{g}, \sigma} = \mathcal{N}_{\mathfrak{g}}^{\mathcal{G}} \times \mathcal{N}_{\mathfrak{g}}^{\mathcal{G}}
\end{cases} \leqno{(184)}
$$

$\Gamma = \Gamma_{\mathfrak{g}, \sigma} = \mathbf{D}_{12}^*$.

En fait, les groupes $\mathcal{S}_0 M_{\mathfrak{g}, \sigma}, \mathcal{S}_0 M_{\mathfrak{g}, \sigma}^{\mathcal{G}}, \mathcal{S}_0 \Gamma_{\mathfrak{g}, \sigma}$ jouent plutôt un rôle analogue, en fait des extensions des groupes $M_{\mathfrak{g}, \sigma}, M_{\mathfrak{g}, \sigma}^{\mathcal{G}}$ en $\mathfrak{S}_{\mathfrak{g}, \sigma}$.
$L_{\mathfrak{M}} = \\operatorname{Ker}(\mathbf{Z}_{\mathfrak{g}, \sigma} \to \mathbf{D}_{12}^*)$ (dérivés de l'action), ainsi $\mathcal{S}_0 \mathfrak{S}_{\mathfrak{g}, \sigma}$ en $\mathfrak{S}_{\mathfrak{g}, \sigma}$ par $L_{\mathfrak{M}} \cong \mathbf{G}_m$.

On a l'homomorphisme $M_{\mathfrak{g}, \sigma} \to \mathfrak{S}_{\mathfrak{g}, \sigma}$ et l'opération $M_{\mathfrak{g}, \sigma}^{\mathcal{G}} \to \mathfrak{S}_{\mathfrak{g}, \sigma}$.
Nous utilisons les notations, en fait des groupes de suites exactes (cf. 89) (79).

\newpage

%% ===== Page 86 =====

(190)
\begin{tikzcd}
1 \ar[r] & \underline{C}^\natural \ar[r] \ar[d] & \underline{M}^\natural_{\rho,\sigma} \ar[r] \ar[d] & \Gamma_{\rho,\sigma} \ar[r] \ar[d, equal] & 1 \\
1 \ar[r] & \underline{C}^+ \ar[r] & \underline{M}_{\rho,\sigma} \ar[r] & \Gamma_{\rho,\sigma} \ar[r] & 1
\end{tikzcd}

et des suites exactes

(191)
\begin{tikzcd}
1 \ar[r] & S N^?_P \times S\mathbb{Z}^?_\rho \times \underline{C}^\natural \ar[r] \ar[d] & \underline{M}^\natural_{\rho,\sigma} \ar[r] \ar[d] & \mathbb{O}_{12}^* \ar[r] \ar[d, equal] & 1 \\
1 \ar[r] & S N^?_P \times S\mathbb{Z}^?_\rho \times \underline{C}^+ \ar[r] & \underline{M}_{\rho,\sigma} \ar[r] & \underset{\gamma_{\rho,\sigma}}{\mathbb{O}_{12}^*} \ar[r] & 1
\end{tikzcd}

(192)
$$1 \to S N^?_P \times S\mathbb{Z}^?_\rho \to \Gamma_{\rho,\sigma} \to \underset{\gamma_{\rho,\sigma}}{\mathbb{O}_{12}^*} \to 1$$

où

(193)
$$ \begin{cases} S N^?_P = \mu_P \simeq \{\pm 1\} & (172) \\ S\mathbb{Z}^?_\rho \simeq C_+ & (179) \end{cases} $$

Ces suites exactes se précisent, grâce
aux décompositions $1/2$ directes :

[MARGIN: crossed out lines]
$??? \simeq \mathbb{O}_{12}^* \cdot C_+$ (150, 151)
$??? \simeq \mathbb{O}_{12}^* \cdot C_+$ (180)
$N^?_P / \mu_P \simeq \mathbb{O}_{12}^*$ (172)

(194)
$$ \begin{cases} \mathbb{Z}_{\rho,\sigma} \simeq \mathbb{O}_{12}^* \cdot C_+ & (150, 151) \\ \underline{C}_{\rho,\sigma} \simeq \mathbb{O}_{12}^* \cdot \underline{C}^+ & (15) \\ \underline{C}^\natural_{\rho,\sigma} \simeq \mathbb{O}_{12}^* \cdot \underline{C}^\natural & (160) \\ N^\natural_S \simeq \mathbb{O}_{12}^* \cdot \underline{C}^\natural & (180) \\ \text{et enfin} \\ N^?_P / \mu_P \simeq \mathbb{O}_{12}^* & (172) \end{cases} $$

\newpage

%% ===== Page 87 =====

On trouve alors des isomorphismes canoniques
$$
\begin{cases} 
M_{g, \nu} / \mu_p \isommap \hat{\mathbb{Q}}_{12}^* \cdot (G_a \times \mathfrak{G}^+) \\ 
M_{g, \nu}' / \mu_p \isommap \hat{\mathbb{Q}}_{12}^* \cdot (G_a \times \mathcal{M}) \\ 
\Gamma_{g, \nu} \isommap \hat{\mathbb{Q}}_{12}^* \cdot G_a \times \mathfrak{G}^+ 
\end{cases} \leqno{(195)}
\marginpar{$\hat{\mathbb{Q}}_{12}^* \cdot (G_a \times \mathfrak{G}^+ \dots)$ (181) et (151). Ses $\mathbb{Q}_{12}^*$ $\mathcal{G}_a \times \mathcal{M} \dots$}
$$
$\underline{M}_{g, \nu}$ et $\underline{M}_{g, \nu}'$ sont, sauf précision, des diviseurs de l'extension. $\underline{M}_{g, \nu} / \mu_p$ et $\underline{M}_{g, \nu}' / \mu_p$ pour $\mu_p$ sous-groupes de diviseurs de l'extension
$$
1 \to \mu_p \to N_g^? \to \hat{\mathbb{Q}}_{12}^* \to 1 \leqno{(196)}
$$
de $\hat{\mathbb{Q}}_{12}^*$ par $\mu_p$, par images inverses des lois canoniques des groupes
$$
(195) \text{ sur la forme } \hat{\mathbb{Q}}_{12}^* \cdot \mathcal{H}
$$
Il faudrait revenir sur l'extension (195), qui est justement l'image inverse d'un cas de $\mathfrak{T}_{0,3} \to H$.
$$
(E_2, E_3)
$$
$$
\hat{\mathbb{Q}}_{12}^* \to (\mu_p) \mathfrak{S}_0
$$
Comparons § 49 V 8°). L'opération de $\hat{\mathbb{Q}}_{12}^*$ sur $G_a \times \mathfrak{G}^+$ ($G_a \times \mathfrak{G}^+ = G_a \times \mathcal{M}$) se fait via des points sur $G_a$ et $\mathfrak{G}^+$ que j'ai explicités dans (181) et (151).

\newpage

%% ===== Page 88 =====

10) Les sous-groupes $\underline{M}_{\rho, \sigma}[0]$, $\underline{M}_{\rho, \sigma}(0)$, $\underline{M}_{\rho, \sigma}(1/2)$, $\underline{M}_{\rho, \sigma}'(-J)$ $[\underline{M}_{\rho, \sigma}'''(-J) \supset \underline{M}_{\rho, \sigma}'(-J)]$.

On utilise le plongement

(195) $\underline{M}_{\rho, \sigma} = N_\rho^? \times_\gamma N_\sigma^? \times_\gamma \underline{G}_{\rho, \sigma} \subset \mathcal{N}_\rho^? \times \mathcal{N}_\sigma^? \times \underline{G}_{\rho, \sigma}$
$$ \Big\| $$
$$ \{(a, b, u) \in \mathcal{N}_\rho^? \times \mathcal{N}_\sigma^? \times \underline{G}_{\rho, \sigma} \mid \gamma(a) = \gamma(b) = \underline{\chi}(u)\} $$

On définit donc

(196) $\underline{M}_{\rho, \sigma}[0] \subset \underline{M}_{\rho, \sigma}$
[MARGIN: NB on aura $\underline{M}_{\rho, \sigma}[0] \subset \underline{M}_{\rho, \sigma}^S$]
$$ \| $$
$$ \{(a, b, u) \in \underline{M}_{\rho, \sigma} \mid u \in \underline{Z}_{\rho, \sigma}\} $$

donc
$$ \underline{M}_{\rho, \sigma}[0] \simeq N_\rho^? \times_\gamma N_\sigma^? \times_\gamma \underline{Z}_{\rho, \sigma} $$

(197) Soit $\underline{Z}_{\rho, \sigma}(0) \subset \underline{Z}_{\rho, \sigma}$
le facteur $1/2$ direct $\mathbb{Q}_{12}^*$ de (150) (151),
alors, on a $\underline{M}_{\rho, \sigma}(0) \simeq N_\rho^? \times_\gamma N_\sigma^? \times_\gamma \underline{Z}_{\rho, \sigma}(0)$

[MARGIN: NB $\underline{Z}_{\rho, \sigma}(0) \subset \underline{Z}_{\rho, \sigma} \subset \underline{G}_{\rho, \sigma}$ est en fait le fact. $1/2$ direct $\mathbb{Q}_{12}^*$ de $\underline{G}_{\rho, \sigma}$ ds (150), il proviendrait la notation $\underline{G}_{\rho, \sigma}(0)$]

(198)
\begin{tikzcd}
\underline{M}_{\rho, \sigma}(0) \ar[r, hook] \ar[d, "\downarrow \simeq"'] & \underline{M}_{\rho, \sigma}[0] \ar[d, "\downarrow \simeq"] \\
S \underline{\Gamma}_{\rho, \sigma}(0) \ar[r, hook] \ar[d] & S_0 \underline{\Gamma}_{\rho, \sigma} \\
\underline{\Gamma}_{\rho, \sigma}(0) &
\end{tikzcd}

(199) $\underline{M}_{\rho, \sigma}(0) = \{(a, b, u) \in \underline{M}_{\rho, \sigma} \mid u \in \underline{Z}_{\rho, \sigma}(0)\}$
[MARGIN: dans la descript. (195) de $\underline{M}_{\rho, \sigma}$ comme prod. $1/2$ direct.]
$$ = \underset{\cap \atop \Gamma_Q}{\underline{\mu}} \cdot ( \underset{\cap \atop \Gamma_{\underline{Q}^+}}{\underline{\lambda}} , 0 ) $$

\newpage

%% ===== Page 89 =====

\noindent 88

\begin{enumerate}
\item[200)] $M_{p, \sigma} [0] \simeq \mathbf{Q}_{12}^* \cdot G_a \hookrightarrow M_{p, \sigma} \simeq \mathbf{Q}_{12}^* \cdot (G_a \times \mathfrak{S}^*)$
\end{enumerate}
(l'inclusion canonique $G_a \hookrightarrow G_a \times \mathfrak{S}^*$ ; d'autre part $M_{p, \sigma} \simeq \dots$)

$$
M_{p, \sigma} [0] \simeq \mathbf{Q}_{12}^* \cdot (G_a \times L_a^*) \simeq \mathbf{Q}_{12}^* \cdot (G_a \times \mathfrak{S}^*) \leqno{(201)}
$$

On définit d'autre part
$$
M_{p, \sigma} (1/2) = \{ (a, b, u) \in M_{p, \sigma} \mid u = b \} \leqno{(202)}
$$
$c$-à-$d$ $M_{p, \sigma} (1/2) \subset M_{p, \sigma}$.

Dans ce cas
$$
\begin{cases}
M_{p, \sigma} (1/2) \isom \Gamma_{p, \sigma} \\
M_{p, \sigma} \simeq M_{p, \sigma} (1/2) \cdot M \\
M_{p, \sigma} \simeq M_{p, \sigma} (1/2) \cdot \mathfrak{S}^+
\end{cases} \leqno{(203)}
$$

Dans la représentation de $M_{p, \sigma} / M_p$ en produits de directs\dots on a la description
$$
\left\{
\begin{aligned}
M_{p, \sigma} (1/2) &= \{ 1 \cdot (\lambda, \lambda (\lambda_1 + \lambda_2)) \mid \lambda \in \Gamma_{p, \sigma} \} \\
&\simeq \mathbf{Q}_{12}^* \cdot G_a \longrightarrow \mathbf{Q}_{12}^* (G_a \times M) \\
&\quad \quad \quad \quad \quad \quad \quad \quad \quad \quad \text{associé à } \Gamma_{p, \sigma} \to G_a \times M \\
&\quad \quad \quad \quad \quad \quad \quad \quad \quad \quad \quad \quad \quad \quad \quad \lambda \to \lambda, \lambda (\lambda_1 + \lambda_2)
\end{aligned}
\right. \leqno{(204)}
$$

\marginpar{$H_p \subset M_{p, \sigma} (1/2)$}

\newpage

%% ===== Page 90 =====

Enfin, on pose
$$
M_{p, \sigma}(\bar{j}) \subset M_{p, \sigma} = \{ (a, b, u) \in \Gamma_{p, \sigma} \mid u = a \} \leqno{(205)}
$$
i.e. on n'en a pas une, $U \in \mathfrak{T}(\mathcal{G}_{p, \sigma}^{\#})$... donc $M_{p, \sigma}(\bar{j}) \neq M_{p, \sigma}$.

On a alors
$$
\begin{aligned}
M_{p, \sigma}(\bar{j}) & \isom \Gamma_{p, \sigma} \\
M_{p, \sigma} & \isom M_{p, \sigma}(\bar{j}) \cdot \mathcal{G}^+
\end{aligned} \leqno{(206)}
$$

Mais on a $H_p \not\subset M_{p, \sigma}(\bar{j})$, plus précisément $H_p \cap M_{p, \sigma}(\bar{j}) = \{1\}$. D'où la relation $M_{p, \sigma}(\bar{j})$ définie par comme $M_{p, \sigma}(\bar{j}) \subset M_{p, \sigma}$ une section $s$. $M_{p, \sigma} / H_p$ s'identifie à $\Gamma_{p, \sigma} / H_p \simeq \mathcal{N}_0^{\#} \simeq \mathcal{D}_{12}^* \cdot \mathcal{G}_{\sigma}$.

On a cependant $M_{p, \sigma}(\bar{j}) \cdot H_p = M_{p, \sigma}(\bar{j}) \cdot \mathcal{G}_{\sigma}$.

$$
M_{p, \sigma}(\bar{j}) \cdot H_p / H_p \isom \operatorname{Im}(M_{p, \sigma}(\bar{j}) \to M_{p, \sigma}/H_p) \simeq \Gamma_{p, \sigma} / H_p \leqno{(207)}
$$
où $\Gamma_{p, \sigma} \simeq \Gamma_{p, \sigma} / H_p$ est isomorphe.

\newpage

%% ===== Page 91 =====

\marginpar{NB : $\Gamma_{p,\sigma} \to \mathfrak{S}_{p,\sigma}$ est épi de noyau $K_{p,\sigma} \cong \mathcal{N}_{p,\sigma} / H_{p,\sigma}$ (cf. dernier \S\ de 89).}

aussi
$$
\mathcal{M}_{\sigma} \hookrightarrow \widetilde{\mathcal{M}}_{p,\sigma}(-\overline{j}) \twoheadrightarrow \Gamma_{p,\sigma} \to 1 \leqno{(208)}
$$
$\widetilde{\mathcal{M}}_{p,\sigma}(-\overline{j})$ a une section de $\mathcal{M}_{\sigma} = \widetilde{\mathcal{M}}_{p,\sigma}(-\overline{j}) / H_{p,\sigma}$ au-dessus de $\Gamma_{p,\sigma}$. Posant $\widetilde{\mathcal{M}}'_{p,\sigma}(-\overline{j}) \defeq \widetilde{\mathcal{M}}_{p,\sigma}(-\overline{j}) \cap \mathcal{M}^{\sim}_{p,\sigma}$ (on a)

$$
\widetilde{\mathcal{M}}'_{p,\sigma}(-\overline{j}) \isom \Gamma'_{p,\sigma} \twoheadrightarrow \Gamma_{p,\sigma} \leqno{(209)}
$$
$\widetilde{\mathcal{M}}'_{p,\sigma}(-\overline{j})$ d'indice 4 dans $\mathcal{M}^{\sim}_{p,\sigma}(-\overline{j})$ sections partielles de $\mathcal{M}_{\sigma}$ au-dessus de $\Gamma_{p,\sigma}$, d'indice 2 des $\Gamma'_{p,\sigma}$. On a donc

$$
\begin{cases}
\widetilde{\mathcal{M}}_{p,\sigma}(-\overline{j}) \defeq \\operatorname{Ker}(\mathcal{M}^{\sim}_{p,\sigma} \to \mathfrak{S}_{\sigma}) \isom \widetilde{\mathcal{M}}'_{p,\sigma}(-\overline{j}) \cdot \mathfrak{S}^+ \\
\widetilde{\mathcal{M}}_{p,\sigma}(-\overline{j}) \defeq \\operatorname{Ker}(\mathcal{M}^{\sim}_{p,\sigma} \to \mathfrak{S}_{p,\sigma}) \isom \mathcal{M}_{p,\sigma}(-\overline{j}) \cdot \mathcal{M}
\end{cases} \leqno{(210)}
$$

Remarques : Je m'arrête là sur les développements de sous-groupes d'indice de $\widetilde{\mathcal{M}}_{p,\sigma}(-\overline{j}) - j'y reviendrai sans doute, dans le cahier des schémas en groupes $\mathcal{M}_{p,\sigma}$ généraux. Je dirai clair qu'il faut choisir...

\newpage

%% ===== Page 92 =====

par la suite, entre les notations possibles
répartition pour les groupes notés ici
$\mathfrak{M}_{g,\nu}$ et $\mathfrak{M}_{g,\nu}/\mathcal{M}_{\mathfrak{p}} \simeq \mathfrak{M}_{g,\nu}$ et idem ... 4)

ci-haut, ci-dessous, il est vrai que les
semblent les plus intéressants, car c'est
\marginpar{Plainte sur la surabondance des tildes...}
$\mathcal{M} = \mathcal{N}_{g,\nu} \subset \hat{\mathfrak{T}}_{g,\nu}(\\operatorname{Sl}(2, \hat{\mathbf{Z}}))$.

Il semblait que ces groupes intéressent
le moins, les plus simples, en l'occurence
$\mathfrak{M}_{g,\nu}$, et les groupes notés précédemment
$\mathfrak{M}_{g,\nu}$ prenaient les tildes ...
plus ambiguës. Le tilde est déjà utilisé
dans un sens différent des $\widetilde{\Gamma}_{g,\nu}$,
$\widetilde{\mathfrak{M}}_{0,3}$ etc — d'autre part, les
exposants '$\cdot$', '$\cdot$', '$\cdot$', '$\cdot$', '$\cdot$' déjà utilisés
dans des sens différents. Le double emploi
des tildes risque d'ailleurs d'être mineur
des conflits de notation, tôt ou tard
il est peut-être préférable de remplacer
le tilde des $\widetilde{\Gamma}_{g,\nu}$ par des lettres telles que $\bar{\Gamma}_{g,\nu}$
(suggérant "étendus") ou "généralisés" ...
Je pense que je résoudrai ceci ...

\newpage

%% ===== Page 93 =====

[Note: This transcription is based on page 92 of the manuscript.]

\noindent lement notationnel, pour l'instant je garde les notations introduites par la bande dans la présente section, qui seront révisées sans doute dans un \S \space spécial ultérieur.

Il faudrait aussi décrire des sous-groupes $\widetilde{M}_{g, \nu}(-J)$, $\widetilde{M}_{g, \nu}(-j)$ dans $\widetilde{M}_{g, \nu}$ écrit comme par une $\frac{1}{2}$ droite, à travers :

$$
\widetilde{M}_{g, \nu}(-J) \defeq \left\{ (\underline{\lambda}, x), \underline{\lambda} \in \mathbb{G}_a^{\nu} \times \mathbb{G}_a^{\nu} \in \mathcal{G}_{g, \nu} \right\} \leqno{(211)}
$$

où $u_{\nu}$ est défini dans (151), en particulier $u_{\nu}(\varphi) = \lambda_i \varphi_i \cdots = \nu(\varphi)$,

$$
(u_{\nu} \varphi)(\varphi) = \chi(u_{\nu}) \varphi(u_{\nu}) \quad \dots \quad \text{avec} \quad (154)
$$

\marginpar{$\delta = \lambda_1 \delta_0^{-1}$}
et posant $\sigma = \omega \tau$, le premier cas s'écrit aussi

$$
\nu = \sigma \nu_3 \delta^{-\nu} \in \Gamma_{\nu}, \quad \chi(u) = \underline{k}
$$
$$
\nu = \omega_3 (\omega_3 \cdot \tau_3 \delta^{-\nu} \mu_{\nu}) \leqno{(212)}
$$
avec $\in \Gamma_{g, \nu}$, car les trois facteurs de $\in \Gamma_{\nu} \mathcal{G}_{g, \nu} = \Gamma_{\nu}^g$

\newpage

%% ===== Page 94 =====

Il en suit que $u \in N_p^2$, dans la forme formant les deux classes de $N_p^2$.
Soit $\kappa = (\nu_1, \nu_3, \nu_0) \in \Gamma_{\mathbf{Q}}^{\alpha}$.

$$ (N_p^2)_\kappa = \{ \omega^i \sigma^{\nu_3} \delta^{-\nu''} \mu \mid i \in \{0, 1\} s_0 \} \leqno{(212)} $$

où $\nu_3 = \nu_3(\mu)$, $\nu'' = \nu''(\mu)$, $\delta = \lambda_1 \lambda_0^{-1}$. Dans la suite, on trouve
\marginpar{$\delta = \lambda_1 \lambda_0^{-1}$}

$$ \mathcal{M}_{p,\sigma}^{\sim}(-j) = \{ (\lambda, x) \cdot \mu \mid x = \omega^i \sigma^{\nu_3} \delta^{-\nu''} \mu, \text{où } i \in \{0, 1\} s_0 \} \leqno{(213)} $$

On aura donc

$$ \mathcal{M}'_{p,\sigma}^{\sim}(-j) = \mathcal{M}_{p,\sigma}^{\sim}(-j) \cap \mathcal{M}_{p,\sigma} = \{ (\lambda, \delta^{-\nu''}) \mu \mid \mu \in \Gamma_{\mathbf{Q}}^{\pm 1} (\text{avec } \varepsilon(\mu)=1), \lambda \in \Gamma_{\mathbf{Q}}, \nu'' = \nu''(\mu) \} \leqno{(214)} $$

\newpage

%% ===== Page 95 =====

11) L'homomorphisme $M \to M_{\rho,\sigma}(\hat{\mathbb{Z}})^{\sim}$, relatif à la tour des courbes de Fermat.

Notons qu'on a pour tout entier $N \in \mathbb{N}^*$

[MARGIN: Le $\Delta$ indique qu'on prend (215) sections de [unclear] qui sont [unclear]]

$$ \mathbb{O}_N^{\Delta}(\hat{\mathbb{Z}}) \simeq \mathbb{Z} \quad \text{d'où} \quad \mathbb{O}_\Gamma^{\Delta}(\hat{\mathbb{Z}})^{\wedge} \simeq \hat{\mathbb{Z}}^* $$

On a donc, compte tenu des décompositions en produits semi-directs (150) et des isomorphismes

(216)
$$ Z_{\rho,\sigma}^\sim \overset{df}{=} Z_{\rho,\sigma}^\Delta(\hat{\mathbb{Z}})^\wedge \simeq \hat{\mathbb{Z}}^* . \hat{L}_0 \quad \text{où} \quad \hat{L}_0 \simeq \hat{\mathbb{Z}} . \lambda_0 \simeq \hat{\mathbb{Z}} $$
$$ G_{\rho,\sigma}^\sim \overset{df}{=} G_{\rho,\sigma}^\Delta(\hat{\mathbb{Z}})^\wedge \simeq \hat{\mathbb{Z}}^* . \underline{G}^{+\wedge} $$
$$ C_{\rho,\sigma}^\sim \overset{df}{=} C_{\rho,\sigma}^\Delta(\hat{\mathbb{Z}})^\wedge \simeq \hat{\mathbb{Z}}^* . M^\wedge $$
$$ N_{\rho}^\sim \overset{df}{=} N_\rho^\sim(\hat{\mathbb{Z}})^\wedge \simeq \mathbb{O}_\Gamma^\Delta(\hat{\mathbb{Z}})^\wedge \simeq \hat{\mathbb{Z}}^* $$
$$ N_{\sigma}^\sim \overset{df}{=} N_{\sigma}^\sim(\hat{\mathbb{Z}})^\wedge \simeq \hat{\mathbb{Z}}^* . \hat{\mathbb{Z}} $$
$$ M_{\rho,\sigma}^\sim \overset{df}{=} M_{\rho,\sigma}^\sim(\hat{\mathbb{Z}})^\wedge \simeq \hat{\mathbb{Z}}^* . (\hat{\mathbb{Z}} \times \underline{G}^{+\wedge}) $$
$$ M_{\bar{\rho},\bar{\sigma}}^\sim \overset{df}{=} M_{\bar{\rho},\bar{\sigma}}^\sim(\hat{\mathbb{Z}})^\wedge \simeq \hat{\mathbb{Z}}^* . (\hat{\mathbb{Z}} \times \hat{M}) $$

où les opérations de $\hat{\mathbb{Z}}^*$ sur les $\hat{\mathbb{Z}}$ facteurs à droite sont explicitées dans les sections précédentes. On notera que, modulo des ajustages de notations et de définitions, les groupes $Z_{\rho,\sigma}^\sim$ etc. précédents sont ceux qui serais construits directement, en termes de $(\hat{\underline{G}}^\wedge, \underline{G}^{+\wedge}, M^\wedge)$ (non indexé de $\rho,\sigma,\tau$), dans §49 VI.

On a des hom. canoniques

(217)
\begin{tikzcd}
M = [unclear] \ar[r] & M_{\rho,\sigma}^\sim \\
M' \ar[r] \ar[u, hook] & M_{\bar{\rho},\bar{\sigma}}^\sim \ar[u, hook]
\end{tikzcd}

qui comme d'habitude se décompose en

\newpage

%% ===== Page 96 =====

trois, compatibles avec ces lois
$$
\mathcal{M} \to \mathcal{G}_{\rho, \sigma} = \gamma_\rho^*(\hat{\mathbf{Z}}) = \hat{\mathbf{Z}} \rtimes \hat{\mathbf{Z}}^*
$$
qui n'est autre que la condition cyclotomique $X$.
$$
\mathcal{M} \to N \isom \hat{\mathbf{Z}}^*
$$
se réduit ici à cette condition, et on passe de anciens questionnements. On a d'autre part
$$
\begin{tikzcd}
\mathcal{M} \arrow[r] & \mathcal{G}_{\rho, \sigma} = \hat{\mathbf{Z}}^* \rtimes \hat{\mathbf{Z}} \\
\mathcal{M}' \arrow[r] & \mathcal{G}_{\rho, \sigma} = \hat{\mathbf{Z}}^* \rtimes \widehat{M}
\end{tikzcd} \leqno (218)
$$
qu'il faudra examiner de plus près, comme la condition cyclotomique est surjectif, et en tout cas que ces deux sont surjectifs. Enfin, on a homomorphisme $\mathcal{M} \to N^{\hat{\sigma}}$, qui s'inscrit en
$$
\Gamma_{\mathbf{Q}} \to N^{\hat{\sigma}} = \hat{\mathbf{Z}}^* \rtimes \hat{\mathbf{Z}} \leqno (219)
$$
(1/2 discret)

qu'il faudra examiner également, notamment voir s'il est surjectif. Bien entendu, je n'ai pas de doute de le fait que par "voie arithmétique", on prouvera que l'action revient sur $\Gamma_{\mathbf{Q}}$.

\newpage

%% ===== Page 97 =====

J'ai envie de mettre ces homomorphismes en
relation avec la "tour" des courbes algébriques
définies sur $\mathbb{Q}$, indexées par $n \in \mathbb{N}^\ast$

(220) $X_n : X_0^n + X_1^n + X_\infty^n = 0$ \qquad \qquad (courbe de
$\subset \mathbb{P}^2_{\mathbb{Q}}$ \qquad \qquad \qquad \qquad Fermat d'indice $n$)

Soit $U_n$ l'ouvert des pts où $X_0 X_1 X_\infty \neq 0$

Sur $X_n$ le groupe fini $\mu_n^{\{0,1,\infty\}} / \mu_n$ (diag.)
opère ($\mu_n \subset \mathbb{G}_m$), et on a aussi
que $\mathfrak{S}_3 = \mathfrak{S}_{\{0,1,\infty\}}$, de sorte que le
produit semi-direct

[MARGIN: Notation $F_n$ pas très heureuse - il est possible qu'il y aura un besoin ultérieur de la notation $F_n$ pour les [unclear] de $U_n, U_\infty$ avec un in-dice $n$? (comparer p. 719)]

(221) $F_n = \mathfrak{S}_3 \cdot \mu_n^{\{0,1,\infty\}} / \mu_n = F_n^0$

opère sur $X_n$.
On a

(222) $X_n / F_n^0 \simeq X_1 \simeq U_{0,3}$ ,

isomorphismes compatibles aux opérations de
$\mathfrak{S}_3$. Plus généralement, si
$n = n'd$, on a

(223)
\begin{tikzcd}
X_n \ar[r] & X_{n'} \\
x_i \ar[r, |->] & x_i^d
\end{tikzcd}

et de cette façon $X_n$ devient un revête-
ment principal (= torseur) sur $X_{n'}$, de
groupe $\mu_d^{\{0,1,\infty\}} / \mu_d$, de façon plus précise,
[unclear], (223) est compatible avec
l'épimorphisme

(224)
\begin{tikzcd}
F_n \ar[r, "\text{épi}"] \ar[d, equal] & F_{n'} \ar[d, equal] & g \ar[r, |->] & g \\
\mathfrak{S}_3 \cdot \mu_n^{\{0,1,\infty\}}/\mu_n \ar[r] & \mathfrak{S}_3 \cdot \mu_{n'}^{\{0,1,\infty\}}/\mu_{n'} & (x_0, x_1, x_\infty) \ar[r, |->] & (x_0^d, x_1^d, x_\infty^d)
\end{tikzcd}

\newpage

%% ===== Page 98 =====

J'aimerais mettre en homomorphisme les actions sur le "tour" des courbes algébriques définies sur $\mathbf{Q}$, indexées par $n \in \mathbf{N}^*$.
$$
X_n : x_0^n + x_1^n + x_\infty^n = 0 \quad (\text{courbe de Fermat d'indice } n) \leqno{(220)}
$$
$\mathbf{P}^1_{\mathbf{Q}} \setminus \{0, 1, \infty\}$ soit l'ouvert des points $x_0, x_1 \neq 0$.
Soit $X_n$ le schéma. Le groupe $\mu_n^{\{0, 1, \infty\}} / \mu_n$ opère (où $\mu_n = \mathbf{G}_m$), et on a ainsi que
$$
\mathcal{G}_n = \mathcal{G}_{\{0, 1, \infty\}} / \mu_n \quad \text{direct} \leqno{(221)}
$$
$F_n = \mathcal{G}_{\{0, 1, \infty\}} / \mu_n$.
On a $X_n \tilde{F}_n \isom X_1 \isom U_{0,3}$,
et des isomorphismes compatibles aux actions de $\mathcal{G}_n$. Plus généralement, si $n = n'd$, on a
$$
X_n \xrightarrow{\pi} X_{n'} \leqno{(223)}
$$
$$
x_i \mapsto x_i^d
$$
et de cette façon $X_n$ devient un revêtement principal (ou torseur) sur $X_{n'}$, de groupe $\mu_d^{\{0, 1, \infty\}} / \mu_d$, de façon plus précise encore, compatible avec...
$$
\begin{tikzcd}
F_n \arrow[r, "\text{épimorphisme}"] \arrow[d] & F_{n'} \arrow[d] \\
\mathcal{G}_{\{0, 1, \infty\}} / \mu_n \arrow[r] & \mathcal{G}_{\{0, 1, \infty\}} / \mu_{n'} \arrow[r, "\text{groupes}"] & (x_0, x_1, x_\infty \mapsto x_0^d, x_1^d, x_\infty^d)
\end{tikzcd} \leqno{(224)}
$$
et $\mu_d^{\{0, 1, \infty\}} / \mu_d$ fonctionnant sur $X_n$, et...
$$
1 \to \mu_d^{\{0, 1, \infty\}} / \mu_d \to F_n \to F_{n'} \to 1 \leqno{(224)}
$$
$(x_0, x_1, x_\infty) \mapsto (x_0, x_1, x_\infty)$
$(x_0, x_1, x_\infty) \mapsto (x_0^d, x_1^d, x_\infty^d)$
pour $\alpha, \beta, \gamma \in \mu_d$, $g \in F_{n'}$.

Et l'isomorphisme (223) est compatible avec (224), et induit un isom. compatible aux actions de $F_{n'}$ :
$$
X_n \isom X_{n'} / (\mu_d^{\{0, 1, \infty\}} / \mu_d) \leqno{(225)}
$$
De plus, l'action de $\mu_d^{\{0, 1, \infty\}} / \mu_d$ sur $X_n$ est libre, et faisant elle un $\mu_d^{\{0, 1, \infty\}} / \mu_d$-torseur, de sorte que $X_n$ est bien un torseur sur $X_{n'}$ de groupe $\mu_d^{\{0, 1, \infty\}} / \mu_d$.
Passant à la limite, on trouve :
$$
X_\infty = \varprojlim_{n \in \mathbf{N}^*} X_n, \quad U_\infty = \varprojlim X_n \leqno{(226)}
$$
$$
F_\infty = \varprojlim F_n = \left(\mathcal{G}_{\{0, 1, \infty\}} / \varprojleftarrow \mu_n \right)
$$

\footnote{Note marginale sur la page 96 : "La notation $F_n$ pour ces homomorphismes est-elle possible ? Elle est bonne. Elle est la suite de la définition de la ramification de la suite des revêtements. La notation $F_n$ est un sous-groupe. (p. 719)"}

\newpage

%% ===== Page 99 =====

\textit{Suite de la page précédente, début de la ligne:}
\dots\ une opération de $F_{\infty}$ sur $M_{\infty} \simeq \hat{\mathbf{Z}}(1)$, et on obtient:

$$
F_{\infty} = \varprojlim F_n \simeq \hat{\mathfrak{T}}(\mathcal{C}_{\infty}) \leqno{(217)}
$$

qui est libre, sur lequel opère $\Gamma$, et $X_{\infty} \simeq U_{0,3}$.

Pour la théorie des Galois et $(U_{0,3}, \mathfrak{S}_{0,3})$, le groupe profini $F_{\infty}$ sur $\mathbf{Q}$ est muni d'une action, et on a un groupe profini agissant, et un groupe d'automorphismes — définis via les \dots

$$
\mathcal{M}_{0,3} = \pi_1(U_{0,3}, \overline{\mathbf{Q}}, \mathfrak{S}_{0,3}, -\overline{j}) \to \mathcal{M}^* \leqno{(217)}
$$

$$
F_{\infty}^{\circ} = F_{\infty}^{\circ}(\overline{\mathbf{Q}}) \simeq \mathfrak{S}_{0,3} \left( \hat{\mathbf{Z}}^{\{0,1, \infty\}} \right) \cong \widehat{\mathfrak{T}}(\mathcal{C}) \leqno{(218)}
$$

\marginpar{NB: L'action de $\mathcal{M}$ sur $M \simeq \hat{\mathbf{Z}}(1)$ est celle qui est donnée par le caractère cyclotomique, et d'ailleurs $\mathcal{M}$ opère sur $\pi_1$ via l'action galoisienne et l'action de $\mathfrak{S}_{0,3}$.}

Sur lequel $\Gamma_{\mathbf{Q}}$ opère. L'action correcte est la suivante:
$\Gamma_{\mathbf{Q}}$ fait opérer $\mathfrak{S}_{0,3}$ sur l'image du groupe, et on a un groupe qui opère sur $U_{0,3}$, $\Gamma_{\mathbf{Q}}$ sur $U_{0,3}$, via $v \cdot \dots \dots$, et on a un groupe d'extensions $\mathcal{M}_{0,3} = \pi_1(U_{0,3}, \overline{\mathbf{Q}}, \mathfrak{S}_{0,3}, -\overline{j})$
Donc $F_{\infty}^{\circ}$ qui en fait est lié aux groupes classiques, $M \simeq \hat{\mathbf{Z}}^*$:

$$
M \to \mathfrak{S}_{0,3}, \quad M \simeq \hat{\mathbf{Z}}^* \leqno{(219)}
$$

$$
\mathcal{M}_{0,3}(-\overline{j}) \text{ un torseur sous } F_{\infty}^{\circ}, \text{ sur lequel } \mathcal{M}_{0,3} = \pi_1(U_{0,3}, \overline{\mathbf{Q}}, \mathfrak{S}_{0,3}, -\overline{j}) \text{ opère}
$$
de façon compatible avec son opération sur $F_{\infty}^{\circ}$.

Le domaine des $(F_{\infty}^{\circ}, \mathcal{M}_{0,3})$-torseurs profinis équivaut à des revêtements principaux de $U_{0,3}$, de groupe $F_{\infty}^{\circ}$.

\newpage

%% ===== Page 100 =====

...sur les opérations de $\mathfrak{S}_3$ rendues compatibles avec le groupe de $\mathfrak{S}_3$ sur $F_\infty$.

D'autre part, la catégorie des torseurs sur des groupes abéliens $M \otimes T$, où $T$ un $\hat{\mathbf{Z}}$-module inversible, est équivalente à celle des torseurs sous le groupe affine $\hat{\mathbf{Z}}^* \cdot M^\wedge$.

Le donnée d'une classe d'isomorphisme de tels torseurs, avec opérations de $\pi_{0,3}$ rendue compatible avec l'action de $M^\wedge$, permet de définir d'autre part la suite exacte
$$
1 \to M \to \operatorname{Aut}(U_{0,3}) \to \mathcal{O}_{12}^* \times \left( \mathfrak{S}_3 \right)_{/S} \to 1 \leqno{(231)}
$$
$$
\text{ou } \left\{ \begin{matrix} H^+ / \mu_2 \end{matrix} \right\}
$$
elle donne, pour $T = (\hat{\mathbf{Z}})^\Delta$, une suite exacte
$$
1 \to M^\wedge \to \operatorname{Aut}(U_{0,3}) \to \hat{\mathbf{Z}}^* \times \mathfrak{S}_3 \to 1 \leqno{(232)}
$$
qui induit l'opération du sous-groupe de $\hat{\mathbf{Z}}^* \times \mathfrak{S}_3$ sur $M^\wedge$. Elle contient la suite exacte
$$
1 \to M^\wedge \to \mathfrak{S}^{+\wedge}_{0,3} \to \mathfrak{S}_3 \to 1 \leqno{(233)}
$$
qui s'inscrit dans une suite exacte de suite exacte.

\newpage

%% ===== Page 101 =====

698

$$ (234) \qquad 1 \to \Pi_1(U_{0,3 \overline{\mathbb{Q}}}, \overline{y}) \to \Pi_1(U_{0,3 \mathbb{Q}}, \mathfrak{S}_3; \overline{y}) \to \mathfrak{S}_3 \to 1 $$

par le sous-groupe invariant des commutateurs
du groupe $\Pi_1(U_{0,3 \overline{\mathbb{Q}}}, \overline{y})$. Cette structure
[MARGIN: si, elle est triviale, cf. plus bas...]
d'extension (233) de $\mathfrak{S}_3$ par $M^\wedge$ n'est pas
triviale, sauf erreur, bien que l'extension
induite de $\{1, \rho, \rho^2\} \simeq \mathbb{Z}/3\mathbb{Z}$ par $M^\wedge$, - ou
$\{1, \sigma\}$ par $\hat{M}$, le soit. A fortiori (232) n'est
pas triviale (bien qu'elle ait des scindages
canoniques au-dessus de $\hat{\mathbb{Z}}$, au-dessus de
$\{1, \rho, \rho^2\}$ et au-dessus de $\{1, \sigma\}$). Donc
(sous réserve de vérifier ces points tantot...)
le groupe $C_{\rho,\sigma}/\mu_a$ ne peut s'identifier
à un sous-groupe du groupe 
$$ \boxed{ Aff_{\hat{\mathbb{Z}}}(\hat{M}) \simeq \operatorname{Aut}_{\hat{\mathbb{Z}}}(\hat{M}). \hat{M} } $$
comprenant le sous-groupe $\hat{\mathbb{Z}}^\times \times \mathfrak{S}_3$ de $\operatorname{Aut}_{\hat{\mathbb{Z}}}(\hat{M})$.

Considérons le sous-groupe
$$ (235) \qquad Aff \simeq (\hat{\mathbb{Z}}^\times \times \mathfrak{S}_3) . \hat{M} $$
On a un diagr.

\begin{tikzcd}
(237) & 1 \ar[r] & \hat{M} \ar[r] & Aff \ar[r] & \hat{\mathbb{Z}}^\times \times \mathfrak{S}_3 \ar[r] & 1 \\
& & & & \mathcal{M}_{0,3} \ar[u]
\end{tikzcd}

: ligne exacte. Le donné de la structure
(230) équivaut à isoler (à un conjugué près, à l' [unclear]

\newpage

%% ===== Page 102 =====

\noindent de-ci de-là d'un s-g de $\hat{M}$-conjugués. Le relèvement de $\mathcal{M}_{0,3} \to \widehat{\mathfrak{S}}_{0,3} \times \hat{\mathbf{Z}}^*$ en
$$
\mathcal{M}_{0,3} \to \operatorname{Aff} \leqno{(238)}
$$
Le choix d'un tel relèvement dépend des choix d'un point de $U_{0,3}(\overline{\mathbf{Q}})$ --- dessus des points $\bar{q}$ de $U_{0,3}(\overline{\mathbf{Q}})$ --- le choix de ce point originel fait par transitivité pour $x \in \hat{M}$, et le relèvement (238) et bon défini par ailleurs par cet élément $x$.

Dans le livre de F. l'on fournit un telle classe de conjugaison de relèvements dont la connaissance pourrait, en principe, aller le livre de F. à bon pris. Bon on doit mettre, si possible, un relèvement en correspondance avec le hom.\ (217), et plus particulièrement son dual ou (218), qui fournit un hom.\ canonique $\mathcal{M}_{0,3} \to \operatorname{Aff}$, mais de façon à ce que le groupe ambiant $\mathcal{G}_{\mathbf{Q}}$ de (232), qui relie le hom.\ $\mathcal{M}_{0,3} \to \widehat{\mathfrak{S}}_{0,3} \times \hat{\mathbf{Z}}^*$ \marginpar{NB. Cet h.m. est défini sur $\pi_1(U_{0,3}) \to \hat{M}$ propriété initiale.}

\newpage

%% ===== Page 103 =====

699

(239)
\begin{tikzcd}
1 \ar[r] & M^\wedge \ar[r] & G_{p\Gamma}/H_G \ar[r] & \hat{\mathbb{Z}}^* \times \mathfrak{S}_3 \ar[r] & 1 \\
& & \hat{\mathbb{Z}}^* \times G^{+ \wedge} / H_G \ar[u, "\simeq"'] & & \\
& & & \mathcal{M}_{0,3} = \pi_1(\mathcal{U}_{0,3 \overline{\mathbb{Q}}, \mathfrak{S}_3} ; -\vec{1}) \ar[uu] \ar[ul, dashed]
\end{tikzcd}

Sur $G_{0,3}^{+ \wedge} \simeq \pi_1(\mathcal{U}_{0,3 \overline{\mathbb{Q}}, \mathfrak{S}_3}, -\vec{1}) \subset \mathcal{M}_{0,3}$ , cet hom. induit l'hom. canonique du passage au quotient de $G_{0,3}^{+ \wedge}$ obtenu en rendant $\Pi_{0,3}^\wedge$ abélien , soit $G_{p\Gamma}/H_G \simeq G^{+ \wedge}/H_G$ , extension de $\mathfrak{S}_3$ par $M^\wedge$.

Reprenons l'hom. $\mathcal{M}_{0,3} \to \text{Aff}$ (238) , il s'insère dans un diag. commutatif

(240)
\begin{tikzcd}
1 \ar[r] & \Pi_{0,3}^\wedge \ar[r] \ar[d] & \mathcal{M}_{0,3} \ar[r] \ar[d] & \Gamma_{\mathbb{Q}} \times \mathfrak{S}_3 \ar[r] \ar[d, "\chi \times \text{id}"] & 1 \\
1 \ar[r] & \hat{M} \ar[r] & \text{Aff} \ar[r] & \hat{\mathbb{Z}}^* \times \mathfrak{S}_3 \ar[r] & 1
\end{tikzcd}

où $\Gamma_{\mathbb{Q}} \times \mathfrak{S}_3 \to \hat{\mathbb{Z}}^* \times \mathfrak{S}_3$ est l'hom. de (237) , et
où $\Pi_{0,3}^\wedge \to \hat{M}$ est l'hom. canonique

(241) $$ \Pi_{0,3}^\wedge \to (\Pi_{0,3}^\wedge)_{ab} \simeq ((\Pi_{0,3})_{ab})^\wedge \simeq \hat{M} \ . $$

Si on désigne par $\mathcal{M}_1$ le quotient de $\mathcal{M}_{0,3}$ par $[\Pi_{0,3}^\wedge, \Pi_{0,3}^\wedge]$ , on trouve de même

\newpage

%% ===== Page 104 =====

un hom. canonique de suites exactes

(242)
\begin{tikzcd}
1 \ar[r] & M^\wedge \ar[r] \ar[d, equal] & M_1 \ar[r] \ar[d] & N_{0,3} \ar[r] \ar[d] & 1 \\
1 \ar[r] & M^\wedge \ar[r] & Aff \ar[r] & \hat{\mathbb{Z}}^\times \times \mathfrak{S}_3 \ar[r] & 1
\end{tikzcd}

(dépends du choix d'un pt de $\mathcal{U}_{0,3}(\overline{\mathbb{Q}})$ au-dessus de $-\vec{j}$)
qui montre que l'extension $M_1$ de $\hat{\Gamma}_Q$ par $M^\wedge$
est [unclear: isome] : l'image inv. de l'extension
triviale $Aff$ de $\hat{\mathbb{Z}}^\times \times \mathfrak{S}_3$ par $M^\wedge$, donc est
triviale. Plus précisément encore, considérons
le quotient $M_2$ de $M_1$ par le noyau de $M_1 \to Aff$,
on trouve un diagramme

(242 bis)
\begin{tikzcd}
1 \ar[r] & M^\wedge \ar[r] \ar[d, "\simeq"'] & M_2 \ar[r] \ar[d, hook] & N_2 \ar[r] \ar[d, hook] & 1 \\
1 \ar[r] & M^\wedge \ar[r] & Aff \ar[r] & \hat{\mathbb{Z}}^\times \times \mathfrak{S}_3 \ar[r] & 1
\end{tikzcd}

(où $N_2 \overset{def}{=} M_2 / \text{Im } M^\wedge$). On voit donc
que l'extension de $N_2$ par $M^\wedge$ est triviale
- [unclear: avec un choix de] $M^\wedge$ - conjugaison -
de scindages de cette extension - les
scindages correspondant aux pts de
$\mathcal{U}_{0,3}(\overline{\mathbb{Q}})$ au-dessus du pt $-\vec{j} \in \mathcal{U}_{0,3}(\mathbb{Q})$

\newpage

%% ===== Page 105 =====

On aurait vraiment envie de montrer que
cet homomorphisme $M_{0,3} \to Aff$ est isomorphe
à l'hom. $M_{0,3} \to G_{\rho,0}/\mu_{\hat{\mathbb{Z}}}$ de (239) , moyennant
un isomorphisme canonique de
$Aff$ avec $G_{\rho,0}/\mu_{\hat{\mathbb{Z}}}$. Cela impliquerait que
cette dernière extension est scindée. Les
isomorphismes de cette extension avec $Aff$
correspondraient aux scindages de cette
extension.

12°) Scindages de l'extension $G_{\rho,0}/\mu_{\hat{\mathbb{Z}}}$ de $\hat{\mathbb{Z}}^* \times \mathfrak{S}_3$ par $M$.
Comme on a déjà un relèvement
$$ \hat{\mathbb{Z}}^* \hookrightarrow G_{\rho,0} $$
(150, 151), je vais regarder les scindages
de l'extension (231) qui prolongent celui-ci -
ce qui revient à, les
scindages de l'extension induite sur $\mathfrak{S}_3$ :

[MARGIN: partout : il faut se donner ce scindage qui $\mathfrak{S}_3 \hookrightarrow \underline{S}^+ / \mu_{\hat{\mathbb{Z}}}$ ($G_{\rho,0} \simeq \hat{\mathbb{Z}}^* \dots$)]
[MARGIN: On regarde les scindages [unclear] correspondent bien uniquement aux repères]

(241) $$ 1 \to M \to \underline{S}^+ / \mu_{\hat{\mathbb{Z}}} \to (\mathfrak{S}_3)_{S_0} \to 1 $$

\newpage

%% ===== Page 106 =====

(244) $(\rho', \sigma') = (\rho, y \rho) \quad y \in \underline{M}$
$\rho'^3 = \sigma'^2 = 1, \sigma'(\rho') = \rho'^{-1}$.

Nous interprétons ces conditions en termes de $\{1, y\}$. La condition $\rho'^3 = 1$ s'écrit comme $\rho^3 = 1$ ($\underline{G}_{0,3}$).
$$ (245) \quad \{ \rho \rho \rho = 1 \quad \text{i.e. additive } (Id + \rho + \rho^2)(y) = 0 \leqno{(245)} $$
Cette condition est très satisfaisante, car
$$ (246) \quad \sigma(y) = 1 \quad \text{i.e. additive } (\operatorname{Aut})(y) = 0 \leqno{(246)} $$
et elle signifie que
$$ (247) \quad y = b\delta \quad \text{avec } b \in \Gamma, \quad \delta = \lambda_1 - \lambda_0 \leqno{(247)} $$

Nous poserons
$$ (248) \quad \zeta = a_0 \dots a_1, \quad a_0, a_1 \in \Gamma \leqno{(248)} $$
et nous allons exprimer en termes de $a_0, a_1, b$ la condition $\sigma'(\rho) = \rho'^{-1}$.
On a
$$ \sigma'(\rho) = i(y)\sigma(\rho) = i(y)[\sigma \rho \sigma^{-1} \dots ] $$
La condition $\sigma'(\rho) = \rho'^{-1}$ s'écrit
$$ \sigma(\zeta) \rho \dots \lambda_1 = \rho^{-1} \dots $$
soit, en multipliant $\dots$ par $\rho$
$$ \rho(\sigma\rho) \cdot \lambda_1 = \rho \dots $$

\newpage

%% ===== Page 107 =====

$$
\rho_M \sigma_M (\xi) + \lambda_1 = \rho_M(\eta) - \xi + \eta \quad \text{ou encore}
$$
$$
(1 + \varepsilon_1 \rho_M)(\xi) + \lambda_1 = (1 - \rho_M)(\eta) \leqno{(249)}
$$
$$
(2a_1 - a_0)\lambda_1 \qquad 3b\lambda_1
$$
$$
\text{ou encore, par un calcul immédiat}
$$
$$
2a_1 - a_0 + 1 = 3b \quad \text{ou} \quad 2a_1 - a_0 - 3b = -1 \leqno{(250)}
$$
(lire pour $g_{M}$),
Ceci montre que pour $g_{M}$ que l'on veut, il existe $y$ bien défini tel que les conditions (244) — i.e. définies par $(\rho^3 = \rho_M, \sigma^2 = \rho_M)$ définissent un sous-groupe de (243) il faut il que l'on ait
$$
2a_1 - a_0 + 1 \equiv 0 \leqno{(251)}
$$
et les $y$ étant déterminés de façon unique par (250), au moins si $S$ est un schéma régulier au-dessus de $S$ (ce qui est le cas si $S = \operatorname{Spec} \mathbf{Z}$). Ainsi les sous-groupes de (243) forment un schéma $S_0$, que l'on définit par les points
$$
S_0 \simeq \{ (a_0, a_1, b) \in \mathbf{Z}^3 \mid 2a_1 - a_0 - 3b = -1 \} \leqno{(252)}
$$
$$
\simeq \{ (\xi, \eta) \in \mathcal{M} \times \mathbf{Z} \mid (1 + \varepsilon_M)(\xi) - \rho_M(\eta) = -\lambda_1 \}
$$
qui est un schéma de façon naturelle.

\newpage

%% ===== Page 108 =====

toujours sur les solutions, on considère
$$ S_0 \simeq K, \text{ où } K \text{ est le noyau de } \mathbf{Z}^2 \to \mathbf{Z} \to 0 \leqno{(253)} $$
\marginpar{$K \simeq \mathbf{Z}^2$}
$$ (\alpha, \beta) \mapsto 2\alpha - \alpha_0 - 3\beta $$

D'où, d'autre part, une opération de $M$ sur $S_0$ par automorphismes intérieurs, que je vais expliciter. $\lambda \in M$ conjugue l'opération

$$ (\xi, \eta) \mapsto (\xi + (1-\rho_M)(\lambda), \eta + (1-\sigma_M)(\lambda)) \leqno{(254)} $$

Ainsi, les solutions en question procèdent de la structure de $S_0$ — toujours sur $S_0$ — et d'un homomorphisme de $M \to S_0$.

$$ \begin{cases} M \to S_0 \\ \lambda \mapsto ((1-\rho_M)(\lambda), (1-\sigma_M)(\lambda)) \end{cases} \leqno{(255)} $$

qui en fait, associé à l'homomorphisme de $\mathbf{Z}$-modules
$$ \begin{cases} M \to K \subset \mathbf{Z}^3 \\ c_0 + c_1\lambda \mapsto (\alpha_0 = c_0 + c_1, \alpha_1 = 2c_1 - c_0, \beta = c_1 - c_0) \end{cases} \leqno{(256)} $$

\newpage

%% ===== Page 109 =====

On voit tout de suite que cet homomorphisme est mis à part injectif (car le composé $\mathcal{M} \to \mathbf{Z}^3$ l'est) donc pour les questions de $\mathbf{Q}$.

\noindent \textbf{Proposition.} L'extension $(243)$ de $\mathfrak{S}_{0,3}$ par $\mathcal{M}$ est scindée sur $S_0$.

--- des scindages de $S_0$ sur $\mathfrak{S}_{0,3}$ (et $\mathfrak{S}^+_{0,3}$) correspondent \marginpar{$1-1$} aux scindages de l'extension de $\mathfrak{S}_{0,3}$ par $\mathcal{M}$ sur un bon quotient $S$, les deux scindages sont conjugués par un unique scindage de $\mathcal{M}$ sur $S$, donc le scindage de $(243)$ est unique.

Notons qu'il est immédiat, sur un bon quotient $S$, que les scindages de $(243)$ correspondent bijectivement aux scindages de l'extension de $\mathfrak{S}_{0,3}$ par $\mathcal{M}$ :
$$ 1 \to \mathcal{M}(S) \to \mathfrak{S}_{0,3}(S) \to \mathfrak{S}_{0,3} \to 1 \leqno (244) $$
qui s'identifie à l'extension...

\newpage

%% ===== Page 110 =====

$\mathfrak{S}_3$ par $M(S) = \Gamma(S, \mathfrak{g}) \otimes_{\mathbf{Z}} M$, déduite de l'extension universelle $\mathfrak{S}^+/\mathfrak{T}_g$ de $\mathfrak{S}_3$ par $M$, via l'homomorphisme $M \to M(S) = M \otimes_{\mathbf{Z}} \Gamma(S, \mathfrak{g})$.

Appliquons ceci au cas $S = \operatorname{Spec}(\hat{\mathbf{Z}})$, dans le cas

Corollaire L'extension $\mathfrak{S}^{+ \wedge}/\mathfrak{T}_g$ de $\mathfrak{S}_3$ par $M^\wedge$ est scindée et ses scindages fixent un torseur sous $M^\wedge$ (ceci, attention !).

Étudions maintenant les scindages de $\mathfrak{S}_{\sigma}/\mathfrak{T}_g$ en termes de déterminations, pour un scindage donné
$$
\begin{cases}
\rho' = \rho + a\lambda_1 \\
\eta' = y\sigma = b\delta = b(\lambda_1 - \lambda_0)
\end{cases} \leqno{(254)}
$$
de $\mathfrak{S}^+/\mathfrak{T}_g$ au-dessus de $(\mathfrak{S}_3)_{\sigma}$, les relevants $\mathfrak{D}_{12} \to \mathfrak{S}_{\sigma}/\mathfrak{T}_g$ compatibles avec les scindages précités, i.e.\ qui correspondent à $\rho', \sigma'$. Pour ceci

\newpage

%% ===== Page 111 =====

\newpage
\setcounter{page}{110}
\noindent Je vais déterminer le centralisateur de $(\rho', \sigma')$ dans $\mathfrak{S}_{g,\nu} / \mathcal{K}_{\mathfrak{q}}$, i.e.\ le centralisateur du sous-groupe section public \leqno{(258)}
$$
H(\rho', \sigma') \subset \mathfrak{S}_{g,\nu} / \mathcal{K}_{\mathfrak{q}}
$$
L'image de ce centralisateur dans $(\mathfrak{S}_{g,\nu})_0 \cong (\mathfrak{S}_{g,\nu} / M) / \mathcal{D}_{12}^*$ est contenue dans le centre de $(\mathfrak{S}_{g,\nu})_0$, donc on conclut
$$
\operatorname{Centr} (H(\rho', \sigma')) \subset \mathcal{D}_{12}^* \cdot M = \mathfrak{S}_{g,\nu} / \mathcal{K}_{\mathfrak{q}} \leqno{(259)}
$$
(où on a posé $\mathcal{D}_{12}^* \cdot M = \left\{ \left. \begin{smallmatrix} \lambda & \mu \\ \nu & \rho \end{smallmatrix} \right| \lambda \in \mathcal{D}_{12}^* \right\}$)

On est défini dans (150, 151). On se donne déterminer le centralisateur de $H(\rho', \sigma')$ en exprimant la condition pour un $\lambda, \mu, \nu, \rho$ d'y être, i.e.\ le commute ; $\rho', \sigma'$ :
On a :
$$
(\lambda\mu\nu\rho)(\rho') = \operatorname{int}(\lambda)(\rho') = \nu(\lambda)\rho'(\nu(\lambda))^{-1} \lambda\rho' = \operatorname{int}(\lambda)(H(\rho', \sigma'))
$$
$$
\text{avec } \rho' = \lambda' \lambda_1 \rho(\lambda)^{-1} \rho
$$
$v = \nu(\lambda)$, d'où la relation $\lambda\mu\nu\rho(\rho') = \rho'$
S'écrit aussi :
$$
\lambda\lambda_1\rho(\lambda)^{-1}\rho = \zeta\rho \quad \text{i.e.} \quad \lambda\lambda_1\rho(\lambda)^{-1} = \zeta
$$
(où $\zeta$ est une action additive $\lambda+\mu\xi + \nu\lambda_1 - \rho(\lambda) = \zeta$, avec $\mu-1=2\nu$)

\newpage

%% ===== Page 112 =====

$$(260) \quad (1-\rho)(\lambda) + 2\nu + \rho\lambda_1 = 0$$

La relation $(\lambda_{\mu, \rho})(\sigma') = \sigma'$ s'explicite, via
$$(\lambda_{\mu, \rho}(\sigma')) = \operatorname{int}(\lambda)\mu_{\rho}(\eta\sigma) = \operatorname{int}(\lambda)(\mu_{\rho}(\eta)\mu_{\rho}(\sigma)) = \lambda\eta^{\mu}\sigma\lambda^{-1}$$
avec $\mu_{\rho}(\sigma) = \sigma$ car $\nu=1$ et $\mu_{\rho}(\sigma') = \sigma'$, d'où la relation s'écrivant dans $\mathcal{G}_{\rho, \sigma}/\mathcal{M}$ :
$$\lambda\eta^{\mu}\sigma\lambda^{-1} = \eta\sigma \quad \text{soit} \quad \lambda\eta^{\mu}\sigma\lambda^{-1} = \eta\sigma$$
et, en écriture additive :
$$\lambda + \mu\eta - \sigma(\lambda) = \eta$$
$$(261) \quad (1-\sigma)\lambda + 2\nu\eta = 0$$

En termes de $a_0, a_1, b$ de $(257)$, ces équations s'expliciteront comme des couples $\lambda = x\lambda_0 + y\lambda_1$ par les équations
$$
\begin{cases}
x+y = -2\nu b \quad \text{(équation (261))} \\
x+y = -2\nu a_0 \\
-x+2y = -\nu(2a_1+1)
\end{cases} \leqno{(262)}
$$
où les quantités $a_0, a_1, b$ figurent dans le deuxième membre par des facteurs de multiplicité $\nu = \nu(\mathcal{K})$ et extraits de $(250)$
$$(263) \quad -a_0 + 2a_1 - 3b + 1 = 0$$
La première et troisième équation $(262)$ nous donne $x, y$ :
$$
\begin{cases}
y = \nu(-2a_1-b-1) \\
x = y + 2\nu b = \nu(-2a_1+b-1)
\end{cases} \leqno{(264)}
$$
et la deuxième équation $(261)$ devient, via

\newpage

%% ===== Page 113 =====

y portant ces valeurs de $x, y$, une relation liant $a_0, a_1, b, \nu$ :
$$ \nu(-4a_1 + 6b - 2) = -2\nu a_0 \quad \text{i.e. } \quad 2\nu(a_0 - 2a_1 + 3b - 1) = 0 $$
qui est une conséquence de (263)! Ainsi on trouve que pour un scindage

\textbf{Proposition.} Pour $H$ scindage de l'extension $\mathfrak{S}^+/\mu$ de $(\mathcal{G})_\mu$ par $M$ (scindage bien d'ailleurs), le commutant (dans ce groupe section) est un sous-groupe section de l'extension $\mathcal{G}_{\rho, \sigma}$ de $\mathcal{D}_{12}$ par $M$.

Le reliement $\mathcal{U}_{\rho', \sigma'} : \mathcal{D}_{12} \to \mathcal{G}_{\rho, \sigma}$ que l'on associe à $\rho', \sigma'$ donnés par (257) :
$$ \begin{cases} \rho' = \rho \\ \sigma' = \gamma \sigma = a_0 \lambda_0 + a_1 \lambda_1 \\ \nu = b \delta = b(\lambda_1 - \lambda_0) \end{cases} \quad a_0, a_1, b \in \Gamma_{\mathbf{Q}} $$
est donné par
$$ u_{\rho', \sigma'}(\lambda) = \lambda u_{\rho, \sigma} \leqno{(265)} $$
où $u_{\rho, \sigma}$ est défini par (150, 151), et $\lambda = x \lambda_0 + y \lambda_1$, avec $x, y \in \Gamma_{\mathbf{Q}}$ d'après (264), et $\nu = \nu(\lambda)$.

\textbf{Corollaire.} Le morphisme commutant
$$ \text{Scindage}(\mathcal{G}_{\rho, \sigma} \to \mathcal{D}_{12} \times \mathcal{G}) \to \text{Scindage}(\mathfrak{S}^+ \dots) \leqno{(266)} $$

\newpage

%% ===== Page 114 =====

et un isomorphisme de faisceaux dans $(Sch)$, et par suite (prop. p. 702) le premier est un $\mathcal{M}$-torseur via co-présentation.

Cet énoncé s'explicite en un énoncé correspondant sur les extensions, comme dans le cas précédent. On trouve en particulier, sur $S = \operatorname{Spec}(\mathbf{Z})$ :

Corollaire : Les scindages de l'extension $\mathcal{G}_{0,3}/\boldsymbol{\Gamma}$ de $\mathbf{Z}^\times \times \mathcal{G}_3$ par $\mathcal{M}$ forment (en restriction) un torseur sous $\mathcal{M}$, isomorphe au torseur des scindages de l'extension $\mathfrak{S}^{+ \wedge}_{0,3}$ de $\mathcal{G}_3$ par $\mathcal{M}$.

Variante de l'étude précédente pour l'extension $\mathcal{G}_{0,3}$ de $\Phi_{12}^\times \times \mathcal{H}_0^+$ par $\mathcal{M}$ (sans division par $\mathcal{H}_0 = \{1, \sigma\}_0$). Considérons d'abord le quotient des scindages de la sous-extension $\mathfrak{S}^+$ de $\mathcal{H}_0^+$ par $\mathcal{M}$ — ils correspondent aux couples.

\newpage

%% ===== Page 115 =====

\begin{cases} \rho'^3 = \sigma'^2 = \gamma, \quad (\gamma \in \Gamma M) \text{ tel que soit} \\ \rho'^3 = \sigma'^{-2} = \omega, \quad \sigma'(\rho') = \rho'^{-1} \end{cases} \leqno{(267)}

On interprète qu'un scindage dont transformation $\omega$, qui précise sur $M$ et sa centralisatrice des $\mathfrak{H}^+$, dans une extension centrale de l'extension
$$
\begin{cases} \operatorname{Aut}(\mathfrak{H}^+) = M \rtimes \operatorname{Aut}(\mathfrak{H}^+) \simeq \mathbb{H}_{1, \omega} \end{cases} \leqno{(268)}
$$
- ce qui montre qu'un scindage de transformation $\omega$ en est, d'où finalement la construction (267) des $\rho', \sigma'$ correspondants ! un scindage. On a
$$
\begin{cases} \rho'^3 = |\rho|^3 = (\epsilon | \rho) \rho^3 (\epsilon | \rho) = (\epsilon | \rho) \gamma \\ \sigma'^2 = (\gamma | \sigma)^2 = (\gamma | \sigma) \sigma^2 (\gamma | \sigma) = (\gamma | \sigma) \gamma \omega \end{cases}
$$
donc les deux équations $\rho'^3 = \omega$ et $\sigma'^2 = \omega$ équivalent respectivement :
$$
\begin{cases} (\epsilon | \rho) \gamma = \omega \quad \text{i.e.} \quad (1 + \operatorname{Id}_M - \rho)(\gamma) = \omega \\ (\gamma | \sigma) \omega = \omega \quad \text{i.e.} \quad (1 + \operatorname{Id}_M - \sigma)(\gamma) = 0 \end{cases}
$$
i.e. cela revient aux équations (245) (246) et à ce que le scindage garantit, un sous-groupe ... $\gamma$ ... $\gamma = \dots$
La condition $\sigma'(\rho') = \rho'^{-1}$ s'explicite d'autre part comme celle du p. 700', ...

\newpage

%% ===== Page 116 =====

sur l'équation (249), qui s'explicite par (250).

Cela montre qu'il existe un sous-groupe $w$ — un sous-groupe $Z$ — qui vérifie : un scindage de l'extension $\mathfrak{S}^+$ est $\simeq \mathfrak{T}_0$, $H^+ = \mathfrak{S}^+ / M$ (avec $M \simeq \hat{\mathbf{Z}}$) — et que les scindages en question fixent un sous-groupe sur $K_0 \simeq \hat{\mathbf{Z}}$ défini dans (253), et — soit $K$ [unclear: comme] $M \simeq K$ (254), et $M \simeq K$, correspondant à l'opération de $M$ sur le relevé des scindages, par conjugaison. D'où en finissant...

\section*{Proposition (Sur le cas général).}

L'extension $\mathfrak{S}^+ \text{ de } H^+ \text{ par } M$ a scindage de sous-groupes — et conjugués par un unique section de $M$ — d'où l'application $\operatorname{Im } g_0$ — qu'on peut par $\mathcal{H}^1$ définit une bijection de l'ensemble des scindages de $\mathfrak{S}^+$ sur $H^+$, vers l'ensemble des scindages de $\mathfrak{S}^+ / \mathfrak{H}^+ \text{ sur } H^+ / \mathfrak{H}^+ \simeq \mathfrak{S}_3$.

\newpage

%% ===== Page 117 =====

Étudions maintenant le centralisateur des $\mathfrak{G}_{g,\sigma}$ d'un sous-groupe section au-dessus de $H^+$, disons $H(\varphi, \sigma)$. On doit avoir maintenant :

$$
\operatorname{Centr}_{\mathfrak{G}_{g,\sigma}}(H(\varphi, \sigma)) = \mathfrak{G}_{12}^* \cdot (M \times \{e\}) \leqno{(269)}
$$
(où $\lambda \omega \mu$ avec $\lambda \in \Gamma_{\mathbf{Q}}^*$, $\mu \in M$, $\omega \in \mathfrak{G}_{12}^*$). D'ailleurs $\lambda \omega \mu$ centralise $H(\varphi, \sigma)$ si et seulement si les deux centralisent. La condition de $\sigma$ s'explicite avec la condition (260) du calcul p. 703, intégrée, ou voisine. Pour la condition de commutation avec $\sigma$, le calcul de p. 703, avec $\omega^2 = 1$, la condition s'écrit
$$
\omega \lambda \eta \sigma(\lambda^{-1}) \sigma = \eta \sigma \quad \text{ou } \quad \omega^2 \lambda \eta \sigma(\lambda^{-1}) = \eta
$$
et équivaut à la conjonction de la condition
$$
v_2(\lambda) = 0 \quad \text{i.e. } \mu \in \Gamma_{\mathbf{Q}}^{\times 11} \leqno{(270)}
$$
et de la relation (261). En poursuivant le calcul de p. 703, 704, s'appliquent finalement:
\textit{Le centralisateur dans $\mathfrak{G}^+$ d'un sous-groupe section $H(\varphi, \sigma)$ des $\mathfrak{G}_{g,\sigma}$ au-dessus de $H^+$...}

\newpage

%% ===== Page 118 =====

On a alors une suite exacte
$$ 1 \to H^+ \to \mathcal{G}_{g,\sigma} \to \mathcal{Q}_{12}^* \to 1 \leqno (278) $$
et pour $\mathcal{T}_{\sigma} \subset \mathcal{G}_{g,\sigma}$,
$$ \mathcal{T}_{\sigma} \cong \mathcal{Q}_{12}^* \leqno (278\text{bis}) $$
i.e. $\mathcal{T}_{\sigma}$ est une section petite de $\mathcal{G}_{g,\sigma}$ au dessus de $\mathcal{Q}_{12}^*$.

\noindent Corollaire. L'extension $\mathcal{G}_{g,\sigma}$ de $\mathcal{Q}_{12}^* \times H^+$ par $M$ n'est scindée sur aucune base $S$ non vide. L'extension $\mathcal{G}_{g,\sigma}$ de $\mathcal{Q}_{12}^* \times H^+$ par $M$ est scindée sur la base $S$, et deux scindages sont conjugués par une unique section de $M$. L'application centrale fournit une bijection des scindages de $\mathcal{G}_{g,\sigma}$ sur $\mathcal{Q}_{12}^* \times H^+$, avec l'un des scindages de $\mathcal{Q}^+_{12}$ sur $H^+$.

\noindent Corollaire. L'ext. $\mathcal{G}_{g,\sigma}$ de $\hat{\mathbf{Z}}^* \times H^+$ par $M$ n'est pas scindée. Celle de $\mathcal{G}_{g,\sigma}$ de $\hat{\mathbf{Z}}^* \times H^+$ par $M$ est scindée, les scindages sont conjugués par un unique élément de $M$. L'application centrale définit une bijection de l'un des scindages de $\mathcal{G}_{g,\sigma}$ sur $\hat{\mathbf{Z}}^* \times H^+$ avec l'un des scindages de $\hat{\mathcal{G}}$ sur $H^+$.

\newpage

%% ===== Page 119 =====

Explication d'un scindage de l'extension $\mathfrak{S}^+_{g,\nu}$ de $H^+$ par $M$. --- Correspondance : un scindage de $\mathfrak{S}^+_{g,\nu}$ de $H^+$ par $M$, i.e. une solution particulière de l'équation
$$
a_0 - 2a_1 + 3b - 1 = 0 \leqno (273)
$$
en posant $a_1 = b = 0$, d'où $a_0 = 1$.

i.e.
$$
\left\{ \begin{aligned} &\xi = \lambda_0, \eta = 0 \text{ (écriture additive)} \\ &\rho' = \lambda \circ \rho, \quad \sigma' = \sigma \end{aligned} \right. \leqno (274)
$$
les $\text{formules fournissent un scindage de } \mathfrak{S}^+_{g,\nu}$. Le scindage correspondant on essaie de $\mathfrak{S}^+_{g,\nu} \cong H^+/\mu_H^+$. --- Explicitation reliante la structure correspondante à celle de $\mathbb{Q}_{12}^*$ dans $\mathfrak{G}_{g,\nu}$, d'où
$$
\left\{ \begin{aligned} &\mu_{|\rho'} : \mathbb{Q}_{12}^* \to \mathfrak{G}_{g,\nu}/\mu_H^+, \text{ resp. } \mathbb{Q}_{12}^* \to \mathfrak{G}_{g,\nu} \\ &\mu \longmapsto \lambda_0 \mu \quad (\nu = \nu(K)) \end{aligned} \right. \leqno (275)
$$
où $\mu$ défini dans (151). (NB (275) est contenu dans rappel p. 104 (265) et (264), qui donne $\mu_{\rho'} = \lambda \mu, \lambda = x\lambda_0 + y\lambda_1, x=y=-\nu \implies \lambda = -\nu(\lambda_0 + \lambda_1) = \nu\lambda_0$).

On arrive donc, posant $u'_H = u_{\rho', \sigma'}(H)$,
$$
\left\{ \begin{aligned} &u'_H(\rho') = \rho', \quad u'_H(\sigma') = \sigma' \\ &u'_H(\lambda) = \mu \lambda (\lambda \in \Gamma) \quad (\mu = \mu(H)) \end{aligned} \right\} \text{(à condition } u'_H) \leqno (276)
$$

\newpage

%% ===== Page 120 =====

et on en déduit
$$
\begin{cases} u'_K(\rho) = \lambda^{-2\nu} \rho \\ u'_K(\sigma) = \omega^{\nu/2} \sigma \end{cases} \quad (= \sigma \text{ dans } \mathfrak{S}_{g, \nu} / \mu_G \dots) \leqno{(277)}
$$

\newpage

%% ===== Page 121 =====

70bis

13°) Retour sur la tour de Fermat.

Commençons par dégager une suite topologique générale, qui était un peu embrouillée entre les lignes dans 110).

Soit $\Gamma$ (ici $\Gamma_{\mathbb{Q}} \times \mathfrak{S}_3 = \mathcal{N}_{0,3}$) opérant sur un topos $\mathcal{X}$ (ici $\underline{M}_{0,3} \simeq M_{0,4 \overline{\mathbb{Q}}}$) conn. loc. connexe, muni d'un "pt" (géométrique) $x_0$ (ici : le point $\vec{1}$ à valeurs dans $\mathbb{Q}$), soient

[MARGIN: en fait on discute [unclear: localement] ($\mathcal{X}$ loc. simplemt [unclear: connexe] ou profini, il faudra y revenir...)]

$$ (278) \begin{cases} \Pi = \Pi_1(\mathcal{X}, x_0), \mathcal{M} = \Pi_1(\mathcal{X}, \Gamma; x_0) \\ M = \Pi_{ab} = \Pi/[\Pi,\Pi] \quad, \tilde{\mathcal{M}} = \mathcal{M}/[\Pi,\Pi] \end{cases} $$

d'où des suites exactes et hom. de suites exactes

(279)
\begin{tikzcd}
1 \ar[r] & \Pi \ar[r] \ar[d, "epi"'] & \mathcal{M} \ar[r] \ar[d, "epi"'] & \Gamma \ar[r] \ar[d, "\sim"'] & 1 \\
1 \ar[r] & M \ar[r] & \tilde{\mathcal{M}} \ar[r] & \Gamma \ar[r] & 1
\end{tikzcd}

(Au lieu de prendre $M = \Pi/[\Pi,\Pi] = \Pi_{ab}$ on aurait pu prendre un quotient de $\Pi_{ab}$ invariant par action de $\Gamma$, mais peu importe...). Considérons "le rev. universel abélien" $\tilde{\mathcal{X}}_{ab}$ de $\mathcal{X}$ en $x_0$, i.e. le quotient par $[\Pi,\Pi]$ du rev. universel de $\mathcal{X}$ en $x_0$ - c'est donc un rev. principal de $\mathcal{X}$, de groupe $M$. On

\newpage

%% ===== Page 122 =====

je propose d'examiner les façons de relever l'action de $\boldsymbol{\Gamma}$ sur $\mathfrak{X}$ (ou $\mathfrak{X}_0$) en un $\boldsymbol{\Gamma}$-torseur compatible avec l'action de $\boldsymbol{\Gamma}$ sur le groupe d'opérateurs $M$ (ou plutôt un groupe de type $M$, cf. compatible avec l'action sur $\mathfrak{X}$).

Plus généralement, si $\mathfrak{X}'$ est un $M$-torseur sur $\mathfrak{X}$ isomorphe à $\mathfrak{X}_0$ déduit par trait; l'acide d'un $M$-torseur $E$ (le fibré de $\mathfrak{X}'_{x_0}$) étudie les relèvements de l'action de $\boldsymbol{\Gamma}$ sur $\mathfrak{X}$ en une action de $\boldsymbol{\Gamma}$ compatible avec l'action de $\boldsymbol{\Gamma}$ sur $\mathfrak{X}$.

Notons que, si $\widetilde{\mathfrak{X}}$ désigne le rev.\ universel de $\mathfrak{X}, x_0$, $M$ opère sur $\widetilde{\mathfrak{X}}$ de façon compatible avec l'action canonique sur $\mathfrak{X}$, et puisque $\pi_1 = [\pi, \pi]$ est distingué dans $M$, l'action de $\boldsymbol{\Gamma}$ sur le quotient canonique $\mathcal{G} = M/[\pi, \pi]$ sur $\mathfrak{X}' = \widetilde{\mathfrak{X}}/[\pi, \pi]$, qui satisfait aux...

\newpage

%% ===== Page 123 =====

\setcounter{page}{122}

\section*{}
\marginpar{709}
conditions analogues. Soit $\theta$ cette application, et $\theta$ une opération quelconque, satisfaisant auxdites conditions.

On veut donc faire de $\mathfrak{X}'$, qui est un revêtement de $\mathfrak{X}$, donné par le fibré $E$ et l'action de $\pi = \pi_1(\mathfrak{X}, x_0)$ sur $E$ (par translations, grâce à une structure de $M$-torseur sur $E$) un revêtement de $(\mathfrak{X}', \mathcal{G})$ --- i.e.\ prolonger l'action donnée de $\pi$ sur le fibré $E$, une action de $M$ sur $E$ --- qui d'ailleurs se factorise nécessairement par $\mathcal{M}$.

Je dis que cette action se fait victorieusement via le groupe affine de torseur $E$ --- de façon précise, $\Gamma$ opère sur $M$, soit $\gamma_M$ l'opération définie par $\Gamma$, et soit $g \in \mathcal{M}$ dessinée par $\gamma \in \Gamma$, --- on a donc, pour $x \in E, \lambda \in M$ :
$$
g(\lambda \cdot x) = g(\lambda) \cdot gx \leqno{(280)}
$$
ce qui ressort trivial par le calcul :
$$
g(\lambda \cdot x) = (g\lambda g^{-1}g)x = (g\lambda g^{-1}) \cdot (gx) \leqno{(281)}
$$
$$\gamma_M(\lambda)$$

\newpage

%% ===== Page 124 =====

Donc si on introduit le groupe $Aff(E)$
on aura un homomorphisme de
suites exactes

(281)
\begin{tikzcd}
1 \ar[r] & M \ar[r] \ar[d, "id_M"] & \tilde{\Gamma} \ar[r] \ar[d, "\theta_{\tilde{\Gamma}}"] & \Gamma \ar[r] \ar[d, "\Theta_\Gamma"] & 1 \\
1 \ar[r] & M \ar[r] & Aff(E) \ar[r] & \operatorname{Aut}(M) \ar[r] & 1
\end{tikzcd}

où les flèches verticales extrêmes $id_M$
et $\Theta_\Gamma$ sont connues, et où la flèche médiane
est à déterminer. On peut dire
que la catégorie des $\tilde{X}'$ au dessus
d'un [unclear: l'é.p.] de $\Gamma$ dessus satisfaisant
aux autres [unclear: axiomes] ) équivaut
à celle des couples $(E, \theta_{\tilde{\Gamma}})$, où $E$
est un $M$-torseur, et $\theta_{\tilde{\Gamma}}$ un
hom. de groupes $\tilde{\Gamma} \to Aff(E)$
compatible avec $id_M$ et $\Theta_\Gamma$ i.e. qui
rend commutatif (281). On
peut dire aussi une fois $E$ choisi,
(et ce choix est unique à isom. [unclear: unique] près -)
qu'on a défini une extension
$Aff(E)$ de $\operatorname{Aut}(M)$ par $M$, d'où
[unclear: tirée par] $\Gamma \to \operatorname{Aut}(M)$ une extension

\newpage

%% ===== Page 125 =====

\marginpar{Les scindages de $\mathcal{A}ut(E)$ déterminés par les points de départ de la scission.}

(710) qui d'ailleurs s'identifie, sur $M$, (ce qui dit, les relèvements de $\Gamma$ dans $\mathcal{A}ut(E)$ qui sont des isomorphismes de cette extension de $\Gamma$ par $M$ sur l'extension $\mathcal{E}$ de $\Gamma$ par $M$). Donc elle existe, c'est l'extension $\mathcal{E}$ de $\Gamma$ par $M$.

$$
c^2(\mathcal{E}) \in H^2(\Gamma, M) \text{ est nul.} \leqno{(282)}
$$

S'il existe un relèvement, c'est-à-dire un isomorphisme, alors on fixe un torseur sous le groupe des centres de l'extension (triviale) $\mathcal{E}$ de $\Gamma$ par $M$, qui s'identifie au groupe $Z^1(\Gamma, M)$ des homomorphismes croisés de $\Gamma$ dans $M$. Les classes de $M$-conjugaison de tels relèvements comme isomorphismes correspondent aux éléments de $H^1(\Gamma, M)$.

Une autre façon de le dire, en posant $E = M$-torseur trivial $= M$, de $\mathcal{A}ut(E)$ sont les relèvements divisés eux-mêmes, c'est-à-dire scindages de l'extension $\mathcal{E}$ de $\Gamma$ par $M$, et les classes de conjugaison de relèvements d'action de $\Gamma$ sur $E$.

\newpage

%% ===== Page 126 =====

correspondant aux classes de $M$-conjugaison de scindages de ladite extension $\mathfrak{S}$ par $M$.

Ceci bien posé, revenons à la situation de $(U_{0,3} \otimes \overline{\mathbf{Q}}, \Gamma = \widehat{\Pi}_{0,3} \times \mathfrak{S}_3)$, les $\Gamma$-revêtements de $U_{0,3} \otimes \overline{\mathbf{Q}}$, i.e.\ des $\mathfrak{S}_3$-revêtements de $(U_{0,3}, \mathfrak{S}_3)$. \marginpar{NB. le groupe $G_{\sigma}$ mentionné ici $M_{0,3}$ devient $\widehat{\Pi}_{0,3}$} \marginpar{en posant $\widehat{\Pi} = \widehat{\Pi}_{0,3}$}

À l'aide des constructions générales du $\S$ VI,
$$
\mathcal{G}^{+\wedge} / \mu_6 \simeq \mathfrak{S}_{0,3}^+ / [\widehat{\pi}_0, \widehat{\pi}_0] \defeq \mathcal{M}
$$
(sous les notations des pages précédentes), le groupe profini $\mathcal{G}$ qui était noté $\mathcal{G}_{\sigma}$ dans les sections précédentes, apparaissait comme un produit semi-direct $\mathcal{G}_{\sigma} = \widehat{\pi}_0 \rtimes \mathcal{M}$, extension de $\widehat{\mathbf{Z}} \times \mathfrak{S}_3$ par $\mathcal{M}^\wedge$ et un homomorphisme canonique
$$
\mathcal{M} = \mathcal{M}_{0,3} / [\widehat{\pi}_0, \widehat{\pi}_0] \to \mathcal{G}_{\sigma}
$$
d'origine d'une suite exacte.

\newpage

%% ===== Page 127 =====

711

(28bis)
\begin{tikzcd}
1 \ar[r] & \hat{M} \ar[r] \ar[d] & \tilde{\mathcal{M}} \ar[r] \ar[d] & \Gamma \ar[r] \ar[d, "N_{0,3}=\Gamma_{\mathbb{Q}} \times \mathfrak{S}_3" description, "\hat{\mathbb{Z}}^\times \times id_{\mathfrak{S}_3}"] & 1 \\
1 \ar[r] & \hat{M} \ar[r] & G_{\tilde{\rho}} \ar[r] & \hat{\mathbb{Z}}^\times \times \mathfrak{S}_3 \ar[r] & 1
\end{tikzcd}

D'ailleurs comme l'on attend - on a trouvé que
l'extension $G_{\tilde{\rho}}$ de $\hat{\mathbb{Z}}^\times \times \mathfrak{S}_3$ par $\hat{M}$ était
scindée, d'où il résulte que
l'ext. $\tilde{\mathcal{M}}$ de $\Gamma$ par $\hat{M}$ l'est aussi
(ce qui au demeurant est aussi conséquence de
l'existence de la tour de Fermat, en
y appliquant les réflexions générales
des pages précédentes). Bien mieux,
on a vu qu'il y a une unique classe
de $\hat{M}$-conj. de scindages de l'extension
$G_{\tilde{\rho}}$, ce qui induit une classe de
$\hat{M}$-conjugaison canonique (mais non
unique!) de scindages de $\tilde{\mathcal{M}} \to$ 
de $\Gamma$. Cela définit donc une
$\hat{M}$-conjugaison canonique de relèvements
sur $\tilde{X} = \widetilde{U_{0,3 \overline{\mathbb{Q}}}}$
de l'action de $\Gamma$ (à $\simeq$ près, ou ce
qui revient au m., un foncteur sur
$(U_{0,3 \mathbb{Q}}, \mathfrak{S}_3)$, de fibre un groupe profini

[MARGIN: NB $H^1(\Gamma, \hat{M})$ est un groupe fort gros! Mais voir cf th. p. 716]

\newpage

%% ===== Page 128 =====

... commutatif ... dessus $\mathcal{M}_{0,3}(\mathbf{Q}, \mathfrak{S}_3)$ qui ... dessus déjà $\mathcal{M}_{0,3}$ soit isomorphe au revêtement $\mathcal{X}$ envisagé plus haut.

Revenons pour un instant à ce qui vient de tout à l'heure. Notons que le donné d'un diagramme (281) important définit un quotient de $\tilde{\mathcal{M}}$ donnant lieu à une suite exacte :
$$
1 \to M \to \tilde{\mathcal{M}} \to \tilde{\Gamma} \to 1 \leqno{(283)}
$$
($\tilde{\mathcal{M}}/I = M$)

$$
\tilde{\Gamma} \to \\operatorname{Aut}(M) \leqno{(284)}
$$
$\tilde{\Gamma}$ s'identifie à l'image de $\tilde{\Gamma}$ dans $\\operatorname{Aut}(M)$ par l'action de $\tilde{\Gamma}$ sur $M$ [préserve] les structures d'extension $\tilde{\mathcal{M}}$ de $\tilde{\Gamma}$ sur $M$. Cette $\mathcal{G}$ qui à $\tilde{\mathcal{M}}$ et $\mathcal{M}$ donne comme on a vu un calcul (283), c'est-à-dire un groupe ... à ce fait que $\tilde{\mathcal{M}}$ s'y envoie injectif, et que $\tilde{\Gamma} = \tilde{\mathcal{M}} / \dots$

\newpage

%% ===== Page 129 =====

\hfill 712

opère fidèlement sur $M$ : en plus, l'extension (283) est triviale, et le fait qu'elle provient d'un $\tilde{\theta}_{\tilde{u}}$ implique une structure supplémentaire sur cette extension ; à savoir la donnée d'une classe de $M$-conjugaison de scindages. Plus précisément, nous avons une application canonique
$$
E \to \operatorname{Scind}(\tilde{M} \text{ sous-groupe ext. de } \tilde{\Gamma} \text{ par } M) \leqno{(285)}
$$
compatible avec les opérations de $M$ sur $E$ (par translation) et sur $\operatorname{Scind}$ (par construction, les opérations). La donnée de $M$-conjugaison, ainsi que le scindage. Pour que $x$ et $y \in E$ soient des images dans $\operatorname{Scind}$, il faut que l'élément $u=y^{-1}x \in M$ soit dans $M^\Gamma = M^{\tilde{\Gamma}}$, ce qui montre que
$$
\text{l'homomorphisme } E \to \operatorname{Scind} \text{ est injectif si et seulement si } M^\Gamma = \{e\} \leqno{(286)}
$$

Dans le cas $M^\Gamma = \{e\}$, on voit donc que $E$ s'identifie exactement à un domaine de $M$-conjugaison de scindages de $\tilde{M}$.

\newpage

%% ===== Page 130 =====

Résumé des considérations ci-dessus

Proposition. Soient $\mathcal{X}$ un topos connexe localement connexe muni d'un point géométrique $x_0$, on suppose de plus les extraits discrètes ($\mathcal{X}$ localement simple, c'est-à-dire [unclear: signifiant] revêtement profini).
On a un groupe $\Gamma$ d'automorphismes de $\mathcal{X}$ ($\Gamma$ profini dans le cas profini — il faudrait expliciter pour bien faire les hypothèses basiques : faire [unclear: bien] compte de la topologie profinie de $\Gamma \dots$). On pose
$$
\pi = \pi_1(\mathcal{X}, x_0), M = \pi_1(\mathcal{X}, \bar{x}_0) \quad \text{d'où} \quad (*) \quad 1 \to M \to \pi \to \Gamma \to 1
$$
où $M$ est un sous-groupe distingué de $\pi$, et que [unclear: aura] une question de $\bar{\pi} = \pi/[\pi, \pi]$ stable par l'action de $\Gamma$, d'où pour $M^{\sim} = M^{\text{ext}}$ quotient de $(*)$
$$
(287) \quad 1 \to M \to \widetilde{M} \to \Gamma \to 1
$$
où $M$ commute avec $\Gamma$ agissant sur $M$.
On s'intéresse à la catégorie des revêtements $\mathcal{X}'$ de $(\mathcal{X}, \Gamma)$ satisfaisant la condition.

\newpage

%% ===== Page 131 =====

suivante : après l'oubli des opérations de $\Gamma$ on trouve un revêtement [unclear: connexe & abélien] galoisien de $\mathfrak{X}$, correspondant à un groupe quotient $M$ de $\pi = \pi_1(\mathfrak{X}, x_0)$ - de sorte que $\mathfrak{X}'$ devient de façon naturelle un torseur sous $M_{\mathfrak{X}}$.

Le foncteur
$$ \mathfrak{X}' \longmapsto \mathfrak{X}'(x_0) \quad (\text{fibre en } x_0) $$
de la catégorie de ces $\mathfrak{X}'$ dans la catégorie des ensembles, s'en va en fait vers les ensembles $E$ munis des structures suivantes

a) $E$ est un $M$-torseur, et en tant que tel donne naissance à une extension
$$ 1 \to M \to \text{Aff}(E) \to \operatorname{Aut}(M) \to 1 $$
d'où une sous-extension induite

(288)
$$ 1 \to M \to \text{Aff}_\Gamma(E) \to \tilde{\Gamma} \to 1 $$

où $\tilde{\Gamma} = \operatorname{Im}(\Gamma \to \operatorname{Aut}(M))$

b) On a un hom d'extensions

(289)
\begin{tikzcd}
1 \ar[r] & M \ar[r] \ar[d, "id"'] & \tilde{M} \ar[r] \ar[d, "\theta_{\tilde{M}}"'] & \Gamma \ar[r] \ar[d, "can"] & 1 \\
1 \ar[r] & M \ar[r] & \text{Aff}_\Gamma(E) \ar[r] & \tilde{\Gamma} \ar[r] & 1
\end{tikzcd}

La catégorie des $\mathfrak{X}'$ (à un iso...) équivaut à celle des couples $(E, \theta_{\tilde{M}})$, où $E$ est un $M$-torseur, et $\theta_{\tilde{M}}$ un hom. d'extensions (289).

\newpage

%% ===== Page 132 =====

Les automorphismes d'un objet $\mathcal{X}'$ correspondent aux « translations » par des éléments de $M$ (en termes de $(E, \Theta_{\widetilde{M}})$) aux translations de $E$ par de tels éléments de $M$ qui sont dans $\operatorname{Aff}(E)$ (ils sont centraux et se conjuguent par ces éléments dans le ( = hom. $\Theta_{\widetilde{M}}$ ...)). La classe d'isomorphie de ces éléments [unclear: (définie par)] ... binaires de la catégorie (287) de $\Gamma$ et $\widetilde{M}$. De façon précise : un $(E, \Theta_{\widetilde{M}})$ ... un homomorphisme de $M$-ensembles
$$
E \to \operatorname{Scind}(\widetilde{M} \text{ ext. de } \Gamma \text{ par } M) \leqno{(290)}
$$
et à l'image ... décrit la classe d'isomorphie ... L'homomorphisme (290) se factorise par $M$
$$
M \to \operatorname{Scind}(\widetilde{M} \text{ ext. de } \Gamma \text{ par } M) \leqno{(291)}
$$
et si $M=0$, on voit que $E$ se constitue canoniquement à partir de la classe de

\newpage

%% ===== Page 133 =====

conjugaison envisagée. L'opération (définie par $\tilde{\Theta}_u$) de $\tilde{M}$ sur $E$ n'est autre que l'opération (ci-après posée) de conjugaison sur les scindages (que bien entendu on fait par section d'une classe de $M$-conjugués, $=$ conjugués par des éléments quelconques de $\tilde{M}$).

Corollaire. Supposons que l'on ait
$$ H^0(\Gamma, M) = H^1(\Gamma, M) = 0 \leqno{(292)} $$
i.e. que tous deux scindages d'une catégorie $F$ par $M$ soient conjugués par un élément unique de $M$. Ce qui, le cas $M=0$ impliquant que la catégorie des $F'$, ou mieux celle des couples $(E, \tilde{\Theta}_u)$, est rigide, alors les classes sont mieux prises à un unique isomorphisme près. Ceci posé, les classes d'isomorphisme en question (i.e. les classes de $M$-conjugués des scindages de l'ext. $\tilde{M}$ par $\Gamma, M$) correspondent biunivoquement aux groupes quotients de $\tilde{M}$.

\newpage

%% ===== Page 134 =====

de $M$ vérifiant les conditions suivantes :
\begin{enumerate}
\item[a)] $M \to \tilde{M}$ (pour $M \hookrightarrow \tilde{M}$) est injectif, d'où l'extension
$$
1 \to M \to \tilde{M} \to \Gamma' \to 1 \leqno{(293)}
$$
\item[b)] $\Gamma' = \tilde{M}/M$ opère fidèlement sur $M$, i.e. $\Gamma' \simeq \tilde{\Gamma} \subset \\operatorname{Aut}(M)$.
\item[c)] L'extension $\tilde{M}$ de $\tilde{\Gamma}$ par $M$ est triviale (i.e. sa classe dans $H^2(\tilde{\Gamma}, M)$ est nulle).
\end{enumerate}

Pour un tel groupe $\tilde{M}$, on désigne $E$ comme l'ensemble de tous les sous-groupes de l'extension $\tilde{M}$ de $\tilde{\Gamma}$ par $M$, sur lequel on fait opérer $\tilde{M}$ via $\tilde{M}$ par conjugaison — la restriction de cette opération : $M \subset \tilde{M}$ fait de $E$ un $E$-tordu (par l'hyp. (292)).

Revenons à nouveau en ce particulier qui nous intéresse. Au principe (120) on a vérifié les relations (292), maintenant on se demande quels groupes extérieurs $\tilde{\Gamma} \simeq \hat{\Gamma} \times \mathcal{G}$ par $M$ (ou $M$ dans les considérations générales).

\newpage

%% ===== Page 135 =====

était triviale, et que deux scindages étaient conjugués par un élément de $\mathcal{M}$ et un seul (c'est la deuxième assertion d'existence et unicité de conjugaisons qui est le contenu géométrique de (292)).

De plus, on a donc construit justement, ici, un homomorphisme surjectif
$$
\mathcal{M} = \mathcal{M}_{0,3} \twoheadrightarrow \mathcal{G}_{\sigma}
$$
faisant donc de $\mathcal{G}_{\sigma}$ une extension de $\mathcal{M}$ et de $\tilde{M}$, avec les propriétés a) b) c) du corollaire précédent. Cela nous définit donc, à isomorphisme près, un objet de la catégorie des $\mathfrak{T}$, i.e. un pro-revêtement de on a que celui-ci, un pro-revêtement de $(U_{0,3}, \mathfrak{G}_3) \cong \mathcal{M}_{1,1, \overline{\mathbf{Q}}}$, que après extension du groupe de $\mathfrak{G}_3$ à extension de la loi de $\mathbf{Q} \to \overline{\mathbf{Q}}$, donne le revêtement universel défini de $U_{0,3, \overline{\mathbf{Q}}}$.

Mais d'autre part le tour des revêtements de $\mathfrak{F}$ nous fournit très exactement un tel objet — et si

\newpage

%% ===== Page 136 =====

je n'avais pas eu à ma disposition
cette tour de Fermat, je n'aurais
certes pas suspecté que [crossed out: l'extension] $\widetilde{C}^{+\wedge} = \widetilde{C}^{+\wedge}_{H}$ de
$H^+ \simeq \widehat{\mathfrak{S}}_3$ par $M^\wedge$ soit scindée !

La conjecture qui se présente tout de
suite, c'est que ces deux revêtements
pro-étales sont isomorphes - l'isomor-
phisme entre les deux est alors
unique.

Pour l'autre voie de détermination
des classes d'isomorphie [unclear: (c.a.d. $\tilde{M}$ sur $M_{0,3} / [\Gamma_{\mathbb{Q}}, \Gamma_{\mathbb{Q}}]$)]
de scindages de $\Gamma = \Gamma_{\mathbb{Q}} \times \mathfrak{S}_3$ par
$M^\wedge$, on a par la suite exacte associée
à la suite spectrale de Hochschild

$$ (294) \quad 0 \to H^1(\Gamma_{\mathbb{Q}}, H^0(\widehat{\mathfrak{S}}_3, \widehat{M})) \to H^1(\Gamma, \widehat{M}) \to H^0(\Gamma_{\mathbb{Q}}, H^1(\widehat{\mathfrak{S}}_3, \widehat{M})) $$

Or on a vu que
$$ (295) \quad H^0(\widehat{\mathfrak{S}}_3, \widehat{M}) = 0 \quad (\text{[unclear: car] } H^0(L_f, \widehat{M}) = 0) $$

[crossed out: Et] d'autre part
$$ H^1(\widehat{\mathfrak{S}}_3, \widehat{M}) = H^1(\mathfrak{S}_3, M) \otimes_{\mathbb{Z}} \hat{\mathbb{Z}} \quad , $$

[MARGIN: donc $H^1(\Gamma, \widehat{M}) \simeq H^0(\Gamma_{\mathbb{Q}}, H^1(\widehat{\mathfrak{S}}_3, \widehat{M}))$]

\newpage

%% ===== Page 137 =====

715

Ce qui nous ramène à calculer
$$
H^1(\mathfrak{S}, \widehat{M}) ?
$$
Mais l'étude faite de l'extension $\mathfrak{S}/\mathfrak{T}_{\hat{\mathbf{Z}}}$ de $\mathfrak{S}$ par $\widehat{M}$ — qui nous a montré que cette extension était scindée, et que les scindages étaient conjugués par un unique scindage de $\widehat{M}$, nous donne un fait
$$
\left. \begin{aligned} H^2(\mathfrak{S}, \widehat{M}) &= H^2(\mathfrak{S}, \widehat{M}) = 0 \\ H^0(\mathfrak{S}, \widehat{M}) &= H^1(\mathfrak{S}, \widehat{M}) = 0 \end{aligned} \right\} \quad \text{d'où} \leqno{(296)}
$$
d'où finalement, grâce à (295), la relation triomphale
$$
H^1(\Gamma_{\mathbf{Q}} \times \mathfrak{S}, \widehat{M}) = 0 \leqno{(297)}
$$
!

On trouve donc — ô merveille — entre tout abouti, je l'avoue ! —
le
\textbf{Théorème.} Il existe un revêtement \smash{\rlap{$\underbrace{\phantom{= \widehat{M}_{\mathbf{Q}}}}_{= \widehat{M}_{\mathbf{Q}}}$}} proétale $U_\infty$ de $(U_{0,3} \otimes_{\mathbf{Q}} \mathfrak{S})$, ayant la propriété suivante : L'image inverse de $U_\infty$ par $U_{0,3} \otimes \overline{\mathbf{Q}}$ (revêtement proétale de $\mathfrak{S}$) est "le" revêtement abélien universel de $U_{0,3} \otimes \overline{\mathbf{Q}}$. Pour deux tels revêtements

\newpage

%% ===== Page 138 =====

mants prostales de $(U_{0,3} \mathbf{Q}, \mathfrak{S}) = M_{1, \mathbf{Q}}$, il existe un unique isomorphisme de l'un sur l'autre.

Nous avons obtenu l'existence de deux façons : par la tour de Frattini et par la théorie générale du §9, appliquée au cas présent, ce qui nous donne un $U_{\infty}$ et un $U'_{\infty}$, chacun décrit à isomorphismes près !

Le théorème d'unicité résulte du petit calcul précédent — j'avais pour des "gens" [unclear: ...],
$$
H^i(\Gamma_{\mathbf{Q}}, \hat{M}) \simeq M \hat{\otimes} H^i(\Gamma_{\mathbf{Q}}, \mathbf{Z}_{\ell}(1))
$$
mais finalement c'est le calcul de $H^i(\mathfrak{S}, M)$ pour $i=0, 1$ qui implique (en lien avec l'action sur $\Gamma_{\mathbf{Q}}$ !) celui de $H^i(\Gamma_{\mathbf{Q}} \times \mathfrak{S}, \hat{M})$, qui donne l'unicité qu'il nous fallait.

\newpage

%% ===== Page 139 =====

717

Le théorème précédent élucide complètement, il me semble, la signification authentico-géométrique (en termes de la tour de Fauré) de l'homomorphisme (1218) p. 695 de
$$
\mathcal{M} = \mathcal{M}_{1,1} \to G_{\Gamma, \sigma} \simeq \hat{\mathbf{Z}}^* \times \hat{\mathfrak{S}}^\wedge
$$
obtenu dans la section 11°) par voie purement algébrique. On en tire celui de l'homomorphisme
$$
\mathcal{M} \to G_{\Gamma, \sigma} / \mu_6 \simeq \hat{\mathbf{Z}}^* \times \hat{\mathfrak{S}}^\wedge / \mu_6
$$
qui se factorise par le sous-groupe
$$
\tilde{\mathcal{M}} = \mathcal{M}/\mu_6 = \mathcal{M}_{\mu_6}\footnote{noté $\tilde{\mathcal{M}}$ des pages précédentes} \to G_{\Gamma, \sigma} \cong \hat{\mathbf{Z}}^* \times \tilde{\mathfrak{S}}_{0,3}
$$
Le groupe de $G_{\Gamma, \sigma} \simeq \tilde{G}_{\Gamma, \sigma}$ des genres 0,1, de un actuel, avait comme rôle essentiel (en particulier!) de fournir une extension par $\mathcal{M}^\wedge$ (contenant la s.g. d'une composante neutre) d'un groupe opérant fidèlement sur $\mathcal{M}^\wedge$, qui n'était pas le cas de l'ext. $G_{\Gamma, \sigma}$ de $\hat{\mathbf{Z}}^* \times \hat{\mathbf{H}}^\wedge$ par $\mathcal{M}^\wedge$, puisque les éléments de $\mu_6 \subset \mathcal{M}^\wedge$ y opéraient trivialement.

\newpage

%% ===== Page 140 =====

Ceci, joint aux relations
$$H^\circ(\hat{\mathbb{Z}}^* \times \mathfrak{S}_3, \hat{M}) = H^1(\hat{\mathbb{Z}}^* \times \mathfrak{S}_3, \hat{M}) = 0 \quad,$$
permettait d'interpréter le groupe
$G_{\rho, \tilde{\sigma}}$ comme un groupe de transformations
sur l'âme $E$ de ce [unclear: pro-jeu singulier],
ce qui permettait d'interpréter
l'hom $\mathcal{M} \to G_{\rho, \tilde{\sigma}}$ (ou $\mathcal{M}^\sim \to G_{\rho, \tilde{\sigma}}$)
comme une spécification de $\mathcal{M}$
(ou $\tilde{\mathcal{M}} = \Pi_1(U_{0,3 \mathbb{Q}}, \underline{\mathfrak{S}}_3) = \Pi_1(M_{1,1 \mathbb{Q}}, \dots)$
comme, [unclear], sur $E$ [unclear: correspondant]
à un revêtement de $(U_{0,3 \mathbb{Q}}, \underline{\mathfrak{S}}_3)$
$\simeq M_{1,1 \mathbb{Q}}$ - et ce revêtement
n'est autre que le prorevêtement
de Fermat.

Il resterait à interpréter de façon
analogue l'hom $\mathcal{M} = \mathcal{N}_{1,1} \to G_{\rho, \sigma}$,
s'insérant dans l'hom. de suites exactes

(298)
\begin{tikzcd}
1 \ar[r] & \mu_M \ar[r] \ar[dd] & \mathcal{M} \ar[r] \ar[dd] & \mathcal{M}^\sim \ar[r] \ar[dd] & 1 \\
& (\pm 1, \omega_f) & (\mathcal{N}_{1,1}) & (\tilde{\mathcal{N}}_{1,1} = \mathcal{M}_{0,3}) & \\
1 \ar[r] & \mu_G \ar[r] & G_{\rho, \sigma} \ar[r] & G_{\rho, \tilde{\sigma}} \ar[r] & 1 \\
& \pm 1, \omega_f & & &
\end{tikzcd}

Pour [unclear] d'interpréter d'ailleurs
$\mathcal{M}$, en tant qu'extension centrale de $\mathcal{M}^\sim$

\newpage

%% ===== Page 141 =====

718

par $\mu$, comme image inverse de l'extension $\mathcal{G}_{\sigma}$ de $\mathcal{G}_{\tilde{\sigma}}$ par $\mu$, par l'homomorphisme $\widetilde{M} \to \mathcal{G}_{\tilde{\sigma}}$ précédent. Il y a un analogue frappant avec ce qui s'était passé, d'ailleurs, quand on était parti au $\S$ 48, et de certains développements qui liaient des groupes de type Lie $\\operatorname{Sl}(2, \hat{\mathbf{Z}})$ de $\mathcal{M}_{0,3}$, qu'on construit un homomorphisme
$$
\mathcal{M}_{0,3} \to \\operatorname{Sl}(2, \hat{\mathbf{Z}})/\pm 1
$$
D'où il m'a été possible de reconstituer l'extension $\mathcal{M} = \mathcal{N}_{1,1}$ et $\mathcal{N}_{0,3}$ par ce jeu comme l'image inverse de l'extension $\\operatorname{Sl}(2, \hat{\mathbf{Z}})$ de $\\operatorname{Sl}(2, \hat{\mathbf{Z}})/\pm 1$ par $\mu = \{\pm 1\}$. Des points de vue actuel, j'essaye d'expliciter autant que possible des structures automorphiques en termes d'automorphismes des groupes à lacets $\pi_{0,3}$. Le fait remarquable, c'est que les deux

\newpage

%% ===== Page 142 =====

extensions de $\tilde{M}$ par $\mu$, sont-ce des calculs en apparence effrayants, soient conséquents [unclear: [?]]\dots Je pense qu'une explicitation géniale de ce fait suffirait, sans effrayer; dans le contexte actuel, du fait que, dans les deux cas, l'image de $\\operatorname{Sl}(2, \hat{\mathbf{Z}})/\pm 1$ sur $\mathcal{G}_{\sigma}$ provenant du fait que j'étais parti de un groupe ($\\operatorname{Sl}(\hat{\mathbf{Z}})$) — $\mathfrak{S}_{0,3}^{+\wedge} \cong \\operatorname{Sl}(2, \hat{\mathbf{Z}})/\pm 1$ — mais $\\operatorname{Sl}(2, \hat{\mathbf{Z}}) \cong \mathcal{N}_{1,1}$ \dots

Si on cherche malgré tout une interprétation (genre « tour de Fermat ») de l'homomorphisme $\mathcal{M} \to \mathcal{G}_{\sigma}$ (comme $\SL_2$ sur $\mathcal{G}_{\sigma}$) il faudrait interpréter $\mathcal{G}_{\sigma}$ (comme un groupe d'automorphismes d'un « fibre-type » « pas trop gros ») — la difficulté est d'inventer quelque chose sur quoi on opère de façon non triviale.

\newpage

%% ===== Page 143 =====

Interprétation $M_{0,1} = N_{1,1}$ comme $\pi_1(M_{1,1})$ où $M_{1,1}$ est le schéma modulaire sur $\mathbf{Q}$ des courbes elliptiques à symétrie près, on voit que la tour de Fermat (de caractéristique 0) ceci: sur tout schéma, ou bien $S$ (disons) et $\tilde{X}$ quasi-courbe elliptique sur $S$, considérons les variétés étales de degré $2$ de $\tilde{X}$ ), posons $\overset{\text{déf}}{=}$
$$ M_0(\tilde{X}/S) = \pi_*(\mathbf{Z}_2)/\mathbf{Z}_2 \leqno{(299)} $$
le $\mathbf{Z}$-module, en deux fois correspondant, et considérons $\overset{\text{déf}}{=}$
$$ M_1(\tilde{X}/S) = M_0(\tilde{X}/S) \otimes_{\mathbf{Z}} \hat{\mathbf{Z}}(1) \leqno{(300)} $$
$\overset{\text{déf}}{=} T_0(G_m)$

qui est un faisceau étale sur $S$, un système de coefficients localement constant de rang $\hat{\mathbf{Z}}^2$. Ceci posé, donc, $M_1(\tilde{X}/S)$ dépend fonctoriellement de $\tilde{X}/S$ pour $D$ et $\mathcal{G}$-isom. (On a un sentiment antérieur d'une cat. fibrée au-dessus de $S$...). Ceci posé;

\newpage

%% ===== Page 144 =====

La donnée de la "tour de Fermat" équivaut à la donnée, pour $H \tilde{X}/S$ comme dessus, d'un torseur
$$ F_\infty(\tilde{X}/S) \text{ torseur sur } S \text{ sous } M_1(\tilde{X}/S) \leqno{(301)} $$
\marginpar{provenance \dots (370) cf. p. 724}
compatible avec isomorphisme et degré de base. Le théorème précédent dit que un tel $F_\infty(\tilde{X}/S)$ existe, et est unique : isomorphisme unique pris, en exigeant de plus que l'homomorphisme
$$ \hat{\pi}_1 \to M^\wedge $$
associé à une telle donnée $F_\infty$ soit l'homomorphisme canonique\dots

Soit d'abord $X$ une courbe elliptique définie sur une base $S$ de caractéristique 0, et soit $S' \subset X$ le sous-schéma fermé, étale de degré 3 sur $S$, complément aux points d'ordre 2. Considérons la factorisation, l'étude des automorphismes, soit $\mathcal{J}^0$ la partie de degré 0, on a alors

\newpage

%% ===== Page 145 =====

720
une suite exacte
$$
1 \to \mathcal{J}^\circ \to \mathcal{J}^\circ \to X \to 1 \leqno{(302)}
$$
$$
\text{\phantom{1} \hspace{0.5cm} \hspace{0.5cm} \phantom{\mathcal{J}^\circ} \hspace{0.5cm} \mathbf{G}_m \otimes_{\mathbf{Z}} M_0}
$$
où $M_0 = M_0(X/S)$ est défini plus haut (299). On en déduit une suite exacte
$$
1 \to M_1 \to \mathcal{T}_\infty(\mathcal{J}^\circ) \to \mathcal{T}_\infty(X) \to 1 \leqno{(303)}
$$
$$
\text{\phantom{1} \hspace{0.5cm} \phantom{M_1} \hspace{0.5cm} \parallel}
$$
$$
\text{\phantom{1} \hspace{0.5cm} \phantom{M_1} \hspace{0.5cm} M_0 \otimes_{\mathbf{Z}} \mathcal{T}_\infty(\mathbf{G}_m)}
$$
$$
\mathcal{T}_\infty(\mathbf{G}_m) = \varprojlim_n \mathbf{G}_m \leqno{(304)}
$$
et le système local des $\pi_1$ des schémas à fibres commu\text{-}
nes, groupe $\varprojlim$ commu. sur $S$. Donc
le donné d'un torseur sous $M_1$, équi\text{-}
vaut à le donné d'un torseur sous
$\mathcal{T}_\infty(\mathcal{J}^\circ)$, muni d'une section du
torseur associé de groupe $\mathcal{T}_\infty(X)$.
Y a-t-il une description remarquable du
torseur de Fermat $\mathcal{F}_\infty(X/S)$, en termes
de la géométrie de la courbe elliptique
relative $X/S$ ?

\newpage

%% ===== Page 146 =====

14°) Relation aux jacobiennes généralisées.

Considérons la « jacobienne généralisée » de $U_{0,3} = \mathbf{P}^1 - \{0, 1, \infty\}$, c'est donc un schéma en groupes $\mathcal{J}_{0,3} = \mathcal{J}^*$, agissant vers $\mathbf{Z}$, d'où dans la suite exacte
$$
1 \to \mathcal{J}^0 \to \mathcal{J}^* \xrightarrow{\deg} \mathbf{Z} \to 1 \leqno{(305)}
$$
\marginpar{NB $S = \operatorname{Spec} \mathbf{Z}$}
\noindent et muni d'un homomorphisme canonique
$$
\varphi: U_{0,3} \to \mathcal{J}^1 \leqno{(306)}
$$
(sous-schéma de $\mathcal{J}^*$ image inverse de la section $1$ de $\mathbf{Z}$).

Les actions de $\mathfrak{S}_{0,3}$ sur $U_{0,3}$ définissent par fonctorialité des actions de $\mathfrak{S}_{0,3}$ sur $\mathcal{J}^*$, compatibles avec la structure de l'extension (305), et telles que $\varphi$ commute aux actions de $\mathfrak{S}_{0,3}$. Notons que $\mathcal{J}^1$ est un $\mathcal{J}^0$-torseur et...

\newpage

%% ===== Page 147 =====

\section*{71}

\noindent muni d'une structure de $\mathfrak{S}_3$-$\mathcal{J}^0$-torseur (en fait, un $\mathcal{J}^0$-torseur sur $S_0$ sur lequel $\mathfrak{S}_3$ opère de façon compatible avec son action sur $\mathcal{J}^0$), donc classifié par un élément de $H^1(S_0, \mathfrak{S}_3; \mathcal{J}^0)$. Or je présume qu'on a
$$
H^0(S_0, \mathfrak{S}_3, \mathcal{J}^0) = H^1(S_0, \mathfrak{S}_3, \mathcal{J}^0) = 0 \leqno{(307)}
$$
$$
S_0 = \operatorname{Spec} \mathbf{Z}, \quad \mathcal{J}^0 = \mathbb{G}_m \times \mathbb{G}_m \simeq \mathbb{G}_m^{\{0,1,\infty\}} / \text{diag} \simeq \\operatorname{Ker}\left(\mathbb{G}_m^{\{0,1,\infty\}} \xrightarrow{(x_0, x_1, x_\infty) \mapsto x_0 x_1 x_\infty} \mathbb{G}_m\right) \leqno{(308)}
$$

\noindent On a en effet
\marginpar{Calcul complet. Estimant que $H^0(\mathfrak{S}_3, \mathcal{J}^0) = 0$.}
$$
H^1(\mathfrak{S}_3, \mathcal{J}^0) = \{ (x_0, x_1, x_\infty) \in \mathbf{Z}^{*\{0,1,\infty\}} \mid x_0 x_1 x_\infty = 1 \}
$$
$$
= \{ (x, x, x) \mid x \in \mathbf{Z}^*, x^3 = 1 \}
$$
$$
= 1
$$

\noindent seule racine cubique de 1 dans $\mathbf{Z}$, d'où $= 1$. Ce qui est la première vérification de (307). Pour la seconde, on utilise la suite spectrale.

\newpage

%% ===== Page 148 =====

$$
H^1(S, \mathfrak{S}_3; \mathcal{J}) \implies H^1(\mathfrak{S}_3, H^0(S, \mathcal{J}))
$$
qui donne la suite exacte
$$
0 \to H^1(\mathfrak{S}_3, H^0(S, \mathcal{J})) \to H^1(S, \mathfrak{S}_3; \mathcal{J}) \to H^1(\mathfrak{S}_3, H^1(S, \mathcal{J})) \leqno{(309)}
$$
(où $\downarrow$ pointe vers $H^2(\mathfrak{S}_3, H^0(S, \mathcal{J}))$).

Or,
\marginpar{NB: on a besoin de la relation $Pic(S)=0$}
$$
H^0(S, \mathcal{J}) \isom H^0(S, \mathbb{G}_m) \isom Pic(S) = 0
$$
donc il vient
$$
H^1(S, \mathfrak{S}_3; \mathcal{J}) \isom H^1(\mathfrak{S}_3, H^0(S, \mathcal{J}))
$$
[unclear: best-guess: (\mathbf{Z}^*)^{\{0,1,\infty\}}/\mathbf{Z}^*]
$$
\isom \mathbf{Z}^* \otimes M
$$

Mais en fait, les calculs faits par ailleurs montrent que
$$
H^1(\mathfrak{S}_3, \mathbf{M} \otimes \mathbf{Z} I) = H^1(\mathfrak{S}_3, \mathbf{M} \otimes \mathbf{Z} I) = 0 \leqno{(309)}
$$
quel que soit le $\mathbf{Z}$-module $I$ (cf. ex. on a $I = H^1(S, \mathbb{G}_m) = Pic(S); I = J^0(S), S$ un schéma pur), d'où finalement,
$$
H^1(S, \mathfrak{S}_3; \mathcal{J}) \isom H^1(S, \mathfrak{S}_3; \mathcal{J}) = 0 \leqno{(310)}
$$
pour le schéma $S$. On trouve ainsi :

\newpage

%% ===== Page 149 =====

\textbf{Proposition.} Soit $S$. $J^0 = \mathbf{G}_m / \mathbf{G}_m \cong \mathbf{G}_m \otimes M$. Alors pour $(\mathcal{G}_S, J^0)$ sur un bon $S$ quelconque, il existe un unique discretisation des abéliens, i.e. une unique section, invariant par les opérations de $\mathcal{G}_S$.

\textbf{Corollaire 1.} Sur tout bon $S$, il existe un unique morphisme
$$
\psi: U_{0,3} S \to J_S^0 \leqno{(311)}
$$
qui soit compatible avec les opérations de $\mathcal{G}_S$, et qui (via l'homomorphisme $J_S^0 \to J_S^0$ que la factorisation, par la propriété universelle des jacobiennes généralisées)\footnote{NB Avec les notations 695, $U_{0,3} \defeq \{ (x_0, x_1, x_2) \dots \}$} induise l'homomorphisme $J_S^0 \to J_S^0$ identique.

\textbf{Corollaire 3.} Soit $Y$ un schéma de \leqno{(312)}

\newpage

%% ===== Page 150 =====

Revêtements de $S$, de degré rationnel 1 (i.e. un faisceau de $\mathbf{P}^1_S$), munis d'une section $S'$ sur $S$, et soit $U = X \setminus S'$, $\mathcal{J}_{U/S}$ le jacobien généralisé, alors on a :
$$ 1 \to \mathcal{J}^0_{U/S} \to \mathcal{J}_{U/S} \to \mathbf{Z} \to 1 \leqno{(313)} $$
avec $\mathcal{J}^0_{U/S} \simeq \mathbf{G}_m \otimes_{\mathbf{Z}} M_1(S'/S)$. On a $M_1(S'/S) \simeq \operatorname{Div}(Z_S) / \bar{Z}_S$ (d.g.), et le $S$-schéma (loc. isom. $Z$), décrit en le tordant (via action de $\mathfrak{S}$) par le fibré principal associé : $S'$.

Il existe alors un unique morphisme
$$ \varphi: U \to \mathcal{J}_{U/S} \leqno{(314)} $$
commutant à l'action de $\\operatorname{Aut}(S'/S) \simeq \\operatorname{Aut}(S, S')$, et induisant l'isomorphisme $\mathcal{J}_{U/S} \isom \mathcal{J}_{U/S}$ qui soit l'identité.

\newpage

%% ===== Page 151 =====

723.

Explicitons l'isomorphisme
$$
U_{0,3} = \mathbb{P}^1 \setminus \{0, 1, \infty\} \isom U_1 \subset \mathbb{P}^2 \leqno{(315)}
$$
où $U_1 = \{ (x_0, x_1, x_\infty) \in \mathbb{P}^2 \mid x_0 + x_1 + x_\infty = 0 \}$.

C'est l'isomorphisme de $U_{0,3}$ avec $U_1$ que l'on connaît :
$$
\begin{cases}
\bar{\varphi}: U_{0,3} \to U_1 \subset \mathbb{P}^2 \\
\bar{\varphi}(0) = (0, 1, -1) \\
\bar{\varphi}(1) = (-1, 0, 1) \\
\bar{\varphi}(\infty) = (1, -1, 0)
\end{cases} \leqno{(314)}
$$

Il suffit de prendre :
$$
\psi(z) = (-z, z-1, 1) = (1, 1/z - 1, -1/z) \leqno{(317)}
$$

Revenons maintenant à la donnée d'une courbe elliptique relative sur $S$, modèle symétrique. C'est-à-dire un morphisme $S \to M_{1,1}$.

\newpage

%% ===== Page 152 =====

On a deux inclusions sur $S$, de sorte que le donné de $X$ corresponde par le type de lignes $X \to (X/ \dots)$ : la donnée d'un $(Y, S')$ comme dessus, avec une section de $U = Y \setminus S'$.

En fait, la donnée de $(Y, S')$ revient à un $\mathcal{G}$-torseur et la donnée de $S' \to X$ et réciproquement.

$\mathcal{T}$ sur $S$ et $Y$ un $\mathcal{G}$-torseur donné comme le schéma $\mathcal{T}_{\mathcal{G}, \mathbb{P}^1 \setminus \{0,1,\infty\}_S} \dots$
$$
U = Y \setminus S' \simeq (\mathcal{T}_{\mathcal{G}} \mathbb{P}^1) \setminus \mathcal{T}_{\mathcal{G}} \{0,1,\infty\}_S
$$
Plongement conséquemment des $\mathcal{T}_{\mathcal{G}, \{0,1,\infty\}_S} \simeq \mathcal{M}_{U/S}$. Le plongement $U \hookrightarrow \mathcal{M}_{U/S}$ puis s'interpose comme un morphisme canonique
$$
X \to \mathcal{M}_{U/S} = \mathcal{M}_{X/S} \leqno{(318)}
$$
qui doit être ainsi : construction canonique

\newpage

%% ===== Page 153 =====

ment, \`a expliciter en termes de ``fonctions elliptiques'' standard. Sans attendre une telle explicitation, et revenant au plongement correspondant
$$
U \hookrightarrow \mathfrak{T}_{1,1/S} \leqno(318)
$$
et \`a la section marquée
$$
e \in U(S) = \mathcal{H}(S, \mathfrak{T}_{1,1/S}) \leqno(319)
$$
--- par utilisation du tore d'isogénies, on utilise le plongement
$$
\mathcal{H}(S, \mathfrak{T}_{1,1/S}) \to \mathcal{H} \leqno(320)
$$
$$
1 \to \mathfrak{T}_{1,1/S} \to \mathfrak{T}_{1,1/S} \xrightarrow{\text{n.i.}_{\mathfrak{T}_{1,1/S}}} \mathfrak{T}_{1,1/S} \to 1 \leqno(321)
$$
(pour $n$ variable), i.e.\ le diagramme universel déduit par passage \`a la limite :
$$
1 \to \mathfrak{T}_{1,1/S} \to \widetilde{\mathfrak{T}}_{1,1/S} \xrightarrow{p} \mathfrak{T}_{1,1/S} \to 1 \leqno(322)
$$
$$
\text{T}_{\infty}(\mathbf{G}_m) \otimes_{\mathbf{Z}} \mathcal{M}(S/S) \quad \varprojlim (\mathfrak{T}_{1,1/S})
$$
pour déduire, par image inverse de $e$, un torseur sous $\text{T}_{\infty}(\mathfrak{T}_{1,1/S}) \isom \text{T}_{\infty}(\mathbf{G}_m) \otimes_{\mathbf{Z}} \mathcal{M}(S/S)$.
Ceci posé, --- avec (cf.\ 051 p.\ 719)
$$
\mathcal{F}_{\infty}(\widetilde{X} | S) \isom p^{-1}(e), \text{ le } \mathfrak{T}_{1,1/S}\text{-torseur déduit de la section } e \text{ de } \mathfrak{T}_{1,1/S} \text{ sur } S \leqno(323)
$$

\newpage

%% ===== Page 154 =====

Remarque. Le choix d'isomorphisme du $T_\infty(\mathcal{J}_{U/S})$-torseur obtenu est un élément de $H^0(S, T_\infty(\mathcal{J}_{U/S}))$, image de l'élément $e$ de $H^0(S, \mathcal{J}_{U/S})$, par le morphisme canonique
$$ H^0(S, \mathcal{J}_{U/S}) \to H^0(S, T_\infty(\mathcal{J}_{U/S})), \leqno{(324)} $$
dont le noyau est formé des sous-groupes de $H^0(S, \mathcal{J}_{U/S})$ des éléments "infiniment divisibles". Dans le cas où $S$ est de type fini sur $\mathbf{Z}$ — p.ex. de type fini sur $\mathbf{Q}$ — (et ce sont les éléments "finis" qui en noyau s'admettent) n'admettent qu'un noyau : zéro. Ce qui signifie donc que la section $e$ (i.e. la "fonction elliptique visée") est donc univoquement déterminée, quand on connaît...

\newpage

%% ===== Page 155 =====

725

\begin{enumerate}
\item[1°)] Les revêtements étales $S'$ de $S$ associés aux points d'ordre 2 de $\widetilde{X}/S$ — d'où $\pi_1(\widetilde{X}/S) \simeq \pi_1(\widetilde{X}, S)$ — et
\item[2°)] Le $T_\infty(\widetilde{X}/S)$ — torseur $F_\infty(\widetilde{X}/S)$.
\end{enumerate}

\newpage

%% ===== Page 156 =====

15°) Interprétation de $\Gamma_{\mathbf{Q}} \to N_{\mathfrak{S}}$.

Au n° 11°, nous avons introduit l'homomorphisme
$$
\begin{cases} \mathcal{M} = N_{1,1} \to \mathcal{M}_{g, \nu} \\ \mathcal{M}^! \longrightarrow \mathcal{M}_{g, \sigma}^! \end{cases} \text{cf. } (217) \text{ p. 69}
$$
que nous avons "décomposé" bien plus tôt (en plus d'un homomorphisme trivial $\mathcal{M} \to N_{\mathfrak{P}} \simeq \hat{\mathbf{Z}}^{\times}$) en deux hom.
$$
\begin{tikzcd}
\mathcal{M} \arrow[d, hook] \arrow[r] & \mathcal{G}_{g, \sigma} \simeq \hat{\mathbf{Z}}^{\times} \cdot \boldsymbol{\Gamma} \\
\mathcal{M}^! \arrow[r] & \mathcal{G}_{g, \nu} \simeq \hat{\mathbf{Z}}^{\times} \cdot \mathcal{M}^{\wedge}
\end{tikzcd} \leqno{(325)} \text{ (cf. } (218))
$$

et
$$
\Gamma_{\mathbf{Q}}^{\text{it}} \longrightarrow N_{\mathfrak{S}}^{\mathcal{G}_{g, \nu}} \simeq \hat{\mathbf{Z}}^{\times} \cdot \hat{\mathbf{Z}} \leqno{(326)} \text{ (cf. } (219))
$$

Le premier hom. (325) a été caractérisé via des termes \textit{arithmétiques} géométriques (pour ce qui est de sa restriction à $N_{1,1} \subset N_{1,1}^{\text{it}}$, lui-même, et un réduction à ... via \marginpar{...})
$$
\mathcal{M}_{g, \nu} \defeq \mathcal{G}_{g, \nu} \simeq \mathcal{G}_{g} / \mathcal{M}_a \leqno{(327)}
$$
dans la section 13°), via la tour des

\newpage

%% ===== Page 157 =====

725
courbes de Fermat. Du point de vue actuel, il s'agit d'expliciter la géométrie algébrique en type fini sur $\mathbf{Q}$ en termes d'automorphismes de groupes profinis, pour dire en somme, inversement, que nous sommes arrivés à "décrire" la tour des courbes de Fermat, via l'homomorphisme (324) construit ici "à la force du poignet".

Il resterait maintenant à trouver une explication arithmético-géométrique de l'homomorphisme $\Gamma_{\mathbf{Q}} \to \hat{\mathbf{Z}}^* \cdot \hat{\mathbf{Z}}$, induit par (325). Mais cet homomorphisme utilise le plongement $\mathcal{N}^{\mathcal{G}_k} \hookrightarrow \mathcal{M}_{0, \sigma} \simeq \hat{\mathbf{Z}}^* \cdot \mathcal{M}$, l'homomorphisme (323) apparaissant comme le composé $\Gamma_{\mathbf{Q}} \xrightarrow{\text{ét}} \mathcal{M}^{\text{ét}} \xrightarrow{(325)} \mathcal{G}_{0, \sigma}^{\mathcal{G}_k}$, de façon plus précise ... un diagramme :

\newpage

%% ===== Page 158 =====

(328)
\begin{tikzcd}
{[\Gamma_\mathbb{Q} \hookrightarrow]} \Gamma_\mathbb{Q}^{et} \ar[r] & \mathcal{N}_\sigma^\S \simeq \hat{\mathbb{Z}}^* \cdot \hat{\mathbb{Z}} \\
{[\underset{\Pi_1(U_{0,3\mathbb{Q}})}{M_{0,3}^!} \hookrightarrow] M_{0,3}^{! et}} \ar[u, hook, "K_{1/2}"] \ar[r, "\text{''Fermat''}"', "(32\bar{5})"] & G_{p,\sigma}^\S \simeq \hat{\mathbb{Z}}^* \cdot \hat{M} \ar[u, hook]
\end{tikzcd}

Cela signifie donc que l'on trouve
l'hom. induit par (326)

(329) $$\Gamma_{\mathbb{Q}} \longrightarrow \mathcal{N}_\sigma^\S \simeq \hat{\mathbb{Z}}^* \cdot \hat{\mathbb{Z}}$$

en regardant l'image de l'hom.

(330) $$\Gamma_{\mathbb{Q}} \longrightarrow \hat{\mathbb{Z}}^* \cdot \hat{M} = Aff(\hat{M}) \ ,$$

dans $\mathcal{N}_\sigma^\S \subset \hat{\mathbb{Z}}^* \cdot \hat{M}$ , déduit de la
tour des courbes de Fermat en
regardant la fibre de la dite tour
au point $1/2$ , ce qui donne un
torseur sous $T_a(\hat{\mathbb{Z}}_{(m)}) \otimes_{\hat{\mathbb{Z}}} \hat{M}$ , s'explici-
tant galoisiennement justement par
l'hom. (326). Notons que le $K_{1/2}$
de $U_{0,3\mathbb{Q}}$ corresp. par l'iso. $\varphi$ (317)
de $U_{0,3\mathbb{Q}}$ avec $U_{1,\mathbb{Q}}$, au point [unclear: $(\xi_{1/2}, 1)$]

\newpage

%% ===== Page 159 =====

727

$$ U_{0,3} \defeq \left\{ (x_0, x_1, x_\infty) \in \mathbb{G}_m^3 / \mathbb{G}_m \mid x_0 + x_1 + x_\infty = 0 \right\} \leqno{(331)} $$

La fibre de $U_{0,3}$ en le dit point s'identifie aux systèmes de relations des équations ci-dessous :

$$
\begin{aligned}
U_{0,3}(s_{1/2}) &\defeq \{ (x_0, x_1, x_\infty) \in \mathbb{G}_m^3 / \mathbb{G}_m \mid x_0 = -1, x_1 = -1, x_\infty = -1 \} \\
&\simeq \{ (x_0, x_1, x_\infty) \in \mathbb{G}_m^3 \mid x_0 = -1, x_1 = -1, x_\infty = -1 \} \\
&\simeq \{ (x_0, x_\infty) \in \mathbb{G}_m^2 \mid x_0 = -1/2, x_\infty = -1/2 \} \\
&\simeq T_0(-1/2) \times T_\infty(-1/2)
\end{aligned} \leqno{(332)}
$$

--- pour tout $\zeta \in \mathbb{G}_m(S)$ \marginpar{Schéma...}

$$ T_u(\zeta) = \dots \leqno{(333)} $$

Posons : la limite ---

$$ U_\infty(s_{1/2}) \simeq T_\infty(1/2) \times T_\infty(-1/2) \leqno{(334)} $$

$$ T_\infty(\zeta) = \varprojlim T_\infty(\zeta) \quad (\text{limite de } T_\infty(\mathbb{G}_m)\text{-tores}) \leqno{(335)} $$

\newpage

%% ===== Page 160 =====

Ainsi, l'hom. (329) doit pouvoir s'interpréter en termes de l'hom.
$$
\Gamma_{\mathbf{Q}} \to \\operatorname{Gal}(\overline{\mathbf{Q}}[\sqrt{-1}] / \mathbf{Q}) \leqno{(335)}
$$
à désigner l'extension cyclotomique universelle de $\mathbf{Q}$
$$
\widetilde{\mathbf{Q}} \simeq \mathbf{Q}[\sqrt{1}] \leqno{(336)}
$$
Pour apprécier la "taille" de l'image de $\Gamma_{\mathbf{Q}}$ dans $\hat{\mathbf{Z}}^* \cdot \hat{\mathbf{Z}} \simeq N_{\hat{\mathbf{Z}}}$, il y a lieu de faire une digression de nature générale sur les groupes de Galois des extensions de la forme $\overline{\mathbf{Q}}[\sqrt[n_1]{\alpha_1}, \dots, \sqrt[n_N]{\alpha_N}]$.

Pour ceci sur un bien bref cas, pour simplifier (on a ici $S = \Spec \mathbf{Q}$) donnons-nous un groupe abélien $A$ (tel que nous prendrons $A = \mathbf{Z}$), et

\newpage

%% ===== Page 161 =====

un homomorphisme
$$
A \to \mathcal{G}_m(S) = H^1(S, \mathcal{G}_m) \leqno (337)
$$
(dans le cas présent, il est donné par $1 \to \dots$)
- ce qui revient, dans le cas
$$
A_S \to \mathcal{G}_m S \leqno (338)
$$
Considérons le type des isogénies
$$
1 \to \mu_n \to \mathcal{G}_m \xrightarrow{n} \mathcal{G}_m \to 1 \leqno (339)
$$
d'où, par $1 \to \mathcal{G}$ iso via (338), une
tour d'extensions des schémas en
groupes étales sur $S$
$$
1 \to \mu_n \to E_n(S) \to A_S \to 1 \leqno (339)
$$
et par passage à la limite
$$
1 \to \widehat{\mathcal{T}}(\mathcal{G}_m) \to E_\infty(S) \to A_S \to 1 \leqno (340)
$$
Supposons $S$ connexe, et prenons un
point géométrique $\bar{s}$ sur $S$, on en déduit
un homomorphisme canonique
$$
\pi_1(S, \bar{s}) \to \\operatorname{Aut}_{gr}(E_\infty(S)_{\bar{s}}) \leqno (341)
$$
\marginpar{pro-finis}

\newpage

%% ===== Page 162 =====

Le fibre $E_\infty(\mathcal{G})_{\bar{s}}$ est muni d'une structure d'extension abélienne
$$
1 \to T_\infty(\bar{k}^*) \to E_\infty(\mathcal{G})_{\bar{s}} \to A \to 1 \leqno{(342)}
$$
L'homomorphisme (341) est compatible avec cette structure d'extension, qui fait opérer $\pi_1(S, \bar{s})$ trivialement sur $A$ ($\mathbf{Z}$-module libre) et sur $T_\infty(\bar{k}^*)$ par le groupe $\pi_1$ par multiplication par les caractères cyclotomiques
$$
\pi_1(S, \bar{s}) \to \pi_1(\mathbf{Q}, \bar{s}) \xrightarrow{\chi} \hat{\mathbf{Z}}^* . \leqno{(343)}
$$
Soit $\mathcal{G}$ le groupe
$$
\mathcal{G} = \\operatorname{Aut}(E_\infty(\mathcal{G})_{\bar{s}}) \leqno{(344)}
$$
qui respecte la structure d'extension et induit l'identité sur $A$ et sur $T_\infty(\bar{k}^*)$, homomorphisme qui est défini par $\hat{\mathbf{Z}}^*$ bien déterminé.
On a un homomorphisme canonique $\mathcal{G} \to \hat{\mathbf{Z}}^*$, d'où un diagramme
$$
1 \to \operatorname{Hom}(A, T_\infty(\bar{k}^*)) \to \mathcal{G} \to \hat{\mathbf{Z}}^* \leqno{(345)}
$$

\newpage

%% ===== Page 163 =====

74

où l'image de $G \to \hat{\mathbb{Z}}^*$ est formée des $\mu \in \hat{\mathbb{Z}}^*$ tels que l'homothétie $\mu_T$ ne change pas la classe de l'ext. (342) de $A$ par $T$. L'hom (341) peut lors s'interpréter comme s'insérant dans un diagramme

(345)
\begin{tikzcd}
1 \ar[r] & \operatorname{Hom}_\mathbb{Z}(A, T_\infty(\bar{k}^*)) \ar[r] & G \ar[r] & \hat{\mathbb{Z}}^* \\
& & \Pi_1 \ar[u] \ar[ur, "\text{"caractère cyclotomique"}"']
\end{tikzcd}

[MARGIN: NB.] Si $S$ est sch. de type fini sur $\mathbb{Q}$, alors l'image de $\Pi_1 \to \hat{\mathbb{Z}}^*$ est un sous-groupe ouvert de $\hat{\mathbb{Z}}^*$, donc il en est a fortiori ainsi de l'image de $G \to \hat{\mathbb{Z}}^*$ - qui sera surjectif si le car. cyclotomique sur $\Pi_1$ est surjectif, i.e. si les seules racines de 1 sur $S$ sont $\pm 1$ (p. ex. si $S = Spec \mathbb{Q} \dots$).

Le diagramme (345) a
une propriété de fonctorialité évidente
pour $A$ variable (disons), i.e. pour un

\newpage

%% ===== Page 164 =====

hom. de groupes abéliens
(347) $A' \xrightarrow{f} A$ ,

l'extension $E_{\infty}(A')_S$ s'identifie
à l'image inverse de l'ext. $E_{\infty}(A)_S$
par (347), d'où un hom.
$G_A \to G_{A'}$ (injectif si $f$ est surjectif)
s'insérant dans un diagramme
commutatif

(349)
\begin{tikzcd}
1 \ar[r] & \operatorname{Hom}_{\mathbb{Z}}(A', T) \ar[r] & G_{A'} \ar[r] & \hat{\mathbb{Z}}^* \\
1 \ar[r] & \operatorname{Hom}_{\mathbb{Z}}(A, T) \ar[r] \ar[u] & G_A \ar[r] \ar[u] & \hat{\mathbb{Z}}^* \ar[u, "1"'] \\
& & \Pi_1 \ar[u, "\gamma_A"'] \ar[uu, bend left=45, "\gamma_{A'}"] \ar[ur, "\text{car. cyclot.}"']
\end{tikzcd}

où la flèche sur les $\operatorname{Hom}_{\mathbb{Z}}$ est trans-
posée de (347). Cela montre que
pour étudier la structure de l'image de $\gamma_A$ (345),
il est commode, par fonctorialité, de se

\newpage

%% ===== Page 165 =====

730

étudier le cas où $j$ est injectif
(350) $$j : A \hookrightarrow \mathbb{G}_m(S) \ .$$

On peut d'ailleurs "diviser" encore un peu le diagramme (345), en introduisant

(351) $$\Pi_1^+ = \Pi_1^+(S, \bar{s}) \overset{df}{=} \operatorname{Ker}\Big(\Pi_1(S, \bar{s}) \overset{cyclot.}{\longrightarrow} \hat{\mathbb{Z}}^*\Big)$$

et l'hom. de suites exactes

(352)
\begin{tikzcd}
1 \ar[r] & \Pi_1^+ \ar[r] \ar[d, "\gamma_j^+"'] & \Pi_1 \ar[r, "\chi_{\Pi}"] \ar[d, "\gamma_j"'] & \hat{\mathbb{Z}}^* \ar[d] \\
1 \ar[r] & \operatorname{Hom}_{\hat{\mathbb{Z}}}(A, T) \ar[r] & G_j \ar[r] & \hat{\mathbb{Z}}^*
\end{tikzcd}
.

On peut considérer que l'étude de la "taille" de $\operatorname{Im}(\gamma_j : \Pi_1 \to G_j)$ est essentiellement ramenée à celle de la taille de $Im \chi_{\Pi}$ et de celle de $Im \gamma_j^+$; de façon précise, si on pose

(353) $$Im \gamma_j = \tilde{\Pi}_1 \ , \ Im \gamma_j^+ = \tilde{\Pi}_1^+ \ , \ Im \chi_{\Pi} = \Gamma$$

on a une suite exacte

(354) $$1 \to \tilde{\Pi}_1^+ \to \tilde{\Pi}_1 \to \Gamma \to 1 \ .$$

Je voudrais essayer de donner des

\newpage

%% ===== Page 166 =====

indications sur la taille de $\tilde{\Pi}_1^+$\footnote{$\simeq \operatorname{Hom}(A, \mathbf{Z})$} dans le cas où :
$$
\left.
\begin{aligned}
1) & \quad A \text{ est un } \mathbf{Z}\text{-module de type fini} \\
2) & \quad S \text{ est de type fini sur } \mathbf{Q} \text{ et normal} \\
3) & \quad j: A \to \operatorname{Gm}(S) \text{ est injectif}
\end{aligned}
\right\} \leqno{(355)}
$$

L'hypothèse 2) implique déjà, par $\Gamma_S \subset \hat{\mathbf{Z}}^*$ un sous-groupe ouvert. L'hypothèse que $S$ est normal implique d'ailleurs que les groupes étendus $\tilde{\Pi}_1^+$ et $\Gamma_S$ ont des points génériques (on suppose pour simplifier que $S$ est connexe, le dit point générique), dans des cas $S = \operatorname{Spec} K$, où $K$ est une extension de type fini de $\mathbf{Q}$. Et on peut [voir] que $S$ est affine, le lien vers $\mathbf{Q} \dots$

Notons enfin, posons
$$
\breve{A} \defeq \operatorname{Hom}_{\mathbf{Z}}(A, \mathbf{Z}) \leqno{(356)}
$$

\newpage

%% ===== Page 167 =====

que $\check{A}$ est un $\mathbf{Z}$-module libre de type fini (de rang $n$ : voir 354), et, on a un isomorphisme canonique
$$
\operatorname{Hom}_{\mathbf{Z}}(A, \hat{\mathbf{Z}}) \cong \check{A} \otimes_{\mathbf{Z}} \hat{\mathbf{Z}} \quad (\cong \hat{\mathbf{Z}}^n \text{ non canonique}) \leqno{(357)}
$$

J'avais envie de prouver que l'image de $\pi_1^+$ dans $\pi_1$ des $\operatorname{Hom}_{\mathbf{Z}}(A, T) \cong \check{A} \otimes_{\mathbf{Z}} T$ est ouverte. J'ai un doute, dans le cas $S$ normal de type fini sur $\mathbf{Z}$ et $A = H^0(S, \mathbf{G}_m)$, pour $S = S_{\mathbf{Q}}$, si l'homomorphisme $\pi_1 = \pi_1(S) \to \operatorname{Hom}(A, \hat{\mathbf{Z}}) \cong \check{A}(\hat{\mathbf{Z}})$ n'est pas nécessairement surjectif ?? (Plus généralement, il en est de même pour $\mathbf{G}_m(S) / \operatorname{Im} A$ et, disons, en ce qui concerne $A$ n'est rien [unclear: de clair], d'après l'hyp. (355), qu'à agrandir $A$ en $A'$ tel que $A'/A$ soit fini ...).

\newpage

%% ===== Page 168 =====

Je vais revenir sur le cas $S = \operatorname{Spec} \mathbf{Q}$, alors
$$
\mathbf{G}_m(S) = \mathbf{G}_m(\mathbf{Q}) = \mathbf{Q}^* \simeq \pm 1 \times \mathbf{Z}^{(P)}
$$
Pour $p$ l'un des nombres premiers,
$$
\pi_1 = \operatorname{Gal}(\overline{\mathbf{Q}}/\mathbf{Q}) = \boldsymbol{\Gamma}_{\mathbf{Q}}.
$$
La question est donc d'étudier les homomorphismes
$$
\zeta_p: \boldsymbol{\Gamma}_{\mathbf{Q}}^+ \to \hat{\mathbf{Z}}(1) \leqno{(358)}
$$
associés à $p$, et de répondre notamment aux questions :
\begin{enumerate}
    \item[a)] Les $\zeta_p$ sont-ils surjectifs, et quelle est l'image, est-elle un sous-groupe ouvert ?
    \item[b)] Les $\zeta_p$ sont-ils linéairement indépendants sur $\hat{\mathbf{Z}}$ ?
\end{enumerate}
Plus précisément, si $p_1, \dots, p_n$ sont des nombres premiers distincts, l'homomorphisme
$$
\boldsymbol{\Gamma}_{\mathbf{Q}}^+ \to \hat{\mathbf{Z}}(1)^n
$$
$( \zeta_{p_1}, \dots, \zeta_{p_n} )$ est-il surjectif, et quelle est l'image de $\boldsymbol{\Gamma}_{\mathbf{Q}}^+$ — est-elle un sous-groupe ouvert ?

Je vais commencer par regarder un peu la question a) (le cas où [unclear: la condition de Fermat] était celui en $p=2$). Pour ceci, considérons
$$
\zeta_p: \boldsymbol{\Gamma}_{\mathbf{Q}}^+ \to \hat{\mathbf{Z}}(1) \quad (p \in P) \leqno{(359)}
$$

\newpage

%% ===== Page 169 =====

732 (au moins!)
l'aimerais pouvoir démontrer que ces homomorphismes surjectifs, ou du moins surjectifs, sont pour $\ell$ fini de $\ell$, et [unclear: surtout] pour $\ell$ quelconque. Étudier $\chi_{p,\ell}$, c'est se restreindre à la tour des isogénies
$$
1 \to \mu_{\ell^n} \to \mathbb{G}_m \xrightarrow{\ell^n} \mathbb{G}_m \to 1 \leqno{(359)}
$$
pour $n$ puissance de $\ell$, et étudier l'image inverse de $\Gamma_{\mathbf{Q}}$ par ces isogénies. Il est naturel alors de regarder le groupe de décomposition
$$
D_\ell \defeq \Pi_1(\mathbf{Q}_\ell, \overline{\mathbf{Q}}_\ell) \to \Gamma_{\mathbf{Q}} \leqno{(360)}
$$
et la caractère cyclotomique $\chi_\ell$ :
$$
D_\ell \to \mathbf{Z}_\ell^* \leqno{(361)}
$$
ou, mieux, le sous-groupe d'inertie
$$
I_\ell \subset D_\ell \leqno{(363)}
$$
et la restriction de $\chi_\ell$ sur $I_\ell$.
On a ainsi
$$
I_\ell^+ \defeq \\operatorname{Ker}(\chi_\ell : I_\ell \to \mathbf{Z}_\ell^*) \subset \Gamma_{\mathbf{Q}}^+ \leqno{(364)}
$$
car les $\chi_\ell$ ($\ell \neq \ell$) sont triviaux sur $I_\ell$ (ils ne sont pas triviaux cependant sur $D_\ell$!).

\newpage

%% ===== Page 170 =====

On va donc regarder la restriction
$$
\chi_{\ell | I_{\ell}^+} : I_{\ell}^+ \to \hat{\mathbf{Z}}_{\ell}(1) \leqno{(365)}
$$
et on peut se demander déjà si cet homomorphisme est surjectif, ou si du moins son image est assez ?

\newpage

%% ===== Page 171 =====

\section*{16°) Réflexions sur le théorème de Fermat.}

Le ``théorème de Fermat'' peut s'énoncer en disant que pour $n \ge 3$, la courbe de Fermat affine sur $\mathbf{Q}$
$$
U_n \defeq \{ (x_0, x_1, x_\infty) \in \mathbf{G}_m^2 \mid x_0, x_1, x_\infty \text{ inv.}, x_0^n + x_1^n + x_\infty^n = 0 \} \leqno (366)
$$
$$
\simeq \{ (x_0, x_1, x_\infty) \in \mathbf{G}_m^{\{0, 1, \infty\}} / \mathbf{G}_m \mid x_0^n + x_1^n + x_\infty^n = 0 \}
$$
n'a pas de points rationnels sur $\mathbf{Q}$. Nous avions introduit cette ``tour'' des courbes de Fermat dans § 110 (pg 696 ff), et dès 140) on écrit la description de $U_n$ comme une image isom. de
$$
U_n \subset \mathbf{G}_m \otimes_{\mathbf{Z}} M \simeq \mathbf{G}_m^{\{0, 1, \infty\}} / \mathbf{G}_m \leqno (367)
$$
par l'isogénie de Kummer $x \mapsto x^n$
$$
1 \to \mu_n \to T \xrightarrow{n} T \to 1 \leqno (368)
$$
Calculons le groupe de $\widehat{U}_n$, pour un revêtement de degré $n^2$ de $U_n \simeq \mathbb{P}_{0,3}^1$, ramifié en les pts $0, 1, \infty$ et...

\newpage

%% ===== Page 172 =====

$$ (2 - 2g_n) - 3n = n^2(2 - 3) $$
\begin{center}
\small \text{(EP de $\widehat{U}_n$)} \hspace{1cm} \text{(EP de $U_n$)} \hspace{1cm} \text{(EP de $U_1$)}
\end{center}
$$ 2(g_n - 1) = n^2 - 3n $$
$$ g_n = 1 + \frac{1}{2}n(n-3) \leqno{(369)} $$
Donc,
$$ g_1 = g_2 = 0, \quad g_3 = 1, \quad g_n > 1 \text{ si } n > 3 \leqno{(370)} $$

En tout état de cause, tous les $U_n$ sont anabéliens, et aussi les $\widehat{U}_n$ pour $n \ge 4$.

Si le yoga développé dans la première partie de ces notes est valable, les existences de points v.b. / $\mathbf{Q}$ d'une courbe $U_n$ s'expriment par les invariants de l'extension en groupes profinis
$$ 1 \to \pi_1(\widehat{U}_n) \to \pi_1(U_n) \to \Gamma_{\mathbf{Q}} \to 1 \leqno{(371)} $$
qui soient distincts des seuls points : l'infini de $U_n$ (localisés en... : l'infini de $U_n$...), ou encore, les seuls tels que $(\Pi^+ = \operatorname{Centr} \dots \hat{\mathbf{Z}}^*)$ dans $\pi_1(U_n)$ soit réduit à $\{1\}$. Donc - si le...

\newpage

%% ===== Page 173 =====

de Faltings et bel et bien valable (pour $U_n$, disons), il devrait se vérifier par l'inexistence d'un tel scindage. D'ailleurs, si nous pouvions prouver l'inexistence de scindages de (371) autres que ceux "provenant des points à l'infini", le théorème de Faltings serait prouvé pour $U_n$, même sans avoir établi le "yoga anabélien" (et une démonstration semble bien hors d'atteinte!). Or il n'est nullement exclu qu'une analyse poussée de $\mathcal{M}_{0,3}$ et $\mathcal{M}_{1,1}$ dans la direction actuelle — impliquant aussi des "approximations" de plus en plus fines de $\pi_1(U_n)$ et de la structure d'extension (371) — permette finalement de prouver l'inexistence de scindages "non triviaux" de (371), pour $n \ge 3$. Chaque nouvel affinement des "équations de Lie" de $\mathcal{M}_{0,3}$ devrait amener un réexamen de la situation (371) et de la question de ses scindages, dans l'optique d'une démonstration éventuelle du théorème de Faltings.

\newpage

%% ===== Page 174 =====

Il s'impose de partir d'approximations de
Lie au moins aussi fines que celle du
présent $\S$, qui apparait comme l'approximation
"minima" dans laquelle on puisse exprimer
"algébriquement" la tour des courbes de
Fermat, ~~via~~ comme on l'a vu dans 13°),
à partir de $M_{0,3} \subset \hat{M}_{0,3}$.

Tenant compte des opérations de $\mathfrak{S}_3$
sur $U_n$, la suite ex.te (371) se prolonge
d'ailleurs dans une suite ex.te plus
riche

$$(372) \quad 1 \longrightarrow \pi_1(U_{n,\overline{\mathbb{Q}}}) \longrightarrow \underset{\overset{\wedge}{\pi_1(U_1,\mathfrak{S}_3) \simeq M_{0,3}}}{\pi_1(U_n,\mathfrak{S}_3)} \longrightarrow \Gamma_{\mathbb{Q}} \times \mathfrak{S}_3 \longrightarrow 1$$

et les calculs faits dans le présent $\S$ semblent
suggérer que cette structure d'extension plus
riche jouera un rôle important.

Du point de vue de l'approximation
de Lie de $M_{0,3}$ développée dans le présent
$\S$, il y a lieu de remplacer la structure
d'extension (372) par une structure quotient
moins fine, en divisant par un sous-groupe
invariant de $\pi_1(U_n, \mathfrak{S}_3)$ contenu dans

\newpage

%% ===== Page 175 =====

$73r$

$\Pi_1(\mathcal{U}_n, \overline{\mathbb{Q}})$, savoir l'intersection dans $M_{0,3}$ de ce
groupe avec $[\overline{\pi_0, \hat{j}}, \overline{\pi_0, \hat{j}}]$. On trouve un
diagramme commutatif

(373)
\begin{tikzcd}
1 \ar[r] & M^\wedge[n] \ar[r] \ar[d, equal] & \widetilde{\Pi_1(\mathcal{U}_n, \mathfrak{S}_3)} \ar[r] \ar[d, "p_n"'] & \Gamma_{\mathbb{Q}} \times \mathfrak{S}_3 \ar[r] \ar[d, "\chi \times id"] & 1 \\
1 \ar[r] & M^\wedge[n] \ar[r] \ar[d] & \widetilde{G}_{\rho, \sigma}[n] \ar[r] \ar[d] & \hat{\mathbb{Z}}^* \times \mathfrak{S}_3 \ar[r] \ar[d, equal] & 1 \\
1 \ar[r] & M^\wedge \ar[r] & \widetilde{G}_{\rho, \sigma} \ar[r] & \hat{\mathbb{Z}}^* \times \mathfrak{S}_3 \ar[r] & 1
\end{tikzcd}

où l'image de $M^\wedge[n] \to M^\wedge$ est $n M^\wedge$.

Voici une façon de construire "algébriquement"
ces diagrammes (373) (pour $n$ variable)
à partir du diagramme

(374)
\begin{tikzcd}
1 \ar[r] & M^\wedge \ar[r] \ar[d, equal] & \overset{\mathcal{E}_1}{\widetilde{\Pi_1(\mathcal{U}_1, \mathfrak{S}_3)}} \ar[r] \ar[d] & \Gamma_{\mathbb{Q}} \times \mathfrak{S}_3 \ar[r] \ar[d, "\chi \times id"] & 1 \\
& & (M_{0,3}/[\overline{\pi_0, \hat{j}}, \overline{\pi_0, \hat{j}}]) \ar[d] & & \\
1 \ar[r] & M^\wedge \ar[r] & \widetilde{G}_{\rho, \sigma} \ar[r] \ar[d, "\simeq"'] & \hat{\mathbb{Z}}^* \times \mathfrak{S}_3 \ar[r] & 1 \\
& & (\hat{\mathbb{Z}}^* \times \mathfrak{S}_3).M^\wedge & & 
\end{tikzcd}

On choisit un scindage de la première,
qui donne [unclear] un scindage de la seconde
(scindage [unclear: par] un [unclear: sous-groupe] fini [unclear]
[unclear: engendré par] un unique élément de $M^\wedge \dots$). [unclear: En prenant]
[unclear] on peut déduire $U_n$ en-dessous de [unclear]

[MARGIN: On a vu qu'en un sens un peu tautologique les scindages de l'ext. $\mathcal{E}_1$ [unclear: équivalent] à ceux de l'ext. [unclear] (dans 12°), i.e. à un [unclear: ensemble] de choix tels [unclear: canoniques]]

\newpage

%% ===== Page 176 =====

pour tout $j \in U_{0, \mathbf{Q}} \simeq U_{1, \mathbf{Q}}$ il définit un point géométrique de chacun des $U_{\mathbf{Q}}$ au-dessus de $j$. En termes de semi-directs de $\mathcal{E}_1$, permet d'écrire
$$
\mathcal{E}_1 \simeq (\Gamma_{\mathbf{Q}} \times \mathfrak{G}_3) \cdot \hat{M} \leqno{(375)}
$$
(pr. $\frac{1}{2}$ direct)

Les $\mathcal{E}_n \subset \mathcal{E}_1$ s'identifiant aux $\pi_1(U_n, \mathfrak{G}_3)$... produits semi-directs
$$
\mathcal{E}_n = (\Gamma_{\mathbf{Q}} \times \mathfrak{G}_3) \cdot (\hat{M}) \leqno{(375)}
$$
$$
\text{ab. } \hat{M}[n]
$$

(Les interprétations sont des niveaux de généralités de la page 712 ff, et méritent de figurer dans la scholie-proposition en question...)

Les points rationnels $/ \mathbf{Q}$ de $U$ définit un semi-groupe de l'extension (à $\Gamma_{\mathbf{Q}}$ où $\hat{M} = \hat{M}$)
$$
\mathcal{E}_n^! \simeq \Gamma_{\mathbf{Q}} \cdot (\hat{M}) \leqno{(377)}
$$
($\frac{1}{2}$ direct)
quotient de $(371)$, de la connaissance

\newpage

%% ===== Page 177 =====

$736$

modèles $\hat{M}$-conjugués impliquent celles des scindages correspondants à $\mathcal{E}_n$, donc impliquent la connaissance des points correspondants à $x_n \in U_n(\mathbf{Q})$, ce qui signifie que $x_n$ comme tel provient de $r_n = \pi_1(\mathbf{Q})^{0,1,03} / \pi_n(\mathbf{Q})$. Or, $r_n = \pi_n(\mathbf{Q})^{0,1,03} / \pi_n(\mathbf{Q}) = 1$
$$
(378)
$$
si $n$ impair
car $\pi_n(\mathbf{Q}) = 1$ si $n$ impair — donc dans ce cas $x_n$ est connu quand on connaît le scindage correspondant, il est induit dans $\hat{M}$-conjugués.

Ceci nous amène à déterminer toutes les classes de $\hat{M}$-conjugués possibles des tels scindages, afin d'examiner (si fait peut) lesquels peuvent raisonnablement correspondre à un $x_n \in U_n(\mathbf{Q})$. Ces classes correspondent aux éléments de
$$
H^1(\Gamma_{\mathbf{Q}}, \hat{M}) \simeq H^1(\Gamma_{\mathbf{Q}}, \hat{M}) \simeq H^1(\Gamma_{\mathbf{Q}}, \hat{\mathbf{Z}}(-1) \otimes_{\mathbf{Z}} M) \simeq \mathcal{G}(\mathbf{Q})^\wedge \quad (\mathbf{Q}^*)^\wedge \leqno{(379)}
$$

\newpage

%% ===== Page 178 =====

736
modulos $\widehat{\mathcal{M}}$ – ce genre implique alors du scindage correspondant à $\mathcal{E}_1$, donc impliquent la connaissance des points correspondants $x_i$ de $U_n(\mathbf{Q})$, ce qui signifie que $x_n$ se connaît par la position de $\gamma_n = \pi_1(\mathbf{Q})^{\{0,1,\infty\}} / \pi_1(\mathbf{Q})$. Or
$$
\gamma_n = \pi_1(\mathbf{Q})^{\{0,1,\infty\}} / \pi_1(\mathbf{Q}) = 1 \leqno{(378)}
$$
si $n$ impair, car $\pi_1(\mathbf{Q})=1$ si $n$ impair – donc dans ce cas $x_n$ se connaît quand on connaît le scindage correspondant, il se réduit à dans $\widehat{\mathcal{M}}$ – ce genre.

Ceci nous amène à déterminer toutes les classes des $\widehat{\mathcal{M}}$ – ce genre possible de tels scindages, afin d'examiner (si faire peut) lesquels peuvent raisonnablement correspondre à un $x \in U_n(\mathbf{Q})$. Ces classes correspondent aux éléments de
$$
H^1(\Gamma_{\mathbf{Q}}, \widehat{\mathcal{M}}) \simeq H^1(\Gamma_{\mathbf{Q}}, \widehat{\mathcal{M}}) \simeq \operatorname{Hat}(\Gamma_{\mathbf{Q}}, \hat{\mathbf{Z}}(-1) \otimes_{\hat{\mathbf{Z}}} \widehat{\mathcal{M}}) \simeq \mathcal{G}(\mathbf{Q})^\wedge \otimes (\mathbf{Q}^*)^{\wedge} \leqno{(379)}
$$

\newpage

$$
\mathcal{G} \defeq \mathcal{G}_{\mathbf{Q}}^{\{0,1,\infty\}} / \mathcal{G}_{\mathbf{Q}}
$$
et le tore de dimension 2 sur $\mathbf{Q}$ dans lequel se trouvent logés les $U_n$. On voit que tout ça passe comme si on étudiait simplement l'existence (!) et la répartition algébrique des points $\mathbf{Q}$-rationnels de $\mathcal{G}$ lui-même — sans nous préoccuper de l'équation $x_0^n + x_1^n + x_\infty^n = 1$ qui définit $U_n$. Cela montre : \underline{vis-à-vis} que l'approximation laquelle nous travaillons est beaucoup trop grossière — en fait, pratiquement nous sommes dans le domaine d'un objet typiquement “dilué”. Il est clair qu'il est nécessaire de poser des approximations plus fines.

\newpage

%% ===== Page 179 =====

\section*{52. Enveloppe schématique des L-groupes}

1°) Notion de L-groupe, de LT-groupe (ceci dit). La notion de L-groupe est la généralisation naturelle de celle de "groupe à lacets" qui s'impose si l'on veut pouvoir englober des groupes fondamentaux mixtes de surfaces orientables (de type fini, i.e. un $X \setminus S$, $X$ surface compacte, $S$ partie finie) avec un groupe fini d'automorphismes.

Nous avons vu que si $\pi$ est un groupe à lacets (discret), alors tout sous-groupe d'indice fini $\pi_0$ hérite une structure de "sous-groupe à lacets" de $\pi_0$, les intersections avec $\pi_0$ des sous-groupes à lacets de $\pi$.\marginpar{NB. On a \dots}

On a par ailleurs des cas $(g, \nu)$ avec $g \ge 0, \nu \ge 0$, où la structure à lacets n'est déterminée par le seul donné d'une famille de sous-groupes discrets ($\simeq \mathbf{Z}$) --- le cas qui ici faisait problème.

\newpage

%% ===== Page 180 =====

est le cas $(0,1)$ --- le cas $(0,0) \to \pi=\{1\}$ et $\pi_0=\pi \dots$ ]. Inversement, la structure à lacets de $\pi$ se réinterprète à partir de celle de $\pi_0$ (quand elle existe!), des sous-groupes à lacets de $\pi$ étant les anti-images des sous-groupes à lacets de $\pi_0$.

Soit $\pi$ un groupe discret. Soit $\mathcal{L}_0(\pi)$ l'ensemble des structures de groupe à lacets sur $\pi$, de type $(g, \nu)$.

Pour les sous-groupes $\pi_i$ d'indice fini, $\mathcal{L}_0(\pi_i)$ forment un système inductif d'ensembles. On pose, si $\pi_i$ est distingué,
$$ \mathcal{L}(\pi) \defeq \varinjlim \mathcal{L}_0(\pi_i) \leqno{(1)} $$
$$ \mathcal{L}_0(\pi_i) = \text{ensemble des structures à lacets sur } \pi_i \text{ invariants par } \mathcal{G}_i = \pi/\pi_i $$

Il existe $\pi_i$ d'indice fini tel que $\mathcal{L}(\pi_i) \neq \emptyset$, alors $\pi$ admet $\mathcal{L}(\pi) \neq \emptyset$, et il s'ensuit que l'ensemble des structures d'indice fini [unclear: ...]. On pose alors.

\newpage

%% ===== Page 181 =====

\setcounter{page}{180}
\section*{738}
\footnote{M.B. Si $\pi$ n'est pas distingué dans $\pi$, on peut définir dans $\mathcal{L}(\pi)$ une $L$-structure pour $\pi' \subset \pi$, (qui est $\pi'$ dans $\pi$ ...)}
On a
$$\mathcal{L}(\pi) = \varinjlim_{\substack{\pi_i \text{ distingué} \\ \text{d'indice fini} \\ \text{dans } \pi}} \mathcal{L}^\natural(\pi_i)$$

On sait que les morphismes canoniques $\mathcal{L}^\natural(\pi_i) \to \mathcal{L}(\pi)$ (pour $\pi_i$ distingué d'indice fini dans $\pi$) sont injectifs. Si $\lambda \in \mathcal{L}(\pi)$, on dit que $(\pi, \lambda)$ est une $L$-structure sur $\pi$ (et si $\lambda \in \operatorname{Im}(\mathcal{L}^\natural(\pi_i) \to \mathcal{L}(\pi))$, on dit que $\pi_i$ est un sous-groupe de définition de la $L$-structure $\lambda$ --- on la considère comme un groupe : doté par l'image $\lambda_i \in \mathcal{L}^\natural(\pi_i)$ de $\lambda$).

Soient $\pi$ un groupe discret, $\pi_0$ un sous-groupe de $\pi$ d'indice fini. Alors, les morphismes canoniques de réduction de $L$-structures
$$
\mathcal{L}(\pi) \to \mathcal{L}(\pi_0) \leqno{(4)}
$$
qui est injectif, et dont l'image est l'ensemble des $L$-structures sur $\pi$ qui...

\newpage

%% ===== Page 182 =====

"invariantes sous $\pi$" dans un sens évident.
En particulier, si $\pi$ est une extension
d'un groupe fini par un groupe
discret $\pi_0$, alors les $\mathcal{L}$-structures sur
$\pi$ correspondent aux $\mathcal{L}$-structures
sur $\pi_0$ invariantes par $\pi / \pi_0 \to$ [unclear: $Out(\pi_0)$]
opérant sur l'ensemble des $\mathcal{L}$-structures de $\pi_0$
de façon évidente. [NB $\pi \mapsto \mathcal{L}(\pi)$ est
un foncteur en les groupes extérieurs
$\pi$, avec les isomorphismes extérieurs
comme flèches ...]

Exemple standard : soient $X$ une
surface 0-conn. "de présentation finie",
$G$ un groupe fini opérant sur $X$,

(5) $\pi = \pi_1(X, G ; -)$

où $-$ est un pt base de $X$, de
sorte qu'on a une suite exacte

(6) 
\begin{tikzcd}[row sep=small]
1 \ar[r] & \pi_0 \ar[r] & \pi \ar[r] & G \ar[r] & 1 \\
& \pi_1(X,-) \ar[u, equal]
\end{tikzcd}

On a donc sur $\pi_0$ une $\mathcal{L}$-structure : [unclear: locale]

\newpage

%% ===== Page 183 =====

invariante par l'op.\ externe de $G$, d'où une $L$-structure sur $\pi$. Il semble que le $L$-groupe soit isomorphe à un groupe obtenu ainsi — mais c'est là un résultat qui semble profond (c'est équivalent à celui des points fixes des sous-groupes finis des groupes de Teichmüller...), et qui n'est pas compris totalement. Il faudrait dégager un "bon énoncé" général, reliant la 2-catégorie des $L$-groupes, ou celle des surfaces (connexes, finies, ... \footnote{Le manuscrit comporte une correction : "à support fini (à valeurs dans $\mathbf{N}^*$)"}) en associant à telles surfaces un $\pi_1(X, d, a)$, relatif à un pt base $a \in X$-support.

À l'exclusion de quelques cas singuliers, on devrait trouver dans le $L$-groupe $\pi$, isomorphes au $\pi_1(X, G, a)$ où $G$ opère fidèlement, ce qui d'ailleurs...

\newpage

%% ===== Page 184 =====

invariante par l'opération extérieure de $G$, d'où une $L$-structure sur $\pi$. Il semble que le $L$-groupe soit isomorphe à un groupe obtenu ainsi — mais c'est là un résultat qui semble profond (c'est équivalent à celui de l'existence des points fixes des sous-groupes finis des groupes de Teichmüller...), et qui n'est pas compris totalement. Il faudrait dégager "bon énoncé" général, reliant la 2-catégorie des $L$-groupes, ou celle des surfaces de revêtement finis (munies de structures finies (à volumes des $\mathbf{N}^*$), en associant à de telles surfaces un $\pi_1(X, d, a)$, relatif à une base $a \in X$-sous-groupe "triviaux". À l'exclusion de quelques cas singuliers, on devrait trouver dans les $L$-groupes $\pi$, isomorphes à $\pi_1(X, G, a)$ où $G$ opère fidèlement, ce qui s'exprime

en terme de $\pi$ par le fait que la centralisation de tout sous-groupe d'indice fini de $\pi$ est réduit à $\{e\}$. Cela devrait donner une description complète, en termes topologiques, des $L$-groupes ayant cette propriété supplémentaire. Il y aurait ensuite encore un travail de raccordage : fait en termes de certaines extensions des $L$-groupes par des groupes finis, pour traiter le cas général. De plus, il faudrait faire le lien avec les "topos $L$-galoisiens" associés, en complétant les développements correspondants sur les topos galoisiens associés aux groupes à lacets.

\newpage

%% ===== Page 185 =====

740

Il ne pouvait que le "bon" candidat général pour définir et étudier des groupes de Teichmüller généraux soit celui des $L$-groupes, plutôt que celui des groupes à lacets.

Prenant p. ex. le groupe
$$
\mathfrak{S}_{0,3}^+ \isom \\operatorname{Sl}(2, \mathbf{Z}) / \{\pm 1\} \isom \pi_1(U_{0,3}(\mathbf{C}), \mathfrak{S}_{0,3}^+) \leqno{(4)}
$$
celui-ci est muni d'une $L$-structure, qui peut être définie en termes de l'image de la classe de conjugaison du sous-groupe engendré par $\epsilon_0 = \left( \begin{smallmatrix} 0 & -1 \\ 1 & 0 \end{smallmatrix} \right)$, ou plus simplement par le sous-groupe engendré par $\epsilon_1 = \left( \begin{smallmatrix} 1 & 1 \\ 0 & 1 \end{smallmatrix} \right)$.

Le $L$-groupe (discret) a la propriété remarquable que tout $L$-automorphisme "direct" (induisant l'identité sur le module d'orientation) de celui-ci est intérieur. Il se pourrait que cette propriété, ajoutée à celle que le centre est trivial (ou mieux, peut-être, que la centralisatrice...

\newpage

%% ===== Page 186 =====

de tout sous-groupe d'indice fini et réduit à $\{1\}$) caractérise $\mathfrak{S}_{0,3}^+$ à isomorphisme (non unique) près.

NB Pour prouver que tout automorphisme de $\mathfrak{S}_{0,3}^+$ est intérieur, le point essentiel me semble de prouver que l'on conserve le sous-groupe $\pi_{0,3}$, noyau de
$$
1 \to \pi_{0,3} \to \mathfrak{S}_{0,3}^+ \to \mathfrak{S}_3 \to 1
$$
En fait, tout automorphisme de la structure de groupe de $\mathfrak{S}_{0,3}^+$ (pas nécessairement un $L$-automorphisme) conserve le sous-groupe distingué $\pi_{0,3}$. Cela provient du lemme : tout sous-groupe distingué $H$ de $\mathfrak{S}_{0,3}^+$ tel que $\mathfrak{S}_{0,3}^+/H \isom \mathfrak{S}_3$, alors $H = \pi_{0,3}$.

Plus précisément, tout tel surjectif $p: \mathfrak{S}_{0,3}^+ \to \mathfrak{S}_3$ est conjugué à l'homomorphisme canonique (dont le noyau est $\pi_{0,3}$).

On utilise la présentation de $\mathfrak{S}_{0,3}^+$ :
$$
\rho^3 = \sigma^2 = 1
$$

\newpage

%% ===== Page 187 =====

741

qui implique qu'un lien par $\mathfrak{S}_{0,3}^+$ dans un autre groupe, dont l'image n'est pas d'ordre 1, 2 ou 3 i.e. d'ordre $\geq 4$, transformant $\rho$ et $\sigma$ resp. en des éléments d'ordre 3 et d'ordre 2. Donc si l'image est $\mathfrak{S}_3$, $\rho$ est transformé en $\rho$ ou $\rho^{-1}$ (i.e. à conjugaison près c'est $\rho$) et $\sigma$ est transformé en $\sigma$ ou $\sigma_1 = \sigma\rho$ (i.e. à conjugaison près c'est $\sigma$), et on gagne.

On peut appeler groupe de Teichmüller d'un $\mathcal{L}$-groupe $\pi$, le groupe de ses $\mathcal{L}$-automorphismes extérieurs ou peut-être, mieux, le groupe des $\mathcal{L}$-automorphismes extérieurs du sous-groupe $\pi^+$ (d'indice 1 ou 2) de $\pi$, fixant des $g \in \pi$ tels que l'action intérieure soit direct. La propriété envisagée de $\mathfrak{S}_{0,3}^+$ est que son groupe de Teichmüller "direct" est réduit à $1$. La propriété

\newpage

%% ===== Page 188 =====

correspondant aux $\mathcal{L}$-groupes $\mathfrak{S}_{0,3} \simeq \\operatorname{Aut}(\hat{\mathbf{Z}})/\pm 1$ est que tout $\mathcal{L}$-automorphisme du sous-groupe $\mathfrak{S}_{0,3}^+$ est induit par un unique élément de $\mathfrak{S}_{0,3}$. Cette propriété est peut-être caractéristique des $\mathcal{L}$-groupes discrets $\mathfrak{S}_{0,3}$.

J'ai envie de présenter sous une forme légèrement différente la notion de $\mathcal{L}$-groupes en termes de la donnée d'un $\mathbf{Z}$-module $T$ libre de rang 1 (le $\mathbf{Z}$-module d'orientations) et d'une famille finie $(k_i)_{i \in I}$ de groupes d'homomorphismes extérieurs
$$ k_i : \mathfrak{T} \longrightarrow \pi \quad (i \in I) \leqno{(5)} $$
(groupes d'homomorphismes extérieurs), où la notion de ``groupes'' est relative à celle de sous-groupes d'indice fini de $T$.

\newpage

%% ===== Page 189 =====

conclure du groupe $\mathfrak{T}_{0,3} \simeq \\operatorname{Aut}(\hat{\mathbf{Z}})/\{\pm 1\}$ est que tout groupe à lacets du sous-groupe $\mathfrak{T}_{0,3}^+$ est induit par un unique élément de $\mathfrak{T}_{0,3}$. Cette propriété est peut-être caractéristique des L-groupes discrets $\mathfrak{T}_{0,3}$.

J'ai envie de présenter sous une forme légèrement différente la notion de L-groupe en termes de la donnée d'un $\mathbf{Z}$-module $T$ (le $\mathbf{Z}$-module d'orientations) et d'une famille finie $(k_i)_{i \in I}$ de germes d'homomorphismes extérieurs
$$
k_i : \mathfrak{T} \to \pi \quad (i \in I) \quad (\text{germes d'homomorphismes extérieurs}) \leqno{(5)}
$$
où la notion de « germes » est relative à celle de sous-groupes d'indice fini de $T$.

Et tout sous-groupe $\pi_0$ de $\pi$ d'indice fini, on trouve un germe de germes d'homomorphismes extérieurs $T \to \pi_0$ induits par les homomorphismes effectifs $\widetilde{\pi} \to \pi$ de la classe $k_i$ (cela passe en un fini de germes d'homomorphismes extérieurs $T \to \pi_0$ pour $i$ fixé, et $i$ varie dans $I$). Ceci dit, l'unification d'une L-structure, c'est que l'on puisse trouver un sous-groupe $\pi_0$ d'indice fini, tel que la « structure induite » provienne d'une structure : lacets de $\pi_0$.

On peut généraliser la notion de L-structure, en s'inspirant de la définition précédente (p. ex.), et en y laissant tomber l'exigence que $\pi_0$ soit d'indice fini — mais en...

\newpage

%% ===== Page 190 =====

prenant soin par contre que $\pi_0$ soit distingué. On exige que, pour $i \in I$, les groupes d'homomorphismes $\kappa_i: \mathcal{G} \to \pi$ de la classe $\mathcal{K}_i$ ne factorisent par un ensemble de groupes d'homomorphismes extérieurs de $\mathcal{T}$ dans $\pi_0$.\marginpar{MB il n'est pas a priori clair que les groupes d'hom. ext. soient finis.}

et on exige que (avec $\pi_0$ convenable) sur $\pi_0$ provienne d'une structure à lacets sur $\pi_0$, et évidemment déterminée une fois $\pi_0$ choisi. On est défini à commensurabilité près — i.e. $\pi_0, \pi'_0$ convenables, des $\pi_0, \pi'_0$ d'indice fini dans $\pi_0$ et $\pi'_0$, et convient également.

S'agissant de la terminologie "substitutions", j'ai une série de suggestions rappelant les groupes $\pi$ munis de telles.

\newpage

%% ===== Page 191 =====

structure liée à un choix de sous-groupes discrets commensurables entre eux — des L-T-groupes. Le L est initial de «lacets», et T de «Teichmüller») — ce sont des sortes de mixtures de groupes de Teichmüller et de L-groupes. Pour illustrer ce propos, partons d'une extension
$$
1 \to \pi \to \mathfrak{S} \to \mathfrak{T} \to 1 \leqno{(6)}
$$
Si $\pi$ est muni d'une L-structure invariante par $\mathfrak{S}$, alors $\mathfrak{S}$ est un L-T-groupe admettant (si $\pi$ est un groupe à lacets) $\pi$ comme sous-groupe de définition. Pour des groupes à lacets $\pi$, on pose
$$
\begin{cases} 
\mathfrak{S} = \\operatorname{Aut}_L(\pi) \\ 
\mathfrak{T} = \Autext_L(\pi) 
\end{cases} \leqno{(7)}
$$

\newpage

%% ===== Page 192 =====

\section*{}
structure liée à un s-groupe de s-groupes discrets commensurables, ou aux L-T-groupes.
$L$ est initial de "lacets", et $T$ de "Teichmüller" — ce sont des sortes de mélanges de groupes de Teichmüller et de L-groupes. Pour illustrer, partons d'une extension
$$
1 \to \pi \to \mathcal{G} \to \mathfrak{T} \to 1 \leqno(6)
$$
Si $\pi$ est muni d'une L-structure invariante par $\mathcal{G}$, alors $\mathcal{G}$ est un L-T-groupe admettant (ch. \ldots) un s-groupe à lacets $\pi$ comme s-groupe de définition. Pour des groupes à lacets $\pi$, on pose pour $p.i.$
$$
\begin{cases}
\mathcal{G} \defeq \\operatorname{Aut}_L(\pi) \\
\mathfrak{T} \defeq \Autext_L(\pi)
\end{cases} \leqno(7)
$$
et ainsi le groupe $\mathcal{G}$, extension du groupe de Teichmüller $\mathfrak{T}$ par le groupe à lacets $\pi$, apparaît comme un L-T-groupe.

\newpage

%% ===== Page 193 =====

194
$2^\circ$) Cas profini

L'extension de la notion de $\pi_1$-groupe et de $\mathfrak{T}$-groupe en termes de germes d'homomorphismes extérieurs
$$
k^\# : \mathfrak{T} \to \pi
\leqno{(8)}
$$
est essentiellement évidente. Voici l'exemple le plus important pour nous.

Soit d'abord $k$ un corps alg. clos et soit $\mathcal{M}$ un $\mu$-groupe multiplicatif algébrique sur $k$, qui soit de type fini sur $k$, lisse géométriquement et lisse en dimension $1$. On doit pouvoir pour lors que $\mathcal{M}$ peut se décrire par une courbe alg. lisse et quasi-projective $U$ sur $k$, un groupe fini $\mathcal{G}$ qui opère sur $U$, et comme (outil)
$$
\mathcal{M} = (U, \mathcal{G}) = \text{Topos muni des faisceaux étales sur } U \text{ où } \mathcal{G} \text{ opère...}
\leqno{(9)}
$$

\newpage

%% ===== Page 194 =====

Plus intrinsèquement, il doit être
vrai qu'il existe un rev. étale fini
de $M$ (qu'on peut supposer galoisien,
de groupe $G$) qui soit un
[crossed out: schéma algébrique] - lequel sera néc.
quasi-projectif, lisse, de dim relative 1
sur $k$. Supposons en tous cas
qu'il en soit ainsi, choisissons
un pt base géom. $a$ de $M$, et
considérons

(10) $$\pi = \pi_1(M, a) \simeq \pi_1(U, G; a) \ .$$

Nous supposerons désormais $k$ de
car. zéro. Supposons d'abord $k$
alg clos. Utilisons la suite exacte

(11) $$1 \longrightarrow \underset{\pi_0}{\pi_1(U_{\tilde{a}})} \longrightarrow \underset{\pi}{\pi_1(M, G; a)} \longrightarrow G \longrightarrow 1$$

et la structure : [unclear: toute] naturelle
sur $\pi_0$, on trouve une $L$-structure
sur $\pi = \pi_1(M, a)$, qui [unclear: manifeste] sur

\newpage

%% ===== Page 195 =====

745

ne dépend pas du choix du revêtement étale fini $U$ de $M$, et des relèvements choisis des points géométriques $a$ de $M$ en points géométriques de $U$. Nous munissons (techniquement) $\pi = \pi_1(M, a)$ de cette L-structure.

Si à nouveau $k$ est un corps quelconque (le cas le plus intéressant étant celui où $k$ est de type fini sur $\mathbf{Q}$), on a une suite exacte
$$
1 \to \pi_1(M_{\bar{k}}, a) \to \pi_1(M, a) \to \Gamma_{\bar{k}/k} \to 1 \leqno{(12)}
$$
et $\pi_1(M_{\bar{k}}, a)$ apparaît comme muni d'une L-structure, inscrite par $\pi_1(M, a) = \pi$, de sorte que $\pi$ apparaît comme un $\Theta$-groupe.

Je suis particulièrement intéressé au cas « universel » de type $(g, \nu)$ — mais la base n'est pas alors la

\newpage

%% ===== Page 196 =====

spectre et des corps $k$, mais les multiplicités $M_{\bar{k}}$. [A voir, pour la structure de $L$-groupes discret $\pi_0$, jouant le rôle de groupe de référence, on doit pouvoir faire introduire une multiplicité sur $\mathbf{Q}$, et en sur $\mathbf{Z}$. Le fait qui risque de faire des difficultés conceptuelles pour une telle construction, est qu'un "modèle" multiplicité (sur $\mathbf{Q}$, ... sur $\mathbf{Z}$...) décrit en principe une catégorie fibrée -- telle celle des schémas, ou des relatifs de type $(g,\nu)$ -- alors que les multiplicités relatives : un "type" $\pi_0$ (un groupoïde), fournit une 2-catégorie fibrée sur $(Sch)$, et pas seulement une catégorie fibrée...]

\newpage

%% ===== Page 197 =====

745

On a donc une suite exacte
$$
1 \to \Pi_{g,\nu} \to \underset{M_{g,\nu}}{\Pi_1(M_{g,\nu})} \to \underset{N_{g,\nu}}{\Pi_1(M_{g,\nu,\mathbf{Q}})} \to 1 \leqno{(13)}
$$
$M_{g,\nu}$ est la "courbe de type $(g,\nu)$ universelle" (en car. 0). Ici c'est $\Pi_1(M_{g,\nu}, \bar{\eta})$ qui apparaît comme un $\mathfrak{T}$-groupe profini.

Je voudrais dégager une propriété commune à ce cas "universel" (13), (qui est d'ailleurs en fait : un corps de base $k$ extension de type fini de $\mathbf{Q}$. Le premier cas est d'ailleurs essentiellement en commun avec le second, en pensant au "point générique" de $M_{g,\nu} \dots$). Puis ceci, considérons un $\mathfrak{T}$-groupe $\mathcal{M}$ et un homomorphisme $k: \mathfrak{T} \to \mathcal{M}$ appartenant à un clan $K_{\mathfrak{T}}$.

\newpage

%% ===== Page 198 =====

p. ex. choisissons un sous-groupe de définition distingué $\pi$ dans $M$, qui sera donc un groupe à lacets, et donnons une suite exacte
$$
1 \to \pi \to M \to N \to 1 \leqno{(14)}
$$
et un homomorphisme
$$
k : T \to \pi, \quad k(T) = L \leqno{(15)}
$$
relatif : la structure à lacets de $\pi$.
Considérons le normalisateur de $L$ dans $M$
$$
\operatorname{Norm}_M(L) = N(L) \leqno{(16)}
$$
et considérons l'action de $N(L)/L$ sur $L$. (NB : On voit aisément que l'image de $N(L)/L$ dans $N$ est un sous-groupe distingué de $N \dots$) Cette action se fait par un « caractère cyclotomique », induit par une section
$$
N(L)/L \to \hat{\mathbf{Z}}^* \leqno{(17)}
$$
L'hypothèse essentielle est que l'image

\newpage

%% ===== Page 199 =====

d'un $N$ pour $X$ l'image de $N(L)/L$ dans $X$ soit un sous-groupe ouvert de $\hat{\mathbf{Z}}^*$.

Cette hypothèse est un caractère intrinsèque du $\ell$-groupe $M$ (bien $= \mathcal{G}$-fondamental des diviseurs de $\pi$) — car le caractère cyclotomique
$$
\chi_M : M \to \hat{\mathbf{Z}}^* \leqno{(18)}
$$
est un caractère intrinsèque : $M$, et l'hypothèse signifie que l'image des diviseurs dans $\hat{\mathbf{Z}}^*$ est ouvert.

Revenant à $L$ et $N(L)$, on voit que l'hypothèse fait impliquer ceci :
$$
\text{Tout sous-groupe ouvert de } N(L)/L \text{ opère non trivialement sur } L. \leqno{(19)}
$$

On sait, par les arguments standards que ceci implique que pour tout corps $K$ (notamment un corps local $\mathbf{Q}_\ell$),

\newpage

%% ===== Page 200 =====

et tout homomorphisme de $\mathcal{M}$ dans un groupe $G(K)$; où $G$ est un schéma en groupe de type fini affine sur $K$; et
$$
u : \mathcal{M} \longrightarrow G(K) \leqno{(20)}
$$
la restriction de $u$ à $L$ est "quasi-unipotente", i.e.
$$
\exists \text{ sous-groupe ouvert } L' \text{ de } L, \text{ tel que } u(L') \subset G(K) \text{ soit formé d'éléments unipotents}, \leqno{(21)}
$$
ce qui revient à $\implies$, si $l_i$ est un générateur de $L = \langle l_i \rangle$; alors
$$
u(l_i) \text{ est un élément quasi-unipotent de } G(K), \text{ i.e. } \exists n \in \mathbf{N}^* \text{ tel que } u(l_i)^n \text{ soit unipotent.} \leqno{(22)}
$$

\newpage

%% ===== Page 201 =====

(30) Enveloppe schématique d'un T-L-groupe profini : les homomorphismes admissibles.

Soit $M$ un T-L-groupe profini, de la condition structurale (18)
$$ M \to \hat{\mathbf{Z}}^\times \leqno{(18)} $$
soit un $i$-groupe ouvert, $\pi$ un sous-groupe de définition (distingué, c'est-à-dire dans le groupe profini), et $\Pi_0 \subset \pi$ un sous-groupe discret de $\pi$, qui soit une "discrétification" --- qui est un groupe à lacets discret, c'est-à-dire dans son complété profini.

Soit $G$ un schéma en groupes sur $\mathbf{Q}$. Soit $\mathbf{A}_{\mathbf{Q}}$ l'anneau des adèles de $\mathbf{Q}$
$$ \mathbf{A}_{\mathbf{Q}} = \hat{\mathbf{Z}} \otimes_{\mathbf{Z}} \mathbf{Q} \subset \prod_{\ell \in P} \mathbf{Q}_\ell \leqno{(23)} $$
et considérons le groupe des points adéliques
$$ G(\mathbf{A}_{\mathbf{Q}}) = \prod_{\ell \in P} G(\mathbf{Q}_\ell). \leqno{(24)} $$
Si $G$ provient d'un schéma en groupes

\newpage

%% ===== Page 202 =====

Schème de type fini $G$ sur $\mathbf{Z}$, de
$$ G(\mathbf{A}_{\mathbf{Q}}) \cong \{(x_\ell)_{\ell \in P} \in \prod_\ell G(\mathbf{Q}_\ell) \mid x_\ell \in G(\mathbf{Z}_\ell) \text{ pour presque tout } \ell \} \leqno{(25)} $$
On dit que
On se donne de plus un homomorphisme
$$ G_{\mathbf{Q}} \xrightarrow{\chi} G_{\mathbf{Q}} \leqno{(26)} $$
On s'intéresse aux homomorphismes
$$ \mathcal{M} \to G(\mathbf{A}_{\mathbf{Q}}) \leqno{(27)} $$
satisfaisant à des propriétés que nous allons énumérer.

1°) Continuité. Ce qui revient à dire que
\begin{enumerate}
    \item[a)] le homomorphisme $\mathcal{M} \to G(\mathbf{A}_{\mathbf{Q}})$ est continu, et
    \item[b)] $\exists S$ partie finie de $P$, telle que pour $\ell \notin S$, on ait
\end{enumerate}
$$ \mathcal{M} \to G(\hat{\mathbf{Z}}_\ell) \quad (\ell \notin S) \leqno{(28)} $$
— le tout qui en effet ne dépend pas des «modèles sur $\mathbf{Z}$» $G$ de $G_{\mathbf{Q}}$ choisi.
NB (28) est un homomorphisme de groupes profinis.

\newpage

%% ===== Page 203 =====

741

ReNB Sauf erreur, la propriété 1) signifie qu'
on peut trouver un schéma en groupes $G$
affine de t.f. sur $\mathbb{Z}$, qui définit $G_\mathbb{Q}$, et
tel que le (27) se factorise par $G(\hat{\mathbb{Z}})$.
(29) $$M \longrightarrow G(\hat{\mathbb{Z}})$$

2°) Compatibilité avec $\chi$ :

(30)
\begin{tikzcd}
M \ar[r] \ar[d, "\chi_M"'] & G(\mathbb{A}_\mathbb{Q}) \ar[d, "\chi_G"] \\
\hat{\mathbb{Z}} \ar[r, hook] & \mathbb{G}_m(\mathbb{A}_\mathbb{Q}) \simeq \mathbb{A}_\mathbb{Q}^* \subset \prod_{l \in P} \mathbb{Q}_l^*
\end{tikzcd}

est commutatif.

3°) Compatibilité avec la [unclear: désintégration] $\pi_0$

(31)
\begin{tikzcd}
\Pi'_0 \ar[r, hook] \ar[d, "\text{factorisation}"'] & \Pi_0 \ar[r, hook] & M \ar[d] \\
G(\mathbb{Q}) \ar[rr, hook] & & G(\mathbb{A}_\mathbb{Q})
\end{tikzcd}
,

pour un sous-groupe d'indice fini
convenable $\Pi'_0$ de $\Pi_0$.

NB Quitte à changer les notations, [unclear: on peut supposer] $\Pi'_0 = \Pi_0$.

4°) Caractère minimal (irrédundance) :
Si $G'$ est un sous-schéma en groupes de

\newpage

%% ===== Page 204 =====

$G_{\mathbf{Q}}$, tel que $\mathcal{M} \to G(\mathbb{A}_{\mathbf{Q}})$ se factorise par $G'(\mathbb{A}_{\mathbf{Q}})$, alors $G=G'$.

NB Si cette hypothèse n'est pas vérifiée, alors pour tous les $G' \subset G_{\mathbf{Q}}$ qui la vérifient, il y en a un unique, plus petit, et en remplaçant $\chi$ par $\chi' = \chi|_{G'}$, les hypothèses 1) à 4) sont satisfaites pour $G'$, et $\mathcal{M} \to G'(\mathbb{A}_{\mathbf{Q}})$.

\medskip

\noindent Remarques. Soit $\mathcal{G}_0 \subset \mathcal{M}$ le sous-groupe de $\mathcal{M}$ qui \marginpar{Existence de $\mathcal{M}$ d'un groupe de Teichmüller discret} soit les groupes $\pi_0$ sont commensurables. Ceci généralise le ``groupe de discrétisation'' $\pi_0$ de $\pi$.
Il est possible qu'il faille renforcer la condition 3), en exigeant que l'on ait $\mathcal{G}_0 \to G(\mathbf{Q})$.

Il convient le décider!

\newpage

%% ===== Page 205 =====

250

Réflexion faite, je préfère poser les choses de
façon légèrement différente, en introduisant,
dans le LT-gr. profini $M$, un ~~[unclear]~~
distingué $\pi$ avec une L-structure (par
ex. une structure : toute pièce,
si par ex $\pi = \operatorname{Sl}(2, \mathbb{Z})$) définissant la
L T-structure de $M$, une discrétisation
$\pi_0$ de $\pi$ (dans un sens évident, pour les
L-structures) un sous-groupe discret
$\mathfrak{S}_0 \supset \pi_0$ normalisant $\pi_0$ (NB. dans mon petit
exemple [unclear] on a $\mathfrak{S}_0 = \pi_0$), tel que
[MARGIN: son adh.] $\mathfrak{S} = \widehat{\mathfrak{S}}_0$ dans $M$ soit distingué dans $M$.

On a donc un diag. de groupes

(32)
\begin{tikzcd}
& \tau \ar[d, no head] & \Gamma \ar[d, no head] \\
1 \ar[r, phantom, "\subset"] & \pi \ar[r, phantom, "\subset"] \ar[d, equal] & \mathfrak{S} \ar[r, phantom, "\subset"] \ar[d, equal] & M \\
& \widehat{\pi}_0 & \widehat{\mathfrak{S}}_0 \\
& \uparrow & \uparrow \\
1 \ar[r, phantom, "\subset"] & \pi_0 \ar[r, phantom, "\subset"] \ar[uu] & \mathfrak{S}_0 \ar[uu]
\end{tikzcd}

où la première ligne est formée de
groupes profinis, la deuxième de groupes
discrets. On notera que $\mathfrak{S}$ s'identifie : au
quotient du complété profini de $\mathfrak{S}_0$, mais

\newpage

%% ===== Page 206 =====

Je reviens pour le moment qu'ils sont isomorphes à ce complété\footnote{Dans le cas $\pi = \operatorname{Sl}(2, \mathbf{Z})$}. Dans ce cas, $\pi$ et $\pi_0$ (le groupe) pour le rôle d'un groupe de Teichmüller "discret", le groupe $\mathfrak{T}$ alors qu'un complété de groupe de Teichmüller, le groupe $\Gamma = \mathcal{M}/\mathcal{G}$ pour le rôle d'un groupe Galois profini. La situation je suis intéressé par deux situations extrêmes :

\begin{enumerate}
\item[a)] $\mathcal{G}_0 = \pi_0$ d'où $\mathcal{G} = \hat{\pi}$ --- situation provenant d'une multiplication lisse relative sur un corps de type fini sur $\mathbf{Q}$, ou plus généralement d'un schéma (comme multiplication) en fait\footnote{C'est en fait le cas fondamental d'existence d'une telle multiplication relative.} de type fini sur $\mathbf{Q}$.
\item[b)] $\mathcal{G}/\pi_0 \simeq$ groupe de Teichmüller ou d'indice fini dans un groupe $\mathfrak{T}$ de $\pi_0$ --- des cas, en fait, semblant à une situation "modulaire".
\end{enumerate}

\newpage

%% ===== Page 207 =====

$\Gamma_{\mathbf{Q}/\mathbf{Q}}$ s'interprète comme le groupe fondamental des lieux multiplicité modulaire convenable, $\mathcal{G}$ comme sa partie "géométrique", $\Pi_0$ comme sa partie "transcendante", et $\Gamma$ comme $\Gamma_{\mathbf{Q}/\mathbf{Q}} \dots$

On fait, au (32), l'hypothèse que la courbe algébrique sur $\mathcal{M}$ (opérant sur $\pi$ pro-abélien, $\ell$-adique) soit livre, une action sur $\pi$ (avec (33) $\pi = \pi^\tau, \mathcal{G} = \mathcal{G}^+$), i.e. qu'il y ait une factorisation
$$
\Gamma \to \mathcal{X}_\Gamma \hat{\mathbf{Z}}^* \leqno{(34)}
$$
dont l'image soit un sous-groupe ouvert.

Si $(G_{\mathcal{G}, a})$ est comme plus haut (cf. § 48), s'intéresse à la (37) $\mathcal{M} \to G(\mathbb{A}_{\mathbf{Q}})$ satisfaisant les conditions 1), 2), 3), 4) comme dessus, mais où 3) est remplacée par la forme un peu plus forte (30 forte) (compatibilité de discrétification de Teichmüller $\mathcal{G}$) : on a une factorisation.

\newpage

%% ===== Page 208 =====

tion $\mathbb{G}_0 \to G(\mathbb{Q})$, rendant commutatif le diagramme

(35)
\begin{tikzcd}
\mathbb{G}_0 \ar[r, hook] \ar[d] & \mathcal{M} \ar[d] \\
G(\mathbb{Q}) \ar[r, hook] & G(\mathbb{A}_{\mathbb{Q}})
\end{tikzcd}

Les structures de la forme (32) 
[MARGIN: Quitte par la suite à introduire des variantes sur ces structures]
méritent sans doute un nom, je vais les appeler - ne serait-ce que provisoirement - structures de Galois-Teichmüller.

Si $G$ est un groupe algébrique affine sur $\mathbb{Q}$, un hom $\mathcal{M} \to G(\mathbb{A}_{\mathbb{Q}})$ satisfaisant les conditions 1) 2) $3^{\circ \text{forte}}$) 4) est appelé un "hom de Lie" de la structure de Galois-Teichmüller.

Je vais désigner par la suite par (en abrégé $M$) $M_{\mathbb{Q}}$ le schéma en groupes sur $\mathbb{Q}$ affine.
Je désigne par $\mathbb{G}_{\mathbb{Q}}$ le sous-schéma en groupes de $M_{\mathbb{Q}}$ engendré par l'

\newpage

%% ===== Page 209 =====

image de $\mathfrak{G}_{\mathbf{Q}} \to G(\mathbf{Q})$, par $\Pi_{\mathbf{Q}}$ celui engendré par l'image de $\Pi \to G(\mathbf{Q})$.

C'est ainsi que des sous-groupes :
$$
\Pi_{\mathbf{Q}} \subset \mathfrak{G}_{\mathbf{Q}} \subset \mathcal{M}_{\mathbf{Q}} \leqno{(36)}
$$
distingués des $\mathcal{M}_{\mathbf{Q}}$, donnant lieu à des schémas de groupes quotients :
$$
\overline{\mathfrak{G}}_{\mathbf{Q}} \cong \mathfrak{G}_{\mathbf{Q}} / \Pi_{\mathbf{Q}}, \quad \Gamma_{\mathbf{Q}} \cong \mathcal{M}_{\mathbf{Q}} / \mathfrak{G}_{\mathbf{Q}}, \quad \mathcal{N}_{\mathbf{Q}} \cong \mathcal{M}_{\mathbf{Q}} / \Pi_{\mathbf{Q}} \leqno{(37)}
$$
donnant lieu à des suites exactes de schémas de groupes, dont la dernière est :
$$
1 \to \overline{\mathfrak{G}}_{\mathbf{Q}} \to \mathcal{N}_{\mathbf{Q}} \to \Gamma_{\mathbf{Q}} \to 1 \leqno{(38)}
$$
Notons enfin le morphisme $\chi_{\mathbf{Q}} : \mathfrak{G}_{\mathbf{Q}} \to \mathcal{G}_{\mathbf{Q}}$ un épi, et sa factorisation par $\Gamma_{\mathbf{Q}}$ :
$$
\Gamma_{\mathbf{Q}} \xrightarrow{\chi_{\mathbf{Q}}} \mathcal{G}_{\mathbf{Q}} \leqno{(39)}
$$

\newpage

%% ===== Page 210 =====

Le système de groupes (36), (37), muni de l'homomorphisme
$$
\pi \to \mathcal{G} \to \mathcal{M}
$$
(40) $$
\pi(\mathbb{A}_{\mathbf{Q}}) \subset \mathcal{G}(\mathbb{A}_{\mathbf{Q}}) \subset \mathcal{M}(\mathbb{A}_{\mathbf{Q}})
$$
induit par $\mathcal{M} \to \mathcal{G}(\mathbb{A}_{\mathbf{Q}})$ satisfaisant aux conditions 1) à 4) ci-dessus, i.e.\ d'un "lieu de Lie" s'appellera une "approximation de Lie" du système de Galois-Teichmüller envisagé.

Considérons un homomorphisme
$$
(41) \quad \mathcal{T}_0 \to \pi_0
$$
\marginpar{NB. \textit{Pour fixer les idées}: Le choix du "système" (comprenant les conditions d'homomorphisme) et la "complétion" $\pi_0 \to \mathcal{G}(\mathbf{Q})$ et des "groupes de homomorphismes" \textit{dérivés} du choix de $\mathcal{T}_0$ (voir (42))}
où $\mathcal{T}_0$ est la structure de Teichmüller, et $\pi_0 \to \mathcal{G}(\mathbf{Q})$ la complétion.
$$
(42) \quad \mathcal{T}_0 \to \mathcal{G}(\mathbf{Q})
$$
(voir (21) p. 747), le unit pour cet homomorphisme un "morphisme" \dots un quotient, et qui est défini par suite d'homomorphismes de schémas en groupes.

\newpage

%% ===== Page 211 =====

(43) $k_i \colon T_0 \xrightarrow{\sim} \Pi_{\mathbf{Q}}$
($\cong \mathbf{Q} \otimes_{\mathbf{Z}} T_0$)
qui définit, par extension par les points $\mathbf{Q}$-rationnels, le groupe d'homomorphismes (42). On trouve ainsi un cas fini d'un schéma de groupes de $\pi_0$ et $\Pi$ qui vient d'ailleurs, satisfaisant (43) de schéma, un groupe. J'ai envie d'appeler l'approximation de Lie un "diagramme" les cas précédents sous tous leurs faciès, de $2$ des mêmes.

(44) $\underline{k}_i \colon \underline{T}_{\mathbf{Q}} \cong \mathbf{G}_{\mathbf{Q}} \otimes_{\mathbf{Z}} T_0 \to \underline{\Pi}_{\mathbf{Q}}$.

Supposons que $\mathcal{M}_{\mathbf{Q}}$ provienne d'un schéma de groupes affine t.f. $\mathcal{M}$ sur $\mathbf{Z}$, tel que l'on ait
(45) $\mathcal{M} \to \mathcal{M}(\hat{\mathbf{Z}})$
(il en existe, on peut toujours trouver un tel $\mathcal{M}$...). On dirigera le par...

\newpage

%% ===== Page 212 =====

$\underline{\Pi}$, $\underline{\mathfrak{S}}$ les adhérences schématiques
dans $\underline{M}$ de $\underline{\Pi}_\mathbb{Q}$, $\underline{\mathfrak{S}}_\mathbb{Q}$, d'où un diagramme
de schémas en groupes sur $\mathbb{Z}$

(46) $$ \underbrace{\underline{\Pi} \subset \underline{\mathfrak{S}}}_{\underline{\Gamma}} \subset \underline{M} $$
[MARGIN: with another brace under $\underline{\mathfrak{S}} \subset \underline{M}$ labeled $\underline{N}$]

(sous réserve de représentabilité des
quotients $\underline{\tau}, \underline{\Gamma}, \underline{N}$ - mais je crois que
pour Raynaud il n'y a jamais de pb...).

[unclear: On a des factorisations] [crossed out]

(47)
\begin{tikzcd}
\underline{\Pi} \ar[r, hook] \ar[d] & \underline{\mathfrak{S}} \ar[r, hook] \ar[d] & \underline{M} \ar[d] \\
\underline{\Pi}(\hat{\mathbb{Z}}) \ar[r, hook] & \underline{\mathfrak{S}}(\hat{\mathbb{Z}}) \ar[r, hook] & \underline{M}(\hat{\mathbb{Z}})
\end{tikzcd}

On suppose bien sûr également que
l'hom $\chi_{\underline{M}} : \underline{M}_\mathbb{Q} \to \mathbb{G}_{m \mathbb{Q}}$ se
prolonge en

(48) $$ \chi_{\underline{M}} : \underline{M} \to \mathbb{G}_{m \mathbb{Z}} $$

lequel nécessairement se factorise par

(49) $$ \chi_{\underline{\Gamma}} : \underline{\Gamma} \to \mathbb{G}_{m \mathbb{Z}} . $$

\newpage

%% ===== Page 213 =====

Il serait intéressant de dégager un $k$-d'existence (voire d'unicité, moyennant des conditions plus fortes) d'un "modèle" $\mathcal{M}$ sur $\mathbf{Z}$ et $\mathbf{Q}$, satisfaisant à nos conditions précédentes ($\mathcal{M} \to \mathcal{M}(\mathbf{Z})$ et $\mathcal{M} \xrightarrow{\chi_\mathcal{M}} G_\mathbf{Z}$).

Je n'ai pas un seul plan en général, je présume que les homomorphismes extérieurs (43) se prolongent tels quels en des homomorphismes $T_k \times T_\mathbf{Z} \to \Pi$. Mais je vais considérer le centralisateur de $k: T_0 \to \Pi_0$, soit $L_{i,0}$, qui est un groupe métabélien : $\operatorname{Centr}_{\Pi_0}(k: T_0 \to \Pi_0)$ est un sous-groupe distingué d'indice fini. Considérons l'homomorphisme
$$
\begin{array}{ccc}
L_{i,0} & \to & \Pi(\mathbf{Z}) \\
\parallel & & \\
\operatorname{Centr}_{\Pi_0}(k: T_0 \to \Pi_0) & &
\end{array} \leqno{(50)}
$$
qui est unipotent — il se factorise donc par un homomorphisme de schémas en groupes
$$
E_\infty(L_{i,0}) \to \Pi \leqno{(51)}
$$
(cf. § 50, 4$^\circ$, p. 658 et suiv.), se factorisant.

\newpage

%% ===== Page 214 =====

par des groupes lisses\footnote{Les $Z_{i,n}$ sont des groupes lisses sur $S_0 = \operatorname{Spec} \mathbf{Z}$.} $E_{\operatorname{niv}}(\widetilde{L}_{i,0})$, et on a une suite :
$$
1 \to \underline{\pi}_0 \to Z_{i,n} \to (\pi_1)_{S_0} \to 1 \leqno (52)
$$
avec $Z_{i,n} \cong \widetilde{L}_{i,0}/\underline{\pi}_0$.
Mais, bien entendu, la restriction de l'homomorphisme
$$
Z_{i,n} \to \underline{\pi}
$$
est $Z_{i,n} \simeq \underline{\pi}_0$.

Ce n'est pas en général un mono — s'il l'est pour $n$ fini, alors $n'|n$, $n \neq n'$; il ne l'est plus pour $n'$. Une condition favorable serait alors de pouvoir trouver un homomorphisme
$$
\underline{\pi}_0 \to \underline{\pi}
$$
qui est un mono. Cela indique que l'adhérence schématique dans $\underline{\pi}$ du "sous-groupe additif" $\widetilde{L}_{i,0} \defeq \\operatorname{Ker}(\text{liste de } \pi)$ est un sous-schéma en groupes de $\underline{\pi}$.
$$
\underline{\pi}_0 \to \underline{\pi} \leqno (53)
$$
\marginpar{cat. et très unique}

\newpage

%% ===== Page 215 =====

Cette condition doit être\marginpar{755} suffisante à ce que cette schématique soit une composante neutre qui soit un isomorphisme : Cas 0.
On trouve alors, un homomorphisme unique (à conjugaison près)
$$ k'_Q \colon \underline{T}_0 \to \underline{\Pi} \leqno{(54)} $$
tel que $k'_Q$ soit de la forme $k_Q \circ \text{id}$, pour $\mathcal{G} \in \mathcal{Q}^*$ convenable. Je ne suis pas sûr que sans cette condition, on puisse trouver $k'$ déduit de $k'_Q$ passant aux $\mathbf{Z}$-points
$$ T_0 \isom \underline{T}_0(\mathbf{Z}) \to \underline{\Pi}(\mathbf{Z}) $$
se factorisant en un
$$ T_0 \xrightarrow{k'_Q} \underline{T}_0 \xrightarrow{\text{comp.}} \underline{\Pi}(\mathbf{Z}) \leqno{(55)} $$
$T_0 \to \underline{T}_0$ m'amène au ``groupe'' $k'_Q$.
Je vais nécessiter le recours comme hypothèse supplémentaire lorsque les $k'_Q$ existent.
J'ai envie d'énoncer quelques autres ``bonnes hypothèses'', qui s'effacent en réalité.

\newpage

%% ===== Page 216 =====

en travaillant avec des approximations de Lie « anti-corps » i.e. sur $\operatorname{Spec} \mathbf{Z}$.

Considérons tous les homomorphismes $k'_i$ possibles (i.e. variables, mais à une conjugaison près $k'_i : \mathfrak{T}_0 \to \mathfrak{T}_0$ en défaut par des éléments quelconques de $\mathfrak{T}_0$), et les homomorphismes composés sur les algèbres de Lie :
$$
\operatorname{Lie}(k'_i) : \mathfrak{T}_0 \to \operatorname{Lie}(\mathfrak{T}) = \operatorname{Lie}(\mathfrak{L}) \leqno{(58)}
$$
qui sont par hypothèse des isomorphismes universels, on voudrait que $\operatorname{Lie}(\mathfrak{T})$ soit $\mathbf{Z}$-module engendré par les images qui seraient les images d'un générateur de $\mathfrak{T}_0$. \marginpar{cette hyp. raisonnable $\mathfrak{T}_0 = \\operatorname{Sl}(2, \dots)$ mais...} Je verrais... que n'est d'autre que $\operatorname{Lie}(\mathfrak{T})$ pour là : $k'_i$, dans l'image déduit générateurs des $\mathfrak{T}_0$, les $I$ éléments obtenus de $\operatorname{Lie}(\mathfrak{T})$, et leurs conjugués par les éléments de $\mathfrak{T}_0$, engendrent $\mathbf{Z}$-module $\operatorname{Lie}(\mathfrak{T})$. Ceci signifierait que

\newpage

%% ===== Page 217 =====

756

Le schéma en groupes $\underline{\Pi}^\circ$ sur $S_0 = Spec \mathbb{Z}$, vu en tant
que faisceau étale sur le "gros topos étale" de
$S_0$, est engendré par ses sous-faisceaux en
sous-groupes
$\underline{k}'_!(\underline{G}_\circ)$ (avec répétition possible d'une même
classe de conjugaison), et encore $\underline{\Pi}$ est
(Et il y a [unclear] pratiquement quant aux
les $\underline{k}'_!(\underline{G}_\circ)$ lors d'un changement de base) et
le groupe discret $(\Pi_0)_{S_0}$. (Et il se pose
la question naturellement d'expliciter les
relations... suffit-il d'écrire les relations de
$\Pi_0$, plus les commutativités de

(57)
\begin{tikzcd}
T_0 \ar[r, "k'_!"] \ar[d, "can"'] & \Pi_0 \ar[d, "can"] \\
\underline{T}_0 \ar[r, "\underline{k}'_!"'] & \underline{\Pi}
\end{tikzcd}
? )

Autres conditions concernant les schémas
en groupes

(58)
$$ \gamma_{\underline{\Pi}} = \underline{\Pi} / \underline{\Pi}^\circ, \gamma_{\underline{\mathfrak{S}}} = \underline{\mathfrak{S}} / \underline{\mathfrak{S}}^\circ, \gamma_{\underline{M}} = \underline{M} / \underline{M}^\circ $$

- on voudrait que ces schémas soient
en groupes constants, correspondant aux
schémas groupes finis discrets qui décrivent les
quotients correspondants des groupes sur $\mathbb{Q}$,
lesquels groupes discrets sont des quotients

\newpage

%% ===== Page 218 =====

(59) $\gamma_{\Pi}$ (resp.\ $\gamma_{\mathfrak{S}}$, resp.\ $\gamma_{\mathcal{M}}$) quotients de $\Pi, \mathfrak{S}, \mathcal{M}$
$$
\gamma_{\Pi} \simeq (\hat{\Pi})_{S_0}, \quad \gamma_{\mathfrak{S}} \simeq (\hat{\mathfrak{S}})_{S_0}, \quad \gamma_{\mathcal{M}} \simeq (\hat{\mathcal{M}})_{S_0} \footnote{d'après le théorème [unclear]}
$$
$$
\gamma_{\Pi} \to \gamma_{\mathfrak{S}} \to \gamma_{\mathcal{M}} \leqno{(60)}
$$
et des suites exactes
$$
\begin{cases}
\gamma_{\Pi} \to \gamma_{\mathfrak{S}} \to \gamma_{\mathfrak{T}} \to 1 \\
\gamma_{\Pi} \to \gamma_{\mathcal{M}} \to \gamma_{\Gamma} \to 1 \\
\gamma_{\mathfrak{S}} \to \gamma_{\mathcal{M}} \to \gamma_{\mathfrak{T}} \to 1 \\
\gamma_{\mathfrak{T}} \to \gamma_{\Gamma} \to \gamma_{\Gamma'} \to 1
\end{cases} \leqno{(61)}
$$

Cette hypothèse sur $\Pi, \mathfrak{S}, \mathcal{M}$ permet de définir les groupes $\Pi^{\wedge}(\hat{\mathbf{Z}}), \mathfrak{S}^{\wedge}(\hat{\mathbf{Z}}), \mathcal{M}^{\wedge}(\hat{\mathbf{Z}})$ et on aura nécessairement les deux premières factorisations de
$$
\begin{cases}
\Pi \to \Pi^{\wedge}(\hat{\mathbf{Z}}) \\
\mathfrak{S} \to \mathfrak{S}^{\wedge}(\hat{\mathbf{Z}}) \\
\mathcal{M} \to \mathcal{M}^{\wedge}(\hat{\mathbf{Z}})
\end{cases} \leqno{(63)}
$$
$( \Pi_0 \to \Pi(\hat{\mathbf{Z}}) \subset \Pi^{\wedge}(\hat{\mathbf{Z}}), \mathfrak{S}_0 \to \mathfrak{S}(\hat{\mathbf{Z}}) \subset \mathfrak{S}^{\wedge}(\hat{\mathbf{Z}}) )$. Quant à la troisième factorisation, elle devrait être incluse dans la définition de l'objet, à l'aide d'une structure de Teichmüller-Galois, sachant que $\hat{\mathfrak{T}}_g / \hat{\mathfrak{T}}_g^{\circ}$ est construit, ce qui permet \marginpar{C'est un [unclear]...}

\newpage

%% ===== Page 219 =====

alors de définir $\mathfrak{M}^A(\mathbf{A}_Q)$\marginpar{757}, et on exige que $M \to \mathfrak{M}^A(\mathbf{A}_Q)$ factorise par
$$
M \to \mathfrak{M}^A_{\hat{\mathbf{Q}}}(\mathbf{A}_Q) \leqno{(63)}
$$
Si de plus on suppose que les images des trois homomorphismes (62) sont des sous-groupes invariants, (conditions qui renforcent considérablement les conditions données dans 4) (p. 149), pouvons-nous espérer la surjectivité des applications
$$
\begin{cases}
\pi \to \pi^A(\hat{\mathbf{Z}}) \to \Gamma_\pi \\
\mathfrak{G} \to \mathfrak{G}^A(\hat{\mathbf{Z}}) \to \Gamma_\mathfrak{G} \\
\mathfrak{M} \to \mathfrak{M}^A(\hat{\mathbf{Z}}) \to \Gamma_\mathfrak{M}
\end{cases} \leqno{(64)}
$$
on sera particulièrement heureux de dire qu'on a une uniformisation de Lie pour $\mathbf{Z}$, dite "entière" parfaite, si toutes les conditions ci-dessus sont satisfaites ...

\newpage

%% ===== Page 220 =====

Remarques. Dans la présentation précédente, le groupe $\Pi_0$ lui-même semble pris en second plan, au profit de $\mathfrak{T}_g$ et de la structure-LT.

Il me semble en fait qu'importent les revêtements de $\Pi_0$ (s'il en existe) qui seraient engendrés par les $\mathfrak{T}_{g, \nu}(\Pi_0)$ et leurs conjugués, soit $\Pi_0'$, et les sous-groupes de lacets exactement. Il est très probable, si un tel $\Pi_0'$ existe, qu'il soit unique (prenons le cas $I = \emptyset$ !). Le cas qui m'avait occupé jusqu'à présent, $\Pi_0 = \mathfrak{T}_{0,3} = \mathfrak{T}_{0,3}^+$, a été un peu "misleading", car -- puis alors $\Pi_0 = \mathfrak{T}_{0,3}$ justement. Si j'étais parti -- dans la vie, de ce fait -- prenant moralement $\Pi_0 = \Pi_{0,3}$, comme le groupe adjoint -- aurait un (comme alors d'ailleurs il est en effet, et) le rôle des $\Pi_{0,3}$ s'amenuise au profit des sous-groupes plus petits (type $\mathfrak{T}_g$, cf. § VI ...).

\newpage

%% ===== Page 221 =====

40) Approximation de $\mathfrak{T}_0(\pi_0, \mathfrak{G}_0, U)$ induite par une approximation de $\mathfrak{G}_0$

Partons d'un système $(\mathfrak{G}_0, \pi_0)$ d'un $LT$-groupe $\mathfrak{G}_0$ et d'un sous-groupe discret\marginpar{Je suppose $\pi_0$ discret. $\pi_0 \hookrightarrow \mathfrak{G}_0$ défini par $\mathfrak{G}_0$ [unclear: l'image] $\pi_0$ [unclear: est] défini.} $\pi_0$ de $\mathfrak{G}_0$. Supposons donné une donnée $\mathfrak{G}_{\mathbf{Q}}$ sur $\mathbf{Q}$, et un homomorphisme
$$
\varphi: \mathfrak{G}_0 \to \mathfrak{G}(\mathbf{Q}) \leqno{(65)}
$$
ayant la propriété que pour $H$ [unclear: un sous-groupe] 
$$
k_i: \mathfrak{T}_0 \to \mathfrak{G}_0 \leqno{(66)}
$$
qui soit partie de la $LT$-structure, l'on complète $\mathfrak{T}_0 \to \mathfrak{G}_0 \to \mathfrak{G}(\mathbf{Q})$ qui soit quasi-surjectif. On suppose de plus que $\mathfrak{G}_{\mathbf{Q}}$ soit engendré par l'image de certains $\mathbf{Q}$-éléments définis par les éléments de $\mathfrak{G}_0$ (65).

On considère $\mathfrak{G}_{\mathbf{Q}}$ comme muni de l'homomorphisme (65) et de la famille des homomorphismes extérieurs
$$
k_i: \mathfrak{T}_{\mathbf{Q}} \to \mathfrak{G}_{\mathbf{Q}} \leqno{(67)}
$$

\newpage

%% ===== Page 222 =====

donnant lieu à des compatibilités

(68)
\begin{tikzcd}
T_0 \ar[r, "K_i"] \ar[d, "can"'] & \underline{\mathfrak{S}}_0 \ar[d, "(65)"] \\
T_{0\,Q}(Q) \ar[r, "\underline{K}_i(Q)"'] & \underline{\mathfrak{S}}(Q) \ .
\end{tikzcd}

Je considère alors le schéma en groupes [MARGIN: extérieurs (resp. à un...)] sur $Q$ des (automorphismes de $\underline{\mathfrak{S}}_Q$, formé des $\mu \in \underline{Aut}(\underline{\mathfrak{S}}_Q)$, tels que pour tt $i \in I$, il existe loc. une section $f_i$ de $\underline{\Pi}_Q$ (engendré par l'image de $\Pi_0$ dans $\underline{\mathfrak{S}}(Q)$), telle que l'on ait

(69) $$ \mu(\underline{K}_i(\lambda)) = int(f_i)^{-1}(\mu \lambda) \ . $$

Je désigne par $\underline{M}_Q^\varphi$ ce schéma en groupes sur $Q$, s'envoyant vers $\underline{G}_Q$ par $(u,\mu) \mapsto \mu$, par un hom. qui est trivial sur $\operatorname{Im}(\underline{\mathfrak{S}}_Q \to \underline{M}_Q^\varphi)$, d'où se factorise par un hom.

(70) $$ \underline{\Gamma}_Q^\varphi \overset{df}{=} \underline{M}_Q^\varphi / Im \underline{\mathfrak{S}}_Q \xrightarrow{\lambda^\varphi} \underline{G}_Q \ . $$

Si maintenant $(\underline{\mathfrak{S}}_0, \Pi_0)$ fait partie d'un système de [unclear]-Teichmüller

\newpage

%% ===== Page 223 =====

$(\Pi_0 \subset \mathbb{G}_0 \subset \mathcal{M})$, et si $\mathbb{G}_0 \xrightarrow{\varphi} \underline{\mathbb{G}}^\varphi_\mathbb{Q}$ provient d'un
hom de Lie de ce système, $\mathcal{M} \to M(\mathbb{A}_\mathbb{Q})$
par la méthode vue la section précédente,
alors on trouve un hom. naturel
via op. de $\underline{M}_\mathbb{Q}$ sur $\underline{\mathbb{G}}_\mathbb{Q}$, et $\lambda_M$
(71) $$ \underline{M}_\mathbb{Q} \xrightarrow{\varphi} \underline{M}^\varphi_\mathbb{Q} $$
rendant commutatif le diagramme

(72)
\begin{tikzcd}
\mathbb{G}_0 \ar[r] \ar[d] & \mathcal{M} \ar[d] & \underline{\mathbb{G}}_\mathbb{Q} \ar[r] \ar[dr, "via\ \text{can.}\ int."'] & \underline{M}_\mathbb{Q} \ar[d, "\varphi"] \\
\underline{\mathbb{G}}(\mathbb{Q}) \ar[r] \ar[d, "int."'] & M(\mathbb{A}_\mathbb{Q}) \ar[d] & & \underline{M}^\varphi_\mathbb{Q} \\
\underline{M}^\varphi_\mathbb{Q}(\mathbb{Q}) \ar[r, hook] & \underline{M}^\varphi_\mathbb{Q}(\mathbb{A}_\mathbb{Q}) & &
\end{tikzcd}

Si le centre de $\underline{\mathbb{G}}_\mathbb{Q}$ est trivial (ce qui est
une hyp. "crédible" surtout si le centre de
$\mathbb{G}_0$ l'est ...) lors $\underline{\mathbb{G}}_\mathbb{Q} \to \underline{M}^\varphi_\mathbb{Q}$ est
injectif, et $\underline{\mathbb{G}}_\mathbb{Q}$ s'interprète comme
"le $\underline{\mathbb{G}}_\mathbb{Q}$" "procréé" par un hom. de
Lie
$$ (\Pi \subset \mathbb{G} \subset \mathcal{M}) \longrightarrow (\Pi \subset \underline{\mathbb{G}}_\mathbb{Q} \subset \underline{M}^\varphi_\mathbb{Q}) , $$
qui a un certain caractère universel
p.r. à l'hom. donné (65) $\varphi : \mathbb{G}_0 \to \underline{\mathbb{G}}(\mathbb{Q})$.

\newpage

%% ===== Page 224 =====

Quand on ne suppose pas d'avance l'exis-
tence d'un tel hom. de Lie dans lequel
s'insère $\underline{\mathfrak{S}}_\mathbb{Q}$ , il n'est pas automatique
qu'on puisse trouver un hom.

(73) $$\varphi_M : \mathcal{M} \longrightarrow \underline{\mathcal{M}}_\mathbb{Q}^\varphi (A_\mathbb{Q})$$

qui rende compte des opérations de
$\mathcal{M}$ sur $\underline{\mathfrak{S}}_\mathbb{Q}$ , par la commutativité
d'un diagramme , pour tout $g \in \mathcal{M}$

(74)
\begin{tikzcd}
\underline{\mathfrak{S}} \ar[r, "int(g)|\underline{\mathfrak{S}}"] \ar[d, "\hat{\varphi}"'] & \underline{\mathfrak{S}} \ar[d, "\hat{\varphi}"] \\
\underline{\mathfrak{S}}(A_\mathbb{Q}) \ar[r, "int(\varphi_M(g))|\underline{\mathfrak{S}}(A_\mathbb{Q})"'] & \underline{\mathfrak{S}}(A_\mathbb{Q})
\end{tikzcd}

On s'intéresse seulement aux $\varphi$ (65)
pour lesquels (73) existe - au moins
quitte à se restreindre à un sous-groupe
ouvert de $\mathcal{M}$...

\newpage

%% ===== Page 225 =====

\section*{53}

1$^\circ$) Action \emph{différentiable} du groupe de Teichmüller spécial.

Soit $X$ une surface connexe, orientable, compacte dès le cas qui nous intéresse, $I$ une partie discrète de $X$. Pour $H \in I$, soit $\mathfrak{S}_i$ l'espace (homéomorphismes à bord...) des demi-droites tangentes de l'espace tangent, et soit $\widetilde{\mathfrak{S}}_i$ un rev. universel de $\mathfrak{S}_i$.

Soit $\phi \in \\operatorname{Aut}_{\text{diff}}(X, S)$, il s'induit un autom. de
$$
\mathfrak{S} \defeq \coprod_{i \in I} \mathfrak{S}_i \leqno{(1)}
$$
noté $\phi_{\mathfrak{S}}$. Considérons l'un des relèvements $\widetilde{\phi}$ de $\phi_{\mathfrak{S}}$ au-dessus de
$$
\widetilde{\mathfrak{S}} = \coprod_{i \in I} \widetilde{\mathfrak{S}}_i \leqno{(2)}
$$
Notons $T_0$ le module d'oscillation de $\mathfrak{S}$, que chaque $T_0$ opère sur chaque $\widetilde{\mathfrak{S}}_i$.

\newpage

%% ===== Page 226 =====

de sorte que $T_0$ opère en fait sur $\tilde{\mathcal{S}}$
$$
\tilde{\mathcal{S}}/T_0 \isom \mathcal{S} \leqno{(3)}
$$
Ceci montre que (hi. et ait. comme ci-dessus) le groupe $T_0^I$ opère à gauche (p. ex.) sur l'un des relèvements $\tilde{u}$ que définit aussi un foncteur vers $T_0^I$.

Soit $G$ un groupe topologique, qui opère effectivement sur $(X, S)$, avec des actions continues sur $\mathcal{S}$. On pose
$$
\tilde{G} = \{ (g, \tilde{u}) \mid u \in G, \tilde{u} \text{ un relèvement continu de } u \text{ sur } \tilde{\mathcal{S}} \} \leqno{(4)}
$$
\marginpar{NB $G$ opère sur $T_0^I$ via la condition d'orientation}
on trouve sur $\tilde{G}$ une topologie (considérée dans et comme un sous-groupe de $\tilde{\mathfrak{T}}$), de telle sorte que l'on ait une extension de $\tilde{G}$ par le groupe $T_0^I$:
$$
1 \to T_0^I \to \tilde{G} \to G \to 1 \leqno{(6)}
$$

Proposition. Supposons que $X$ soit compact,

\newpage

%% ===== Page 227 =====

761

et considérons l'hom. canonique
$$ G \longrightarrow G/G^\circ \xrightarrow{\sim} \mathcal{T}(X,S) \quad \text{(groupe de Teichmüller)} $$
on a alors un hom.
canonique d'extensions

[MARGIN: Pour la démonstration de la dite proposition voir [unclear] (p. 755 ff, p. 764 ff)]

(7)
\begin{tikzcd}
1 \ar[r] & T_0^I \ar[r] \ar[d] & \tilde{G} \ar[r] \ar[d] & G \ar[r] \ar[d] & 1 \\
1 \ar[r] & T_0^I \ar[r] & Sp\mathcal{T}(X,S) \ar[r] & \mathcal{T}(X,S) \ar[r] & 1 \ ,
\end{tikzcd}

$Sp\mathcal{T}(X,S)$ est le groupe de Teichmüller
spécial de $(X,S)$.

NB On notera que si pour tout $i \in I$ on
fixe une origine $e_i$ sur $S_i$, il
résulte de la propriété universelle de
$\tilde{S}$ que la donnée d'un relèvement
$\tilde{u}$ de $u \in G$ à $\tilde{S}$ revient à la donnée,
pour tout $i \in I$, d'un relèvement de
$u(e_i)$ en un point $\tilde{a}_i$ de $\tilde{S}_{\sigma(i)}$ (i.e., de $\tilde{S}_{u(i)}$).

\newpage

%% ===== Page 228 =====

2°) Le groupe $\operatorname{Gl}(2, \mathbf{R})^{\pm 1}$ : On pose
$$
\operatorname{Gl}(2, \mathbf{R})^{\pm 1} \defeq \{u \in \operatorname{Gl}(2, \mathbf{R}) \mid \det u \in \{\pm 1\}\} \leqno(8)
$$
Ce groupe, qui est identifiable au groupe des transformations affines sur $\mathbf{C} \cong \mathbf{R}^2$ qui laissent fixée l'origine, on peut lui appliquer la construction des $\mathfrak{T}$, dont $\mathfrak{T}_0$ s'identifie au cercle unité
$$
\mathcal{U} \defeq \{z \in \mathbf{C} \mid |z|=1\} \longrightarrow S^1_0
$$
$$
z \longmapsto z^2
$$
et pour $u \in \operatorname{Gl}(2, \mathbf{R})^{\pm 1}$, son action sur $S^1_0$ s'identifie à l'action sur $\mathcal{U}$ donnée par
$$
z \longmapsto \frac{u(z)}{|z|} \leqno(10)
$$
Une action de cette nature sur $S^1_0$ s'identifie à une action de $S^1_0$ sur $\mathcal{U}$ (au moyen de $z \mapsto z^2=1$ au niveau des $\mathcal{U}$-constructions\dots) choisissons un point base $x_0$ de $S^1_0$, correspondant à $\dots$

\newpage

%% ===== Page 229 =====

point $1$ de $\mathcal{U}$, de sorte que le revêtement universel de $\mathbb{S}^1$ partant au-dessus de $1$ s'identifie au revêtement exponentiel
$$
\mathbf{R} \xrightarrow{\exp} \mathbb{S}^1 \quad (\simeq \mathbb{S}^1) \leqno{(11)}
$$
$$
t \mapsto \exp(2i\pi t)
$$
Ainsi, un relèvement de l'action de $G$ se traduit par le choix de $t \in \mathbf{R}$ tel que
$$
\exp(2i\pi t) = \frac{u(1)}{|u(1)|} \leqno{(12)}
$$
i.e. le choix d'un argument $t$ de $u(1)$.

On applique la construction de (9) au cas $G = \operatorname{Gl}(2, \mathbf{R})^{\pm 1}$, et donne une structure d'extension
$$
1 \to \mathbf{Z} \to \operatorname{Gl}(2, \mathbf{R})^{\sim} \to \operatorname{Gl}(2, \mathbf{R})^{\pm 1} \to 1 \leqno{(13)}
$$
qui induit au-dessus de la composante connexe $\operatorname{Sl}(2, \mathbf{R})$ de $\operatorname{Gl}(2, \mathbf{R})^{\pm 1}$ une structure d'extension
$$
1 \to \mathbf{Z} \to \operatorname{Sl}(2, \mathbf{R})^{\sim} \to \operatorname{Sl}(2, \mathbf{R}) \to 1 \leqno{(14)}
$$
Proposition Le groupe $\operatorname{Sl}(2, \mathbf{R})^{\sim}$ est connexe et un revêtement universel : $\mathbf{R}^3$.

\newpage

%% ===== Page 230 =====

Dim. Comme l'inclusion
$$
\mathcal{U} \hookrightarrow \\operatorname{Sl}(2, \mathbf{R}) \leqno{(15)}
$$
est une équivalence d'homotopie, qui s'identifie à l'inclusion $\mathcal{U} \hookrightarrow \mathcal{U} \times \mathbf{R}^2$, et pourvu que le revêtement $\mathcal{U} \to \\operatorname{Sl}(2, \mathbf{R})^{\sim}$ s'identifie au revêtement universel (de base $1$ ou $0 \dots$) Or si $u \in \mathcal{U}$, le dernier choix d'un $\tilde{u}$ au-dessus de $u$, i.e. d'un relèvement de $u(1)/|u(1)|$, en $\mathbf{R}$ par la formule (12), s'identifie (puisque $u(1)/|u(1)| = u(1) = u$) au choix d'un $t \in \mathbf{R}$ tel que $\exp(2i\pi t) = u$. O.K.

Cor. 1 Le sous-groupe d'un sous-groupe discret de $\\operatorname{Sl}(2, \mathbf{R})^{\sim}$ est isomorphe à $\mathbf{Z}$ ($\exists$ générateur $w_0$ tel que $w_0^n = w_0$, i.e. $w_0^{n-1} = 1$ dans $\\operatorname{Ker}(\\operatorname{Sl}(2, \mathbf{R})^{\sim} \to \\operatorname{Sl}(2, \mathbf{R}))$)\footnote{Expliquez à ce sujet que dans les $\\operatorname{Sl}(2, \mathbf{R})^{\sim}$ on tend à confondre $w_0$ et $w_0^{-1}$}.

Cor. 2 L'action des $g \in \mathcal{GL}(2, \mathbf{R})^{\pm 1}$ sur $\\operatorname{Sl}(2, \mathbf{R})^{\sim}$ (ou sur le centre $w_0$ de $\\operatorname{Sl}(2, \mathbf{R})^{\sim}$) va via $w_0 \mapsto w_0^{-1}$ (resp. $w_0 \mapsto w_0^{-1}$) $\dots$

\newpage

%% ===== Page 231 =====

\section*{3.) Le groupe $\\operatorname{Gl}(2, \mathbf{Z})^\sim$.}

Soit $\\operatorname{Gl}(2, \mathbf{Z})^\sim$ l'image inverse de $\\operatorname{Gl}(2, \mathbf{Z})$ dans $\\operatorname{Gl}(2, \mathbf{R})^\sim$, de sorte que nous avons une structure d'extension :

$$
1 \to \mathbf{Z} \xrightarrow{n \mapsto \omega^n} \\operatorname{Gl}(2, \mathbf{Z})^\sim \to \\operatorname{Gl}(2, \mathbf{Z}) \to 1 \leqno{(16)}
$$

Je voudrais décrire $\\operatorname{Gl}(2, \mathbf{Z})^\sim$ par générateurs et relations. Pour ceci, il m'est utile :

\textbf{Proposition.} L'image inverse des groupes cycliques $\mathfrak{L}_\rho \cong \mathbf{Z}/6\mathbf{Z}$ et $\mathfrak{L}_\sigma \cong \mathbf{Z}/4\mathbf{Z}$ engendrés par $\rho$ et $\sigma$ sont isomorphes à $\mathbf{Z}$.

Ceci résulte du

\textbf{Corollaire} (à la page p. 762) : Pour tout sous-groupe cyclique fini $\mathfrak{Z} \subset \\operatorname{Sl}(2, \mathbf{R})$, l'image inverse de $\mathfrak{Z}$ dans $\\operatorname{Sl}(2, \mathbf{R})^\sim$ est isomorphe à $\mathbf{Z}$.

En effet, pour un sous-groupe $\mathfrak{Z} \subset \mathcal{U}$, i.e. $\mathfrak{Z} \cong \mathcal{U}_1$, mais alors $\mathfrak{Z}^\sim$ s'identifie à l'un des $\mathfrak{L}_\theta$ où $(\exp 2i\pi t) = 1$ i.e. $\exp 2i\pi t = 1$ i.e.

\newpage

%% ===== Page 232 =====

$t \in \frac{1}{n} \mathbf{Z}$,

le groupe des centres $\simeq \mathbf{Z}$, q.f.d.

Corollaire. Il existe $\tilde{\rho} \in \\operatorname{Sl}(2, \mathbf{Z})^{\sim}$ et $\tilde{\sigma} \in \\operatorname{Sl}(2, \mathbf{Z})^{\sim}$ \marginpar{Il serait plus simple de poser des relations, tel qu'on a $\tilde{\rho}^3 = \tilde{\sigma}^2$ pour la question}
(14) \quad \quad $\tilde{\rho}^3 = \tilde{\sigma}^2$

L'un de ces systèmes $(\tilde{\rho}, \tilde{\sigma})$ est le torseur sous $\mathbf{Z}$, par l'opération
(18) \quad \quad $(\tilde{\rho}, \tilde{\sigma}) \longmapsto (\tilde{\rho} w_0^{2n}, \tilde{\sigma} w_0^{3n}) \quad ,$
où $w_0$ est le générateur canonique de $\\operatorname{Ker}(\\operatorname{Sl}(2, \mathbf{Z})^{\sim} \to \\operatorname{Sl}(2, \mathbf{Z}))$.

On peut dire\footnote{Le centre de l'élément est le groupe $\mathbf{Z}$, et le choix d'un relèvement unique générateur $w_0$ satisfaisant $w_0^2 = 2\text{id}$ (générateur canonique de $\\operatorname{Sl}(2, \mathbf{Z})^{\sim} \to \\operatorname{Sl}(2, \mathbf{Z})$)} que l'image de l'élément central est le groupe $\mathbf{Z}$, et le choix d'un unique générateur $w_0$ satisfaisant
(19) \quad \quad $w_0^2 = 2\text{id}$

et ceci étant il est possible de relever $\rho$ en $\tilde{\rho}$, $\sigma$ en $\tilde{\sigma}$ satisfaisant
(20) \quad \quad $\tilde{\rho}^3 = w_0^{-1}, \tilde{\sigma}^2 = w_0$
Et cette fois-ci, $\tilde{\rho}$ et $\tilde{\sigma}$ sont uniquement définis. \marginpar{Faire la vérification}

\newpage

%% ===== Page 233 =====

minis. Ceci résulte d'une conséquence de la prop. précédente.

Considérons l'image de $\{1, \pm 1\} \subset \\operatorname{Sl}(2, \mathbf{Z})^\sim$, cette extension de $\{1, \pm 1\} \simeq \mathbf{Z}/2\mathbf{Z}$ par $\mathbf{Z}$ — je dis que cette extension est scindée. En effet, l'action de $\tau$ sur $\mathcal{H}_0$ se relie à l'action de $\tilde{\tau}$ d'ordre 2 sur $\mathcal{H}_0 \simeq \mathbb{H}$, qui se définit par
$$
\tilde{\tau}(t) = \frac{1}{2} - t \leqno{(21)}
$$
En effet, pour $z \in \mathcal{H}, z = \exp 2i\pi t$
$$
\tau z = -\bar{z} = -\exp 2i\pi \bar{t} = \exp 2i\pi \left(\frac{1}{2} - t\right) = \exp 2i\pi \tilde{\tau}(t) \leqno{(22)}
$$
Je me suis un peu trompé, et travaillons pour $\tau' = \tau_0 = \begin{pmatrix} 0 & 1 \\ 1 & 0 \end{pmatrix}$, qui est, dis-je, le produit de $e_0 = 1$ et $e_1 = i$, i.e. la symétrie par rapport à la diagonale principale. Cet automorphisme stabilise $\mathcal{H}_0$, et fixe le point $\exp \frac{\pi i}{4} = \frac{\sqrt{2}}{2}(1+i)$, il se relie en
$$
\tilde{\tau}'(t) = \frac{1}{4} - t \leqno{(23)}
$$
Notons que l'action de $\tilde{\tau}$ \marginpar{NB $\tau = \tau_{\infty} = \begin{pmatrix} 0 & -1 \\ 1 & 0 \end{pmatrix}$ i.e. $\tau(z) = -1/z$ (symétrie p.u.)} comm...

\newpage

%% ===== Page 234 =====

\noindent d $\tilde{\tau}$) où $\tilde{\tau}$ est la symétrie (cf.\ fin du § précédent). Je dis que
$$ \tilde{\tau}(\tilde{\rho}) = \tilde{\rho}^{-1}, \quad \tilde{\tau}(\tilde{\sigma}) = \tilde{\sigma}^{-1} \leqno{(24)} $$
Les formules s'écrivent ainsi :
$$ \tilde{\tau}(\tilde{\rho}^{-1}) = \tilde{\rho}, \quad \tilde{\tau}(\tilde{\sigma}^{-1}) = \tilde{\sigma} $$

Elles sont valables dans $\\operatorname{Sl}(2, \mathbf{Z})$, donc pour voir qu'elles sont valables dans $\\operatorname{Sl}(2, \mathbf{Z})^{\sim}$, il suffit de montrer que les premiers mots sont les autres, (20), qui s'écrivent
$$ \tilde{\tau}(\tilde{\rho}^{-3}) = \tilde{\omega}_0^{-1}, \quad \tilde{\tau}(\tilde{\sigma}^{-2}) = \tilde{\omega}_0^{-1} \leqno{(20)} $$
i.e.\ $\tilde{\tau}(\tilde{\omega}_0) = \tilde{\omega}_0^{-1}$, OK.

\noindent Théorème. Le groupe $\\operatorname{Sl}(2, \mathbf{Z})^{\sim}$ admet la présentation (à deux générateurs $\tilde{\rho}, \tilde{\sigma}$)
$$ \tilde{\rho}^3 = \tilde{\sigma}^2 = 1 \leqno{(25)} $$

\noindent Démonstration. Soit $\bar{\pi}$ le groupe défini par cette présentation, et considérons l'homomorphisme
$$ \bar{\pi} \to \\operatorname{Sl}(2, \mathbf{Z}) / \pm 1 = \{ \tilde{\rho}, \tilde{\sigma} \mid \tilde{\rho}^3, \tilde{\sigma}^2 = 1 \} \leqno{(26)} $$
$$ (\tilde{\rho}, \tilde{\sigma}) \mapsto (\rho, \sigma) \pmod{\pm 1} $$

\newpage

%% ===== Page 235 =====

a $\tilde{\tau}$) en ce qui est la symétrie $\omega_0$ et ceci fin de $\S$ précédent). Je dis que :
$$ \tilde{\tau}(\tilde{\rho}) = \tilde{\rho}^{-1}, \tilde{\tau}(\tilde{\sigma}) = \tilde{\sigma}^{-1} \leqno{(24)} $$
Les formules s'écrivent ainsi :
$$ \tilde{\tau}(\tilde{\rho}^{-1}) = \tilde{\rho}, \tilde{\tau}(\tilde{\sigma}^{-1}) = \tilde{\sigma} $$
Elles sont valables dans $\\operatorname{Sl}(2, \mathbf{Z})$, donc pour ceux qui s'intéressent aux tables de $\\operatorname{Sl}(2, \mathbf{Z})$, il suffit de montrer que les premières sont bien celles caractéristiques :
$$ \tilde{\tau}(\tilde{\rho}^{-3}) = \omega_0^{-1}, \tilde{\tau}(\tilde{\sigma}^{-2}) = \omega_0^{-1} \leqno{(25)} $$
i.e. $\tilde{\tau}(\omega_0) = \omega_0^{-1}$, ok.

Théorème Le groupe $\\operatorname{Sl}(2, \mathbf{Z})^\sim$ admet la présentation à deux générateurs $\tilde{\rho}, \tilde{\sigma}$
$$ \tilde{\rho}^3 = \tilde{\sigma}^2 = 1 \leqno{(26)} $$
Dém. Soit $\Pi$ le groupe défini par cette présentation, et considérons l'homomorphisme
$$ \Pi \to \\operatorname{Sl}(2, \mathbf{Z})/\pm 1 \leqno{(27)} $$
$$ \tilde{\rho}, \tilde{\sigma} \mapsto (\rho, \sigma) $$

Le groupe $\\operatorname{Sl}(2, \mathbf{Z})/\pm 1$ s'identifie au quotient de $\Pi$ par le sous-groupe distingué engendré par l'élément $\omega_0 = \tilde{\rho}^3 = \tilde{\sigma}^2$. Mais cet élément $\omega_0$ est déjà distingué. Dans une suite exacte
$$ \mathbf{Z} \xrightarrow{\omega_0} \Pi \to \\operatorname{Sl}(2, \mathbf{Z})/\pm 1 \to 1 \leqno{(28)} $$
qui s'envoie dans la suite exacte
$$ 1 \to \mathbf{Z} \to \\operatorname{Sl}(2, \mathbf{Z})^\sim \to \\operatorname{Sl}(2, \mathbf{Z})/\pm 1 \to 1 $$
en induisant des isomorphismes sur les facteurs extérieurs — donc par le lemme des 5, $\Pi \to \\operatorname{Sl}(2, \mathbf{Z})^\sim$ est iso, cqfd.

Corollaire ! Le groupe $\\operatorname{Sl}(2, \mathbf{Z})^\sim$ admet la présentation à trois générateurs $\tilde{\rho}, \tilde{\sigma}, \tilde{\tau}$
$$ \tilde{\tau}^2 = \omega_0, \tilde{\rho}^3 = \tilde{\sigma}^2 = 1, \tilde{\tau}\tilde{\rho} = \tilde{\rho}^{-1}, \tilde{\tau}\tilde{\sigma} = \tilde{\sigma}^{-1} \leqno{(29)} $$
Je vais essayer de présenter les générateurs de $\\operatorname{Sl}(2, \mathbf{Z})^\sim$ en trois générateurs explicites, soient $\tilde{\rho}, \tilde{\sigma}, \tilde{\omega}_0$ des éléments de $\\operatorname{Sl}(2, \mathbf{Z})^\sim$, en s'inspirant des générateurs de $\\operatorname{Sl}(2, \mathbf{Z})$ :

\newpage

%% ===== Page 236 =====

$$
(29) \quad \begin{cases}
\Gamma_\infty = \tau_\infty = \begin{pmatrix} -1 & 0 \\ 0 & 1 \end{pmatrix} \\
\Gamma_0 = \tau'_\infty = \begin{pmatrix} 0 & 1 \\ 1 & 0 \end{pmatrix} \\
\Gamma_1 = -\tau'_0 = \begin{pmatrix} -1 & 1 \\ 0 & 1 \end{pmatrix}
\end{cases}
\qquad
(30) \quad \begin{cases}
\Gamma_0(\rho) = \rho^{-1}, \Gamma_0(\sigma) = \sigma^{-1} \\
\Gamma_1 = \rho^{-1} \Gamma_0 = \Gamma_0 \rho \\
\Gamma_\infty = \sigma \Gamma_0 = \Gamma_0 \sigma^{-1}
\end{cases}
$$

Notons que $(\tilde{\rho}, \tilde{\tau})$ engendre un groupe
isomorphe à $D_\infty$, que je note $D_{\tilde{\rho}}$,
et que $(\tilde{\sigma}, \tilde{\tau})$ engendre un groupe $D_{\tilde{\sigma}}$,
isomorphe à $D_\infty$. Enfin $(\omega_0, \tilde{\tau})$ enge-
ndre un groupe $D_{\omega_0}$, isomorphe encore
à $D_\infty$. On a un diagramme

$$
(31) \quad 
\begin{tikzcd}
& D_{\omega_0} \ar[dl, "i"'] \ar[dr, "j"] & \\
D_{\tilde{\rho}} & & D_{\tilde{\sigma}}
\end{tikzcd}
\qquad
\begin{array}{l}
\begin{cases} i(\omega_0) = \tilde{\rho}^3 \\ i(\tilde{\tau}) = \tilde{\tau} \end{cases} \\
\begin{cases} j(\omega_0) = \tilde{\sigma}^2 \\ j(\tilde{\tau}) = \tilde{\tau} \end{cases}
\end{array}
$$

et le cor. 1 s'écrit sous la forme

$$
(32) \quad \operatorname{Gl}(2, \mathbb{Z})^\sim \simeq D_{\tilde{\rho}} *_{D_{\omega_0}} D_{\tilde{\sigma}}
$$

(somme amalgamée de groupes).

Or on peut prendre dans $D_{\tilde{\rho}}$, $D_{\tilde{\sigma}}$ les
générateurs involutifs

$$
(33) \quad \begin{cases}
\tilde{\Gamma}_0 = \tilde{\tau} = \tilde{\tau}'_\infty \\
\tilde{\Gamma}_\rho = .......
\end{cases}
\text{et} \quad
\begin{matrix}
\tilde{\Gamma}_1 \overset{df}{=} \tilde{\rho}^{-1} \tilde{\Gamma}_0 = \tilde{\Gamma}_0 \tilde{\rho} \text{ dans } D_{\tilde{\rho}} \\
\text{et } \tilde{\Gamma}_\infty \overset{df}{=} \tilde{\sigma} \tilde{\Gamma}_0 = \tilde{\Gamma}_0 \tilde{\sigma}^{-1}
\end{matrix}
$$

\newpage

%% ===== Page 237 =====

et on trouve alors \marginpar{755}

Corollaire 2. $\mathcal{G}(\mathbf{Z})^{\sim}$ admet une présentation à trois générateurs involutifs $\tilde{r}_0, \tilde{r}_1, \tilde{r}_{\infty}$ :

$$
\boxed{
\begin{array}{l}
\tilde{r}_0^2 = \tilde{r}_1^2 = \tilde{r}_{\infty}^2 = 1, \\
\underbrace{(\tilde{r}_0 \tilde{r}_1)^3}_{\tilde{\rho}} \cdot \underbrace{(\tilde{r}_{\infty} \tilde{r}_0)^2}_{\tilde{\sigma}} = 1
\end{array}
} \leqno{(34)}
$$

\newpage

%% ===== Page 238 =====

4°) L'isomorphisme $\operatorname{Gl}(2, \mathbf{Z}) \isom \mathfrak{S}_{2,1}$

Posons
$$
X_1 = \mathbf{R}^2 / \mathbf{Z}^2 = (\mathbf{R}/\mathbf{Z})^2 \isom \mathcal{U} \times \mathcal{U} \leqno{(35)}
$$
(surface de type de genre 1), notons additivement sa loi de composition, et soit
$$
a_1 = 0 \in X_1, \quad X_1 \setminus \{a_1\} = X_1 \setminus \{0\} = \mathcal{U}_{1,1} \leqno{(36)}
$$
surface de type et type $(g, \nu) = (1, 1)$.
Le groupe $\operatorname{Gl}(2, \mathbf{Z})$ s'identifie au groupe des automorphismes du groupe topologique sur $X_1$, et opère dans sur $X_1$ en laissant fixe l'origine $a_1=0$, d'où il opère sur $\mathcal{U}_{1,1}$, d'où un homomorphisme
$$
\operatorname{Gl}(2, \mathbf{Z}) \isom \mathcal{G}_{1,1} = \mathfrak{T}(M_{1,1}) \leqno{(37)}
$$
\marginpar{groupe des classes d'isotopie d'homéomorphismes de $\mathcal{U}_{1,1}$}

dont on sait que c'est un isomorphisme. D'ailleurs, $\mathcal{G} = \operatorname{Gl}(2, \mathbf{Z})$ qui définit fidèlement sur $X_1$, et l'action sur l'espace tangent s'identifie à l'action topologique de $\operatorname{Gl}(2, \mathbf{Z}) \subset \operatorname{Gl}(2, \mathbf{R})$ sur $\mathbf{R}^2$. Donc le groupe noté $\widehat{\mathcal{G}}$ dans

\newpage

%% ===== Page 239 =====

19) $\mathcal{G}$ est canoniquement isomorphe au groupe d'extensions de $G$ par $T_0 = \mathbf{Z} \oplus \mathbf{Z}$ ($\mathbf{R}^2$ et sa structure $X_1$ étant canoniquement orientés) au groupe $\operatorname{Gl}(2, \mathbf{Z}) \ltimes \mathbf{Z}^2$.

\smallskip
Théorème. L'image inverse par l'isomorphisme (37) de l'extension $\mathfrak{S}_{1,1}$ de $\pi_{1,1}$ par $\mathbf{Z}$ s'identifie canoniquement à $\operatorname{Gl}(2, \mathbf{Z}) \ltimes \mathbf{Z}^2$ dont il admet les présentations explicites :
\begin{enumerate}
\item[a)] $\tilde{\sigma}_0^3 = 1, \tilde{\sigma}'(\tilde{\rho}) = \tilde{\rho}^{-1}, \tilde{\sigma}'(\tilde{\tau}) = \tilde{\tau}^{-1}, \tilde{\tau}^2 = 1$
\item[b)] $\tilde{\rho}_0^2 = \tilde{\rho}_1^2 = \tilde{\rho}_{\infty}^2 = 1, (\tilde{\rho}_0 \tilde{\rho}_1)^3 (\tilde{\rho}_0 \tilde{\rho}_{\infty})^2 = 1$
\end{enumerate}

Interprétons maintenant $\mathfrak{T}_{1,1}$ et $\mathfrak{S}_{1,1}$ en termes du groupe fondamental de $U_{1,1}$, en choisissant (judicieusement ?) un point base $a_{1,1}$ de $U_{1,1}$ :
$$
\pi_{1,1} = \pi_1(U_{1,1}, a_{1,1}) \leqno{(38)}
$$
$$
\mathfrak{T}_{1,1} = \Autext(\pi_{1,1}) \leqno{(39)}
$$

\newpage

%% ===== Page 240 =====

\marginpar{justification de la relation de la désignation de $\operatorname{Gl}(2, \mathbf{Z})$!}
Nous avons une présentation de $\Pi_{1,1}$ par
$$
\Pi_{1,1} = \{x, y, l_0 \mid [x,y]l_0 = 1\} \leqno{(40)}
$$
et désignant par $L$ le sous-groupe à lacets de $\Pi_{1,1}$ engendré par $l_0$, on aura
$$
\mathfrak{S}_{1,1} \simeq \operatorname{Norm}_{\operatorname{Autext}(\Pi_{1,1})}(L). \leqno{(41)}
$$
Donc il faudrait expliciter l'action
$$
\mathfrak{S}_{1,1} \simeq \\operatorname{Sl}(2, \hat{\mathbf{Z}}) \to \operatorname{Autext}(\Pi_{1,1}),
$$
par l'action des générateurs $\tilde{\sigma}, \tilde{\tau}$ --- action qui doit centraliser $L$, et de façon plus précise, transformer $l_0$ en $l_0^{-1}$ (pour $\tilde{\sigma}$) ou en $l_0^{-1}$ (pour $\tilde{\tau}, \tilde{\sigma}$). Il faut donc expliciter l'action des générateurs sur $x$ et $y$.

Soit $U_{1,1}$ le groupe commun. Considérons le revêtement principal $U'_{1,1}$ de groupe $\mathbf{Z}^2$ de $U_{1,1} = X \setminus \{0\}$, à savoir le groupe principal $\mathbb{R}^2$ de $X = \mathbb{R}^2 / \mathbb{Z}^2$ --- l'image inverse de $U_{1,1}$ est $\mathbb{R}^2 \setminus \mathbb{Z}^2$.
Choisissons un point base...

\newpage

%% ===== Page 241 =====

[MARGIN: 268 crossed out]

$a \in \mathbb{R}^2 \setminus \mathbb{Z}^2$,
au dessus de $a_{1,1}$ , on a donc un iso. can.

(42) $$ \Pi_{1,1} \simeq \Pi_1(\mathbb{R}^2 \setminus \mathbb{Z}^2, \mathbb{Z}^2; a) $$

d'où une suite exacte

(43)
\begin{tikzcd}
1 \ar[r] & \Pi_1(\mathbb{R}^2 \setminus \mathbb{Z}^2) \ar[r] & \Pi_{1,1} \ar[r] \ar[d, equal] & \mathbb{Z}^2 \ar[r] \ar[d, "\simeq"] & 1 \\
& & \Pi_1(U_{1,1}=X_1 \setminus \{a_{1,1}\}) & \Pi_1(X_1) \ar[d, "\simeq"] \\
& & & (\Pi_{1,1})_{ab}
\end{tikzcd}

Les éléments du $\Pi_1$ mixte s'explicitent par le calcul de Jean Malgoire à l'aide de couples

(44) $$ (\gamma, l) \quad (\gamma \in \mathbb{Z}^2, l \text{ classe de chemin de } \mathbb{R}^2 \setminus \mathbb{Z}^2 \text{ de } a \text{ à } a+\gamma) $$

qu'on "compose comme les fonctions".

Nous prendrons, au dessus des générateurs $e_0, e_1$ de $\mathbb{Z}^2$, les éléments $x, y$ définis par les chemins rectilignes de $a$ à $a+e_0$, resp. $a+e_1$, à partir de l'origine

(45) $$ a = \tilde{a}_{1,1} = (+\frac{1}{2}, +\frac{1}{2}) $$

Alors l'élément $xyx^{-1}y^{-1}$ est figuré par le chemin ci-dessus, qui est visiblement opposé au lacet "rectiligne" direct de $a$ autour de l'origine $0$.

\newpage

%% ===== Page 242 =====

5°) Digression sur les germes d'espaces.

\subsection*{A) Revêtements de germes}
\label{sec:rev-germes}

Un « germe d'espace »\footnote{ici on définit un filtre $\Phi$ pour des germes d'espaces} est défini par un espace topologique $S$ et une base de filtre $\Phi$ sur $S$. Un morphisme de $(S, \Phi)$ dans $(S', \Phi')$ est un élément de
$$
\operatorname{Hom}_{\Phi, \Phi'}(S, S') \leqno{(46)}
$$
i.e.\ désignant pour $S_i \in \Phi$, $\Phi_i$ le filtre induit par $\Phi$ sur $S_i$, et $\operatorname{Hom}_{\Phi_i, \Phi'}$ désignant l'ensemble des applications continues bien sûr compatibles avec $\Phi_S, \Phi'$, i.e.\ telles que $\forall S'_j \in \Phi'$, $f^{-1}(S'_j) \in \Phi_S$. La composition des morphismes de germes est évidente.

Un revêtement étalé d'un germe $(S, \Phi)$ est un objet de la catégorie
$$
\operatorname{Rev}(S, \Phi) \defeq \varinjlim_{S_i \in \Phi} \operatorname{Rev}(S_i) \leqno{(47)}
$$
ces revêtements forment une catégorie (sans doute connexe des « bons » revts) que nous pouvons définir comme la catégorie des revêtements d'un type convenable (savoir le type limite...)

\newpage

%% ===== Page 243 =====

\section*{5. --- DIGRESSION SUR LES GERMES D'ESPACES}
\addcontentsline{toc}{section}{5. Digression sur les germes d'espaces}

Un ``germe d'espace'' est défini par un espace $S$ et une base de filtre $\Phi$ sur $S$, un morphisme de $(S, \Phi)$ dans $(S', \Phi')$
$$
\varinjlim_{S_i \in \Phi} \\operatorname{Hom}(S_i, S') \leqno{(46)}
$$
où pour $S_i \in \Phi$, $\Phi_i$ est le filtre induit par $\Phi$ sur $S_i$, et $\\operatorname{Hom}_{\Phi_i, \Phi'}$ désigne l'ensemble des applications continues et finies compatibles avec $\Phi_i, \Phi'$, i.e. telles que $\forall S'_j \in \Phi'$, $f^{-1}(S'_j) \in \Phi_i$. La composition des morphismes de germes est évidente.

Les revêtements (étales) d'un germe $(S, \Phi)$ de base topologique, un objet de la catégorie
$$
\underline{\operatorname{Rev}}(S, \Phi) = \varprojlim_{S_i \in \Phi} \operatorname{Rev}(S_i) \leqno{(47)}
$$
les revêtements forment une catégorie (l'analogue des ``bornés''), que nous pouvons dire finie (car le type de revêtements est conservable, à savoir le type limite, projection des topos étales des $S_i \in \Phi$).

\marginpar{en fait par Olivier Levy ?}
Cette catégorie, ``dans les bons cas'', a tendance à être multigaloisienne, voire galoisienne.

On dit qu'un germe d'espace est 1-connexe, $\pi_1$-connexe, si le foncteur fibre
$$
(Ens) \to \underline{\operatorname{Rev}}(S, \Phi) \leqno{(48)}
$$
est une équivalence de catégories. Cela signifie que:

\begin{enumerate}
    \item[1)] Ce foncteur est pleinement fidèle (le dit que $(S, \Phi) \to$ 0-connexe), ce qui signifie : $S$ loc. connexe, que $\forall S_i \in \Phi, \exists S_j \in \Phi$ tel que $(\pi_0(S_j) \to \pi_0(S_i))$ soit une application constante.\footnote{donc ceci implique que le germe est ...}

    \item[2)] Le foncteur est de plus essentiellement surjectif, ce qui signifie que pour tout $S_i \in \Phi$, $\exists S_j \in \Phi, S_j \subset S_i$, tel que pour tout revêtement $S_a$ de $S_j$, l'homomorphisme $\pi_1(S_j, a) \to \pi_1(S_i)$ soit trivial.
\end{enumerate}

Les revêtements de germes d'espaces forment une catégorie fibrée sur elle.

\newpage

%% ===== Page 244 =====

des germes d'espaces. La catégorie des espaces topologiques peut être considérée comme une catégorie pleine de celle des germes d'espaces topologiques, en associant à un espace $S$, le germe $(S, \Phi)$, où $\Phi$ est le filtre trivial sur $S$. Ceci étend la catégorie fibrée des revêtements de germes d'espaces 1-connexes à celle des revêtements d'espaces topologiques.

Soit $\mathcal{X}$ un germe d'espace multigaloisien tel que $\underline{\text{Rev}} (\mathcal{X})$ soit une catégorie galoisienne.
Alors soit $S$ un germe d'espace 1-connexe, $S \to \mathcal{X}$ un morphisme.
Alors
$$
\underline{\text{Rev}} (\mathcal{X}) \to \underline{\text{Rev}} (S) \simeq (Ens) \leqno (49)
$$
foncteur fibre sur $\underline{\text{Rev}} (\mathcal{X})$, ce qui permet de définir
$$
\pi_1 (\mathcal{X}, S) \leqno (50)
$$
comme le $\pi_1$ relatif à ce foncteur fibre. Si $\mathcal{X}$ est 0-connexe et galoisienne, on a alors $\pi \simeq \pi_1(\mathcal{X}, S)$ une équivalence de catégories.

\newpage

%% ===== Page 245 =====

$$ \operatorname{Rev}(\mathcal{X}) \isom (Ens) \leqno{(51)} $$

Le groupe $\pi_1(\mathcal{X}, \mathcal{S})$ dépend fonctoriellement des triplets $(\mathcal{S}, \mathcal{X}, \mathcal{S} \to \mathcal{X})$ i.e. de la fibre $\mathcal{S} \to \mathcal{X}$ dans la catégorie des germes d'espace, où $\mathcal{S}$ est un germe 1-connexe.

Soit $\mathcal{X}$ un espace topologique, $Y$ une partie de $\mathcal{X}$. Le germe de germe $\mathcal{X}$ autour de $Y$ est le germe de germe des voisinages de $Y$ dans $\mathcal{X}$, germe inductif de $\mathcal{X}$ selon $\mathcal{X} \setminus Y$ le filtre (il est vide des $Y$-voisinages). Nous utiliserons cette notion dans le cas où $Y=\{x\}$, germe autour d'un point et germe voisinage de $\mathcal{X}$ en $x$. (NB le germe de $\mathcal{X}$ en $x$ est très 1-connexe, mais pas le germe voisinage). Pour les $(\mathcal{S}, \mathcal{X})$ germes d'espace pointés comme germe voisinage d'un espace $\mathcal{X}$, $\mathcal{X}=\mathcal{S} \cup \{x\}$ un point connexe...

\begin{tikzcd}
\mathcal{X} \arrow[leftarrow]{d} \\
\operatorname{Rev}(\mathcal{X}) \arrow[leftarrow]{u}
\end{tikzcd} \leqno{(52)}

\marginpar{le germe est le filtre des voisinages de $Y$ dans $\mathcal{X}$, germe inductif de $\mathcal{X}$ selon $Y$ (il est vide des $Y$-voisinages).}

\newpage

%% ===== Page 246 =====

J'introduirai en $x \in X$ un voisinage homéomorphe à $\mathbf{R}^2$ --- $X$ est une surface en $x$. Je distinguerai par $\star$ $D_{x,n}$ et $D_n^*$
\leqno{(53)}

Le groupe résiduel de $X$ en $x$ (contenu dans le groupe $D_{x,n}$ de $X$ en $x$).

Je reviendrai plus tard sur un autre... $D_n^* \sim D_n^*$ --- mais pour qu'il en soit fait une question sur la nature d'un groupe... $\pi_1(U)$.

Soit $(S, \mathfrak{G})$ un groupe d'appui. Un revêtement de $(S, \mathfrak{G})$ [unclear: provient] d'un $S'$ dit... Si $S_i \in \mathfrak{G}$, et définit un groupe d'appui $(S_i, \mathfrak{G}_{S_i})$, $\mathfrak{G}_{S_i}$ l'image inverse de $\mathfrak{G}$ dans $S_i$. Au-dessus de $S' = (S, \mathfrak{G})$, un groupe d'appui $\widetilde{\mathfrak{G}}$ isomorphe au groupe... dépend des choix de $S_i, S_i'$, mais seulement des éléments de $\pi_1(S)$. On trouve ainsi un foncteur...

\newpage

%% ===== Page 247 =====

$$
\mathcal{R}w(S) \to (\mathcal{G}roup)/S \leqno{(71)}
$$

qui est pleinement fidèle. Un groupe $S' \to S$ est dit un revêtement étale de $S$, s'il est de l'image de celle-ci. Ainsi la catégorie $\mathcal{R}w(S)$ des revêtements étales du groupe topologique $S$ équivaut à la catégorie pleine de la catégorie des groupes étalés au-dessus du groupe $S$, qui fournit des revêtements étales du groupe $S$.

Supposons que ce cat.\ $\mathcal{R}w(S)$ est multigaloisienne. Les objets 1-connexes $\to S$ appelés des rev.\ universels de $S$ $\to \mathcal{R}w(S)$ galvanisent \marginpar{Les groupes $\Pi_1(S, \tilde{S})$ sont les $\pi_1$ de la cat.\ galoisienne $\mathcal{R}w(S)$, calculés en l'objet 1-connexe $\tilde{S}$} le groupe $\Pi_1(S, \tilde{S})$ qui galvanise le groupe $S \to$ connexe. Les rev.\ universels de $S$ peuvent s'interpréter comme des groupes d'opérateurs. Soit $\tilde{S}$ un tel rev.\ universel, alors $\tilde{S} \to S$ étant un morphisme dominant, on a $\Pi_1(S, \tilde{S}) \isom \\operatorname{Aut}_{\mathcal{R}w(S)}(\tilde{S})$ (comparer avec (55)).

\newpage

%% ===== Page 248 =====

\marginpar{B) Cas des surfaces pointées}
Ceci dit, revenons à la situation (53), d'un espace topologique $X$, et d'un $x \in X$ tel que $X$ soit une surface en $x$.
Alors la catégorie des $\mathrm{Rev}(D_x^*)$ est une catégorie galoisienne, dont le groupe fondamental s'identifie canoniquement au groupe d'orientation $\mathfrak{T}_x$ de $X$ en $x$.

$$ \pi_1(D_x^*) \defeq \pi_1(\mathrm{Rev}(D_x^*)) \isom \mathfrak{T}_x \leqno{(56)} $$

Ceci provient du fait que le germe $D_x^*$ est un système projectif de disques en $x$, et les $\pi_1$ des dits disques forment un système projectif connexe...
Ainsi les revêtements universels $\widetilde{D_x^*}$ et $D_x^*$ forment un groupoïde connexe, le groupoïde fondamental $\mathfrak{T}_x$.
Choisissons donc un $\widetilde{D_x^*}$, et considérons le morphisme pointé...

\newpage

%% ===== Page 249 =====

$$ D_x^{*\sim} \longrightarrow D_x^* \longrightarrow X \setminus \{x\} $$

---

Plus généralement, soit $S$ une partie de $X$, contenant $x$ et comme $x$ isolée (i.e. une partie discrète de $X$), de sorte que
$$ D_x^* \hookrightarrow U \defeq X \setminus S $$
et considérons le composé
$$ D_x^{*\sim} \longrightarrow D_x^* \longrightarrow U. \leqno{(58)} $$
Si $\text{Rev}(U)$ est la catégorie des revêtements (locaux) $U$, on a par des conditions les groupes fondamentaux correspondants de $U$
$$ \pi_1(U, D_x^*) \leqno{(59)} $$
je vais définir un homomorphisme canonique
$$ \kappa_x: T_x \longrightarrow \pi_1(U, D_x^*), \leqno{(60)} $$
$$\approx \Aut \text{ des foncteurs fibre } \text{Rev}(U) \to \text{Rev}(D_x^*) \text{ sur } \text{Rev}(U)$$
on peut le déduire p.ex. pour
$$ \pi_1(D_x^*, D_x^*) \longrightarrow \pi_1(U, D_x^*) \leqno{(61)} $$
$$ T_x $$

\newpage

%% ===== Page 250 =====

induit par l'inclusion
$$D_x^* \hookrightarrow \mathcal{U}$$
ou encore, pour en termes d'automorph. d'un
foncteur fibre, comme l'hom. opération de
$T_x$ sur $\Gamma(D_x^{*\sim}, \mathcal{U}' \times_\mathcal{U} D_x^{*\sim})$ provenant

\begin{tikzcd}
& \mathcal{U}'_{D_x^{*\sim}} \ar[dl] \ar[d] \\
D_x^{*\sim} \ar[dr] & \mathcal{U}' \ar[d] \\
& \mathcal{U}
\end{tikzcd}

de l'opération de $T_x$ sur
$D_x^{*\sim}$. [unclear: text crossed out]

Supposons qu'on ait déjà
[MARGIN: $\pi_1(\mathcal{U},\tilde{\mathcal{U}}) = \pi$]
choisi (par ailleurs) un [unclear: foncteur fibre sur] $\mathcal{U}$,
par la donnée d' un pt base, ..., ou plus
généralement, d'un rev. universel $\tilde{\mathcal{U}}$ de $\mathcal{U}$

Pour le comparer avec $\pi_1(\mathcal{U}, D_x^{*\sim})$,
- il suffira de donner une "classe de chemins"
de $D_x^{*\sim} / \mathcal{U}$ à $\tilde{\mathcal{U}} / \mathcal{U}$ , i.e. un
isomorphisme

(62) $$ \tilde{\mathcal{U}} \xrightarrow{\sim} [unclear: Rev](\mathcal{U}, D_x^{*\sim}) $$

[MARGIN: l'image dans $\mathcal{U}$ du [unclear] de (62) donne dans $\mathcal{U}$ un point $y$ : [unclear]]

[unclear: crossed out paragraph at the bottom]
NB si $D_x^{*\sim}$ défini
par un rev. universel $\widetilde{D_x^*}$
de $D_x^*$, et si $\tilde{\mathcal{U}}$
défini par un pt base $y \in \mathcal{U}$,
alors [unclear]

\newpage

%% ===== Page 251 =====

773

( l'indétermination est un autom.
int. de $\Pi_1$, ou de $\Pi_1(\mathcal{U}, D_x^{* \sim})$ ) qui
induit un isomorphisme

(63) $$ \pi = \Pi_1(\mathcal{U}, \tilde{\mathcal{U}}) \overset{\sim}{\longleftrightarrow} \Pi_1(\mathcal{U}, D_x^{* \sim}) $$

(qui est changé par un autom. intérieur,
quand on change (62) par un él. de $\pi$).
Composant (63) et (60), on trouve donc
un hom.

(64) $$ T_n \longrightarrow \pi = \Pi_1(\mathcal{U}, \tilde{\mathcal{U}}) , $$

qui dépend du choix de (62) — et
qui, indépendamment de ce choix, définit
un hom. ext. bien déterminé —
qui n'est autre que l'hom. can.

(65) 
\begin{tikzcd}
\Pi_1(D_x^*) \ar[r] \ar[d, "\simeq"'] & \Pi_1(\mathcal{U}) \ar[d, equal] \\
T_n & \pi
\end{tikzcd}

[MARGIN: C) Relation avec
germes d'arcs (approche
originale "par pts base")]
Je vais décrire une façon commode
de fixer un rev. universel $\tilde{D}_x^*$ de $D_x^*$,
par un germe d'arc de courbe [unclear: analyt.]

\newpage

%% ===== Page 252 =====

de $x$. De façon générale, soit $I=[0,1]$, et $[I]$ le groupe résiduel de $I$ en $0$.\footnote{Plutôt? Je préfère dire : le groupe résiduel de $I$ en $0$ est un groupe 1-connexe.} Soit $X$ un espace topologique et $x \in X$, et $\pi_1(X, x)$ le groupe à lacets des courbes de $X$ en $x$. On a le morphisme
$$
[I] \to \operatorname{germes}_x(X) \defeq D^*_x \leqno{(66)}
$$
des germes $[I]$ dans le groupe résiduel de $X$ en $x$. Un tel germe est défini par une application continue d'un intervalle $[0, \epsilon]$ ($0 < \epsilon \le 1$)
$$
[0, \epsilon] \xrightarrow{\varphi} X \leqno{(67)}
$$
telles que
$$
\begin{cases} \varphi(0) = x \\ \varphi(t) \neq x \text{ si } t \in ]0, \epsilon] \end{cases} \leqno{(68)}
$$
Dans le cas où l'on s'intéresse à un système de voisinages de $x$, $\operatorname{Rev}(D^*_x)$ (revêtements abéliens de groupe $T_n$), on peut utiliser cela pour définir un foncteur fibre sur $D^*_x$
$$
\operatorname{Rev}(D^*_x) \to \operatorname{Rev}([I]) \simeq (Ens) \leqno{(69)}
$$

\newpage

%% ===== Page 253 =====

774

d'où un rev. universel
(70) $$D_x^{*\sim} \overset{\text{déf}}{=} Rev_{univ}(D_x^*, [I], \varphi) \ .$$

Bien entendu, on aura une factorisation

\begin{tikzcd}
{[I]} \ar[r] \ar[rr, bend right, "\varphi"'] & D_x^{*\sim} \ar[r] & D_x^* \ (\hookrightarrow U)
\end{tikzcd}

donc un iso can

(71) $$Rev_{univ}(U, D_x^{*\sim}) \simeq Rev_{univ}(U, [I]) \ ,$$

où $I$ est considéré comme un germe
1-connexe sur $U$ par le morphisme composé
$$[I] \longrightarrow D_x^* \hookrightarrow U \ .$$

Donc . . .

(72) $$\pi_1(U, D_x^{*\sim}) = \pi_1(U, [I]) \overset{K_x}{\longleftarrow} T_x \ .$$

[MARGIN: D) Relation avec structure différentiable]
Lorsque $X$ est muni d'une structure
différentiable au voisi-nage de $x$, on dit
qu'un germe de courbe de $X$ sortant
de $x$ est différentiable, si l'application
continue
$$I \longrightarrow X$$
qui le définit est différentiable, et de plus

\newpage

%% ===== Page 254 =====

l'application tangente : l'origine, nous
conclut que la
$$
R \to \Tangent(X) \leqno{(73)}
$$
est injective. Alors l'image de $R$ est
un semi-droit homogène des l'espaces
tangents : $X$ en $x$, qu'on appelle les
semi-tangentes du germe d'arc de courbe
envisagé en $X$.

Revenons pour un moment au cas
topologique, nous considérons
l'un des Connex(X) des germes d'arcs
en $X$ dans $X$, comme l'un des
objets d'un groupoïde (le groupoïde
fondamental local de $X$ en $x$), on prend
une morphisme de $\gamma$ en $\gamma'$
les isomorphismes de Revêtements($D^*_x, \widetilde{D^*_x}$)
$D^*_{\gamma} \leadsto D^*_{\gamma'}$
$$
\Iso(\gamma, \gamma') \defeq \Iso_{D^*}(\widetilde{D^*_{\gamma}}, \widetilde{D^*_{\gamma'}}) \leqno{(74)}
$$
Par construction
on trouve donc une équivalence de
catégories des groupoïdes connexes
$$
\Connex(X) \to \Revêtements(D^*) \leqno{(75)}
$$

\newpage

%% ===== Page 255 =====

Soit d'autre part, dans le cas où $X$ est diff.
en $x$,

(76) $$ \underline{Couronne}_{diff}(X,-) $$

la sous-catégorie pleine de $\underline{Couronne}(X,x)$
formée des germes d'arcs différentiables. Soit
$S_x$
l'espace des demi-droites de
l'espace tangent $Tang_x(X)$

(77) $$ S_x = \text{espace des demi-droites de } Tang_x(X) \simeq \mathbb{S}^1 $$

et considérons le groupoïde fondamental

(78) $$ \Pi_1(S_x) \overset{\simeq}{\longrightarrow} \underline{Revunis}(S_x) $$ .

Je vais définir une équivalence de
catégories

(79)
\begin{tikzcd}
\underline{Couronne}_{diff}(X,x) \ar[r, "\simeq"] \ar[d, "\simeq"'] & \Pi_1(S_x) \ar[d, "\simeq"] \\
\underline{Revunis}(D_x^{*\sim}) & \underline{Revunis}(S_x)
\end{tikzcd}

qui, sur les ensembles d'objets, soit
l'application associant à tout germe d'arc
différentiable en $x$ sa demi-tangente. Il
s'agit donc, si $\gamma, \gamma'$ sont deux tels
arcs, correspondants aux demi-tangentes $\Delta, \Delta'$
de définir une bijection canonique

\newpage

%% ===== Page 256 =====

(80) 
\begin{tikzcd}
Isom(\gamma, \gamma') \ar[r, "\sim"] \ar[d, equal, "\text{déf}"'] & Isom_{\Pi_1 \mathbb{S}_n}(\Delta, \Delta') \ar[d, equal, "\text{déf}"] \\
Isom(D^{*\sim}_{n,\gamma}, D^{*\sim}_{n,\gamma'}) & \text{classes de chemins dans } \mathbb{S}_n \text{ de } \Delta \text{ à } \Delta'
\end{tikzcd}

Au lieu de définir directement la bijection
(80), je vais définir plutôt un
foncteur (qui sera une équivalence de
catégories)

(81) $$ \Phi : \Pi_1(\mathbb{S}_n) \xrightarrow{(\approx)} Revouv(D^*_n) \ , $$

Soit

(82) $A_n$ ens des parties fermées (1-connexes) de $\mathbb{S}_n$
($i.e.$ des $\underline{arcs}$ de cercle de $\mathbb{S}_n$,
mais $\neq \emptyset$,)

c'est un ensemble ordonné, $\underline{donc}$
lequel définit une catégorie ordonnée $\underline{A}_n$
(les objets sont les $a \in A_n$, les flèches
sont les inclusions) -- et on sait que
le groupoïde enveloppant $Gr(\underline{A}_n)$ de cette catégorie
ordonnée est canoniquement équivalent à $\Pi_1(\mathbb{S}_n)$,

\newpage

%% ===== Page 257 =====

de façon précise, soit $\underline{Gr}(A_n)$ la catégorie fibre au-dessus de $Gr(A_n)$ (d'où l'on déduit les objets sur $A_n$) dont les objets sont les points, et les flèches sont les chemins,
$$
\underline{Gr}(A_n) \isom \Pi_1(S_n) \leqno{(83)}
$$
qui sur les objets se réduit à $\{x\} \mapsto x$, et qui, pour toute flèche $A$ d'extrémités $a, b$, donne un chemin de $Gr(A_n)$ dans $\underline{Gr}(A_n)$ :
$$
\left\{ \begin{aligned} a &\sim A \leftarrow b \\ a &\sim b \text{ dans } Gr(A_n) \end{aligned} \right. \leqno{(84)}
$$
associe : à toute flèche $a \sim b$ de $Gr(A_n)$, la flèche de $\Pi_1(S_n)$, définie par l'arc $A$ considéré comme chemin de $a$ à $b$.

Ainsi, pour la groupoïde $g$, le demi-foncteur
$$
\Psi: \Pi_1(S_n) \to g \leqno{(85)}
$$
équivalent à $g$ demi-foncteur.

\newpage

%% ===== Page 258 =====

(86) $\psi: \mathcal{A}_x \to \mathcal{G}$

et (85) est l'équivalence de catégories si $\mathcal{G}$ est encore connexe et $\pi_1(X) \simeq \mathbf{Z}$, et si, pour des points $a, b$ distincts de $\mathcal{S}_x$, définis dans $\mathcal{A}_x$ par des arcs $A$ et $B$ d'extrémités $a, b$, définissant l'élément de $\pi_1(X, a)$ dans $\mathcal{G}(x)$, i.e. dans $\mathcal{G}(\mathcal{A}_x)$, on a
$$
\{a\} \xrightarrow{A} \{b\} \xrightarrow{B} \{a\} \leqno{(86bis)}
$$
l'élément correspondant de $\psi(\{a\})$ dans $\mathcal{G}$ est une équivalence de $\\operatorname{Aut}(\{a\})$.

Nous allons donc définir un foncteur
$$
\mathcal{A}_x \to \operatorname{Rev}_{\text{ét}}(D_x^*) \leqno{(87)}
$$
définissant le revêtement élémentaire. Pour définir un tel foncteur il suffit de définir un foncteur
$$
\mathcal{A}_x \to \text{catégorie des germes } 1\text{-connexes d'espaces au-dessus du germe } D_x^* \leqno{(88)}
$$

\newpage

%% ===== Page 259 =====

$$
\dots \text{ foncteurs sur } \dots
$$

$$
(89) \quad \text{cat. des germes à lacets} \to \underline{\text{Rev}}_\text{alg}(D^*) \text{ au-dessus du germe } D^*
$$

En fait, $(88)$ se traduit : le foncteur $\mathcal{A} \in \mathcal{A}_n \mapsto \text{germe en } x \in X \text{ tendant vers } x \text{, de } \dots \text{ germe de } X_{\mathcal{A}}$ \dots

-- prendre le groupe résiduel

Le foncteur $\mathcal{A} \in \mathcal{A}_n$ s'identifie idéalement aux classes de $\pi_1^\text{top}(X)$, distinctes de $\{0\}$ et des lacets non triviaux sur $\mathcal{C}_{\mathcal{A}}$ l'un de ces cônes. Considérons un isom. local
$$
(90) \quad f : (\pi_1^\text{top}(X), 0) \isom (X, x), \text{ de l'appl. } f \text{ dont l'impl. l'identité. Le groupe } \mathcal{C} \in \mathcal{A}_n \text{, son image par } \dots \text{ tel isom. dépend du choix de l'iso. local } \dots \text{ bien défini ainsi. Mais considérons la } \dots \mathcal{C} \subset \mathcal{A}_n \text{, avec } \mathcal{C} \text{ strictes } \dots
$$
Car, correspondant : des arcs $A'$ qui n'ont des voisinages de $\mathcal{A}$, on a encore tel que $\mathcal{C} - \{0\}$ sont un voisinage de $\mathcal{C} - \{0\}$, dans le

\newpage

%% ===== Page 260 =====

L'image de (90) définit pour le filtre sur $\mathcal{T}_{g,\nu}(X)$ formé des $V_n(C \setminus \{0\})$ : $V$ formant les voisinages de $0$ dans $\mathcal{T}_{g,\nu}(X)$, et un filtre sur $X \setminus \{0\}$ qui ne dépend plus des choix de $f$, et qui est donc défini un groupe $\widetilde{A} \cong \widetilde{C}$ d'espaces de recouvrement de $\mathbb{D}^*$. Ce groupe d'espaces de recouvrement est défini par $\mathcal{T}_{g,\nu}(X)$ muni des filtres et bien des $V_n(C' \setminus \{0\})$ : pour des $V$ assez fins si $V_n(C' \setminus \{0\})$ contient donc, de... Puis $A \subset C$ on ordonne... Puis $A_1 \subset A_2$ i.e. $C_1 \subset C_2$, $\tilde{C}_1 \to \tilde{C}_2$ et la relation correspondante d'inclusion sur les groupes d'espaces de $\mathbb{D}^*$. On \leqno{(81)} trouve donc bien (88), d'où (87), il vérifie que ce dernier a bien un l'équivalence des catégories groupoïdes, via la vérification par l'explicitation plus haut ((86 bis)).

\newpage

%% ===== Page 261 =====

\vspace*{1em}
778

Considérons maintenant, à nouveau, un difféomorphisme $\gamma$ de $X$ en $x$, et sa demi-droite $\Delta$. Alors $\{x\} \in A_x$, et on peut considérer le germe d'espace $\{x\}^\sim$ au-dessus de $\mathbf{D}^*$. On a des morphismes d'induction canoniques
$$ \gamma \to \{\Delta\} \to \mathbf{D}^* \leqno{(91)} $$
et on en déduit un isomorphisme canonique
$$ \operatorname{Rev}(\mathbf{D}^*_x, \gamma) \simeq \operatorname{Rev}(\mathbf{D}^*, \{x\}^\sim) \leqno{(92)} $$
$$ \parallel \operatorname{déf} \qquad \qquad \parallel \operatorname{diff} $$
$$ \mathbf{D}^*_x \qquad \qquad \Theta(\Delta) $$

Si $\gamma'$ est un autre diff.\ en $x$, avec une demi-droite $\Delta'$, on a
$$ \mathbf{D}^*_{x, \gamma} = \Theta(\Delta), \quad \mathbf{D}^*_{x, \gamma'} = \Theta(\Delta') $$

d'où
$$
\begin{tikzcd}
\operatorname{Isom}(\mathbf{D}^*_{x, \gamma}, \mathbf{D}^*_{x, \gamma'}) \arrow[r, "\simeq"] & \operatorname{Isom}(\Theta(\Delta), \Theta(\Delta')) \\
\parallel \operatorname{diff} \arrow[u] & \uparrow \Theta \arrow[u] \\
\operatorname{Isom}(\gamma, \gamma') & \operatorname{Isom}(\Delta, \Delta')_{\pi_1(S_n)}
\end{tikzcd}
$$
ce qui est la bijection (80) cherchée.

Soit $E$ niveau, un diagramme d'équivalences de Galois

\newpage

%% ===== Page 262 =====

(93)
\begin{tikzcd}
germ\, revdiff(X,x) \ar[r, "\simeq"', "\gamma \mapsto \Delta(\gamma)"] \ar[d, "\gamma \mapsto Revunis(D_x^*,\gamma) = D_{x,\gamma}^{*\sim}"'] & \Pi_1(S_x) \ar[d, "\simeq", "a \mapsto Revunis(S_x,-)"'] \\
Rev.unis(D_x^*) \ar[r, "\simeq"'] \ar[ur, dashed, "\Theta \simeq"'] & Revunis(S_x)
\end{tikzcd}

[MARGIN: Bien sûr, la difficulté en fait provenait de ce qu'on n'avait point d'interpr. géométrique explicite entre [unclear] l'espace $S_x$ ...]

Le seul foncteur qui n'ait pas une
explication est le foncteur
(94)
\begin{tikzcd}
Rev.unis.(D_x^*) \ar[r, "\simeq"] \ar[dr, "\Theta^{-1}"'] & Revunis(S_x) \\
& \Pi_1(S_x) \ar[u, "can"']
\end{tikzcd}

défini comme composé de $\Theta^{-1}$ avec
le foncteur can. $\Pi_1(S_x) \to Revunis(S_x)$.

Finalement, dans ces considérations
de géométrie topologico-différentielle, la
construction-clef est celle du foncteur
$\Theta$ (81). Nous avons pris quelques peines
à la définir de façon invariante. Dans
le cas où on dispose d'une iso locale
(95) $$f \colon (T_x(X),0) \to (X,x),$$ la
construction de $\Theta$ peut s'expliciter comme
une construction canonique dans un
espace vectoriel $V$. Toute demi-droite [unclear].

\newpage

%% ===== Page 263 =====

$\Delta$ de $V$ définit un germe d'arc (diff) sur $V$ au point $x$, d'où un revêtement universel $\widetilde{D}_{V,x}$ de $D_{x}^*$, donc un revêtement universel $\widetilde{D}_{x, \Delta}$ de $D_x^*$.

Il faut voir comment $\widetilde{D}_{x, \Delta}$ dépend "continûment" de $\Delta$, et en quoi cela permet de relever $V^*_{\Delta}$ de $V^* = V \setminus \{x\}$ dépendant continûment de $\Delta$, ce qui est en fait évident.

Soit maintenant $U$ un morphisme de $X$, fixant $x$, un difféomorphisme en $x$. On a donc des isomorphismes de $^u D_x^*$ vers $D_x^*$, et $^u S_x$ vers $S_x$.
$$
^u D_x^* \cong D_x^*, \quad ^u S_x \cong S_x \leqno (95)
$$
Considérons d'autre part un revêtement universel $\widetilde{S}_x$ de $S_x$, et soit $\widetilde{D}_x^*$ un revêtement universel de $D_x^*$ associé, qu'on a l'habitude (94). Je dis qu'en choisissant un revêtement universel $\widetilde{D}_x^*$ de $D_x^*$, celui-ci définit un germe de revêtement universel $\widetilde{D}_{V,x}$ de $D_x^*$ (et donc de $S_x$) au voisinage de $x$.

\newpage

%% ===== Page 264 =====

\marginpar{La construction fait suite de la suite de (94) (cf. p. 162, p. 163)}

\noindent bijection canonique $\dots$ des relèvements $\stackrel{\sim}{\to}$ des relèvements de $u_{D_n^*}$ de $u_{\mathfrak{S}_n}$ des groupes $D_n^{*\sim} \dots \mathfrak{S}_n^\sim$
$$
\leqno{(95)}
$$

\noindent bijection qui suit un isomorphisme de $\mathfrak{T}_n \to \dots$ (où $\mathfrak{T}_n$ est le $\mathbf{Z}$-module d'automorphismes de $X$, n. [unclear: discr.]).

\marginpar{E) Axiomatisation : de la construction avec un $\tilde{D}_n^*$ (c'est un groupe qui est le filtre des voisinages en $X$)}

\noindent J'ai envie de donner une description axiomatique de (94) et (96).

\begin{enumerate}
\item[1)] Soit $(X, n)$ un système de surfaces affinitables, $V_{x, n} = V_n = \mathfrak{T} = \mathfrak{T}_{g, n}(X)$, et $\mathfrak{S}_n = V^*/R^* (\dots)$ ainsi que $D_{x, n}^* = D^*$. Soit $B_n$ (déf. par le $\tilde{D}_n^*$) des rev. universels de $D_n^*$, en fait $\in \mathcal{A}_{x, n}$ pour l'un des arcs de $\mathfrak{S}_n$, un rev. bien choisi (c'est le gros [unclear: fini]?) des rev. bien choisis, $\dots$
\end{enumerate}

\noindent $(+ V \text{ et } \{0\})$. Posons $A \in \mathcal{A}_n \dots C \in \mathcal{C}$, disons par $\dots$ l'équi. $\dots$ de $D^*$. Si $\tilde{S}$ est un rev. univ. de $\tilde{D}_n^*$, on considère $\dots$

\newpage

%% ===== Page 265 =====

(97) 
\begin{tikzcd}
\Gamma(A, \tilde{S}_x/S_x) \ar[d, equal] & , & \Gamma(E_p(A), \tilde{D}_x^*/D_x^*) \ar[d, equal] \\
\operatorname{Hom}_{S_x}(A, \tilde{S}_x) \ar[d, equal, "\text{déf}"'] & & \operatorname{Hom}_{D_x^*}(E_p(A), \tilde{D}_x^*) \ar[d, equal, "\text{déf}"'] \\
\tilde{S}_x(A) & & \tilde{D}_x^*(E_p(A))
\end{tikzcd}

qui sont l'un et l'autre des Torseurs sur le groupe $T_{X,x} = T_x$ le $\mathbb{Z}$-module des orientations de $X$ en $x$ (pour $\sim$). Considérons l'ens

(98) $$H(A) = Isom_{T_x}(\tilde{S}_x(A), \tilde{D}_x^*(E_p(A)))$$

qui est lui-même un $T_x$-torseur. Soit en effet une inclusion

$$A \subset A'$$

on en déduit, par restriction, des isomorph. de $T_x$-torseurs

(99) $$\tilde{S}_x(A') \simeq \tilde{S}_x(A), \quad \tilde{D}_x^*(E_p(A')) \simeq \tilde{D}_x^*(E_p(A))$$

d'où un isom. canonique

(100) $$H(A) \simeq H(A')$$

de sorte que 

(101) $$A \mapsto H(A), \text{ foncteur } A_x \xrightarrow{H} \underline{Tors}(T_x)$$

où $A_x$ est la catégorie ordonnée [unclear] associée à l'ens. ordonné $A_x$. On trouve donc par factorisation un foncteur de la catégorie groupoïde associé $Gr(A_x)$ de $A_x$ dans $\underline{Tors}(T_x)$,

(102) 
\begin{tikzcd}
A_x \ar[d] \ar[r, "H"] & \underline{Tors}(T_x) \\
Gr(A_x) \ar[ru, "Gr(H)"'] & 
\end{tikzcd}

\newpage

%% ===== Page 266 =====

Je dis que $\mathcal{G}(\mathcal{A})$ est isom. un $T_{\pi}$-torsor, de sorte, de façon précise le foncteur
$$
\begin{cases}
\mathcal{A} \to \underline{\text{Rev.univ.}}(\mathfrak{S}_{\alpha}) \defeq \mathfrak{S}_{\alpha}\mathcal{A} \\
\mathcal{A} \to \underline{\text{Rev.univ.}}(\mathfrak{S}, \mathcal{A})
\end{cases} \leqno{(103)}
$$
qui factorise un foncteur, je précise
$$ \mathcal{G}(\mathcal{A}) \isom \underline{\text{Rev.univ.}}(\mathfrak{S}_{\alpha}) \leqno{(104)} $$
\marginpar{ici $\mathcal{A}, \mathcal{A}' \in \mathcal{C}$, $\operatorname{Hom}_{\mathcal{G}(\mathcal{A})}(\mathcal{A}, \mathcal{A}') \isom \underline{\text{Isom}}_{\text{Rev}(\mathfrak{S}_{\alpha})}(\mathfrak{S}, \mathcal{A})$}
qui est équivalence de catégories, (104), de...
$$ \operatorname{Hom}_{\mathcal{G}(\mathcal{A})}(\mathcal{A}, \mathcal{A}') \isom \operatorname{Hom}(\mathfrak{S}_{\alpha}\mathcal{A}, \mathfrak{S}_{\alpha}\mathcal{A}') \leqno{(105)} $$
Ceci rappelle ce que le foncteur $\mathcal{G}(\mathcal{A})$ du (102) a établi ; il définit sur $\pi_1$ l'image et $T_{\pi}$... qui vient compléter les isom. (100) en $H(\mathcal{A})$, un système transitif d'isom.

Plusieurs... On désigne par le «colim» du $H(\mathcal{A})$ qui dépend des deux $\mathfrak{S}_{\alpha}$ et $D_{\pi}^{\sim}$, par $H(\mathfrak{S}_{\alpha}, D_{\pi}^{\sim})$, structure définie par des isom. de $\text{Tors}(T_{\pi})$.
$$ i_{\mathcal{A}}: H(\mathcal{A}) \to H(\mathfrak{S}_{\alpha}, D_{\pi}^{\sim}) \leqno{(106)} $$
satisfaisant la condition de compatibilité, s'exprimant...

\newpage

%% ===== Page 267 =====

781

-ment par la commutation des diagrammes

(107)
\begin{tikzcd}
H(A) \ar[r, "i_A", "\sim"'] \ar[d, "(100)"'] & H(\tilde{S}_\alpha, \tilde{D}_\alpha^*) \\
H(A') \ar[ur, "\sim"', "i_{A'}"] &
\end{tikzcd}

Pour vérifier le fait que le foncteur $Gr(H)$
de (102) est [unclear: bien consistant], je considère les
deux foncteurs (dépendant du choix de
$\tilde{S}_\alpha, \tilde{D}_\alpha^*$)

(108)
$$ \begin{cases} 
A \longmapsto \tilde{S}_{\alpha,A} \quad , \quad A \longmapsto \tilde{D}_{\alpha,A}^* \\ 
\underline{A}_\alpha \longrightarrow \underline{Rev}_{univ}(S_\alpha) \quad \underline{A}_\alpha \longrightarrow \underline{Rev}_{univ}(D_\alpha^*) 
\end{cases} $$

qui se factorisent en des équivalences

(109)
$$ Gr(\underline{A}_\alpha) \xrightarrow{\sim} \underline{Rev}_{univ}(S_\alpha) \quad , \quad Gr(\underline{A}_\alpha) \xrightarrow{\sim} \underline{Rev}_{univ}(D_\alpha^*) $$
utilisant $\tilde{S}_\alpha, \tilde{D}_\alpha^*$.

Ceci dit, pour un $A \in Ob \, \underline{A}_\alpha = A_\alpha$ [unclear: donné],
des isom. fonctoriels en $A$

(110)
$$ \begin{cases} 
[unclear](A) \simeq \underline{Isom}_{\underline{Rev}_{univ}(S_\alpha)}(\underbrace{F(A)}_{\tilde{S}_{\alpha,A}}, \tilde{S}) \\ 
[unclear](E,A) \simeq \underline{Isom}_{\underline{Rev}_{univ}(D_\alpha^*)}(\underbrace{G(A)}_{\tilde{D}_{\alpha,A}^*}, \tilde{D}_\alpha^*) 
\end{cases} $$

et par la propriété univ. de $Gr(\underline{A}_\alpha)$,
ces isom. s'étendent sur la
catégorie $Gr(\underline{A}_\alpha)$. Dans la situation où
elle se réduit à un groupe profini $T (= T_\alpha)$
trois $T$-groupoïdes connexes $g, g', g''$

\newpage

%% ===== Page 268 =====

(suite des équivalences des $T$-groupoïdes)
$$ g \cong g', \quad g \cong g'' \leqno{(111)} $$
et des objets $X' \in \mathcal{G}', X'' \in \mathcal{G}''$ (à savoir $\mathbb{S}_n$ et $\mathbb{D}_n^*$), et en considérant les foncteurs
$$ \begin{cases} g \xrightarrow{f=f_n} \mathcal{T}ors(T), \quad & g \xrightarrow{g=g_n} \mathcal{T}ors(T) \\ X \mapsto \underline{\mathcal{H}om}(F(X), X'), \quad & X \mapsto \underline{\mathcal{H}om}(G(X), X'') \end{cases} \leqno{(112)} $$
et le foncteur
$$ \begin{cases} g \to \mathcal{T}ors(T) \\ X \mapsto \underline{\mathcal{H}om}_{\mathcal{T}ors(T)}(f(X), g(X)) \end{cases} \leqno{(113)} $$
à prouver que ce dernier foncteur est constant. On a, si $F', G'$ désignent
$$ F' \cong g, \quad G' \cong g \leqno{(114)} $$
diagrammes des quasi-isomorphismes $F, G$, donc
$$ \begin{cases} f(X) \defeq \underline{\mathcal{H}om}_g(F(X), X') \cong \underline{\mathcal{H}om}_g(X, F_1(X')) \\ g(X) \defeq \underline{\mathcal{H}om}_g(G(X), X'') \cong \underline{\mathcal{H}om}_g(X, F_2(X'')) \end{cases} \leqno{(115)} $$
d'où l'existence des isomorphismes de $T$-torsors
$$ \underline{\mathcal{I}som}(f(X), g(X)) \cong \underline{\mathcal{I}som}(F_1(X'), F_2(X'')) \leqno{(116)} $$
d'où la vérification.

\newpage

%% ===== Page 269 =====

7P2

Définition. — On dit qu'on a une corrélation entre un rev. univ. $\tilde{S}_x$ de $S_x$, et un rev. univ. $\tilde{D}_x^*$ de $D_x^*$, quand on s'est donné un élément de $H(\tilde{S}_x, \tilde{D}_x^*)$.

Donc, si $A \in A_a$, la donnée d'une corrélation revient à la donnée d'un élément de $H(A)$, i.e. d'un isomorphisme entre les $T_a$-torseurs $\tilde{S}_x^\sim(A)$ et $\tilde{D}_x^{*\sim}(E_p(A))$. Cette identification est compatible avec les flèches d'inclusion entre $A \subset A'$. Les corrélations entre $\tilde{S}_x$ et $\tilde{D}_x^*$ forment un torseur sous $T_a$. Pour $\tilde{S}_x$ fixé

Soit $Corr_x(X) = Cour_x(X)$ la catégorie des triplets $(\tilde{S}_x, \tilde{D}_x^*, h)$ de rev. univ. corrélés de $S_x$ et de $D_x^*$, on a des foncteurs $(\tilde{S}_x, \tilde{D}_x^*, h) \rightrightarrows {\tilde{S}_x \atop \tilde{D}_x^*}$

(117)
\begin{tikzcd}
& Corr_x(X) \ar[dl, "\simeq"'] \ar[dr, "\simeq"] & \\
Revuniv(S_x) & & Revuniv(D_x^*)
\end{tikzcd}

qui sont des équivalences de catégories.

\newpage

%% ===== Page 270 =====

Toutes ces constructions ont été faites de façon « intrinsèque » ou simplement par paires de germes de surfaces différentiables $(X, \alpha) \to (X', \alpha')$. Supposons choisis, en plus des $(\tilde{S}_x, \tilde{D}_x^*)$ relatifs à $(X, \alpha)$, les $(\tilde{S}'_{x'}, \tilde{D}'^*_{x'})$ relatifs à $(X', \alpha')$.

À l'iso. de germes est donné, il induit donc :
$$
\begin{cases} u_S : \tilde{S}_x \isom \tilde{S}_{x'} & \text{iso. d'esp. top.} \\ u_{D^*} : \tilde{D}_x^* \isom \tilde{D}_{x'}^* & \text{iso. des germes d'esp. top.} \\ u_T : T_x \isom T_{x'} & \text{iso. des } \mathbf{Z}\text{-modules libres de rk } 1 \end{cases} \leqno{(118)}
$$
considérons les deux ensembles
$$
\begin{cases} F(u) \defeq \text{ens. des iso. } \tilde{S}_x \isom \tilde{S}_{x'} \text{ au-dessus de } u_S \\ G(u) \defeq \text{ens. des iso. } \tilde{D}_x^* \isom \tilde{D}_{x'}^* \text{ au-dessus de } u_{D^*} \end{cases} \leqno{(119)}
$$
ce sont l'un et l'autre des $T_x$-torseurs.

Ceci étant, je dis qu'on a un iso. can. de $T_x$-torseurs
$$
F(u) \isom G(u) \leqno{(120)}
$$
— i.e. "il existe" un relevé $u_S \in F(u)$ et un relevé $u_{D^*} \in G(u)$. Pour le voir, on utilise que par fonctorialité des constructions les...

\newpage

%% ===== Page 271 =====

78s

images inverses par $u_S$ et par $u_{D^*}$ de $\tilde{S}_{x'}^\sim$ et de $\tilde{D}_{x'}^{* \sim}$
soient $\tilde{S}_x^{\prime \sim}$, $\tilde{D}_x^{* \prime \sim}$, sont munis d'une covariation
que je désigne par $h_x^\prime$. Ceci posé, les theoremes
d'iso. de $T_x$-torseurs canoniques

$$ (121) \qquad F(u) \simeq Isom_{\tilde{S}_x}(\tilde{S}_x^\sim, \tilde{S}_x^{\prime \sim}) $$
$$ G(u) \simeq Isom_{\tilde{D}_x^*}(\tilde{D}_x^{*\sim}, \tilde{D}_x^{* \prime \sim}) $$

Or par le diagramme d'équivalence (117) on
a bien un iso de $T_x$-torseurs
$$ Isom((\tilde{S}_x^\sim, \tilde{D}_x^{*\sim}, h_x), (\tilde{S}_x^{\prime \sim}, \tilde{D}_x^{* \prime \sim}, h_x^\prime)) $$

(122)
\begin{tikzcd}
& Isom((\tilde{S}_x^\sim, \tilde{D}_x^{*\sim}, h_x), (\tilde{S}_x^{\prime \sim}, \tilde{D}_x^{* \prime \sim}, h_x^\prime)) \ar[dl] \ar[dr] & \\
F(u) \simeq Isom(\tilde{S}_x^\sim, \tilde{S}_x^{\prime \sim}) & & G(u) \simeq Isom(\tilde{D}_x^{* \sim}, \tilde{D}_x^{* \prime \sim})
\end{tikzcd}

d'où la bijection canonique cherchée (120) —
[unclear: des qui donne un justification (96).]
La bijection canonique (120) a la
propriété $\pm$ évidente suivante, qui la caractérise:

(123) $\begin{cases}
\text{Soient } u_{\tilde{S}} \text{ et } u_{\tilde{D}^*} \text{ des prolongements se correspondants} \\
\tilde{S}_x^\sim \xrightarrow{\sim} \tilde{S}_{x'}^\sim, \ \tilde{D}_x^{*\sim} \xrightarrow{\sim} \tilde{D}_{x'}^{* \sim}, \text{ associés par (120), soit } A \\
\in S_x, \ A' \in S_{x'} \text{ son image par } u_S \text{, soit } \alpha \text{ une} \\
\text{section de } \tilde{S}_x^\sim \text{ sur } A, \ \beta \text{ la section de } \tilde{D}_x^{*\sim} \text{ sur } Ep(A) \\
\text{associé par la covariation } h_x, \text{ soit } \alpha' \text{ la section} \\
\text{de } \tilde{S}_{x'}^\sim \text{ sur } A' \text{ déduite de } \alpha \text{ en appliquant } u_{\tilde{S}}, \\
\beta' \text{ la section de } \tilde{D}_{x'}^{* \sim} \text{ sur } Ep(A') \text{ déduite de } \beta \\
\text{en appliquant } u_{\tilde{D}^*}, \text{ alors } \alpha' \text{ et } \beta' \text{ sont associés} \\
\text{par la covariation } h'.
\end{cases}$

\newpage

%% ===== Page 272 =====

Ceci confirme que $A$ (fixé) contient l'application $(120)$, car pour $u_{\alpha}$ fixé, il détermine $u_{\beta}$ ainsi : $\alpha \to \beta$ par $h_{\ast}$, $\alpha'$ provient de $\alpha$ par $u_{\alpha}$, $\beta'$ déduit de $\alpha'$ par $h_{\ast}$, on $\beta'$ doit être égal à la section déduite de $\beta$ par l'application $u_{\beta} \sim \mathcal{D}^{\sim} \to \mathcal{D}^{\sim}$, et ce détermine bien $u_{\beta}$ de façon unique...

NB : bien que l'utilisation de cette construction, pour $A$ fixé, vienne définir l'application $(120)$, à travers un $\mathbb{T}$-morphisme — tracer que $\dots = \text{conj.}$, qu'elle est celle de $(120)$ et, qu'elle dépend de $\alpha$. Mais a priori il dépend peut-être des choix de $A$, et le point qui demande vérification est que cette bijection est en fait indépendante des choix de $A$. On voit en effet que pour $A \subset A'$, elle est $\mathbb{T}$-funtorielle pour $A \to A'$ d'où... tot le résultat, ou la collection d'équivalences dû à $A$ par le rel. d'équivalence de la collection....

\newpage

%% ===== Page 273 =====

784

F) Application: interprétation différentielle des groupes de Teichmüller spéciaux.

Soit $X$ un espace topologique, muni d'une partie $I$ discrète, telle que $X$ soit un espace de voisinages de $I$. En plus de la donnée $(C_I)$, on se donne une structure différentielle sur les voisinages de $I$. Soit $G$ un groupe topologique, opérant continûment sur $(X, I)$, et respectant la structure différentielle de $X$ et voisinages de $I$. Soient
$$
\begin{cases}
\mathfrak{S}_{X, I} = \coprod_{x \in I} \mathfrak{S}_{X, x} \\
\mathfrak{D}_{X, I} = \coprod_{x \in I} \mathfrak{D}_{X, x}
\end{cases} \leqno{(124)}
$$
On suppose que l'opération de $G$ sur $\mathfrak{S}_X$ induite par l'action de $G$ soit égale et continue.

Soit, pour $x \in I$, $\mathfrak{D}_x^*$ un sous-ensemble de $\mathfrak{D}_x^*$, et posons
$$
\mathfrak{D}_{X, I}^* = \coprod_{x \in I} \mathfrak{D}_x^* \leqno{(125)}
$$
Considérons pour l'instant les données différentiables, et considérons les groupes
$$
\tilde{G} = \{ (u, \tilde{u}) \mid u \in G, \tilde{u} \text{ un relèvement de } u \text{ sur } \mathfrak{D}_{X, I}^* \} \leqno{(126)}
$$
On a un homomorphisme $(\tilde{u}, u) \to \tilde{G}$, puis...

\newpage

%% ===== Page 274 =====

\newpage
\setcounter{page}{273}

\noindent et, pour
$$
\mathfrak{T}_{\mathcal{I}} = \prod_{x \in \mathcal{I}} \mathfrak{T}_x \leqno{(127)}
$$
on a une suite exacte
$$
1 \to \mathfrak{T}_{\mathcal{I}} \to \tilde{G} \to G \to 1. \leqno{(128)}
$$
Je voudrais introduire sur $\tilde{G}$ une topologie telle que (128) devienne une suite exacte de groupes topologiques.
Quitte à remplacer $G$ par le sous-groupe $G'$ d'indice fini moyen de l'homomorphisme continu
$$
G \to \mathfrak{S}_{\mathcal{I}}, \leqno{(129)}
$$
définir $\tilde{G}$ par l'image inverse $\tilde{G}'$ de $G'$,
on voit que, car $G$ opère finiment sur $\mathcal{I}$, et quitte à moyenner,
on peut, pour le groupe $G$ agissant sur les fibres,
considérer les extensions
$$
1 \to \mathfrak{T}_x \to \tilde{G}_x \to G \to 1 \defeq I_x=\{x\}. \leqno{(130)}
$$
Dans ce cas, on a un cas où $\mathcal{I} = 1$.

\newpage

%% ===== Page 275 =====

(131) Soit $\mathcal{D}_{x,0}$ une base discrète de voisinages de $x$ dans $X$, qui soit un filtre des voisinages de $x$ dans $X$, et soit on pose un revêtement universel $\widetilde{\mathcal{D}}_{x,0}^*$ de $\mathcal{D}_{x,0}^* = \mathcal{D}_{x,0} \setminus \{x\}$ qui induit $\widetilde{\mathcal{D}}_{x,0}^*$ sur $\mathcal{D}_{x,0}^*$. Par hypothèse de continuité, il existe un voisinage $U$ de $e$ dans $G$, et un voisinage $\Delta_{x,1}$ de $x$, tels que l'on ait
$$
\begin{tikzcd}
U \times \Delta_{x,1} \arrow[r] \arrow[d] & \Delta_{x,0} \arrow[d] \\
G \times X \arrow[r] & X
\end{tikzcd}
$$
(132) et tel que $\Delta_{x,1}$ soit un disque ouvert. Ceci permet, pour $u \in U$, de définir le relèvement $\tilde{u}$ de $u$ :
$$
\begin{tikzcd}
\widetilde{\Delta}_{x,1} \arrow[r] & \widetilde{\Delta}_{x,0}
\end{tikzcd}
$$
(133) compatible avec $\Delta_{x,2}^* \to \Delta_{x,0}^*$.
Cela dit, que l'image inverse $\widetilde{U}$ des voisinages $U$ de $e$ dans $G$, ...

\newpage

%% ===== Page 276 =====

1-traduire la topologie de la convergence compacte de $\widetilde{U}$ et de ... les applications :
$$
\begin{cases}
p: \widetilde{U} \to U \\
q: \widetilde{U} \to \operatorname{Cont}(\Delta_{x,1}^{\sim}, \Delta_{x,0}^{\sim}), \quad (u, \tilde{u}) \mapsto \tilde{u}_{0,1}
\end{cases}
\leqno{(134)}
$$
$$ q' = q \circ \sigma, \quad \sigma(u,\tilde{u}) = (u,\tilde{u})^{-1} : \widetilde{U} \to \operatorname{Cont}(\Delta_{x,1}^{\sim}, \Delta_{x,0}^{\sim}) $$

On peut, puis on peut prendre $\widetilde{U}$ stable par $u \mapsto u^{-1}$, ce qui implique qu'il est un ouvert de $\widetilde{U}$.

On trouve donc : à l'aide de cette topologie sur $\widetilde{U}$, le filtre des voisinages de 1 dans $\widetilde{U}$, le filtre des voisinages de 1 dans $\widetilde{G}$. Il faut vérifier que ce dernier est bien le filtre des voisinages de 1 pour une topologie de groupe, et que pour celle-ci, $\widetilde{G} \to G$ est un revêtement de $G$, et que ... que tout marche.

Quand $X$ est compact, orientable, connexe, et que $\mathfrak{T}$ désigne un modèle des...

\newpage

%% ===== Page 277 =====

orientations, alors tous les $T_\alpha$ sont canoniquement isomorphes : $T$
$$
T_\alpha \isom T \quad \forall \alpha \in \Gamma \leqno(135)
$$
et la suite exacte (128) pour le groupe
$$
1 \to T^I \to \widetilde{G} \to G \to 1 \leqno(136)
$$
Le cas ``universel'' est celui où
$$
G = \\operatorname{Aut}_{\text{top}}(X, S) \leqno(137)
$$
(avec la topologie de la convergence uniforme et de $u$ et de $u^{-1}$). Alors dans les cas anabéliens (type de $(X, S)$ tel que $2g-2+\nu > 0$, i.e., $2g+\nu \geq 3$, i.e., $(g,\nu) \neq (0,0), (0,1), (0,2), (1,0)$), on sait que $A^0$ est contractible, donc $A$ est un groupe discret, et l'extension $\widetilde{A}$ de $G$ provient d'une extension
$$
1 \to T \to \widetilde{A} \to G \to 1 \leqno(139)
$$
du groupe de Teichmüller
$$
\mathfrak{T} = \mathfrak{T}(X,S) = A^0 = \pi_0(\\operatorname{Aut}_{\text{top}}(X, S)). \leqno(140)
$$

\newpage

%% ===== Page 278 =====

qui prend ainsi le nom de groupe de
Teichmüller spécial de $(X,S)$. Pour
$G$ quelconque et action de $G$ sur
$(X,S)$ orientée : la donnée d'un hom.
de groupes topologiques
$$G \to \mathcal{A}$$
s'intègre dans un diagramme de
groupes top.

(141)
\begin{tikzcd}
G \ar[r] \ar[d] & \mathcal{A} \ar[d] \\
\pi_0(G) \ar[r] & \tau = \pi_0(\mathcal{A})
\end{tikzcd}

et on trouve alors, par fonctorialité,
que le $\tilde{G}$ construit précédemment
est l'image inverse de l'ext. $\tilde{\tau}$ de $\tau$
par $\mathbb{Z}^I$, via $G \to \tau$.

Ces préliminaires étant posés,
revenons à la situation où $X$
admet une structure diff. $C^\infty$ au voisin-
age de $S$, et soit pour tout $x \in S$,

\newpage

%% ===== Page 279 =====

$\tilde{\mathfrak{S}}_x$ «lieu» ou universel de $\mathfrak{S}_x$ [unclear: attaché] au rev. universel $D_x^\sim$ de $D_x^*$. Posons
$$
\tilde{\mathfrak{S}}_I = \bigsqcup_{x \in I} \tilde{\mathfrak{S}}_x \leqno{(142)}
$$
et notons que l'on a, pour $u \in \mathcal{G}$, une correspondance biunivoque entre l'ens. des relèvements $\tilde{D}_x^\sim$ de $D_x^*$ en ceux de $\tilde{D}_x^\sim$, et l'ens. des relèvements $\tilde{\mathfrak{U}}_I$ de $u$ en autom. de $\tilde{\mathfrak{S}}_I$, qui sont donc des $\mathfrak{T}$-tomes. Cette correspondance est compatible avec les compositions, et on a ainsi une isom. canonique
$$
\tilde{\mathcal{G}} \simeq \{ u \in \mathcal{G}, \tilde{u}_I \text{ rel. de } u \text{ sur } \tilde{\mathfrak{S}}_I \} \leqno{(143)}
$$
On vérifie que l'on a bien, top. en $\mathcal{G}$, sur la deuxième part le top. de la réunion finie ouverte. Continuons les applications
$$
\begin{aligned}
&\varphi_1: \tilde{\mathcal{G}} \to \mathcal{G} \\
&\varphi_2: \tilde{\mathcal{G}} \to \\operatorname{Aut}_{\text{top}}(\tilde{\mathfrak{S}}_I) \quad (u, \tilde{u}) \mapsto \tilde{u}
\end{aligned} \leqno{(144)}
$$

\newpage

%% ===== Page 280 =====

\dots \\operatorname{Aut}_{top}(\tilde{\mathfrak{S}}_I) et non de la composante compacte.

Dans le cas où $X$ est une surface orientable, on trouve donc une interprétation des groupes précédents explicite en $\mathfrak{S}_{g, \nu}$, une effectivité (comme image du groupe de Teichmüller spécial $\mathfrak{S} \to \mathfrak{T}$, par l'homomorphisme canonique de $G \to \mathfrak{T}$ des (141)).
